\documentclass[../main.tex]{subfiles}
\begin{document}
%==========================================TIÊU ĐỀ TẬP========================================
\begin{table}[h]
    \begin{adjustwidth}{-1cm}{}
        \begin{tabular}{>{\centering\arraybackslash}p{6cm}|>{\raggedright\arraybackslash}p{14cm}}
            \multirow{5}{6cm}
            % Cột bên trái: ÉPISODE và số tập
            {\\[-40pt]
                \flaregothic\fontsize{20pt}{20pt}\selectfont \centering
                \color{eptype}HISTOIRE\color{black}\\
                \brian\fontsize{72pt}{72pt}\selectfont
                \color{epnum}5
                \\[18pt]
                % \flaregothic\fontsize{14pt}{14pt}\selectfont
                % \color{epattr}BIENVENUE, \uline{\color{Banpire} Mel} !
                % \flaregothic\fontsize{18pt}{18pt}\selectfont
                % \color{epattr}ÉPISODE PERDU
                \color{black}
            }
             &
            % Dòng thứ nhất: tên tập tiếng Nhật
            {\fotskip\fontsize{18pt}{18pt}\selectfont お\ruby{誕}{たん}\ruby{生}{じょう}\ruby{日}{び}おめでとう、\ruby{火}{ひ}\ruby{川}{かわ}\ruby{久}{く}\ruby{遠}{おん}\ruby{先}{せん}\ruby{輩}{ぱい}！} \\[6pt]
            % Dòng thứ hai: tên tập tiếng Anh
             & {\cafeteria\fontsize{18pt}{18pt}\selectfont Happy Birthday, Dear Kuon-senapi!}                                                                             \\[4pt]
            % Dòng thứ ba: tên tập tiếng Việt
             & {\vnmsans\fontsize{18pt}{18pt}\selectfont Chúc mừng sinh nhật, Kuon-senpai!}                                                                               \\[8pt]
            % Dòng thứ tư: Nhân vật xuất hiện
             & {\sen Track Used: The 4\textsuperscript{th} Anniversary  (\ruby{四}{よん}\ruby{周}{しゅう}\ruby{年}{ねん})}
        \end{tabular}
    \end{adjustwidth}
\end{table}
%===========================================NỘI DUNG==========================================
\color{black}

\pov{Hikawa Kuon}

\begin{center}
    \uline{Thư viện - Đại học Xây dựng Kawachi\\Bốn giờ chiều.}
\end{center}

Trời mùa hè giữa tháng Bảy đúng là nóng thật.

Tôi chống cằm nghe mọi người trong nhóm đồ án bàn bạc về các bài tập sắp tới, tay cầm chai nước suối gần như đã hết sạch, trong khi mồ hôi bắt đầu rịn ra dưới lớp áo phông đỏ nhăn nhúm. Nắng giữa mùa hè thi thoảng lấp ló chiếu thằng vào phòng như thể đang trút giận lên đám sinh viên còn vất vưởng nơi đây - những kẻ như tôi, vì một lý do nào đó, vẫn phải mò mặt lên trường giữa kỳ nghỉ.

Lý do là gì ư? Chỉ đơn giản là\dots\ mấy môn học kỳ tới đang đe dọa giấc mơ tốt nghiệp của cả nhóm. Không phải là giảng viên kêu chúng tôi lên, mà chúng tôi đã tự sắp xếp thời gian để lên trường bàn bạc nhằm chuẩn bị kỹ lưỡng.

``\vie{Vậy rốt cuộc mấy môn chọn tự do kia có bị tính điểm GPA không?}'' - Houjou Tamachi (\ruby{北}{ほう}\ruby{条}{じょう}\ruby{田}{たま}\ruby{町}{ち}, \ruby{Hâu}{Bắc}-\ruby{giâu}{Điều}\ \ruby{Ta-ma}{Điền}-\ruby{chi}{Đinh}), một người bạn trong nhóm đồ án tôi hỏi, vừa lật trang giáo trình in sẵn mà cậu mang theo.

Trong số bốn đứa trong nhóm đồ án làm, cậu bạn này có sức học yếu nhất. Nhưng không phải là vì cậu học kém hay lười biếng gì, mà vì cậu ta quá ham làm để lấy tiền đóng học phí và mua công cụ học tập, đến mức không còn thời gian để học. Nhưng may mắn thay, cậu ta lại chưa trượt một môn học nào. Tôi vẫn không hiểu vì sao lại như thế, nhưng có lẽ là do cậu ta có một chút may mắn trong việc học hành.

Tamachi luôn đeo kính cận, mặc một bộ quần áo bình thường dù có đôi phần hơi xộc xệch và đi dép lê, như tôi và Yamazu.

``\vie{Ủa? T-Tất nhiên là vẫn t-tính chứ!}'' Yamazu đáp lại quả quyết. ``\vie{T-Tôi tưởng đã phổ biến với anh em từ trước rồi mà?}''

``\vie{Nghe đồn là môn Lin*x đánh trượt lên tới 60\% lận\dots}'' tôi thở dài ngán ngẩm. ``\vie{Giáo viên mà chúng ta đã từng học mấy kỳ trước hóa ra cũng chằng dễ chịu gì cho cam cả.}''

``\vie{Chưa kể đến mấy môn nâng cao nữa\dots}'' cô bạn nữ trong nhóm vừa ngán ngẩm lên tiếng kia là Furou Yoisai (\ruby{不}{ふ}\ruby{老}{ろう}\ruby{宵}{よい}\ruby{先}{さい}, \ruby{Phư}{Bất}-\ruby{rô-ư}{Lão}\ \ruby{Gioi}{Tiêu}-\ruby{xai}{Tiên}), cũng đang chống cằm, ngón trỏ gõ gõ xuống bàn như đang thể hiện sự khó chịu.

Trái ngược với cậu bạn Đeo kính, cô ấy lại là một trong những trụ cột chính của nhóm đồ án khi là một ``MVP'' cho các bài đồ án môn học và bài tập lớn cùng với Yamazu mỗi khi tôi và Tamachi bận việc. Cũng bởi vì tôi quá dồn việc ở công ty nên nhiều lúc hai người đó phải ``gánh'' phần lớn công việc trên lớp, trong khi hai người còn lại chúng tôi chỉ có thể hỗ trợ một phần việc nào đó.

Yoisai cũng là một người ``bốn mắt'' (mà nói thực, trong cả nhóm đồ án, có mỗi tôi là kẻ không đeo kính), đồng thời cũng mặc trang phục bình dân như chúng tôi: một bộ áo phông và quần dài, đi giày thể thao, nhưng điều mà tôi cảm thấy khó hiểu là cô ấy luôn đeo khẩu trang, mà tôi đoán là đã hình thành thói quen từ hồi dịch.

``\vie{Quả đúng là Yoisai-san nói nhỉ,}'' cậu bạn Đeo kính thở dài. ``\vie{Năm nhất thì yêu mến trường lắm, năm hai đến năm ba thì chửi trường um sùm lên, để rồi năm cuối lại yêu vì sắp ra trường còn gì.}''

Cả bốn đứa cùng bật cười. Nhưng ngay sau đó, cả bốn lại cùng thở dài. Những tiếng thở dài đầy chua chát khi lời nhận xét nửa đùa nửa thật của Tamachi lại không thể đúng hơn.

``\vie{Đại học có như mơ đâu mà\dots}'' tôi thở ra, tự giễu, rồi uống ngụm nước còn dang dở trong chai nước lạnh đã nguội ngắt từ khi nào. ``\vie{Cũng chính vì kiểu dở dở ương ương nên mới cất công vác cái xác này lên trường để bàn bạc nhau chứ.}''

Ngoài trời, ve sầu kêu ran, tiếng quạt trần quay kẽo kẹt như thể nó cũng đang kêu cứu. Tôi ước rằng chiếc quạt trần đó không hề chạy, vì nó dường như đang phả xuống bầu không khí còn bí bức hơn.

Nhưng, chúng tôi vẫn ngồi lại, vây quanh chiếc bàn trong thư viện tầng hai, chắp vá những bản kế hoạch, lịch trình học, danh sách tài liệu cần đọc, và những lời cảnh báo đỏ rực trên cổng thông tin trường.

Và rồi, như một thói quen, tôi lại nhìn về phía cửa sổ. Cảnh vật bên ngoài vẫn như cũ, những tán cây xanh rì rào trong gió lặng, những chiếc xe đạp đậu bên đường, và những người bạn học đang đi lại trên sân trường.

Ngó xuống nhóm bạn, họ vẫn đang bàn bạc rất sôi nổi về kỳ học khó khằn sắp tới, nơi mà chúng tôi sẽ phải đối đầu với nhiều bài tập lớn, nhiều đồ án hơn, lại còn mang tính tổng hợp hơn.

Và dĩ nhiên, không một ai nhắc đến sinh nhật tôi. Mà, đó cũng chẳng phải là điều gì quá ghê gớm. Họ quên? Tôi cũng chẳng thèm quan tâm.

Với tôi, hôm nay cũng như mọi ngày khác. Nếu có một điều gì đáng bận tâm hơn chuyện sinh nhật, thì đó là cái môn học khỉ ho cò gáy có tỷ lệ đánh trượt cao ngất ngưởng kia.

Nhìn Yamazu hí hoáy ghi chú, tôi đành thở ra lần nữa. Mắt liếc nhanh qua các đề mục đã gạch đầu dòng trong sổ tay, tôi gập nắp bút lại.

Đúng vậy. Với tôi, sinh nhật chỉ là một cái ngày đánh dấu việc mình lại già thêm một tuổi. Và\dots\ chấm hết. Tôi chẳng còn là một đứa trẻ con háo hức mong quà hay mong bánh nữa, mà tôi là sinh viên đại học, và có hàng đống thứ đáng lo hơn.

``\vie{Vậy thì anh em mình chốt nhé,}''  Tamachi nói, giọng hơi trầm nhưng dứt khoát. ``\vie{Môn Lin*x thì chia nhau làm outline sớm, sau đó luyện phần cú pháp code thôi. Ăn được điểm bài thi giữa với cuối kỳ thì chắc còn có suất qua môn.}''

``\vie{T-Tôi cũng sẽ c-chuẩn bị cho anh em giáo trình t-từ các anh chị khóa trước nữa,}'' Yamazu tiếp lời, dù với tính nói lắp của mình nhưng vẫn bày tỏ sự quyết tâm cao độ.

``\vie{Ờ, mà đừng quên mấy cái môn nâng cao nữa,}'' Yoisai chen vào, tay gõ nhanh trên điện thoại. ``\vie{Nghe nói các lớp khóa trước bị bắt thuyết trình hoàn toàn bằng tiếng Anh, không có hỗ trợ slide đấy.}''

``\vie{Thế thì để tôi với Kuon lo phần chính. Hai đứa chúng tôi giỏi tiếng Anh nhất mà,}'' cậu bạn Đạo diễn nhìn sang tôi, nheo mắt.

Tôi chỉ gật đầu. ``\vie{Miễn là không phải thuyết trình vào buổi sáng sớm.}''

``\vie{Buổi sáng?}'' - cậu bạn Đeo kính bật cười. ``\vie{Ông là cú đêm à?}''

``\vie{Không. Làm đêm nhiều quá nên mới thế đấy,}'' tôi nói nửa đùa nửa thật, rồi nghiêng người, tựa vào lưng ghế nhựa ọp ẹp. Mắt nhìn ra sân trường loá nắng.

Một phần nhỏ trong đầu tôi bỗng nghĩ về ngày sinh nhật. Hôm nay. Là hôm nay đấy.

Nếu là một người khác, chắc giờ này họ đang nhận tin nhắn chúc mừng từ bạn bè và gia đình, hay đang tận hưởng một bữa tối đầy ắp tiếng cười, hoặc chí ít cũng đăng cái gì đó lên mạng xã hội để khoe bàn dân thiên hạ.

Còn tôi? Tôi đang ngồi đây, giữa phòng thư viện nóng nực, với ba đứa bạn cùng nhóm, lên chiến lược sống sót qua một học kỳ khắc nghiệt sắp tới.

Một giọng nói mơ hồ trong đầu thì thào: ``\vie{\textit{Ít nhất mày cũng nên tự chúc mình một câu gì đó đi chứ.}}''

Nhưng rồi, tôi lại lắc đầu. Không. Không cần thiết. Không phải bây giờ.

``\vie{Vậy cuối tuần sau học nhóm nhé? Cũng phải tranh thủ, chứ tháng sau là không kịp nữa đâu,}'' Yamazu lên tiếng, kéo tôi ra khỏi mớ suy nghĩ vẩn vơ.

``\vie{Ừ. Chốt thứ Bảy tuần sau đi,}'' tôi đáp, rồi lại mở sổ, viết thêm một dòng vào lịch trình. ``\vie{Để tôi nói chuyện với A-san bảo chị ấy sắp xếp lại lịch cho.}''

``\vie{Với cả anh em nếu ngày nào rảnh thì cứ báo nhé,}'' Yoisai nói, xách chiếc ba-lô lên. ``\vie{Tôi cũng cần anh em giúp một số môn đấy.}''

Hầy. Sinh nhật. Chẳng biết liệu năm nay tôi có thể đón sinh nhật một cách bình thường không nữa.

\ct

Năm giờ ba mươi chiều.

Ngẫm nghĩ lại, chúng tôi đã bàn bạc với nhau lâu đến vậy à. Chậc.

Trường đã vơi người. Trên nền trời cuối chiều, ánh nắng vẫn oi ả bám dính lấy từng bóng cây ven lối đi, như thể không chịu buông tha chút nào cho đám sinh viên chúng tôi còn sót lại. Tôi cùng cả nhóm đi bộ trên vỉa hè trải nhựa nóng, hướng về trạm xe buýt phía đông bắc cổng trường chính.

Cả bọn vẫn tiếp tục bàn bạc về bài tập lớn và mấy môn học kỳ tới, như thể cả mùa hè này chỉ để dành cho chuyện đó. Nhưng rồi, câu chuyện lại dần trượt về những mẩu tin tức lặt vặt, những chuyện vặt vãnh trong trường mà tôi chẳng buồn để tâm. Tôi chỉ đi bên cạnh họ, im lặng, đầu óc thì mải trôi dạt đâu đó.

Tôi ngước mắt lên nhìn bầu trời đang chuyển dần sang sắc cam nhạt. Có chút gì đó lạ lẫm len nhẹ trong lòng. Phải chăng\dots\ đây là cảm giác lạc lõng? Hay đơn giản là tôi đang hơi mệt vì đống đồ án và bài tập lớn cùng đống việc bù đầu sắp tới?

Tôi thực sự cũng chẳng cần biết nữa. Thời gian sắp tới sẽ là một giai đoạn cực khó nhằn đây.

Trong lúc tôi còn đang nghĩ tới những điều xa xăm, bỗng chốc Yamazu lên tiếng, như thể vừa nhớ ra điều gì: ``\vie{M-Mà nhắc mới nhớ\dots\ hôm nay là sinh nhật của Kuon à?}''

Tôi khựng bước nửa giây. Không hẳn là bất ngờ, nhưng vẫn hơi ngỡ ngàng khi có người thực sự nhắc đến.

Yoisai liền tiếp lời, gật đầu: ``\vie{Ừ nhỉ. Giờ tôi mới nhớ. Z**o vừa nhắc là hôm nay sinh nhật ông ấy.}''

Tôi khẽ nhếch môi cười. ``\vie{Tưởng hai người quên rồi chứ?}''

Tamachi chỉ nhún vai, giọng tỉnh bơ: ``\vie{Không. Chúng tôi quên thật. Thấy thông báo hiện lên mới nhớ.}''

``\vie{Vậy là cả ba người nhờ Z**o mới nhớ ra à\dots}'' Tôi bật cười khẽ, rồi buông một tiếng thở dài nhẹ như gió thoảng. Nụ cười của tôi có chút chua chát, nhưng tôi không hề giấu nhẹm nó đi.

Mà thực ra, tôi cũng chẳng trách họ được. Tôi vẫn để chế độ hiển thị ngày sinh, tự tôi để như thế thì có gì phải than phiền? Nếu ai đó biết cũng chỉ nhờ vào cái thông báo hệ thống cứng nhắc ấy, và dĩ nhiên, cũng chẳng mang lại giá trị gì.

Thậm chí, dù có biết, chắc gì đã có ai trong lớp buồn gửi lời chúc mừng tôi? Không phải bởi tôi giấu, tôi đâu có. Chỉ đơn giản là tôi không thích bị làm phiền bởi những lời chúc xã giao, và cũng không đáng để khiến họ phải làm như vậy.

\dots Ừ thì, nếu họ có chúc, tôi vẫn sẽ thấy vui chứ. Nhưng tôi cũng không mong người ta để tâm đến những chuyện đó.

Cậu bạn Đeo kính đột ngột hỏi: ``\vie{Vậy nhà ông tính ăn gì cho sinh nhật hôm nay?}''

Tôi hơi nhướng mày: ``\vie{Hỏi làm gì?}''

``\vie{Thì\dots\ tò mò thôi.}''

Tôi khẽ chẹp miệng. ``\vie{Cũng chỉ là một món gì đó đặc biệt hơn ngày thường. Có thể là lẩu, cũng có thể là nướng. Có năm mát tay thì ra ngoài ăn đồ Tây. Có bánh kem, có đồ ăn, vậy là đủ.}''

Cô bạn cùng nhóm cũng tò mò: ``\vie{Nhà ông ăn với nhau à?}''

``\vie{Ừ. Có mỗi bốn người ở nhà thì ăn với nhau thôi. Có ai khác quái đâu.}''

``\vie{Ô-Ông không r-rủ bạn bè à?}'' cậu bạn Đạo diễn lắp bắp, có vẻ ngạc nhiên thấy rõ.

``\vie{Tôi cũng chẳng mời ai. Ai đến thì đến, không đến thì thôi,}'' tôi nhún vai. ``\vie{Sinh nhật thì cả nhà bên nhau là quá đủ rồi.}''

Tôi thoáng dừng lại, một chút ký ức mờ nhạt bỗng trồi lên như bong bóng khí nổi trong nước.

Tôi nhớ hồi còn nhỏ, chừng cấp một, cấp hai gì đó, tôi từng có một sinh nhật được tổ chức trong ba ngày liền. Không đùa đâu, đúng là BA ngày.

Một ngày trước ngày sinh nhật chính thức ở cái xưởng mà bố tôi từng làm việc, nơi mà những cộng sự - nếu không muốn nói là cấp dưới của ông sếp, vốn cũng là đàn em của bố tôi - đã gom góp lại tổ chức tiệc nhỏ cho tôi.

Một ngày ở ngay chính căn nhà của chúng tôi, nơi trước khi bị biến thành một ``ký túc xá'' đầy bất đắc dĩ. Chi là một bữa lẩu với gia đình, với bác Hijaki-san, với ông bà nội và ông bà ngoại.

Và một ngày nữa sau ngày chính thức ở dưới quê, nhà ông bà ngoại. Cũng chỉ từng đó người, và cũng là món lẩu đó.

Hồi ấy, tôi cứ tưởng đó là điều bình thường. Nhưng bây giờ nhìn lại, tôi cũng không hiểu vì sao người lớn lại ``chịu chơi'' đến vậy.

Một sinh nhật ba ngày liền. Giờ nghĩ lại thấy buồn cười. Và có gì đó\dots\ xa xôi tới không tưởng.

Nhưng đó cũng chỉ là quá khứ đã vào dĩ vãng. Giờ chẳng ai rảnh đến mức làm sinh nhật rình rang cho một sinh viên đại học vừa bước sang tuổi hai mươi lăm kiêm tổng quản cho một công ty giải trí (cũng cực kỳ loạn lạc vì thái độ của các tài năng) cả.

Chúng tôi tiếp tục bước đi dọc theo lối mòn dẫn tới trạm xe buýt gần cổng sau của trường, những bước chân giẫm nhẹ lên nền đất khô cằn như thể tránh khuấy động bầu không khí oi nồng vẫn còn vương vất giữa trưa hè. Mỗi người một chút lặng im, có lẽ ai cũng bắt đầu thấm mệt sau hơn một giờ bàn bạc về đống học phần chết tiệt kia.

``\vie{Kuon đang làm việc cho \hll\ đúng không nhỉ?}'' Tamachi chợt lên tiếng, phá vỡ luồng suy nghĩ miên man trong đầu tôi như một nhát cắt dứt khoát.

Tôi quay sang, thoáng rướn mày. Câu hỏi này không mới với tôi, nhưng cách mà cậu ta thốt ra lại khiến tôi hơi bất ngờ. Chắc hôm nay Tamachi cũng lười nghĩ đề tài nói chuyện nên mới moi ra chủ đề đó.

``\vie{Hình như tôi cũng nhớ mang máng như vậy\dots}'' Yoisai gật gù, mắt nhìn đâu đó vào khoảng không như đang cố lục lại trí nhớ. ``\vie{Hình như ông ấy làm chức vụ gì đó quan trọng lắm.}''

``\textit{Trời đất, bạn bè học với nhau hơn hai năm rồi mà bây giờ mới biết à?}'' tôi nghĩ thầm, cười gượng. Cũng chẳng thể mắng họ được, vì tôi chủ động không nói chuyện này ra mà. Những người duy nhất trong lớp biết được vụ này có lẽ chỉ có nhóm Miharu-Tokue, nhóm Jotaro và Yamazu thôi.

Chưa kịp để tôi nói đỡ một câu cho nhẹ chuyện, thì Yamazu đã nhảy vào trước, vấp chữ như một thói quen: ``\vie{Đ-Đúng là như v-vậy mà! Tôi có v-vài lần lên văn p-phòng của ông ấy rồi, cũng g-gặp các thành v-viên ở đó nữa! Ông ấy là tổng quản lý á!}''

Tôi thở dài trong lòng. Tên này đúng là không biết giữ mồm giữ miệng khi cần. Dù cũng chẳng phải bí mật gì to tát, nhưng tôi vốn không thích phô trương chuyện công việc với người ngoài ngành.

``\vie{Tokue có bảo như vậy hôm tôi nói chuyện với ông ấy nên tôi biết,}'' cậu bạn Đeo kính thêm vào, như thể xác nhận lại điều mà cậu Đạo diễn vừa nói.

``\vie{Vẫn loạn như nhà trẻ mẫu giáo vậy thôi,}'' tôi nhún vai đáp khi Yamazu hỏi tình hình công ty. ``\vie{Chẳng có ngày nào yên ổn cả.}''

``\vie{Làm tổng quản cho công ty chắc là cũng cực lắm nhỉ?}'' cô bạn Lập trình nghiêng đầu hỏi, vẻ tò mò xen lẫn cảm phục.

``\vie{Không cực mới ngạc nhiên đấy. Gần như tuần quái nào cũng có chuyện ngớ ngẩn xảy ra,}'' tôi cười nhạt, chống tay lên trán, đầu lắc nhẹ. ``\vie{Tôi mà kể ra thì chắc cả ngày không xong mất. Nói chung là nhiều lắm.}''

``\vie{Thử kể vài chuyện đi,}'' Tamachi hứng thú chen vào, ``\vie{kiểu mấy vụ khôi hài ấy.}''

Tôi thoáng im lặng một lúc, cố lục lọi trong đống hỗn độn ký ức gần đây. ``\vie{Ừm, xem nào\dots\ Kiểu như, mấy tuần trước Peko-chan tự dưng muốn làm tư vấn viên, cuối cùng lại chẳng đâu vào đâu mà còn bị đàn em tư vấn ngược lại\footnote{Xem lại {\flaregothic ÉPISODE 111}, quyển số Chín.}. Rồi cái lần mà Ayame-chan bày trò tập thể dục thể thao chỉ để chữa bệnh trầm cảm tháng Năm nữa\dots\ Chính tôi với Miharu quay phim chứ còn ai\footnote{Xem lại {\flaregothic ÉPISODE 104}, quyển số Tám.}!}''

``\vie{R-Rồi! Vậy ra ở c-công ty ông vẫn như thế à\dots?}'' cậu Đạo diễn tròn mắt, mặt tái đi như thể tôi đang mô tả hiện trường án mạng.

``\vie{Chưa kể đến vụ Subaru bị nhầm là vịt để nấu lẩu\footnote{Xem lại {\flaregothic ÉPISODE 96}, quyển số Tám.} nữa chứ,}'' tôi nói tiếp, khoé miệng nhếch lên đầy bất lực. ``\vie{May mà họ tìm thấy vịt cổ xanh với tỏi tây chứ không là cũng rách việc lắm!}''

``\vie{Vãi cả hồn\dots}'' Yoisai cười không ra tiếng. ``\vie{Nghe như ông đang làm ở chỗ nào đó ảo lắm luôn á.}''

``\vie{Đôi khi tôi cũng tự hỏi liệu mình đang làm việc cho công ty truyền thông hay là một rạp xiếc trá hình nữa,}'' tôi buông một câu mỉa mai nhẹ nhàng, rồi khẽ bật cười. ``\vie{Nhưng mà nghĩ lại thì\dots\ cũng vui.}''

Ba người kia dường như đều gật đầu, đồng cảm phần nào. Dù mỗi người một ngành, nhưng mùi khói tỏa ra từ đống áp lực và công việc chẳng khác nhau là mấy.

Chúng tôi tiếp tục bước đi thêm vài trăm mét nữa. Trạm xe buýt đã hiện ra ở xa, nấp dưới bóng một cây bàng già. Tôi cảm thấy những mạch nghĩ vụn vặt lại bắt đầu lấn chiếm trong đầu, kéo tôi trở lại với một cảm giác trống vắng.

Tamachi nhìn đồng hồ, khẽ gật đầu chào chúng tôi rồi băng qua đường, khuất dần sau làn xe cộ đang nườm nượp khi giờ tan tầm đã gần kề. Ba người chúng tôi còn lại vẫn đứng yên ở trạm xe buýt, lặng lẽ quan sát những chiếc xe chạy qua, những người đi bộ vội vã, và những chiếc xe đạp đậu bên đường.

``\vie{Chắc là mọi người ở trong công ty quý ông lắm nhỉ?}'' Yamazu lại cất tiếng, giọng vẫn giữ cái vẻ rụt rè nhưng lần này nghe có phần chân thành hơn mọi khi.

Tôi ngẩng mặt lên, nhìn bầu trời trong vắt đã ngả dần màu vàng. Gió nhẹ phả vào mặt, mùi mồ hôi xen lẫn chút bụi đường khiến lòng tôi chợt lạnh đi một chút.

``\vie{Tôi cũng chẳng biết nữa,}'' tôi đáp, thở nhẹ ra một hơi.

``\vie{Sao lại thế?}'' Yoisai hỏi, hơi nghiêng đầu.

Tôi nhún vai, quay đầu lại nhìn hai người bạn. ``\vie{Tôi chỉ là một tổng quản lý thôi. Làm công ăn lương như bao người. Có khi còn khổ hơn nữa là đằng khác.}''

Tôi thoáng cười, kiểu cười không vui, không buồn, chỉ đơn giản là để dập tắt câu chuyện đang bị đẩy đi quá xa. ``\vie{Chức vụ Tổng quản thoạt nghe có vẻ oai thật, nhưng so với các thành viên ở đó - là các idol ấy - tôi cách họ cả một trời một vực.}''

Tôi bước tới chiếc ghế bành dưới trạm chờ xe buýt, rồi tiếp tục trải lòng: ``\vie{Họ được hàng triệu người ngưỡng mộ, muốn được xin ký tặng, muốn được chụp ảnh, rồi còn đòi làm vợ. Còn tôi? Chỉ là một thằng con trai tầm thường đi làm văn phòng, viết báo cáo, ký giấy tờ và thi thoảng lắm quát mắng vài người khi cần thiết.}''

Cậu Đạo diễn xị mặt, có vẻ không đồng tình, liền phản bác lại bằng chất giọng có phần hấp tấp: ``\vie{N-Nhưng mà\dots\ đ-đó là vị trí mà b-bao nhiêu người đàn ông con trai muốn m-mơ ước còn gì?}''

``\vie{Có thể,}'' tôi gật đầu hờ hững, ``\vie{nhưng mơ ước với thực tế là hai thứ khác nhau. Người ta chỉ nhìn thấy ánh đèn sân khấu chứ đâu có thấy được chỗ tôi đứng phía sau tấm màn ấy đã cật lực như thế nào.}''

Cả ba chúng tôi lại im lặng thêm một lúc. Những tiếng xe cộ chạy rừm rừm trên phố, tiếng còi kêu inh ỏi trong giờ đã gần tan tầm, và tiếng gió Lào thổi qua cành cây, như thể muốn lấp đầy sự im ắng ấy.

Yoisai chợt lên tiếng, hỏi một câu có phần nhẹ nhàng hơn: ``\vie{Vậy\dots\ hôm nay sinh nhật ông, có ai trong công ty biết không?}''

Tôi bật cười khẽ, lắc đầu. ``\vie{Chỉ có đồng nghiệp của tôi gọi điện tối qua. Gọi đúng nửa đêm như thường lệ, chúc vài câu rồi dặn ngủ sớm, nhớ ăn uống đầy đủ. Và\dots\ chấm hết. Tôi không nói với ai khác cả. Kiểu, nó không đáng để nói lắm.}''

Thực ra trong điều tôi nói cũng có phần đúng. Nhưng đến bây giờ, tôi vẫn không hiểu vì sao A-san lại muốn tôi tạm nghỉ một ngày vào ngày hôm nay. Có lẽ là vì tôi đã làm việc quá sức trong suốt thời gian qua, hoặc đơn giản là chị ấy - hoặc Tanigo-san - muốn tôi có một ngày nghỉ ngơi thoải mái.

``\vie{Không đáng à?}'' Yamazu tròn mắt.

``\vie{Không. Tôi chưa bao giờ nghĩ sinh nhật mình là chuyện quan trọng với ai khác. Cũng không muốn người khác phải bận tâm. Thật đấy.}''

Cậu bạn gãi đầu, ngập ngừng một lúc rồi hỏi: ``\vie{N-Nếu mà\dots\ giả sử họ biết thì sao? Nếu họ làm cho ông một b-bữa sinh nhật hay cái gì đó?}''

Tôi phẩy tay, như xua tan đi ý nghĩ có phần không tưởng ấy: ``\vie{Đừng có ngớ ngẩn chứ. Tôi đâu xứng đáng được như vậy.}''

``\vie{Tại sao lại không?}'' cô bạn Lập trình hỏi, có phần bối rối.

``\vie{Thì\dots\ tôi bảo rồi, tôi không phải là người nổi tiếng. Tôi chỉ là một thằng con trai bình thường thôi. Đừng có nói mấy thứ viển vông như thế chứ.}''

``\vie{Nhưng mà\dots\ ông cũng là một trong những người làm việc cho công ty lớn mà! Ông còn là một người chỉ huy các cô gái ấy nữa chứ!}''

``\vie{Thôi nào,}'' tôi cười nhẹ, ``\vie{đừng có nói mấy cái đó nữa. Tôi không thích đâu.}''

Yoisai định nói gì đó để tiếp tục chất vấn, nhưng rồi khựng lại. Hai người nhìn tôi một lúc lâu, đôi mắt có gì đó là lạ, có lẽ là chút tiếc nuối, chút ngạc nhiên, hoặc là\dots\ thương cảm.

Tôi không rõ. Có thể là tôi tưởng tượng.

Có điều tôi biết chắc một điều: họ hiểu. Hiểu tình cảm mà các thành viên \hll\ dành cho tôi, là sự tôn trọng, sự gắn bó, sự tin tưởng. Nhưng bản thân tôi thì\dots\ tôi không biết. Hoặc đúng hơn là tôi chưa bao giờ thực sự để tâm đến nó.

Tôi chỉ thấy mình như người giữ cửa sau sân khấu. Một người cần mẫn, âm thầm, và hoàn toàn có thể bị quên lãng ở một xó nào đó vào bất kỳ ngày nào.

Và tôi cũng quen với điều đó rồi.

``\vie{À, xe của tôi tới rồi,}'' tôi buông một câu nhẹ tênh, rồi gật đầu chào hai người bạn trước khi bước vội lên chiếc xe buýt vừa đến, tiếng thắng kêu cọt kẹt bên vệ đường như nhấn nhá thêm chút lặng lẽ cuối buổi chiều.

Khi ngồi vào chỗ cạnh cửa sổ, tôi thoáng ngoái lại nhìn. Qua lớp kính mờ phủ bụi, tôi thấy họ vẫn đang đứng ở đó. Họ đang nói gì đó với nhau, thỉnh thoảng lại liếc về phía chiếc xe tôi vừa lên, gương mặt xen lẫn nét đăm chiêu và một chút\dots\ tôi không chắc, nghi hoặc? Lo lắng?

Tôi cũng không để tâm nhiều. Có thể họ chỉ đang bàn xem chiều nay ăn gì, hoặc chê tôi già đời khó gần. Tôi toàn bị trêu như thế ở trên trường, nên nếu vế sau là thật thì cũng không có gì làm lạ.

Gió điều hoà mát lạnh phả lên trán, tôi khẽ tựa đầu vào lưng ghế, để mặc tâm trí trôi dạt trở lại cái ngày này hàng năm.

Sinh nhật.

Một từ nghe có vẻ long trọng, đầy hoa và nến, nhưng với tôi, kể từ khi tôi lên cấp ba, tôi ngẫm nghĩ lại, nó cũng không thực sự đặc biệt đến như vậy.

Tôi đoán năm nay cũng vậy thôi. Một bữa cơm gia đình khác biệt hơn thường ngày đôi chút, có thể là một nồi lẩu đặt giữa bàn, mẹ tôi sẽ lại cằn nhằn tôi ăn ít rau, cha thì vừa ăn vừa xem tin tức, còn Ryuuhou\dots\ chắc con bé sẽ bám lấy tôi đòi cắt bánh sớm.

Có lẽ vẫn là chiếc bánh kem đủ cho sáu người ăn mẹ tôi mua ở cửa hàng quen thuộc, không tên tôi trên đó, chỉ có dòng chữ ``Chúc mừng sinh nhật con trai\~{}'' đơn giản. Không quá phô trương, không quá cầu kỳ, nhưng cũng đủ để khiến tôi cảm thấy ấm lòng.

Và như mọi năm, tôi cũng sẽ thổi nến, cười, nói vài câu xã giao cho đúng vai người được mừng sinh nhật. Để rồi ngày hôm sau lại phải đi làm và quay về guồng làm việc như thường lệ.

Như thể chẳng có gì xảy ra. Có lẽ hôm nay cũng sẽ trôi qua chóng vánh thôi.

\pov{-}

Khi chiếc xe buýt chậm rãi lăn bánh, Yoisai và Yamazu vẫn đứng lặng trước trạm. Ánh hoàng hôn đã nhuốm vàng lên vỉa hè, khiến bóng hai người đổ dài theo mặt đường.

Cậu Đao diễn đưa mắt nhìn theo bóng Kuon đã khuất sau lớp kính. ``\vie{Ông ấy l-l-lúc nào cũng như vậy,}'' cậu lẩm bẩm, ``\vie{tỏ vẻ cam chịu. Cứ như thể\dots\ chuyện b-bị người khác quên l-lãng là điều hiển nhiên.}''

Cô nàng Lập trình nhíu mày, đồng tình: ``\vie{Ừ. Tôi cũng thấy Kuon có gì đó là lạ. Không phải kiểu khó gần, nhưng mà kiểu, ông ấy như đang tự kéo ranh giới xung quanh mình ấy.}''

``\vie{Chắc là ông ấy không muốn bị thất vọng,}'' Yamazu chậm rãi đáp. \vie{Nếu không mong gì thì sẽ không buồn. N-Nhưng mà, nghĩ Kuon không đáng được tôn trọng\dots\ c-cái đó thì quá khiêm tốn rồi.}''

``\vie{Quá mức ấy chứ,}'' Yoisai đồng tình. ``\vie{Một người làm đến chức đó trong \hll\ mà lại tự xem mình chỉ là `thằng văn phòng bình thường'? Nói thật nhá, dám cá là cả công ty có khi còn loạn nữa nếu ông ấy không làm ở đó đấy.}''

Cậu bạn bật cười khẽ, nhưng mắt vẫn không rời khỏi con đường xe buýt vừa chạy qua. ``\vie{Mà\dots\ bà nghĩ ông ấy có buồn không?}''

``\vie{Buồn chứ,}'' Yoisai đáp, hơi khẽ giọng. ``\vie{Tôi nghĩ là ông ấy giấu giỏi thôi.}''

``\vie{Ừ\dots}'' Yamazu gật đầu, lặng thinh một lúc. ``\vie{Hy vọng là bữa tiệc sinh nhật của ông ấy năm nay không chỉ có mỗi nồi lẩu và cái bánh kem.}''

Yoisai không đáp. Cô nàng chỉ nhìn lên trên bầu trời bắt đầu chập chừng tối, nơi mặt trời đang dần lặn sau những tán cây phía xa, ánh sáng nhạt dần như tâm trạng của một người bạn học mà họ không thể với tới.

\pov{Hikawa Kuon}

Sáu giờ ba mươi chiều.

Xe buýt thả tôi xuống khu phố quen thuộc, cách nhà chỉ chừng một cây số đi bộ. Ánh đèn đường đã lên, chen chúc giữa những bảng hiệu nhấp nháy của hàng quán ven đường. Tiếng còi xe, tiếng rao vặt, tiếng cười nói vang vọng từ những quán ăn bên kia đường như một bản hợp xướng hỗn loạn mà quen thuộc.

Tôi bước từng bước chậm rãi trên vỉa hè, lách qua những người đi ngược chiều, băng qua những cung đường nườm nượp xe cộ giờ cao điểm. Gió chiều thổi nhẹ, mang theo mùi khói xe pha lẫn chút thơm thơm của món nướng từ một quán gần đó. Trong cái náo nhiệt đó, đầu óc tôi lại không ngừng vang lên những lời vừa rồi.

``\vie{Chắc là mọi người ở trong công ty quý ông lắm nhỉ?}''

``\vie{Làm tổng quản cho công ty chắc là cũng cực lắm nhỉ?}''

``\vie{Ông còn là một người chỉ huy các cô gái ấy nữa chứ!}''

Tôi thở ra một hơi thật dài. Chắc là họ nói thật lòng. Nhưng tôi lại không biết phải tiếp nhận điều đó thế nào. Họ không ở trong ngành, không biết những gì tôi đã trải qua, không biết những gì tôi đã làm. Họ chỉ thấy một phần của tôi, một phần mà tôi đã cố tình để họ thấy.

``Tổng quản lý.''

Ba chữ. Nghe thì to tát thật, và đúng là có vẻ cũng ra oai. Là người điều phối toàn bộ công việc giữa các phòng ban, giữa các nhân viên và các tài năng thần tượng, là người giữ cho lịch hoạt động không chồng chéo, là người xử lý khủng hoảng khi có sự tranh cãi nổi lên, là người ngồi nghe ông CEO phân tích báo cáo tài chính hàng quý, và cũng là người chịu trận khi một ai đó của Gen 3 rủ nhau đi họp online - mà thực chất chỉ là dựng khung lên và đứng vào các ô đó, để rồi Sui-chan nhìn chúng tôi với đầy sự ghen tị\footnote{Xem lại {\flaregothic ÉPISODE 56}, quyển số Bốn.}, cũng chính là người phải tự bỏ tiền túi chỉ để đi đền cửa sổ mà các cô gái đã phá tung trong khoảng thời gian đầu tôi làm việc.

Tôi làm tất cả những việc đó. Ừ. Làm mỗi ngày. Nhưng tôi vẫn không thể nào cảm thấy mình thực sự xứng đáng với ba chữ ``sự tôn trọng.''

Bởi lẽ đơn giản thôi, tôi đâu có tạo ra những sân khấu rực rỡ. Tôi đâu có khiến hàng nghìn người phải khóc khi xem concert. Tôi đâu có khiến những người hâm mộ từ khắp nơi trên thế giới phải viết thư tay cảm ơn, hay đến Nhật chỉ để được vẫy lightstick (đèn) cho một người.

Tôi chỉ là người đứng sau hậu trường. Và tôi quen với chuyện đó. Tôi thấy mình phù hợp với bóng tối hơn là ánh đèn.

Nhưng vậy\dots\ liệu có đang tự phủ nhận chính mình không?

Tamachi không nói gì lúc chúng tôi tách nhóm, nhưng ánh mắt của cậu ta\dots\ vẫn khiến tôi bận tâm. Có lẽ cậu ấy cũng đang thắc mắc tại sao tôi lại không nhận ra giá trị bản thân.

Mà biết đâu đấy, có thể tôi có nhận ra. Chỉ là tôi không muốn tin vào nó. Vì nếu tin, tôi sẽ mong đợi. Và nếu mong đợi, tôi sẽ thất vọng. Và cứ thế, tôi sẽ lại tự nhốt mình vào cái vòng luẩn quẩn đó.

Vậy nên tốt nhất là cứ bước đi như thế này, trong cái dòng người tấp nập, lặng lẽ như bao người khác. Là một nhân viên công sở, một người con, một người anh cả trong gia đình. Chỉ vậy thôi.

Tôi mải chìm trong suy nghĩ của mình, không nhận ra tôi đã về gần đến nhà từ khi nào.

Tôi liếc nhìn đồng hồ trên điện thoại. Bảy giờ tối.

Bóng đèn vàng nơi khu phố nhỏ trước nhà tôi nhấp nháy nhẹ như chào đón. Tôi rẽ vào đó, dừng lại trước cổng.

Cửa hàng viễn thông vẫn còn đang làm việc. Mọi thứ vẫn diễn ra rất đỗi bình thường.

Chắc giờ này mẹ đang chuẩn bị lẩu rồi.

Chắc có lẽ bố tôi đang nằm ở nhà phụ giúp mẹ tôi.

Chắc Ryuuhou vẫn đang còn vừa xem tivi vừa chơi điện tử.

Nhưng một điều mà tôi biết chắc chắn: nó đang chờ cơm tôi về. Sinh nhật năm nay của tôi chắc cũng sẽ chỉ diễn ra một cách bình thường.

Không hơn, không kém.

\ct
Tôi khẽ mở cửa, bước vào nhà. Cảm giác đầu tiên không phải là ấm áp, mà là\dots\ sự yên ắng. Một sự tĩnh lặng lạ lùng. Tôi ngước nhìn về các cô nhân viên trong cửa hàng rồi bắt đầu leo cầu thang ở bên trái. Từng tầng, từng tầng một.

Tầng hai, cánh cửa phòng của Aloe-chan vẫn đóng im lìm như hồi chiều. Tôi dừng lại một chút, ngó qua ô kính mờ - vẫn không có ánh sáng. Chắc chắn em ấy sẽ không thể nào về kip cho tiệc sinh nhật của tôi rồi, nhưng tôi cũng không trách em ấy. Tôi lắc đầu cười nhạt rồi tiếp tục bước lên.

Tầng ba, căn phòng trống, nơi từng được tổ chức vài sự kiện Tết Nguyên Đán với các thành viên \hll\ năm ngoái. Nhưng giờ, đèn hành lang tắt ngóm, hầu như không có một tiếng động gì. Cái không khí lạnh lạnh của cầu thang hắt lên làm tôi bỗng nghĩ đến một ngôi nhà bỏ hoang. Cứ mỗi lần đi qua tầng này là lại có cảm giác như đang bước vào một khoảng không vô hình nào đó. Một sự lặng thinh, vô chủ, và có chút gì đó buốt sống lưng. Không hiểu sao tôi lại rùng mình nhẹ.

Tầng bốn, phòng của bác Hijaki và bà nội. Cửa khép, bên trong không có tiếng động gì cả. Tôi nghiêng đầu nhìn vào ô kính nhỏ bên hông. Một màu tối đen. Không có tiếng vô tuyến, cũng chẳng có tiếng radio đang phát chương trình tôn giáo như mọi khi. Chắc bác Hijaki có việc ra ngoài, còn bà nội có lẽ đã đi khám sức khỏe định kỳ. Bình thường giờ này thể nào bà cũng đang ngồi ở bàn bếp, vẫy tay gọi tôi xuống ăn trái cây.

Tầng năm, phòng của bố con nhà tôi. Tôi mở cửa, đặt chiếc cặp nặng trĩu công việc xuống đất, rồi cất tiếng: ``\vie{Con về rồi đây\~{}}''

\dots

\dots

Không có tiếng trả lời.

Thứ duy nhất lên tiếng\dots\ là tiếng móng vuốt cào nhẹ lên sàn gỗ, rồi một bóng lông trắng từ cuối hành lang lao ra. Mi sủa lên một tiếng nhỏ như phản ứng lại lời chào ấy, rồi sà vào chân tôi. Tôi cúi xuống xoa đầu nó, tay vuốt nhẹ dọc sống lưng mềm mại.

``Có mỗi em ở nhà thôi à\dots''

Mi rít lên một tiếng khẽ như thể đồng tình. Tôi thở dài.

Không hiểu sao, lòng tôi có chút hụt hẫng. Mỗi năm sinh nhật đều như vậy cả, tôi vẫn biết. Nhưng ít ra, những năm trước, khi tôi về đến nhà, vẫn có người chờ. Một tiếng chào, một câu hỏi: ``Sao, hôm nay ở trường mệt không?'' hay đơn giản là: ``Rửa tay đi rồi xuống ăn cơm con!''

Nhưng hôm nay thì rất khác. Ngoài con chó cảnh của nhà ra, không có lấy một ai.

Tôi đứng dậy, hướng lên tầng sáu, là căn bếp chung của gia đình tôi.

Leo lên từng bậc, tôi vẫn không nghe thấy bất kỳ âm thanh nào: không tiếng dao thớt, không mùi lẩu, không tiếng bát đũa va vào nhau. Khi tôi mở cửa bếp, thứ đón chào tôi là một khoảng trống lạnh lẽo. Đèn chưa bật. Không một bóng người.

Tôi ngẩng đầu nhìn lên tầng bảy. Cửa phòng Coco-chan đóng chặt. Then cũng đã cài. Em ấy vẫn không có nhà.

Tôi chống tay vào lan can cầu thang, nhìn ngược xuống các tầng phía dưới. Căn nhà bảy tầng hôm nay không giống mọi khi.

Không một bóng người.

Chỉ còn tôi và Mi.

Tôi cá là bố mẹ tôi lại đi ra ngoài có việc, còn em gái tôi lại đi chơi với lũ bạn rồi.

Tôi bước chậm lại, ngồi xuống chiếc đệm, để balô xuống cạnh chân. Mi vẫn lẽo đẽo nằm cuộn dưới chân tôi, gác cằm lên đùi tôi như thể đang nói ``Có em đây, chủ đừng buồn nữa.''

Tôi cười nhẹ. ``\vie{Ừ, thôi thì tối nay làm bát mì ăn tạm\dots\ rồi coi như xong một sinh nhật nữa vậy.}''

Tôi nói bâng quơ, như một câu độc thoại ngắn ngủi gửi vào không khí loãng. Không phải tôi mong chờ gì, chỉ là có chút cảm giác lạ. Trống rỗng mà tĩnh lặng.

Đây có lẽ sẽ là lần đầu tiên tôi thực sự cảm nhận được một sinh nhật cô đơn.

Không có lời chúc. Không có bánh kem. Không có tiếng gọi nhau í ới hay ánh đèn màu mè.

Chỉ là một bữa mì, một con chó nhỏ, và tôi.

Một con chó hơn mười năm tuổi còn hiểu được tôi hơn cả em gái tôi. Thật trớ trêu.

Tôi dựa lưng vào tường, tay với lấy điện thoại trong túi áo. Màn hình hiện lên, không có thông báo nào\dots\ Tôi buông ra một tiếng thở dài.

Mà tại sao tôi lại mong chờ có một ai đó gửi lời chúc mừng sinh nhật tôi chứ?

Ngay khi tôi định tắt máy đi, bỗng chốc\dots *TINH!*

Một tin nhắn hiện lên ở thanh thông báo.

Từ Bố. ``\vie{Con xuống tầng ba đi nhé.}''

\dots Gì cơ?

Tôi nhíu mày. Tầng ba? Cái nơi tối om như nhà ma lúc nãy tôi vừa đi qua?

Cái nơi tôi còn đang bỏ trống đó ư?

Tôi ngẩng đầu, nhìn xuống lối cầu thang dẫn xuống dưới. Mi cũng ngẩng đầu theo như thể đang đọc biểu cảm trên mặt tôi.

Tôi vẫn chưa hiểu chuyện gì đang diễn ra.

Nhưng bằng bản năng, hay có lẽ là trực giác, tôi cảm thấy\dots\ thứ gì đó sắp xảy ra.

Một cái gì đó rất khác.

*TINH!*

Lại một tin nhắn nữa vang lên.

Tôi lười nhác mở điện thoại ra xem. Vẫn là từ bố.

``Lấy cho bố thêm hai hộp đá trên tủ lạnh nhé!''

\dots Ờm? Gì cơ?

Tôi chớp mắt vài lần, chưa kịp tiêu hóa với tin nhắn mới nhận. Ủa? Nếu ông muốn lấy đá, thì lên tầng sáu tự lấy đi chứ? Tại sao lại phải nhắn tin cho tôi khi tôi vừa mới đặt mông ngồi xuống chưa được năm phút?

\dots Thôi được rồi.

Dù sao thì cũng đâu phải việc gì quá to tát. Cùng lắm là lấy đá uống bia, hay nước ngọt\dots\ đúng không? Không thể là gì hơn được. Tôi bật dậy, lững thững bước vào khu bếp, mở tủ lạnh.

Dù sao thì cũng đâu phải việc gì lớn. Cùng lắm là lấy đá uống bia, hay nước ngọt\dots\ đúng không? Không thể là gì hơn được. Tôi bật dậy, lững thững bước vào khu bếp, mở tủ lạnh.

Và tất nhiên, đến đoạn này thì mới nhớ, cơ tay tôi yếu. Rất yếu.

Mỗi lần lấy đá là một trận chiến. Hoặc là lấy nước đổ lên khay rồi bẻ từng viên như một ông già lụm cụm, hoặc là đập mạnh cái khay lên nóc tủ lạnh để đá bung ra. Tôi thì lúc nào cũng chọn cách thứ hai. Nhanh gọn mà đỡ phiền phức.

Kể ra thì hơi ồn, vì nghe như có ai đang dùng búa đập vào tủ vậy. Chắc người dưới tầng nghe còn tưởng trên này đang xây nhà.

*CẠCH!* *CẠCH!* \textbf{*RẦM!*}

``Được rồi!''

Tôi cho tất cả đá vào hai hộp nhựa rồi đóng nắp lại. Cũng vừa khít.

``\vie{Mi, ở yên đây. Anh xuống nhà một lúc, tí cho em ăn,}'' tôi cúi xuống nói nhỏ, xoa đầu con chó trắng đang nằm trên thảm như một ông cụ non đang canh nhà.

Mi chỉ vẫy vẫy đuôi, như thể ra dấu ``Vâng ạ. Lên nhanh nhanh nha\~{}''

Tôi bước xuống cầu thang. Cầu thang gỗ vang tiếng két két dưới từng bước chân.

Tầng ba tối om, vẫn y như khi tôi đi ngang hồi chiều.

\dots Nhưng mà khoan. Khi tôi mở cánh cửa khẽ kêu cót két ấy ra, thứ hiện ra trước mắt khiến tôi khựng lại.

Không phải một chiếc bàn ăn đơn giản như mọi năm, mà nếu có, chắc chắn nó phải ở tầng bốn, hoặc tầng năm.

Nhưng, thứ mà tôi ngạc nhiên hơn cả, lại chính là\dots

Một loạt các nồi lẩu. Hơn bốn cái, thậm chí năm cái. Một cái to nằm giữa phòng, vài ba, bốn cái nhỏ nằm rải rác xung quanh. Đĩa thì nhiều không đếm xuể, bát đũa cũng phải hai chục bộ là ít.

Tôi đứng đờ người mất vài giây, mắt nhìn quanh. Đèn trong phòng chưa bật, chỉ có ánh sáng hắt từ cửa sổ bên kia rọi qua rèm, để lộ ra những vật dụng được bày sẵn. Như thể chuẩn bị cho một bữa tiệc lớn.

``\textit{Quái, nhà chỉ có bốn, năm người ăn với nhau thôi mà sao lắm nồi thế?}'' tôi thầm nghĩ, mày khẽ nhíu lại.

Một cảm giác gì đó sai sai.

Rất sai.

Không khí im ắng đến khó hiểu. Không một tiếng động. Không một lời chào đón. Không ai trong bếp. Không ai trong phòng khách.

Chỉ có một căn phòng với hàng chục bộ bát đũa đang đợi sẵn, dường như dành cho một đội quân nhỏ.

Và tôi đây, người vừa vác đá xuống, giờ vẫn đang đứng trước cửa, hộp đá trên tay chưa kịp đặt xuống.

``\vie{Bố? Mẹ?}'' tôi cất tiếng gọi.

Không một tiếng trả lời.

``Ryuuhou?'' tôi thêm một lần nữa, lớn hơn, hy vọng em gái mình ở đâu đó trong căn phòng trống này.

Vẫn không có tiếng đáp.

Tôi khẽ nhăn mặt. Kỳ lạ thật. Mọi người đâu hết cả rồi?

Tôi đặt hai hộp đá xuống sàn. Âm thanh phát ra khi nắp hộp chạm mặt sàn ốp lát nghe lạch cạch một cách trống vắng. Cảm giác như tôi vừa mang lễ vật đến giữa một căn phòng hoang.

Cánh cửa sau lưng tôi khẽ xoạch khi tôi đóng lại. Và rồi ngay thời điểm chiếc cửa đóng lại\dots

\begin{center}
    ``\textbf{Và giờ, tiết mục mừng ngày sinh nhật của Hikawa Kuon xin được bắt đầu!}''
\end{center}

Tôi giật bắn mình.

``\textbf{Cái gì vậy!?}''

Tim tôi đập thình thịch. Tôi quay phắt lại, đảo mắt khắp căn phòng, như thể vừa bị ai đó nhảy ra hù dọa. Không thấy ai. Không một bóng người. Chỉ có giọng nói của một người con gái dội lên, vang vọng trong không gian tối om.

Một luồng sáng đột ngột bật lên, rọi thẳng về phía bức tường đối diện chiếc rèm cửa.

Tôi hơi nheo mắt lại vì chói. ``Cái\dots\ quái gì!?''

Tôi tiến lại gần. Dưới ánh đèn trắng nhợt, tôi thấy một thứ trông như một nút bấm tròn trịa dính trên tường. Mà không phải kiểu nút công tắc đèn bình thường. Cái này nhìn như thể vừa rơi ra từ một cái phim trường hạng hai.

Tôi trừng mắt, càu nhàu trong cổ họng: ``Ông thần thánh sứa nào lại để cái nút bấm kỳ quái giữa tường thế này hả!?''

Lại không có tiếng trả lời. Tuyệt lắm.

Tôi chờ một chút, vẫn không có ai nhảy ra đính chính hay giải thích. Quá quen với cái kiểu ``một mình đối mặt định mệnh,'' tôi thở dài, nhún vai.

``Được thôi\dots\ Bấm thì bấm, sợ đếch gì.''

Tôi giơ tay bấm nút tròn đó.

*CẠCH*

Từ đâu đó trên trần nhà, một tấm màn chiếu trượt xuống, kéo theo tiếng rè rè khe khẽ. Tôi ngước lên, thấy một chiếc máy chiếu đang khởi động. Bên trong, đèn máy lóe sáng.

Kỳ thực thì, chiếc máy chiếu đó ở đâu ra vậy? Tôi đâu có nhớ tầng ba có lắp bất cứ thứ gì ngoài cái điều hòa đâu!

Tôi hướng về bức tường đối diện. Trên nền trắng mập mờ của chiếc máy chiếu, những dòng chữ xuất hiện:

\begin{center}
    ``Những thước phim đáng giá về Hikawa Kuon \\ bởi các thành viên của \hll.''
\end{center}

Tôi há hốc miệng. Chuyện quái gì đang xảy ra vậy?

Tôi tròn mắt nhìn dòng chữ vẫn nhấp nháy. Mọi thứ vẫn tối xung quanh, chỉ có ánh sáng từ màn chiếu hắt lên mặt tôi.

Tôi nuốt nước bọt. Dù chẳng hiểu gì, cũng chẳng có ai giải thích cho tôi biết tôi đang rơi vào trò mèo gì\dots

\dots nhưng có lẽ, nếu đã đến nước này, tôi không thể không xem.

\il

Trên màn hình, nền đen mờ dần, thay vào đó là hình ảnh một căn phòng quen thuộc. Không gì khác, đó là căn bếp tầng sáu nhà tôi.

Và đứng giữa căn bếp đó\dots\ là Tokino Sora-chan.

Em ấy nhìn thẳng vào máy quay, ánh mắt dịu dàng như mọi khi, mái tóc đen dài khẽ đung đưa. Tay vẫn còn đeo tạp dề, chứng tỏ cô vừa mới làm bếp xong.

``Kuon-senpai,'' Sora-chan mỉm cười, tiếp tục dân dắt, ``chúng em biết anh sẽ bất ngờ khi xem tới đây, nên hãy để bọn em  làm rõ một điều nho nhỏ: từ giờ, anh không có quyền tự cô lập bản thân nữa đâu đấy.''

Tôi nhíu mày khó hiểu. ``Cô lập?'' Ý em ấy là gì chứ?

Nhưng tôi chưa kịp hiểu chuyện gì xảy ra, màn hình trên máy chiếu liền đổi sang người tiếp theo.

\ct

Akai Haato xuất hiện, tay vẫn đang cầm một củ hành tây đã cắt dở. Mắt em ấy đỏ hoe, không biết là do hành hay do xúc động.

``Kuon-senpai! Đừng nghĩ là anh có thể lặng lẽ già đi như một ông già lẩm cẩm nữa!'' em ấy làm mặt nghiêm rồi khịt mũi. ``Năm nay sinh nhật của anh, bọn em đã cất công chuẩn bị rồi đó!''

\dots Tôi không biết tôi nên cảm thấy buồn cười hay nên thấy bị xúc phạm nữa. Hơn nữa, sao Haato-chan lại nói ``lặng lẽ già đi một tuổi?''

Cắt cảnh.

\ct

Nakiri Ayame-chan. Vẫn là trong trang phục kimono hiện đại như mọi khi, nhưng lần này em ấy đang vắt chân trên ghế cao, phía sau là một dãy đĩa thức ăn trưng bày như tiệc cưới.

``Yo, Kuon-senpai đó sao?'' Ojou-san hớn hở nói. Em ấy nhảy xuống, đá chiếc ghế cao về phía trước, rồi dang rộng hai tay, nháy mắt như một tên nhóc nghịch ngợm: ``Bọn này không chỉ đơn thuần tổ chức cho anh đâu, mà còn muốn cho anh thấy một điều, rằng \hll\ không thể thiếu em được, yo!''

\dots Đúng, vẫn là Ayame-chan như thường ngày. Vẫn với phong cách lăng xăng như mọi khi.

Nhưng, tại sao?

Họ làm sao biết được tâm trạng của tôi lúc này?

\ct

Đoạn tiếp theo, tôi thấy Ookami Mio-chan, người đang đứng giữa một đống hộp bánh và chai nước giải khát. Trông em ấy hơi mệt, nhưng vẫn cười tươi: ``Kuon-senpai, sinh nhật là lúc để ăn mừng, nhưng quan trọng hơn, là lúc để nhắc cho người ta nhớ rằng họ không bao giờ cô độc.''

Nàng Sói đen bóc ra một chai nước, hướng về phía tôi, như thể đang mời vị khách quý: ``Tại sao anh không bọn em tận hưởng bữa tiệc sinh nhật này nhỉ?''

Một cô nàng Sói đen hiền hậu, tốt bụng và ngay thẳng. Cũng là một thành viên mang dáng dấp của người mẹ trong \hll. Không thể lẫn vào đâu được.

\ct

Tiếp đó là Usada Pekora-chan đang\dots\ thổi bóng bay? Cô nàng nhìn thẳng vào máy quay, tóc mái rủ xuống một bên mắt, nhưng thần thái vẫn cực kỳ ``chính chủ''.

Peko-chan buộc chiếc bóng vừa thổi đó, thả lên trên cao, rồi ngẩng mặt về máy quay trong khi tay lấy một chùm bong bóng khác: ``Chỉ vì anh không lên sân khấu thường xuyên như bọn này, không có nghĩa là anh ở ngoài cuộc đời tụi chị, peko. Nhớ kỹ điều đó, Kuon-senpai!''

\ct

Amane Kanata-chan hiện lên trong đoạn kế tiếp, đứng cạnh một bảng kế hoạch tổ chức tiệc chi chít chữ. Nàng Thiên sứ mỉm cười nhẹ, hướng mặt về phía máy quay: ``Nếu không có senpai, bọn em có thể sẽ không trở thành những người như thế này.''

Em ấy tiến về phía chiếc máy quay, rướn thân về phía trước và nở một nụ cười rạng rỡ: ``Kuon-senpai không đơn thuần là một tổng quản lý, mà thậm chí còn là một phần cốt lõi không thể thiếu, đúng chứ?''

\ct

Và cuối cùng, Flare-chan đang đứng chỉnh máy quay dành cho người tiếp theo. ``Rồi, Polka-chan chuẩn bị lên nào\dots''

Nhưng ngay khi em ấy vừa dứt câu, Omaru Polka-chan chồm từ phía sau, phấn khích thốt lên: ``Polka ở đâu đây\~{}? Polka ở đây này\~{}!\footnote{Nguyên gốc là câu chào khán giả của nàng Cáo sa mạc: ``ポルカおるか？ポルカおるよ！ (Polka oruka? Polka oruyo!)''.}'' Em ấy chồm lấy chiếc máy quay, phấn khích như một nghệ sĩ xiếc quá khích: ``Nếu không có Kuon-senpai, có khi Polka đây sẽ không thể hòa đồng với mọi người được! Anh đã là một phần trong gia đình của \hll\ rồi đó! Hiểu chưa nào!?''

Dần dà,  lần lượt từng người một bước vào khung hình, cho đến khi cả ba mươi thành viên còn lại mà công ty chúng tôi đang có, tất cả những người đã cùng tôi đi qua bao thăng trầm, hiện diện đầy đủ trong khung hình.

Mỗi người đứng ở một vị trí. Có người vẫy tay. Có người giơ biểu ngữ. Có người còn đội cả nón giấy sinh nhật.

Ở giữa, Sora-chan là người dẫn lời: ``Và giờ, Hikawa Kuon-senpai\dots\ hãy để bọn em dẫn anh nhìn lại những khoảnh khắc đáng nhớ nhất trong hành trình vừa rồi.''

Peko-chan tiếp lời: ``Bọn em muốn anh thấy rằng, một điều quan trọng nhất đó chính là: anh chưa từng, và sẽ không bao giờ, cô đơn trong đại gia đình này.''

Màn hình dần mờ đi một màu trắng xóa.

Tôi nhìn chằm chằm vào màn hình, bất giác nuốt khan.

Tôi không thốt được lời nào.

Trong đầu tôi vẫn chưa thôi băn khoăn. Tại sao\dots\ tại sao họ đều nghĩ giống nhau?

Yoisai, rồi Yamazu\dots\ và giờ lại đến tất cả những cô gái này nữa. Sao họ lại cùng có một suy nghĩ: rằng tôi đang cô đơn?

Chẳng lẽ tôi\dots\ đã để lộ điều gì đó? Hay họ chỉ đang nghĩ nhiều quá? Không biết nữa. Tôi cứ tưởng mình đã che giấu khá ổn rồi cơ mà\dots

Và tại sao họ lại biết hôm nay là sinh nhật tôi? Tại sao lại dành cho tôi nhiều tình cảm như vậy?

Tại sao tôi lại trở thành trung tâm của mọi chuyện? Được họ đối xử như thể tôi là người xứng đáng nhận sự tôn trọng ấy?

Tôi chưa từng nghĩ mình cần những thứ đó. Cũng không nghĩ rằng bản thân đủ đặc biệt để lọt vào mắt xanh của các cô gái ấy.

Nhưng họ lại cứ khẳng định như thế. Như thể\dots\ sự tồn tại của tôi thật sự quan trọng với họ.

Hàng tá dấu hỏi cứ hiện lên trong đầu tôi như bong bóng nổi lềnh bềnh trên mặt nước, và chưa kịp vỡ ra câu trả lời nào, thì đoạn phim nhìn lại những kỷ niệm xưa bắt đầu chiếu.

\ct

Ánh sáng dịu nhẹ từ máy chiếu phủ lên màn, và một cảnh tượng quen thuộc hiện ra.

\begin{center}
    \uline{(Mùa xuân năm đầu tiên\dots)}
\end{center}

Hikawa Kuon - là tôi - trong bộ quần áo tạm bợ và có đôi phần luộm thuộm, cùng với đôi dép lào kẹp chân, cũng là người đang cầm chính cái dép ấy đập túi bụi vào một con gián trên sàn. Trong góc khuất, Fubuki-chan, Matsuri-chan và Roboco-chan đang rúm ró, run như cầy sấy\footnote{Xem lại phần {\abv Truyện phụ}, {\flaregothic ÉPISODE 0}, quyển số Không.}.

Tôi bật cười. Hóa ra ngày đầu tiên đến công ty của tôi lại được mọi người lưu giữ như vậy.

\ct

\begin{center}
    \uline{(Mùa hè năm đó\dots)}
\end{center}

Tôi còn nhớ hồi tôi cùng với Robo-PON và Haato-chan đang tạo dáng kỳ lạ trong buổi chụp ảnh quảng bá. Sau đó có cả Ayame-chan và Aki-chan nữa. Nhìn cái cảnh họ cưỡi con gấu giả rồi làm mặt nghiêm túc trước ống kính\footnote{Xem lại phần {\abv Truyện phụ}, {\flaregothic ÉPISODE 6}, quyển số Một.}, tôi không nhịn được, khẽ thốt lên: ``Quả đó đúng là đỉnh cao của sự vô nghĩa\dots\ mà cũng đáng nhớ thật.''

\begin{center}
    \uline{(Một tháng sau\dots)}
\end{center}

Màn hình chuyển đến lúc tôi với Aki-chan và Miko đang lênh đênh trên dòng nước đen ngòm, mà vốn được tiết ra từ dưới váy của nàng Dị giới\footnote{Xem lại {\flaregothic ÉPISODE 10}, quyển số Một.}.

Bây giờ nhớ lại, tôi gần như đỏ bừng mặt khi nhìn lại cảnh khó đỡ ấy. Đến bây giờ, tôi vẫn muốn đặt câu hỏi: ``Thế quái nào em ấy lại có cơ chế ấy!? Và thứ nước đen ngòm đó ở đâu ra!?''

Trời đất, sao hồi đó tôi lại rơi vào cảnh oái oăm như thế được nhỉ?

\begin{center}
    \uline{(Đêm kinh dị đầu tiên, khi mùa thu tới\dots)}
\end{center}

Tôi cùng Sora-chan, Mio-chan, Subaru-chan và Haato-chan ngồi nghe kể chuyện ma trong ánh nến mờ ảo\footnote{Xem lại {\flaregothic ÉPISODE 17}, quyển số Hai.}.

Và rồi\dots\ tôi nhớ chứ.

Tôi và Sora-chan cùng Shion-chan đã lên kế hoạch hù họ một vố. Hiệu quả thì chắc các bạn cũng thấy rồi, đến mức độ nàng Sói đen còn\dots\ ``vãi ra quần''. Tôi khẽ cười khì. Có lẽ đó là lần đầu tiên tôi thấy ba người họ la lên trong tuyệt vọng như vậy.

Nhất là vẻ mặt ngượng ngùng của Mio-chan nữa. Tôi đã phải rối rít xin lỗi em ấy cả ngày hôm sau mới chịu tha.

\begin{center}
    \uline{(Lần đầu khi gặp sinh vật ngoài hành tinh\dots)}
\end{center}

Cảnh quay đến Koro-chan và Okayu-chan chơi đùa với một thứ đen sì, trông chẳng giống sinh vật ở Trái Đất. Tôi nhớ rõ cái khoảnh khắc Korone cố gắng tiếp cận đến cô bạn của mình để giúp đỡ, để rồi em ấy bị nàng Khuyển vô tình đấm thẳng lên trời, phá tung cửa sổ và tiếp cận một thiên thạch\footnote{Xem lại {\flaregothic ÉPISODE 24}, quyển số Hai.}.

Dĩ nhiên tất cả mọi thứ đều chỉ là dựng lại thôi, kể cả cái thứ đen sì đó, vốn được tạo bởi ma thuật của Shion-chan.

Đó là lần đầu tiên tôi vừa vào một vai diễn bất đắc dĩ, lại vừa đứng sau hậu trường\dots\ dựng cả một ``câu chuyện'' từ đống hỗn loạn đó theo kịch bản của Hokuro.

\begin{center}
    \uline{(Kế hoạch điên rồ của Marine\dots)}
\end{center}

Lúc đó, tôi cùng Fubu-chan và Mel-chan đang chôn chân đứng nhìn Marine-chan bay ra ngoài không gian cùng với con rắn xanh mà chính em ấy đã rước về bằng chiếc tên lửa\footnote{Xem lại {\flaregothic ÉPISODE 42}, quyển số Ba.}.

Sẽ không có gì đáng nói nếu loại tên lửa đó lại chính là\dots\ bộ đồ chơi Râu Đen được hóa lớn. Nghĩ lại thì đó hẳn không phải là một kế hoạch quá bài bản\dots

\dots vì sau cùng, em ấy cũng rơi xuống ngay tại vị trí ba đứa bạn chúng tôi đang ngắm sao. Tôi lẩm bẩm, khẽ cười chát: ``Đó cũng là lúc mình thấy em ấy\dots\ đáp xuống sân thượng trường mình lần đầu tiên.''

\begin{center}
    \uline{Cuộc họp online giả qua ``khung hình''\dots}
\end{center}

Tôi cùng với Ru-chan, Noel-chan, Peko-chan và Koro-chan đang giả vờ họp trực tuyến qua một cái khung tự làm, có lẽ là hưởng ứng phong trào làm việc tại nhà hồi dịch. Chúng tôi đã có những trò đùa khá vui nhộn, từ việc ăn bánh uống trà, đến việc cuộc thi vẽ vời nữa chứ\footnote{Xem lại {\flaregothic ÉPISODE 56}, quyển số Bốn.}.

Suisei sau đó bước vào và tỏ vẻ hờn dỗi vì không được chúng tôi mời, và rồi sau đó chúng tôi đã có những trò chơi vui nhộn. Tôi mỉm cười, nhớ lại âm thanh hỗn tạp và tiếng cười vang vọng trong căn phòng vào buổi chiều cuối mùa xuân hồi đó.

\begin{center}
    \uline{Mùi hôi thối ám ảnh\dots}
\end{center}

Đó là lần mà Flare-chan đến văn phòng lần đầu. Tôi chưa từng thấy cảnh tượng mà mọi người cùng nhau đồng loạt đuổi mùi, bịt mũi và quýnh quáng như vậy\footnote{Xem lại {\flaregothic ÉPISODE 58}, quyển số Năm.}. Lúc đó, chỉ có tôi là cảm thấy mùi cá thối đó là bình thường, có lẽ bởi vì tôi đã quen với những mùi khó chịu như vậy sao?

Mà, cũng phải thôi, mùi cá khô dù đã chế biến qua hay là mùi mắm tôm chẳng phải thứ dễ ngửi. Họ còn đổ rạp xuống đất ngay khi biết được chiếc lọ mắm tồn đó là đủ hiểu rồi.

``Họ vẫn ngốc như thế nhỉ\dots'' tôi khẽ thở dài, nhưng trong lòng lại thấy ấm áp đến lạ.

\begin{center}
    \uline{Tiệc mì trôi ống trúc trên bãi biển xan\dots}
\end{center}

Tôi cùng với các cô gái đã cùng nhau tổ chức tiệc mì trôi ống tre\footnote{Xem lại {\flaregothic ÉPISODE 69}, quyển số Năm}, mà sau này tôi mới biết đó là từ kịch bản ``tiêu dùng thông minh'' do Coco-chan và Haato-chan (dù khá miễn cưỡng) bày ra.

Sau bữa tiệc mì trôi đó, các cô nàng nhào ra biển như thể đã lên kế hoạch từ trước và có khoảng thời gian bên nhau. Dưới ánh chiều hoàng hôn ngày đó, tôi cứ đứng nhìn họ cười đùa, và rồi trong phút chốc, tôi thoáng nghĩ: ``Mình thật sự đang sống trong một thứ gì đó rất đẹp.''

\begin{center}
    \uline{Trận chiến sức chịu đựng ly kỳ nhất lịch sử\dots}
\end{center}

Đó là lần mà Sui-chan và Senchou tranh cãi nhau xem ai là người chịu đựng giỏi hơn, dẫn đến việc họ tự tổ chức một cuộc thi, và tôi cùng Fubu-chan làm trọng tài và bình luận viên, trong khi Kana-chan làm khán giả bất đắc dĩ\footnote{Xem lại {\flaregothic ÉPISODE 72}, quyển số Sáu.}.

Kỳ thực, đó quả là một cuộc thi khá vô nghĩa, mà thậm chí tôi còn nghe thoảng đâu đó rằng Marine-chan đã bị gãy lưng và phải nằm nghỉ ở nhà suốt một tuần liền.

Nhưng cũng chính từ đó, tôi bắt đầu thấy mọi người dần gắn kết\dots\ theo cái cách rất riêng của họ.

\begin{center}
    \uline{MV cà ri\dots}
\end{center}

Một đoạn hậu trường quay thước phim quảng bá cho mì cà ri. Họ quay MV như những idol thực thụ. Những động tác vũ đạo dù đẹp mắt nhưng trông ngớ ngẩn, nụ cười rạng rỡ\dots\ và tiếng hát vang vọng.

Tất cả bắt đầu từ việc một tên mì cà ri đột ngột xuất hiện ở xó nào đã thôi miên tất cả chúng tôi\footnote{Xem lại {\flaregothic ÉPISODE 75/SPÉCIAL {\ab XII}}, quyển số Sáu.}.

Và phải đến khi Subaru-chan và tôi lao vào giải cứu, mọi chuyện mới trở lại bình thường. Nhưng cũng chính sự việc hỗn loạn đó, tất cả mọi người đều nghĩ ra ý tưởng cho thước phim quảng bá đó.

Tôi không thể phủ nhận được nữa. Dù gì thì họ cũng là những thần tượng chuyên nghiệp, và\dots\ tôi thật sự thấy tự hào vì được đứng cùng họ, dù chỉ là đứng vai trò quản lý và không tham gia trực tiếp.

\begin{center}
    \uline{(Roboco bị vỡ vụn\dots)}
\end{center}

Thời điểm đó là vào thu, khi gia đình chúng tôi đang phải chịu tang cho người anh trai của bố tôi - một người mà đến cả bác Hijaki-san và bố tôi còn phải nói rằng ``ông ấy đếch có phẩm chất nào tốt cả.'' Nhưng mà ông ấy không phải là nhân vật chính.

Mà đó lại là cô nàng Rô-bốt vụng về kia cơ.

Tôi thấy mình hồi đó đang cắm cúi sửa Roboco-chan, đầu vẫn đeo khăn tang, trong khi Ru-chan và Choco-chan còn đang cố giữ lấy Noel-chan đang hưng phấn quá đà\footnote{Xem lại {\flaregothic ÉPISODE 84}, quyển số Bảy.}.

Khi tôi mắng Roboco vì tội lười bảo dưỡng, em ấy chỉ cười trừ, thậm chí sau đó còn phải nhờ tôi làm bác sĩ tâm lý bất đắc dĩ nữa.

Nhìn lại, tôi khẽ lắc đầu: ``Đúng là đại PON\dots\ không bao giờ thay đổi.''

\begin{center}
    \uline{Khi Giáng Sinh gặp mặt Halloween\dots}
\end{center}

Bộ ba ``Ác ma'' (tên tạm) gồm Mel-chan, Choco-chan và Shion-chan cứ khóe nhau ngày hôm đó là Halloween, trong khi ngày đó đã xa rời tận gần một tháng rồi, và tôi nhớ Ru-chan đập bàn ghê lắm\footnote{Xem lại {\flaregothic ÉPISODE 85}, quyển số Bảy.}.

Nhưng đến cuối cùng, chúng tôi đã cùng nhau hát ngân vang bài ``Jingle Bells'', chờ đón vào một năm mới, và cũng là buổi Fes thứ hai của công ty nữa.

Tôi bất giác mỉm cười khi nhìn lại cảnh mọi người trang trí căn phòng cho hợp với ngày lễ đó: ``Đó là cái ngày lễ hội\dots\ kỳ quái nhất nhưng cũng đáng nhớ nhất trong đời mình.''

\begin{center}
    \uline{Khi những con người thuộc thế hệ mới lần đầu gặp mặt\dots}
\end{center}

Tiếp nối những hình ảnh thân quen trước đó, màn chiếu lại sáng lên. Tôi thấy lần lượt những cảnh quay trong quá khứ tiếp tục hiện ra, đó là những khoảnh khắc không thể quên.

Đó là hồi mà tôi lần đầu chạm trán với ``NePoLaBo'', bốn cô gái đến từ Thế hệ thứ Năm.

Tôi nhớ lại cái cách mà Polka-chan bày ra chiến dịch giải cứu Fubu-chan khỏi cặp đôi Chó-Mèo nghịch ngợm\footnote{Xem lại {\flaregothic ÉPISODE 90}, quyển số Bảy.}, để rồi sau đó em ấy cùng với tôi thực hiện phi vụ thử thách cho cặp đôi đó\dots

\dots hay cái cách Lamy-chan đấm thùm thụp con gấu Bắc Cực\footnote{Xem lại {\flaregothic ÉPISODE 91}, quyển số Bảy.} - là Daifuku nếu tôi nhớ không nhầm - vì phản chủ, khiến chúng tôi cảm thấy vừa thương vừa buồn cười\dots

\dots hay việc Nene-chan thì hớn hở với ``sức mạnh mùa xuân'' được Thần linh (Miko-chi) ban cho\footnote{Xem lại {\flaregothic ÉPISODE 106}, quyển số Tám.}, mà cuối cùng tất cả chỉ diễn ra trong mơ\dots

\dots thậm chí là Botan-chan cũng không thoát khỏi những sự ngớ ngẩn tại căn phòng \hll\ này, khi em ấy bất đắc dĩ phải ngồi trên chiếc xe ba bánh theo hướng dẫn của Luna-chan\footnote{Xem lại {\flaregothic ÉPISODE 92}, quyển số Bảy.}, tất cả cũng chỉ vì những ``bài học'' nghe xàm lơ.

Tôi đã từng nghĩ đó chỉ là một ngày làm việc như bao ngày khác. Nhưng khi nhìn lại\dots\ có điều gì đó trong tim tôi như thắt lại. Những người mới này đã trở thành một phần không thể thiếu của đại gia đình \hll, và tôi, tôi đã may mắn được chứng kiến họ từ những bước đi đầu tiên.

\begin{center}
    \uline{Khi chia tay một trong những thành viên yêu quý nhất\dots}
\end{center}

\dots Rồi đến đoạn tiếp theo, khi mà chúng tôi tổ chức buổi ``lễ tốt nghiệp'' cho Kiryu Coco-chan, người mà sau cùng cũng ``đóng quân'' ngay tại nhà chúng tôi.

Tôi bất giác nín thở khi thấy lại ánh mắt rưng rưng của các thành viên, nhất là khi nàng Rồng sải cánh bay ra khỏi văn phòng, mà chúng tôi tưởng đó là lần cuối được gặp em ấy\footnote{Xem lại {\flaregothic ÉPISODE 112/SPÉCIAL {\ab XX}}, quyển số Chín.}. Lúc đó, mọi người đều cố tỏ ra mạnh mẽ, vui vẻ, như thể đó là một buổi chia tay nhẹ nhàng\dots\ nhưng tôi biết, tận sâu trong tim, ai cũng đang buồn.

Tôi thấy lại những cô gái đứng trên sân khấu ấy, họ gắng gượng mỉm cười, cổ họng nghẹn lại vì biết rằng mình chẳng thể giữ em ấy ở lại, dù đã cố gắng rất nhiều.

``Ngẫm nghĩ lại thì quả là tiếc cho em ấy thật\dots\ Nhưng thôi, dù sao em ấy cũng đã ở đây với chúng ta rồi,'' tôi thầm nhủ, cố gắng an ủi chính mình bằng ký ức tốt đẹp.

\begin{center}
    \uline{Ngày Tết Nguyên Đán đầu tiên của \hll\dots}
\end{center}

Đó là cái Tết đầu tiên trong suốt hơn hai mươi nồi bánh chưng, tôi được ăn một ngày lễ truyền thống xứ Etsunan cùng với các cô gái đến từ xứ Phù tang\footnote{Xem lại trong quyển T.}.

Đó cũng là lần đầu họ tụ họp đông đủ tại nhà Hikawa (theo một cách không ai ngờ tới) để cùng nhau ăn Tết, vốn dĩ là một phong tục của riêng quê tôi, nhưng tất cả đều hưởng ứng như thể đó là chuyện thường ngày. Tôi nhớ tất cả chúng tôi đã dựng những ngày Tết đó làm thành một bộ phim, và nó đã rất thành công ngoài rạp.

Từ đó, những mảnh đời thường của các cô gái càng được bộc lộ rõ hơn.

Làm sao tôi có thể quên được cả ngày hai mươi chín dọn nhà cùng gia đình chúng tôi, buổi chiều ngày ba mươi Tết gói bánh chưng với mỗi người ra một hình thù khác nhau, cuộc đua xe về quê nội tôi vào ngày mùng Một Tết, buổi du hành về quê ngoại Neihei vào mùng Hai và những sự kiện ``uống nước nhớ nguồn'' vào ngày mùng Ba chứ?

Dù tất cả chỉ để quay phim, nhưng niềm vui, tiếng cười hôm đó là thật. Những con người từ một thế giới khác, ngồi lại bên nhau như một gia đình, và tôi, là người - dù cách này cách khác - đã đưa họ đến đó. Tôi thấy mình nhỏ bé\dots\ nhưng cũng thấy mình thật may mắn.

\begin{center}
    \uline{Và rồi, Hikawa Kuon đã tạo cho chúng ta nhiều kỷ niệm khác nữa\dots}
\end{center}

Màn hình bắt đầu tua nhanh. Từng khoảnh khắc, từng đoạn phim chớp hiện, từng thước phim lướt qua trong chớp mắt.

Nhưng điều đặc biệt nằm ở chỗ, là âm thanh đi kèm.

\il

``Những kỷ niệm đáng nhớ này\dots''

``\dots nó đáng trân trọng, đúng chứ?''

``Dù có vui, có buồn, có giận, có nhây đến cỡ nào\dots''

``\dots chúng cũng là một phần trong dòng kỷ niệm đó.''

\il

Tôi sững người lại. Đó là giọng nói của họ.

Những giọng nói của ba mươi con người đã làm nên điều kỳ diệu cho \hll, và cho chính cả gia đình Hikawa chúng tôi.

\il

``Có thể với mọi người, chúng em là những idol chuyên nghiệp, được nhiều người khen ngợi\dots''

``\dots nhưng trong con mắt của anh, bọn em cũng chỉ là một đám hề thôi, đúng không?''

\il

Tôi bật cười, một nụ cười xen lẫn nước mắt vì có phần hài hước, nhưng đâu đó chan chứa một nỗi xúc động. Phải rồi, tôi từng gọi họ như vậy, và đến bây giờ tôi vẫn sẽ gọi như thế. Một cách trìu mến. Một cách chân thành.

\il

``Bọn em không giận anh điều đó. Bọn em biết mình là ai mà.''

``Nhưng cũng nhờ những trò điên rồ đó\dots''

``\dots mà đã tạo nên \hll\ chúng ta\dots''

``\dots và anh cũng góp phần vào những trò `mất não' đó đấy, Hikawa Kuon-senpai à.''

\il

``Không thể nào\dots'' tôi thầm nghĩ, tim đập nhanh hơn. Họ thật sự đã chuẩn bị cả điều này cho tôi sao?

\il

``Nếu không có anh ở đây\dots''

``\dots chúng em còn không biết mọi chuyện sẽ điên rồ đến mức nào\dots''

``\dots và có thể bọn em sẽ không được như ngày hôm nay.''

``Anh đã giúp bọn em hòa nhập với thế giới này\dots''

``\dots hy sinh vì bọn em rất nhiều thứ\dots''

``\dots anh thậm chí còn lắng nghe những lời tâm sự của chúng em nữa.''

\il

Tôi không dám thở mạnh. Mỗi câu nói như đánh thẳng vào lồng ngực, để lại một cảm giác ấm nóng khó tả.

\il

``Anh là người đã truyền cảm hứng cho chúng em đó, Kuon-senpai à.''

``Có thể anh cho rằng đó là việc nhỏ\dots''

``\dots nhưng với \hll\ này, nó là một việc rất ý nghĩa.''

``Chúng em muốn cảm ơn anh\dots''

``\dots vì những gì anh đã làm với chúng em\dots''

``\dots để có thể trở nên như thế này ngày hôm nay.''

``Chúng em chỉ muốn nói\dots''

``\dots cảm ơn anh rất nhiều, Hikawa Kuon-senpai\dots''

``\dots và bọn em muốn chúc mừng sinh nhật anh, ngay bây giờ!''

\il

Tôi đứng chết lặng. Mọi thứ quanh tôi như ngưng đọng. Những lời đó\dots\ là thật sao? Họ tuôn ra những câu nói đó bằng cảm xúc thật của mình ư?

Tôi cảm thấy như có điều gì đó nghẹn lại ở cổ họng. Tay tôi siết chặt lấy vạt áo, tim đập dồn dập.

Một thứ cảm xúc dào dạt trong con người tôi. Một thứ gì đó đang thôi thúc tôi, rằng những lời đó là thật.

``Không lẽ\dots\ mình đã đóng góp nhiều như vậy trong suốt hơn hai năm sao?'' tôi tự hỏi, đầu óc quay cuồng.

``Nhưng\dots\ mình chưa xứng đáng với sự tôn trọng này\dots''

Tôi vẫn còn chưa thể tin được. Chưa thể hiểu nổi tại sao tôi, vốn chỉ là một người bình thường, một cậu quản lý sống sót nhờ may mắn, và có được công việc này từ một sự tình cờ, lại được mọi người yêu quý đến vậy.

Giữa cơn xúc động ấy\dots\ màn chiếu vụt tắt.

Toàn bộ căn phòng trở lại màu đen kịt.

Tôi đứng đó, giữa bóng tối\dots\ với một trái tim vẫn còn đang run rẩy, không dám tin những gì vào sự thật.

\dots

\dots

Bất chợt, một bài nhạc quen thuộc vang lên. Giai điệu của nó nhẹ nhàng, êm ái, có chút hưởng ghi-ta.

Tôi đã từng nghe nó rất nhiều lần trong các buổi sinh nhật ở quê tôi, từ những dịp sinh nhật đơn sơ trong gia đình, tới những buổi sinh nhật của vài người bạn mà cậu thân trong lớp hồi cấp hai, cáp ba, hay thậm chí là cả trong những bản tin tức, những chương trình trên truyền hình quê nhà.

Không gì khác, đó là bài hát ``sinh nhật quốc dân'' của xứ Etsunan chúng tôi - ``Khúc hát mừng sinh nhật'' của một người có đầu trọc lốc\footnote{Ám chỉ đến Phan Đinh Tùng, người này sau đó vào 8/2024 bị dính phải lùm xùm liên quan đến chính trị (mặc dù theo nam ca sĩ, đây là một sự vô tình) và ở thời điểm tập truyện này được mở rộng (5/2025), vụ việc này đã được lắng xuống. Bài hát ``\textit{Khúc hát mừng sinh nhật}'' là một trong những bài hát phổ biến nhất trong các buổi sinh nhật, thậm chí có phần nổi hơn so với bài ``\textit{Happy Birthday}'' gốc.}. Một bản nhạc tưởng như bình dị, nhưng khi được cất lên vào đúng khoảnh khắc này\dots\ nó khiến trái tim của tôi đập mạnh hơn bao giờ hết.

Nhưng hôm nay, bài hát đó vang lên không phải từ một chiếc loa điện tử nào hay bất cứ ai.

Mà\dots\ nó lại vang lên từ người thân của tôi.

Trong lúc tôi còn đang ngỡ ngàng với bản nhạc quen thuộc kia, một ánh sáng phụt lên. Cây pháo sinh nhật nhỏ lóe sáng từ phía sau cánh cửa.

Bố, mẹ, và cô em gái Ryuuhou của tôi. Những người cầm một miếng bánh có cắm nếm.

Từng người một, họ bắt đầu ngân vang bài hát sinh nhật ``quốc dân'' bừng tiếng Etsunan mẹ đẻ của tôi.

``\vie{Ngày hôm nay ta cùng hân hoan nơi đây\dots}''

``\vie{Mọi người bên nhau, ta hát mừng sinh nhật\dots}''

``\vie{Một, hai, ba, ta cùng thổi tắt nến\dots}''

Và rồi, cả ba người cùng nhau đồng thanh, nở một nụ cười tươi và hoàn thành bài hát: ``Happy birthday, happy birthday to you! (Chúc mừng sinh nhật, chúc mừng sinh nhật anh/con trai!)''

Nhưng những câu hát của bài vẫn chưa kết thúc.

Khi tiếng hát của họ vừa dứt, những ngọn nến khác dần được thắp sáng. Và rồi từng khuôn mặt quen thuộc lần lượt xuất hiện, như thể tôi đang ở trong một giấc mơ.

Sora-chan, Sui-chan, Miko-chi, Roboco-chan, và cả A-san.

Họ không nói gì cả, chỉ tiếp tục hát phần tiếp theo của bài hát bằng tiếng Anh, dù cho việc họ phát âm không hoàn hảo, nhưng đủ để tôi hiểu tấm lòng.

``\eng{On this day, all together we'll be,}''

Tiếp đến là Mel-chan, Fubu-chan, Matsuri-chan, Aki-chan và Haato-chan, những thành viên hát tiếng Anh sõi hơn đôi chút:

``\eng{And we'll all sing for your birthday,}''

Rồi đến hội Thế hệ thứ Hai, gồm Aku-tan, Shion-chan, Ayame-chan, Choco-san và Suba-chan:

``\eng{One, two, three, we'll blow out the candles,}''

Mio-chan, Okayu-chan, Koro-chan, Peko-chan và Flare-chan tiếp lời:

``\eng{Happy birthday, happy birthday to you!}''

Ru-chan, Noel-chan, Marine-chan, Kan-chan và Towa-chan nhắc lại câu hát đó: ``\eng{Happy birthday, happy birthday to you!}''

Cuối cùng, sáu người gồm Watame-chan, Luna-chan, Lamy-chan, Nene-chan, Botan-chan và Pol-pol kết thúc bằng câu:

``\eng{And happy birthday, happy birthday to you!}''

Và ngay sau câu hát đó, đèn bật sáng.

Trước mặt tôi là ba mươi gương mặt quen thuộc của nhà Tam Giác Xanh, đứng cạnh gia đình tôi. Tất cả đang cùng vỗ tay, cùng reo hò: ``\textbf{Chúc mừng sinh nhật Kuon(sa)-senpai/san/anh/con trai!}''

Tôi không nói được gì.

Ngực tôi như bị thắt lại. Mắt tôi cay xè, rưng rưng nước mắt. Tôi chưa từng nghĩ chỉ một bài hát thôi lại có thể khiến mình xúc động đến thế.

Có lẽ bởi vì\dots\ chưa bao giờ tôi được nghe người thân mình hát bài này trong một dịp đông đủ như vậy, ở ngay tại Etsunan này, giữa một đại gia đình, mà tôi chưa bao giờ nghĩ là sẽ có.

Miệng mở ra nhưng không nói được câu nào. Trái tim tôi như đang bị nhấn chìm trong một thứ cảm xúc không gọi thành tên.

Vui mừng? Biết ơn? Xúc động?

Tất cả đều đang hòa lẫn vào nhau, một sự giằng xé nội tâm khi tôi không biết phải phản ứng ra sao. Tôi chỉ biết đứng đó, nhìn từng người trong số họ, những người đã âm thầm sắp đặt mọi thứ này mà tôi không hề hay biết, và rồi đã khiến tôi hoàn toàn bất ngờ.

Lâu nay, tôi luôn tổ chức sinh nhật theo cách riêng của mình: kín đáo, yên tĩnh, chỉ cùng gia đình. Dần dần nó trở thành thói quen, và tôi cũng quen với việc đó.

Thậm chí, tôi cũng đã nghĩ đến việc tôi sẽ phải dần làm quen với việc quên những bữa tiệc sinh nhật đi vì chúng còn không đáng quan tâm nữa.

Vậy mà giờ đây, ngay tại thời khắc này, lần đầu tiên trong đời, tôi được mọi người - không chỉ mỗi gia đình, mà còn có các thành viên \hll\ - tổ chức sinh nhật như thế này.

Một giọng nói vang lên trong lúc tôi còn chưa hết sững sờ. Đó là A-san, người mà tôi không ngờ nhất. Chị ấy hớn hở: ``Em thấy bất ngờ chứ?''

Tôi chỉ kịp ấp úng, miệng vẫn chưa ngậm lại: ``C-Có ạ- Mà A-A-san\dots!? Sao chị lại\dots\ \eng{Cái quái gì vậy?}''

Tiếng cười rộ lên khắp căn phòng với phản ứng ngỡ ngàng này.

Tôi gãi đầu, đỏ mặt. Tôi hỏi vội: ``S-Sao mấy đứa biết hôm nay là sinh nhật anh? Anh còn chưa nói với ai mà\dots''

Fubuki-chan vẫy vẫy tay, cười rạng rỡ: ``Có gì đâu. Là A-chan đã tiết lộ cho bọn em đó!''

Tôi tròn mắt: ``Vậy mấy cái này là\dots''

Koro-chan chen vào, tươi như chưa từng thấy: ``Nhờ A-chan đã bày ra cho bọn em á\~{}''

Tôi sững lại với câu nói đầy ngây thơ và tươi tỉnh từ hai người.

Lúc đó, tôi vẫn còn đứng ngơ ngác, không thể tin vào những gì đang xảy ra. Trái tim tôi đập thình thịch, một phần vì sự bất ngờ, một phần vì tôi cảm nhận được tình cảm chân thành từ mọi người xung quanh.

Họ đã làm tất cả vì tôi. Nhưng trong sâu thẳm lòng mình, tôi vẫn cảm thấy một chút gì đó ngần ngại. Cảm giác như mình không xứng đáng với điều này, như thể mình không thực sự cần một sinh nhật to lớn như vậy.

Có lẽ các bạn sẽ không tin, nhưng từ khi tôi biết được mùi ``đời'' là gì khi vào cấp ba, tôi đã bắt đầu không còn mặn mà với cái gọi là sinh nhật nữa.

Tôi luôn giữ thói quen âm thầm, chỉ mừng tuổi mình trong không gian riêng tư cùng với gia đình. Không cần bạn bè, cũng không cần với những người có quan hệ, chỉ quây quanh với những người thân thuộc nhất là đủ.

Tôi không muốn phiền đến ai, càng không muốn mọi người cảm thấy có trách nhiệm phải nhớ đến ngày này. Có những lúc tôi nghĩ, nếu cứ để mọi thứ diễn ra như vậy, tôi sẽ không phải gánh lấy những sự chú ý mà mình không muốn.

Nhưng sao hôm nay lại khác?

Một phần trong tôi vẫn chưa chấp nhận được sự bát nháo này. Cảm giác lạ lẫm, bị thu hút vào tâm điểm chú ý của cả đám đông khiến tôi chỉ muốn lùi lại một bước, thậm chí chỉ muốn tìm một cái lỗ để chui xuống.

Vậy mà, khi nhìn thấy ánh mắt của A-san, tôi lại không thể làm thế. Chị ấy nhìn tôi, dịu dàng và kiên định, như thể biết rõ tôi đang suy nghĩ gì.

``Hikawa-san,'' giọng nói đó bất chợt vang lên, như đang cắt đứt mạch cảm xúc đang dâng trào trong lòng tôi. ``Cậu có biết rằng cậu cũng xứng đáng có được một ngày mà mọi người yêu thương không?''

Chị ấy muốn tôi hiểu được cảm giác hai chữ ``quan tâm'' và ``coi trọng'' như thế nào.

Cả nhóm xung quanh đều im lặng, như thể họ đang chờ đợi tôi trả lời.

Tôi nuốt khan, thở ra một hơi dài, rồi quay sang người đồng nghiệp của mình: ``Nhưng mà chị\dots\ em không thích chuyện này. Được mọi người chú ý, em cảm thấy\dots''

Tôi không thể nói hết câu. Tôi rất muốn phủ định chuyện đó, nhưng không thể tìm được lời nào để nói.

A-san mỉm cười, nụ cười ấm áp như bao lần tôi từng thấy, nhưng lần này nó khác. Chị ấy khẽ lắc đầu, rồi nắm lấy tay tôi: ``Cậu không cần lúc nào cũng cần phải chịu đựng một mình đâu, Hikawa-san. Cậu cũng xứng đáng được người khác quan tâm lắm chứ. Những điều này không phải để làm phiền cậu, mà là để cậu cảm nhận được tình cảm của mọi người.''

Tôi im bặt, không biết phải phản biện như thế nào.

ời nói ấy như một làn sóng, cuốn đi những suy nghĩ tiêu cực trong đầu tôi. Các cô gái khác cũng đồng loạt đến gần, vây quanh tôi, tất cả đều cười, và những cái ôm ấm áp bắt đầu kéo tôi vào. Tôi cảm thấy mình như bị đắm chìm trong một đại dương tình yêu và sự quan tâm từ tất cả những người xung quanh.

Một thành viên trong nhóm mà tôi không để ý ôm chầm tôi từ phía sau, một người khác vỗ nhẹ vào lưng tôi.  Tôi cảm nhận được cơ thể mình dần dần nhẹ nhõm đi, nhưng cũng không thể không thừa nhận rằng trong cái cảm giác ấy, có sự xúc động dâng lên mạnh mẽ.

Tôi đã từng nghĩ mình không cần điều này, nhưng giờ tôi lại nhận ra rằng có những khoảnh khắc mà mọi người muốn dành cho mình, và trong khoảnh khắc đó, tôi biết mình không nỡ từ chối tấm lòng đó.

``Đừng bao giờ cảm thấy mình không xứng đáng đâu đó, Kuon-senpai!'' đôi mắt long lanh của Haato-chan nhìn về phía tôi. ``Anh cũng là một phần trong gia đình \hll\ mà, đúng không?''

Tôi không nói gì nữa. Không một từ ngữ nào có thể miêu tả cảm giác lâng lâng này.

Tôi chỉ có thể đứng đó, nghẹn ngào trong cảm xúc lạ lùng khi tất cả mọi người, từ Sora-chan, Miko-chi, rồi dần dà toàn thể ba mươi người ôm tôi vào lòng. Họ cùng nhau, không một lời nào, chỉ trao cho tôi sự ấm áp và tình cảm chân thành. Một cái ôm đầy sự thấu hiểu, không ai ép buộc, không ai làm tôi cảm thấy áp lực, chỉ đơn giản là họ muốn tôi cảm nhận.

Nhưng khi tôi ngẩng lên, tôi thấy từ xa, ngoài A-san, bố mẹ cùng với em gái Ryuuhou cũng đang đứng nhìn tôi. Họ không nói gì, chỉ khẽ mỉm cười. Đó là nụ cười của sự chấp nhận, của tình cảm gia đình. Một nụ cười ấm áp, nhưng cũng chứa đựng một chút gì đó lặng lẽ, khi họ là những người thân cận nhất với tôi.

Những khoảnh khắc đó đọng lại trong tâm trí tôi mãi mãi. Tôi vẫn nhớ cái cảm giác ấy, không vì cái bánh kem hay những lời hát vang lên, mà là cảm giác của một gia đình thật sự.

Một gia đình mà tôi đã dần dần tìm thấy ở đây, nơi tôi không còn phải chịu đựng một mình, nơi tôi được yêu thương theo cách mà tôi chưa bao giờ nghĩ tới.

Và trong thoáng chốc, tôi cũng đã tự nhủ: ``\textit{Thật may mắn khi mình làm công việc này.}'' Thành thực mà nói, đây đúng là một công việc mà đám đàn ông con trai nào cũng mơ ước được chạm vào, theo lời của Yamazu.

Dù biết rằng, chức Tổng quản này là một chức vụ cần sự trách nhiệm rất cao, và cũng vì thế mà áp lực nhiều hơn\dots

\dots nhưng đổi lại, tôi lại nhận được sự tôn trọng của tất cả mọi người. Tôi đã hy sinh rất nhiều vì họ, làm việc một cách lặng lẽ mà không cần ai công nhận\dots

\dots và có lẽ, đến lúc một ngày nào đó, họ sẽ trả ơn tôi. Họ thật sự tín nhiệm tôi, tin tưởng tôi đến mức này đây.

\pov{-}

A-chan đứng từ xa, ánh mắt nhẹ nhàng quan sát khi các cô gái nhà Tam Giác Xanh quây quần quanh Kuon. Họ ôm chầm lấy cậu, trao cho cậu Tổng sự ấm áp và tình cảm như một món quà không thể nào đong đếm được. Tiếng cười rộ lên khắp phòng, hòa vào tiếng vỗ tay, và trong không gian ấy, cậu dường như trở nên nhỏ bé hơn bao giờ hết, như một đứa trẻ không thể từ chối tình yêu thương.

Cô mỉm cười, nhưng trong lòng lại cảm thấy có chút gì đó không thể diễn tả thành lời, như thể cô vẫn chưa thực sự hiểu, bởi lẽ\dots\

Đây là một dịp đặc biệt, một khoảnh khắc hiếm hoi mà cô đã cùng tất cả mọi người dốc sức chuẩn bị cho người quản lý mến thương của họ. Nhưng điều làm cô suy nghĩ không phải là những lời chúc mừng hay những cái ôm ấm áp, mà là việc cô không thể tin nổi làm sao mà một người như Kuon, luôn kín đáo và ít chia sẻ về bản thân mình, thậm chí lại giữ im lặng về ngày sinh nhật của mình.

Cô hồi tưởng lại khoảnh khắc khi cô cùng với Flare và Mio hỏi gia đình nhà Kuon tại sao lại như thế, để rồi họ nhận về câu trả lời không thể nào tội nghiệp hơn.

\begin{center}
    \uline{Hồi tưởng\\Bốn giờ hai mươi chiều.}
\end{center}

Khi tất cả các thành viên của nhóm trang trí phòng đang thực hiện những công đoạn cuối, Flare gật đầu, tỏ vẻ khá ưng ý: ``Vậy là tốt rồi đấy! Chẳng mấy chốc ta sẽ xong sớm thôi.''

Thế nhưng, A-chan lại không nói gì, chỉ đứng khoanh tay đầy suy tư. ``A-chan?''

Bất chợt, nàng Đeo kính quay sang cô, nét mặt lo lắng nhưng cũng hiện đôi chút đầy tò mò: ``Cô có nghĩ là Hikawa-san sẽ bất ngờ như thế này không? Tôi thật sự không hiểu tại sao cậu ấy lại không nói gì về sinh nhật của mình.''

Đó là điều mà cô cảm thấy lạ, mặc dù cô biết rõ ngày sinh của cậu Tổng - và đó cũng là lý do vì sao cô gọi cho cậu từ đêm qua và thông báo cho mọi người từ sáng nay.

Nhưng câu hỏi là, tại sao lại như thế? Đó không phải là điều mà ai cũng biết.

Và như dự đoán, nàng Dị giới chỉ nhún vai, khuôn mặt tỏ chút bối rối: ``Em cũng chẳng hiểu nữa. Anh ấy luôn giữ mọi thứ cho riêng mình, chẳng bao giờ chia sẻ gì nhiều về bản thân. Nhưng chắc chắn là anh ấy phải muốn một cái gì đó đặc biệt, không phải sao?''

Mio đứng gần đó, đôi tai sói dựng lên, đôi mắt cô cũng thể hiện đồng quan điểm: ``Flare-chan nói đúng đó, chuyện này khá kỳ lạ. Chúng ta chuẩn bị mọi thứ, nhưng chẳng ai biết sinh nhật của anh ấy có nghĩa là gì với anh ấy nữa. Có khi nào Kuon-senpai không muốn chúng ta làm vậy không?''

A-chan nhìn về phía gia đình Kuon, nơi Ryujin, Kaien và Ryuuhou đang đứng từ xa, chăm chú nhìn cả nhóm với nụ cười đầy ấm áp, trò chuyện với nhau khá rôm rả.

Cô tiến lại gần họ, nhẹ nhàng mở lời: ``Xin lỗi mọi người, cháu có thể hỏi một chút không? Tại sao Hikawa-san lại không bao giờ nói gì về sinh nhật của mình ạ?''

Flare và Mio cũng cùng lúc tiếp bước, ánh mắt đầy mong đợi. Họ muốn biết tại sao cậu Tổng của họ lại không bao giờ chia sẻ một điều gì liên quan đến ngày sinh nhật của mình, trong khi cả nhóm lại cố gắng chuẩn bị tất cả mọi thứ cho cậu (mà chưa chắc họ đã hiểu hoàn toàn lý do vì sao, nhưng vẫn làm vì theo đà).

Câu hỏi của A-chan, Flare và Mio dường như đã chạm vào một điều gì đó sâu sắc mà gia đình Kuon hiểu rõ hơn bất kỳ ai.

Ryuuhou khẽ trố mắt ngạc nhiên, ngẩng mặt lên bố mẹ, như thể muốn nhận một tín hiệu: ``\vie{Bố mẹ ơi, mình nên giải thích họ sao đây?}''

Sau một thoáng im lặng, Kaien khẽ mỉm cười rồi quay lại nói với họ, giọng nhẹ nhàng nhưng cũng đầy sự thấu hiểu: ``Chả là\dots\ cô chú không muốn làm con trai cảm thấy áp lực. Cậu ấy là người rất kín đáo và không muốn tạo thêm sự chú ý vào bản thân mình.''

Ryujin tiếp lời: ``Bản thân con nhà chú khi lớn lên cũng cảm thấy ngày này không có gì quan trọng mà. Nó luôn cảm thấy sự chú ý của mọi người có thể khiến nó khó chịu, nhất là khi nó làm thay đổi không khí bình yên mà cậu ấy muốn.''

Cả ba cô gái đều ngạc nhiên với câu trả lời của hai người lớn, nhưng rồi cũng gật gù hiểu được vấn đề. Nàng Dị giới hỏi tiếp: ``Vậy\dots\ ý cô chú là anh áy không thích nhận sự chú ý như vậy ạ?''

Người phụ nữ nội trợ kiêm chủ kho thở dài, ánh mắt đồng cảm hướng về ba cô gái: ``Nói thế cũng không hẳn. Nó\dots\ nói sao nhỉ, cảm thấy sự quan tâm quá mức đôi khi làm mọi thứ trở nên không tự nhiên ấy. Kuon không muốn bất kỳ ai cảm thấy lo lắng cho nó, mà chỉ muốn mọi thứ diễn ra thật yên bình và như lẽ thường thôi.''

Cả ba cô gái nghiền ngẫm lời giải thích của bà, rồi quay sang nhìn nhau, tựa như đã hiểu. Nàng Sói đen lên tiếng: ``Vậy là từ trước đến giờ, Kuon-senpai luôn tự mình âm thầm tổ chức sinh nhật mà không muốn ai biết?''

Cô em gái của cậu Tổng quản lên tiếng: ``Đúng là như vậy ạ. Onii-san thường chỉ đón sinh nhật với gia đình thôi. Em không thấy anh ấy mời bạn anh ấy, hay bất kỳ ai mà anh ấy thân quen tới. Những năm trước nhà em còn mời nhiều người từ bên nhà ngoại nữa, nhưng vài năm trở lại đây những bữa tiệc đã trở nên lặng lẽ hơn rồi ạ\dots''

Người đàn ông cơ khí kết luận: ``Đúng như Ryuuhou nói đấy. Nó dường như muốn làm mọi thứ một mình. Mặc dù chú thấy đôi lúc nó sẽ hơi buồn, nhưng thực sự nó không muốn làm phiền tới ai.''

Mio gật đầu, cô đã bắt đầu hiểu rõ hơn về cậu Tổng của họ. Cô nói nhỏ về phía Flare và A-chan: ``\hl{Giờ em mới để ý, Kuon-senpai luôn làm mọi thứ một mình. Dường như anh ấy không muốn làm phiền ai.}''

``\hl{Kể cả đôi lúc em cũng muốn lên tiếng giúp, nhưng anh ấy lúc nào cũng nói `anh ổn'\dots}'' nàng Dị giới tiếp lời.

``\hl{Cậu ta đúng là luôn tự đùn việc về cho mình thật\dots}'' A-chan thở dài.

Bầu không khí bỗng chốc trở nên lặng lẽ, như thể một điều gì đó sâu sắc đã được hé lộ. Mọi người im lặng lắng nghe những lời giải thích, như thể họ đang dần hiểu ra một phần tính cách của Kuon mà trước giờ chưa từng nhận ra.

Dường như tất cả mọi người cũng đã nghe lỏm được đoạn hội thoại của mọi người. Một số cô gái tỏ vẻ thương cảm tới cậu Tổng của họ, một số người khác thì tỏ vẻ lo lắng với những gì họ nhận được, một số khác thì ngỡ ngàng với suy nghĩ có phần khác biệt của cậu.

``\textit{Chẳng trách mà anh ấy toàn né tránh sự kiện của chúng ta\dots}'' Kanata thầm nghĩ.

``\textit{Không thể ngờ Kuon-senpai rụt rè đến thế đấy ne\dots}'' Miko đang gói lại quà, vẻ mặt đầy sự lo âu.

``\textit{Liệu anh ấy có cảm thấy vui với sự chuẩn bị này không đây\dots}'' Suisei lẩm bẩm.

``\textit{Có vẻ anh Tổng nhà ta không phải là một người thích tiệc tùng nhỉ,}'' Botan vẫn trơ vẻ mặt lạnh lùng, nhưng trong thâm tâm cô là một sự vừa khó hiểu, vừa cảm thông cho cậu.

Và rồi\dots

Dưới ánh sáng chiều muộn vàng óng, sau khi nghe lời giải thích từ gia đình Kuon, ba người phụ trách nhóm trang trí trao nhau một cái gật đầu đầy kiên quyết.

Nàng Đeo kính siết chặt tay, dứt khoát tuyên bố với giọng lắng lại: ``Vậy thì\dots\ chúng ta càng phải làm cho ngày hôm nay trở nên thật đáng nhớ hơn nữa!''

Nàng Sói đen chống hông, nở một nụ cười tươi rói đồng tình: ``Phải rồi! Nếu Kuon-senpai trước giờ luôn âm thầm như vậy, thì lần này chúng ta sẽ làm cho anh ấy phải thật bất ngờ tới mức không thể nào quên được!''

``Làm cho cậu ấy phải đỏ mặt vì hạnh phúc luôn mới chịu nha!'' Flare lém lỉnh tiếp lời.

Không tốn thêm thời gian, cả ba người quay về phía các nhóm đang hoàn thành những công đoạn quyết định cho các bức tường. A-chan đưa tay vỗ nhẹ hai nhịp như hiệu lệnh, gọi sự chú ý của mọi người trong nhóm. ``Nào mọi người! Tôi có điều muốn nói đây!''

Khi tất cả đã quay lại nhìn, cô cất giọng: ``Mọi người, tôi biết là chúng ta không còn có nhiều thời gian nữa, nhưng tôi muốn tất cả hãy dồn hết tâm huyết vào khoảnh khắc này!''

Lời kêu gọi của A-chan vừa dứt, một vài ánh mắt trong nhóm thoáng chần chừ. Matsuri nhún vai, nhăn mặt: ``Nhưng mà, A-chan\dots\ không phải Kuon-senpai ghét bị làm phiền à? Lỡ anh ấy thấy mấy thứ này rồi nổi cáu thì sao?''

Marine khoanh tay, gật đầu đồng tình. ``Phải đó\dots\ em còn nhớ lần trước anh ấy từ chối luôn cái event mừng 1 triệu người đăng ký kênh \hll\ chúng ta vì sợ  `tốn thời gian của mọi người' cơ mà. Anh ấy không giống kiểu người muốn được làm trung tâm đâu.''

Sora lúc này bước tới, nét mặt hiếm khi nghiêm túc đến vậy. Cô nhìn từng người một, rồi nói chậm rãi: ``Nhưng chẳng phải vì vậy mà chúng ta mới nên làm điều này hay sao?''

Korone nghiêng đầu, tai cụp xuống đầy vẻ khó hiểu: ``Ý của Sora-chan là\dots?''

Mio nhẹ nhàng tiếp lời: ``Là vì anh ấy luôn sống cho người khác. Lúc nào cũng để ý cảm xúc của chúng ta, luôn chọn cách yên lặng để tránh gây rắc rối\dots\ Dù là sinh nhật, Kuon-senpai cũng tự gạt mình ra khỏi niềm vui đó. Mọi người không thấy thế sao?''

Tất cả mọi người đều cùng nhìn nhau với lời khẳng định này, rồi ai nấy đều gật gù đồng tình. Một số giọng nói bất bình cũng thốt lên.

``Phải đó, phải đó!'' đầu tiên là Kanata, ``Anh ấy đã sủi năm bọn mình tổ chức holoFes đấy thôi!''

``Không phải chỉ một, mà đến tận hai lần nữa chứ!'' tới lượt Fubuki cũng lên tiếng. ``Có vẻ như anh ấy chủ động không muốn nhập hội với bọn mình hay sao ấy!''

``Dẫu biết rằng anh ấy không thể tham gia vì có việc riêng,'' Haato cũng góp lời, tỏ vẻ không đồng tình và đôi chút nghi ngờ dù hiểu được tình cảnh của cậu, ``lần đầu bọn mình còn hiểu chứ đây là lần thứ n rồi, anh ấy không tham gia ngày vui của bọn mình nữa chứ! Chị là chị thấy nghi rồi đó!''

Rồi, tất cả gần hai mươi người thành viên cũng chia sẻ quan điểm của mình về việc họ cho rằng Kuon ``trốn'' khỏi bữa tiệc của công ty. Không chỉ những lần đại hội công ty, họ cũng ``kể tội'' những lần mà cậu chủ ý không tham dự những sự kiện cá nhân của họ.

``Vậy thì lần này, mình sẽ kéo anh ấy lại,'' Flare nói, ánh mắt ánh lên một tia quả quyết. ``Hãy cùng đưa Kuon-senpai trở về đúng nơi thuộc về: giữa vòng tay của tụi mình!''

Nene ngập ngừng, siết chặt sợi dây trang trí trong tay. ``Nếu là vì vậy\dots\ thì bọn mình làm tiếp thôi.''

Marine chép miệng, nhưng môi đã mỉm cười. ``Anh ấy mà nổi đóa lên thì em không chịu trách nhiệm đâu nhé!''

Matsuri bật cười: ``Thế càng tốt! Tui đợi cái mặt ngượng chín của Kuon-senpai từ lâu rồi!''

A-chan nhìn tất cả, lòng nhẹ đi một chút. ``Vậy chúng ta tiếp tục nhé. Không chỉ là trang trí. Chúng ta đang chuẩn bị một lời cảm ơn. Một món quà. Một điều mà Kuon-senpai chưa bao giờ dám mơ đến.!"'

Nàng Dị giới hùng hồn: ``\textbf{Hãy tạo nên một bữa tiệc sinh nhật thật sự đáng nhớ cho Kuon-senpai nhé!}''

Mio lập tức tiếp lời: ``Chúng ta sẽ biến nơi này thành thiên đường sinh nhật cho Kuon-senpai của chúng ta!''

Tiếng reo hò đồng loạt vang lên từ khắp căn phòng. Ai nấy đều trở lại công việc, tinh thần như được tiếp thêm lửa. Mỗi nhánh dây kim tuyến, mỗi bóng bay, mỗi chi tiết nhỏ đều được chăm chút kỹ lưỡng hơn hẳn.

Trong những giây phút quyết định đó, ai nấy đều hối hả chăm chút cho từng chi tiết của mình, bức tường trang trí cũng đã gần như hoàn tất, các món ăn phục vụ bữa tiệc sinh nhật cũng đã sẵn sàng, trước khi cánh cửa của buổi tiệc được mở ra chào đón người mà họ đã dốc lòng chuẩn bị.

\begin{center}
    \uline{Kết thúc hồi tưởng.}
\end{center}

Khung cảnh trở lại thực tại, khi căn phòng giờ đây tràn ngập tiếng cười, những cái ôm siết chặt, và ánh mắt lấp lánh hạnh phúc. Kuon vẫn đứng đó, như hóa đá giữa vòng vây tình cảm, vừa ngượng ngùng vừa xúc động đến mức không thốt nên lời. Tay cậu run nhẹ, ánh mắt thì không biết nên nhìn ai cho phải.

Từng cái ôm, từng tiếng chúc mừng, từng khuôn mặt quen thuộc\dots\ tất cả hòa thành một cơn sóng lớn ập vào trái tim cậu, thứ cảm giác mà lâu lắm rồi cậu mới dám đối diện.

Bất chợt, chuông điện thoại vang lên, làm gián đoạn bầu không khí đang ấm cúng mà cũng đầy xúc cảm ấy. ``Ồ, có người gọi.''

Cậu Tổng lấy ra chiếc điện thoại trong túi quần mình, trên đó có ghi ``\emph{YAGOO}'' đang gọi.

``\dots Ông sếp của chúng ta gọi này.''

Cô nàng Đeo kính nghe vậy, liền cười: ``Có vẻ như ông ấy cũng có điều gì đó muốn nói với cậu đấy.''

Cậu nhấn nút nghe, tay vẫn chưa hết run vì cảm xúc hiện tại: ``A-lô ạ?''

Ngay từ đầu dây bên kia, một giọng nói của người đàn ông quen thuộc đầy vui vẻ vang lên: ``A, Hikawa-san! Chúc mừng sinh nhật cậu!''

Kuon ngỡ ngàng, tới mức suýt đánh rơi cả điện thoại: ``C-Chú Tanigo-san!? Sao chú lại gọi cháu giờ này\dots!?''

Vị chủ tịch bật cười, tiếp tục: ``Thì nghe nói hôm nay là ngày đặc biệt mà. Không gọi mới là lạ ấy chứ. Sức khỏe vẫn ổn chứ? Không bị sốc vì bất ngờ này chứ hả?''

``Dạ, cháu\dots\ vẫn đang hoang mang đây ạ. Mấy đứa\dots\ cháu vẫn không hiểu vì sao họ tụ tập hết ở nhà cháu nữa\dots''

``Ồ? Vậy là mấy cô gái nhà ta xem chừng cũng quý mến cậu lắm đấy. Được tổ chức tiệc sinh nhật tận nơi, thế này là ``lên hương'' rồi còn gì? Dám chắc là họ chuẩn bị hẳn từ A đến Z luôn nhỉ?''

Kuon buông điện thoại xuống, ngước nhìn quanh căn phòng. Dưới chân cậu, ngoài hàng chục nồi lẩu và hàng tá đia đồ ăn, cậu còn nhìn thấy hàng đống quà ở một bên góc phòng và một thứ khổng lồ được che lấp bởi tấm rèm.

Cậu áp điện thoại lên, ngập ngừng đáp: ``Cháu\dots\ cũng không ngờ tới ạ\dots''

Ông Tanigo bật cười, tiếp tục nói: ``Nghe nói các cô gái nhà ta đã tự lên kế hoạch, phân công nhiệm vụ, chạy đôn chạy đáo cả tuần để làm điều này cho cậu đấy. Tôi mà ở đó chắc tôi sẽ khen nức nở luôn à!''

Do cậu bật chế độ loa ngoài, tất cả mọi người đều nghe thấy những gì ông ấy nói. Có người đỏ mặt, có người nở một nụ cười gian mãnh, người thì cười khúc khích với lời khen này.

Suisei tự hào tiếp lời: ``Đúng đó, Kuon-senpai à! Bọn em lên kế hoạch cực kỳ chỉn chu nghiêm túc luôn nhé!''

Sora thì khiêm tốn hơn một chút: ``A-chan là người chịu trách nhiệm điều phối chính, em chỉ phụ một phần nhỏ thôi.''

Luna cũng loi nhoi chêm vào: ``Còn có cả bố mẹ anh nữ á-nanora\~{} Từ việc tra hỏi thông tin-'' để rồi bị Kanata bịt miệng lại, gắt lên: ``Đừng có nói mấy cái thừa thãi như thế!''

Cậu quay phắt lại, thấy bố mẹ và em gái đang vẫy tay về phía mình. ``\vie{Vậy là hai người cũng biết chuyện này từ đầu rồi ạ!?}''

Người bố của cậu bật cười: ``\vie{Con thì bận bịu như thế, làm sao mà bố nỡ phá bất ngờ chứ.}''

Người mẹ cũng từ đó, tiếp nối: ``\vie{Với cả nếu nói sớm thì còn gì là thú vị nữa! Mẹ phải công nhận là mấy đứa này khéo léo thật đấy.}''

``Chả bù cho Onii-san chúng ta,'' cô em gái của cậu cất giọng móc mỉa, ``suốt ngày lười biếng nằm trong phòng, không biết tán gái, lại còn\dots''

Tới đây, cậu không thể cãi nổi cô bé, vì Ryuuhou nói\dots\ không hề sai. Tất cả mọi người trong phòng cười ồ lên, còn cậu thì ngượng đến chín mặt, tới mức chỉ muốn trốn lủi khỏi nơi này.

Cậu quay lại với điện thoại, thở dài một hơi, rồi đáp: ``Cháu\dots\ thật sự thấy vui ạ.''

``Tốt. Thỉnh thoảng được làm người được yêu thương cũng đâu có gì sai, đúng không?''

Ánh mắt Kuon lướt qua những nồi lẩu đang sôi nghi ngút, những thứ trong căn phòng được sửa gọn gàng, những món ăn bày biện tươm tất khắp phòng\dots\ và rồi, cậu lẩm bẩm như đang tự chửi mình: ``Dạ\dots\ đúng là còn hơn cả mong đợi\dots\ Còn hơn cả một thằng lười như cháu luôn.''

Một lần nữa, tiếng cười phá lên khắp phòng, ai nấy đều cười sặc sụa khi nghe cậu tự tấu hài bản thân không báo trước.

``Cậu đúng là thằng thẳng thắn. Nhưng mà\dots\ hãy nhìn quanh cậu đi, và tự hiểu lý do vì sao mọi người lại làm thế.''

Lúc này, những cô gái xung quanh mới đồng loạt lên tiếng.

Fubuki mỉm cười, hớn hở nói: ``Bởi vì anh là Kuon-senpai đấy! Không quý anh thì bọn này quý ai!''

Marine thì nắm tay cậu thật chặt: ``Cũng không phải ngày nào bọn em cũng thấy anh đỏ mặt nha\~{} Nên phải tranh thủ chứ!''

Polka thì đấm yêu vào vai cậu, nhưng lực của cô dường như hơi mạnh khiến cậu suýt thì ngã ra trước và giẫm vào một trong những nồi lẩu đang sôi lên nghi ngút: ``Thấy chưa! Đây là tình yêu đấy, đừng có trốn nữa!''

``Em làm anh suýt nữa thì chết phỏng rồi đó\dots'' cậu Tổng lầm bầm.

``À, xin lỗi.''

Nhưng ngay khi cậu còn đang lúng túng trước vòng tay vây kín và bầu không khí ấm áp ấy, A-chan chỉ tay về phía một góc tường, nơi những tấm màn đen to bản vẫn còn đang che phủ điều gì đó ``Nhưng mà này, Hikawa-san. Bọn tôi còn một bất ngờ nữa dành cho cậu.''

Mãi đến khi cô nàng ``Tăng ca'' nhắc, Kuon mới để ý thấy xung quanh bốn phía căn phòng là những tấm rèm đen được giăng lên rất khéo léo. Chúng gần như hòa lẫn vào màu sắc trang trí và ánh đèn, khiến cậu chẳng hề để tâm từ trước đó.

Cậu nghiêng đầu chỉ về chúng: ``Những cái này là\dots?''

``Giờ mới hỏi thì muộn rồi, ne!'' Miko tinh nghịch đáp lại, rồi hướng về phía mọi người: ``\textbf{Sẵn sàng chưa nào?}''

Như kiểu được tập dượt từ trước, toàn bộ ba mươi thành viên túm vào một phần của tấm rèm, rồi cùng gật đầu như đã chuẩn bị vào vị trí.

Và rồi, tất cả mọi người cùng nhau đồng thanh: ``\textbf{Một, hai, ba\dots\ A-lê-hấp!}''

*ROẠT!*

Toàn bộ các tấm màn được kéo xuống trong tích tắc, để lộ ra bốn bức tường tràn ngập sắc màu.

Những bức tranh trang trí trên chúng dần hiện ra, với những chi tiết rõ nét và rực rỡ. Đó là một bức họa khổng lồ được ghép từ hơn ba mươi bức tranh nhỏ, mỗi bức phản ánh một mảng ký ức, một lát cắt cảm xúc, một phần con người Hikawa Kuon, qua ánh mắt của từng thành viên \hll. Từ ruy-băng đến hạt cườm, từ sơn màu đến hình dán, tất cả đều được kết hợp kỳ công như một triển lãm nghệ thuật sống động.

Trong đó: Sora là những nốt nhạc, Roboco là mạch điện, Miko là đền chùa cổ kính, Suisei là gạch Tetris rực rỡ.

Tiếp đến là Mel là đôi cánh dơi, Fubuki là cơn tuyết trắng, Matsuri là lễ hội pháo hoa, Aki là cảnh vật từ một thế giới dị giới với những cánh đồng ánh sáng kỳ ảo.

Rồi đến là Haato là những trái tim đỏ thắm, Aqua là đại dương, Shion là phép thuật tím, Ayame là quỷ vương xinh đẹp, Choco là ống nghe và lọ thuốc, Subaru là dải ngân hà thể thao điện tử.

Mio là mặt trăng đen hình sói, Korone là đôi tay đang vẫy, Okayu là những viên cơm nắm dễ thương.

Pekora  thì tất nhiên là cà rốt, Flare là khu rừng cổ, Noel là chiến giáp thời Trung cổ, Marine là con tàu vượt sóng, Kanata là trời cao, Watame là đồng cỏ, Towa là tinh vân vũ trụ, Luna là lâu đài kẹo ngọt\dots

Và cuối cùng, Lamy là bức tường băng trong suốt, Nene là bức tranh ngộ nghĩnh như trẻ thơ, Botan là súng ống pixel, còn Polka là một rạp xiếc lung linh đầy màu sắc.

Mỗi mảng bức tranh, họ đều vẽ dáng dấp hình ảnh Kuon trong góc nhìn của họ: có người thì vẽ toàn bộ người, người thì vẽ theo cách hoán dụ, người thì trang trí lấp ló.

Nhưng, đặc biệt nhất là ở bức tường trung tâm. Đây là một bức tranh lớn vẽ cảnh Kuon đang ôm lấy toàn bộ các thành viên \hll, cùng với bố mẹ và cô em gái cậu đứng hai bên, như một đại gia đình sum vầy.

``\dots'' Im lặng.

Cậu đứng đó, há hốc miệng, không nói thành lời.

Mắt cậu liếc về từng khoảng tranh trên bức tường. Cậu nhìn thấy chính mình trong đó.

Một Hikawa Kuon tốt bụng, dễ mến với mọi người\dots

Rồi đến một Hikawa Kuon nghiêm khắc, có phần cầu toàn với công việc\dots

Rồi còn một Hikawa Kuon đôi lúc mộng mơ, đầu óc bay bổng tới một suy nghĩ sâu xa nào khác\dots

Từng phần trong bức họa đó khắc họa nên một phần làm nên con người của cậu, đồng thời cũng nói lên tính cách của chính các cô gái và cách họ nhìn nhận về phía cậu.

Những bức họa này có phần được chỉn chu rất cẩn thận, có phần thì hơi xộc xệch vì đôi phần vội vã\dots\ nhưng tất cả đã tạo nên một bản kiệt tác hùng vĩ mà ``nhân loại có thể chờ mong cho thế giới nghệ thuật'' (theo lời của Kuon).

Dường như ở trong đời, cậu chưa bao giờ được chiêm ngưỡng một tác phẩm nghệ thuật nào lại đẹp đẽ, khắc sâu vào trái tim cậu như vậy.

``Đ-Đây là\dots''

Cậu thì thầm, miệng vẫn còn há hốc.

Cậu không ngờ rằng, cả đời cậu sẽ xứng đáng có được một khoảnh khắc như thế.

Những sự hy sinh cậu phải đánh đổi chỉ để cho các cô gái có được một tương lai xán lạn hơn\dots

Những sự khó khăn cậu gặp phải khi phải đương đầu với những trò nhí nhố của họ, để rồi chính cậu đã tự tay giải quyết\dots

\dots và rồi giờ đây, chính cậu là người được hưởng tất cả, ngay tại giây phút này.

``\dots Mọi người\dots\ đã vẽ cái này sao?''

Tất cả ba mươi người đều cùng ``Ưm!!'' một tiếng đầy hùng hồn.

Ryujin cũng tự hào nói trước con trai: ``Bố mẹ cũng phụ một tay đó, nhất là mấy đoạn dán hạt cườm, mẹ con làm hết.''

Kaien húc nhẹ vào vai của người chồng, bổ sung: ``Bố con chỉ giúp dán keo thôi, nhưng cũng hăng lắm.''

Cậu ngẩn người, nhìn bức họa trên tường một lần nữa, há hốc miệng, nhưng vẫn lắp bắp hỏi: ``M-Mà\dots\ dày công như vậy, chắc là mọi người mất nguyên cả tuần luôn đấy ạ\dots!?''

Shion lắc đầu: ``Không. Mất đâu đó có ba tiếng rưỡi thôi.''

``\textbf{N-Ngắn đến thế sao!?}''

Những tiếng cười lại rộ lên trong văn phòng. Đúng là trong mắt cậu, các thành viên \hll\ có thể có những ý tưởng điên rồ, nhưng đến mức hợp sức để làm nên một bức tranh trong thời gian ngắn như thế là điều cậu không ngờ tới.

``Cái đó gọi là \emph{nghệ thuật vẽ nhanh} (speed art) đó\~{}'' Okayu ngân lên.

``Nhưng quan trọng nhất là Ogayu với bọn em đã làm ra nó bằng cả trái tim đó, gâu gâu!'' Korone tiếp lời.

``Và tấm lòng của chúng em thì lớn lao lắm luôn!'' Haato và Matsuri đồng thanh, dang rộng hai tay về phía bức họa.

Aki chỉ về phía khoảng tranh của mình, khoe: ``Hình vẽ của em còn có hiệu ứng phát sáng á!''

Luna cũng không kém cạnh: ``Luna thì zẽ hình kẹo siu to bự lun á-nora\~{}!''

Lần lượt, các thành viên khác cũng bắt đầu thuyết trình bức vẽ tường của mình với những giọng nói hãnh diện. Từ những lời nói đơn giản, cho tới nguyên một bài giới thiệu dài như sớ Táo quân.

Lúc đó, tiếng điện thoại trong túi cậu vẫn chưa cúp, một giọng nói từ đầu dây kia vang lên, dường như vẫn theo dõi toàn bộ diễn biến qua điện thoại: ``Thế nào? Cậu cảm nhận được tình cảm của họ dành cho cậu chứ?''

``Cháu\dots\ không biết phải nói gì hơn ạ\dots'' cậu Tổng vẫn đang hoang mang.

``Sốc đến nghẹn họng à? Haha. Tôi bảo mà, họ quý cậu nhiều lắm.''

``\dots''

``Dù sao thì, đừng quên nói lời cảm ơn. Món quà này không đơn thuần là tranh đâu\dots\ mà là tình cảm của ba mươi mấy cô gái dành cho một người duy nhất. Không ai khác là chính cậu.''

``\dots Cháu cảm ơn chú\dots'' cậu thều thào, nhưng rồi vị giám đốc cắt ngang: ``Không, người cậu nên cảm ơn là họ cơ.''

Cậu ngẩng lên, thấy ba mươi tư người - toàn bộ các tài năng của nhà Tam Giác Xanh và ba thành viên trong gia đình cậu cùng A-chan - đang nhìn cậu với ánh mắt dịu dàng và lấp lánh, như chờ một điều gì đó.

Cậu cúi gập người thật sâu, hét to: ``\textbf{Tôi xin cảm ơn tất cả mọi người về món quà này!!!}''

Ryujin và Kaien cùng Ryuuhou chạy tới đỡ cậu đứng dậy. Cậu thở một tiếng, đôi mắt rưng rưng, rồi nói với mọi người bằng sự cảm phục: ``Cảm ơn mọi người vì đã cho anh một món quà ý nghĩa đến thế\dots\ Trần đời này anh chưa từng được tặng món quà nào lớn thế này\dots''

Tất cả đều mỉm cười, cùng nhau nụ cười chân thành nhất, ấm áp nhất. Một nụ cười giống như chính cách mà Kuon đã luôn đối xử với họ.

Cậu áp máy lên tai, nghe lời chúc cuối cùng vang lên từ cấp trên của cậu: ``Chúc mừng sinh nhật cậu nha, Hikawa-san!''

``\dots Cảm ơn chú\dots\ nhiều lắm ạ\dots''

Cậu cúp máy.

Và chính lúc đó, như một phản xạ tự nhiên, những giọt nước mắt lặng lẽ rơi từ lúc nào không hay.

Toàn thể các tài năng cùng A-chan nhanh chóng ôm chầm lấy cậu, có người xoa đầu, có người vỗ về dỗ dành. Trông chẳng khác gì một đứa trẻ đang được hàng tá người lớn thi nhau dỗ dành.

Những người còn lại đứng xung quanh cậu đồng loạt cất giọng và bắt nhịp một lần nữa: ``Nào, lại hát lần nữa nhé! Hai, ba\dots''

``Khúc hát mừng sinh nhật'' một lần nữa lại vang lên, và toàn thể tất cả mọi người trong phòng, kể cả nhân vật chính của bữa tiệc, cùng nhau hòa giọng:

\tabs[4cm] ``Ngày hôm nay ta cùng hợp hoan nơi đây, \\
\tabs[4cm] Mọi người bên nhau ta hát mừng sinh nhật, \\
\tabs[4cm] Một, hai, ba, ta cùng thổi tắt nến, \\
\tabs[4cm] Happy birthday, happy birthday to you\dots''

Tiếng pháo tay rộ lên như sấm, và đó chính là phần mở đầu cho một buổi tiệc sinh nhật mà cậu sẽ không bao giờ quên.

\ct

Sau khi mọi người đã ổn định chỗ ngồi quanh chiếc mâm lớn đặt giữa phòng khách, không khí bắt đầu lắng xuống một chútđể chuẩn bị cho phần tiếp theo. Mâm cỗ này chiếm gần nửa căn phòng, vừa đủ cho tất cả ba mươi lăm người ngồi, gồm hơn chục nồi lẩu bốc khói nghi ngút, dãy đĩa thịt bò, thịt lợn, thịt gà được xếp theo từng loại, có cả loại được tẩm ướp, loại thì để riêng cho người ăn nhạt. Rau nhúng thì phong phú khỏi bàn, từ cải thảo, cải xanh, nấm kim châm cho tới cà rốt thái hoa. Bên cạnh đó là hàng đĩa rau củ để nướng, và một phần nhỏ là cơm nắm ba màu để ăn kèm hoặc chống đói trong lúc chờ đồ chín.

Mọi thứ đều được sắp xếp rất nghệ thuật, đẹp như thể vừa bước ra từ một tạp chí ẩm thực.

Trong lúc mọi người đang cười nói và tranh thủ chụp ảnh cho ``bữa tiệc lịch sử'' do chín họ làm nên, thì Sora đứng dậy, ánh đèn vàng từ trần chiếu lên mái tóc nâu của cô ấy khiến cả gian phòng như dịu lại.

Rồi, nàng Thân tượng hắng giọng, cầm chiếc míc-rô lên (mượn từ cây loa kẹo kéo của Aloe mà một lần nữa, Ryuuhou nói dối rằng đó là của cô bé) và bắt đầu buổi dẫn: ``Mọi người cũng biết hôm nay là ngày gì rồi, đúng không nào?''

Tất cả mọi người nói chen lẫn nhau, người thì tỏ vẻ hứng thú, người thì nói với giọng bình bình, người thì trả lời qua loa vì đã biết rõ từ trước: ``Sinh nhật của Kuon-senpai!'' / ``Sinh nhật của anh Tổng của chúng ta chứ gì\dots''

Nhưng cô nàng vẫn không để tâm lắm tới những tiếng nói đã hạ nhiệt, tiếp tục: ``Đúng rồi! Vậy có nghĩa là chúng ta phải làm gì nào?''

Một nhịp lặng vừa đủ để tạo hiệu ứng. Ngay sau đó, phần lớn mọi người đã mở miệng và đang định đồng thanh ``\textbf{Tặng quà!}''\dots

\dots thì bất ngờ bị chặn ngang bởi ba tiếng hét từ một góc phòng, từ Marie, Rushia và Watame: ``\textbf{Rót rượu tháp ly!!!}''

Không khí im bặt trong một giây. Tất cả các cô gái ngoảnh về phía ba người đó, rồi cùng đồng loạt hoảng hốt: ``\textbf{C-Cái gì cơ!?}''

Trước mặt tất cả mọi người là một tháp ly lấp lánh cao bốn tầng, được đặt trên một chiếc bàn di động phủ khăn trắng, và kéo ra từ sau ghế bành tự bao giờ không ai hay.

Mọi người không ai trả lời mà thay vào đó lại di chuyển rất\dots\ có tổ chức để triển khai kế hoạch bí mật tiếp theo, mặc dù bản thân họ còn không hiểu bộ ba người kia đã chuẩn bị từ khi nào.

Nàng Thuyền trưởng, với sự hỗ trợ của nàng Cừu và nàng Pháp sư, đang được cõng lên cao, tay cầm một chai sâm-banh (champange) lớn, dáng vẻ không khác gì hải tặc thật sự sắp khai trương chiến hạm mới cứng.

Nàng Sói đen đứng bật dậy, thảng thốt: ``Mấy đứa đang làm cái trò gì vậy!?''

Cậu Tổng thì bình tĩnh hơn, nhưng vẫn giấu nổi sự bất ngờ: ``Marine-chan\dots\ Em thật sự muốn làm như vậy từ lâu rồi sao? Mấy đứa đang nhầm đấy\dots''

``Anh không thấy sao!? Bọn em đang thực hiện nghi lễ rượu sâm-banh thôi!''

``Cái đó dùng cho lễ hạ thủy lần đầu mà!!'' cả nàng Sói đen và cậu Tổng đồng thanh.

Dẫu vậy, Senchou vẫn thản nhiên thả chai rượu về phía tháp ly như một nghi thức truyền thống. Chai rượu được buộc vào trần nhà, phóng vèo về phía bức tường ngay bên trên chồng ly, và\dots

*COONG!*

Không vỡ. Nhưng thay vào đó, một mảnh giấy dán tranh tường của Mio bị bong ra, làm cho đôi tai sói của chủ nhân bức tranh giật mình. Cô hét lớn: ``\textbf{Ma-rine!!!}''

Kuon thở phào nhẹ nhõm: ``Trời ạ, mấy đứa đúng là\dots\ Giờ Mio-chan quạu rồi kìa.''

Lúc này, Miko, hoàn toàn không để ý tình hình, vô tình nhặt lấy chai sâm-banh từ tay Marine trong khi vẫn đang ngẫm nghĩ điều gì đó. Với sức mạnh ``thần thông'' đặc trưng của mình, cô\dots\ mở chai bằng tay không, như cách mà cô đánh bại đám người ngoài hành tinh trong khi vẫn còn đang suy nghĩ\footnote{Xem lại {\flaregothic ÉPISODE 105}, quyển số Tám.}.

*PẶC!*

Tiếng nút chai bật vang lên trong căn phòng đầy người. Bản thân cô vẫn không hiểu chuyện gì đang xảy ra.

``\textit{Mà mình đang định làm gì đấy nhỉ\dots?}'' cô thầm nghĩ, đồng thời vô thức đưa chai rượu cho Luna.

Nàng Công chúa xứ Kẹo lập tức hai tay nâng lấy, nhe răng cười: ``Zâng, mời anh đó, Kuonsa-senpai\~{}''

Cậu quản lý ngập ngừng cầm lấy, nghiêng nghiêng đầu rồi gãi đầu gãi tai: ``Sao anh cảm thấy có gì đó không đúng nhỉ\dots''

``Không đúng là cái chắc!!!'' nàng Sói đen trừng mắt, tay vẫn đang dán lại giấy.

Trái lại, nàng Cừu cười tươi, đáp lại: ``Không sai đâu, Kuon-senpai. Giờ thì anh rót nó lên đỉnh tháp thôi!''

Không còn lựa chọn nào khác, cậu cười trừ rồi làm theo. Vừa rót rượu trên đỉnh tháp, cậu vừa cằn nhằn: ``Nhưng mà anh tưởng là tiệc cưới người ta mới làm thế này chứ? Đây là sinh nhật mà?''

Mio thở dài, trát keo lên tờ giấy trang trí rồi giải thích: ``Em đã giải thích cho họ rồi đấy, nhưng mà bọn họ lại không nghĩ như thế\dots\ Kể cả có nhờ A-chan cũng không xoay chuyển được họ cơ\dots''

``Thậm chí Mel-senpai với Matsuri-senpai còn ủng hộ ý tưởng này nữa đó-nanodesu!'' Rushia đáp lại.

Nàng Sói đen và cậu Tổng liếc về hai cô nàng vừa gọi tên, những người đó cười toe, cùng lên tiếng:

\begin{dialogue}
    \speak{Mel} Vì bọn em nghĩ rằng đây là cách chúc hạnh phúc cho anh đó\~{}
    \speak{Matsuri} Đúng đó! Cho nên dù là đám cưới hay sinh nhật thì cũng là chúc phúc cả!
\end{dialogue}

Nghe những lời biện minh của họ, cậu chỉ có thể thở dài: ``Mấy đứa lại râu ông nọ cắm cằm bà kia rồi\dots''

Cậu rót toàn bộ số rượu vào đỉnh tháp. Toàn bộ năm tầng ly rượu đều được đổ đủ, mỗi ly lấp được nửa ly.

Nhưng rồi, ngay khi cậu đặt cái chai xuống, đôi mắt cậu chợt dừng lại ở phần chân tháp.

Một khối đá khô đang bốc khói.

``\dots Vậy đống đá khô ở chân tháp là thế nào đây?''

Nàng Vu nữ bất giác tỉnh lại, phấn khích như vừa nhớ ra điều gì quan trọng: ``A! Anh nhắc em mới nhớ! Cái đó để tạo hiệu ứng chứ sao nữa, ne!''

Cậu nhìn về phía A-chan, người đang toát mồ hôi nhẹ với lời tuyên bố của cô nàng. Cô chỉ nhún vai thở dài: ``Chịu. Tự dưng Sakura-san đòi như vậy. Bọn tôi đã cố ngăn rồi, nhưng không được.''

Ngay khi nàng Đeo kính dứt lời, nàng Cừu bất chợt reo lên: ``Rót thêm rượu vào thôi!!''

Và chỉ chờ có thế, không kịp để ai phản ứng gì, Watame mở một cái *PẶC!* cái nút chai và rót thẳng vào cái ly đỉnh đầu.

``K-Khoan! Một chai là đủ rồi đó, Watame-chan!'' Mel thốt lên định ngăn lại.

Nhưng nàng Cừu vẫn để ngoài tai.

Như một lẽ tất yếu, lớp rượu ở tầng dưới cùng tràn ra, rơi xuống khay đá khô dưới đó, tạo thành một làn khói mờ mịt.

Dưới lớp ly thuỷ tinh trong suốt, những giọt rượu đầu tiên nhỏ xuống khay đá khô bên dưới, lập tức phản ứng lại bằng một làn khói mờ mịt bốc lên. Mọi người xúm lại gần hơn, nín thở quan sát khung cảnh kỳ lạ ấy như đang chứng kiến một nghi thức thần bí.

Towa đứng ngay sát, đôi mắt xanh lá mạ long lanh thích thú. Tính cách quỷ nghịch ngợm của cô trỗi dậy, cô bật cười khoái chí, rồi không biết từ đâu rút ra một chai nước và reo lên: ``Cho thêm tí nước vô nữa cho nó xì khói đã hơn nữa!''

Chưa kịp ai kêu lên cản lại, nàng Ác ma đã hất mạnh nguyên chai nước nhỏ vào giữa khay đá khô. Hiệu ứng xảy ra tức thì: tiếng xèo xèo vang lên dữ dội, khói trắng phun trào lên cao rồi tràn ra xung quanh như một cơn lũ. Một phần nước không kịp bốc hơi kịp, tràn cả ra ngoài, lênh láng khắp sàn nhà gỗ bóng loáng.

Một tiếng ``\textbf{Ối trời ơi!}'' bật ra từ đám đông khi lớp khói trắng đặc quánh bắt đầu bao phủ toàn bộ căn phòng, nhấn chìm cả mâm tiệc trong làn sương mù như cảnh phim kinh dị. Chỉ còn thấp thoáng những chiếc bóng nhòe nhoẹt di chuyển trong màn khói. Ai đó vấp phải chân bàn, suýt nữa làm đạp vào nồi lẩu nghi ngút khói thật đang bốc hơi kế bên.

Towa phấn khích dang hai tay ra như thể đang đứng giữa sân khấu biểu diễn, không để ý rằng cả phòng tiệc đang biến thành một buồng xông hơi đầy khí cacbonic cấp tốc mà chẳng hề ai yêu cầu.

``Mọi người thấy khung cảnh thế nào? Có hùng vĩ không nào!?''

Đáp lại cô chỉ là những tiếng ho khù khụ từ mọi người, hiển nhiên không hề vui vẻ gì. Thoang thoảng trong màn sương, những tiếng la thất thanh vang lên, như ``Ai đó đuổi hết khói ra ngoài đi!!'' hay ``Khó thở quá!!'' hoặc thậm chí ``Chẳng nhìn thấy gì cả!!''

Một tiếng hốt hoảng vang lên: ``Chết thật! Mở tung cửa ra hết đi!''

Miko trố mắt ngạc nhiên với yêu cầu từ phía cậu Tổng: ``Ể!? Sao lại thế!?''

``\textbf{Em định giết chúng ta sao!?}'' một lần nữa, giọng của Kuon cất lên, trước khi chới với được tới nàng Vu nữ.

Tới đây, Miko mới sửng sốt thật sự. Cô chưa kịp mở miệng hỏi tiếp thì người bố của cậu Tổng hét lớn: ``\textbf{Cháu bị mất trí à!? Khói này là khói độc đấy!}''

Nàng Sói đen than vãn: ``Em biết ngay là đống đá khô đó là sai lầm mà! Mọi người, đuổi khói ra ngoài đi!''

Lập tức, tất cả mọi người vừa dò đường trong căn phòng đầy khói, vừa cẩn thận không dẫm phải bất kỳ nồi lẩu hay đĩa đồ ăn nào, vừa cố gắng mở cửa sổ thật nhanh.

Ba mươi lăm người hốt hoảng mở cửa, quạt bằng nắp nồi, bằng tay, thậm chí có người dùng cả đĩa không để làm quạt tạm thời. Khung cảnh giờ đây chẳng khác gì một bữa tiệc trên tiên khí, nhưng lại là một loại khí mà không một ai trong phòng muốn tận hưởng cả.

\ct

Sau vài phút náo loạn, khói mới dần tan, trả lại không gian thoáng đãng. Rất may, nhờ lớp bọc kỹ lưỡng mà đồ ăn không bị ám khói.

Nhìn khung cảnh mọi người thở dốc ngay khi tất cả đống khói độc đó bay ra ngoài, cậu cười chát, châm biếm nhẹ: ``Thật đúng là một nghi lễ mở đầu\dots\ ấn tượng.''

Khi làn khói cuối cùng bị thổi bay ra khỏi cửa sổ và cửa chính, cả căn phòng như vừa trải qua một trận lũ sương. Sàn nhà ướt đẫm, vài chiếc dép đi trong nhà trượt vẹt để lại dấu, còn những chiếc khăn được mọi người vội vã mang ra lau khô góc này góc nọ thì nằm rải rác khắp nơi. Một vài người đang giúp Ryujin gom đá khô còn sót lại vào xô, trong khi Kaien đi vòng quanh xem xét đồ ăn, thở phào vì lớp bọc ni lông đã cứu vớt tất cả.

Trong một góc phòng, Miko ngồi bệt xuống, hai tay ôm đầu gối, mặt nghệt ra như mèo bị tạt nước. Cô lí nhí: ``Em xin lỗi\dots\ Miko không biết là nó có chứa khói cacbonic trong đó\dots''

Towa đứng phía sau, hai tay khoanh lại, miệng cười khì như thể đang tự nhủ: ``\textit{Ừ thì\dots\ hơi quá tay thật.}'' Không ai mắng cô dẫu mọi người đều biết cô nàng là một ác ma, nhưng ánh mắt mọi người đổ dồn tới một cô gái đùa không đúng lúc cũng đủ để khiến một ác quỷ cũng phải thấy toát mồ hôi hột.

Matsuri và Mel cũng không khá hơn. Họ đứng cạnh nhau, hai người cùng gãi đầu gãi tai, miệng cười méo xệch. Nàng Lễ hội lên tiếng trước, lí nhí: ``Bọn em chỉ nghĩ là sẽ vui vui chút thôi mà\dots''

Nàng Ma cà rồng thêm vào, giọng như thỏ con bị bắt quả tang: ``Không ngờ là khói nó dày như vậy\dots''

Kuon bước tới, ánh mắt bất lực nhưng không giấu được sự buồn cười. Cậu khẽ thở dài, rồi buông một câu khiến cả phòng phải cười phá lên: ``Đúng kiểu `nhiệt tình + ngu dốt = phá hoại' luôn!''

Miko giật mình, ngẩng lên đầy lo sợ, nhưng chưa kịp phản ứng thì cậu đã tiến tới, nhẹ nhàng đặt tay lên đầu cô, xoa xoa như an ủi một đứa em nhỏ vừa làm vỡ bình hoa: ``Thôi, nhưng thế thì mới gọi là \hll\ chứ. Loạn xí ngầu thế này là anh cũng vui rồi. Hôm nay cũng là sinh nhật anh, nên\dots\ thôi, anh cho qua.''

Cô nàng ``Ưu tú'' tự xưng chớp mắt, rồi bĩu môi lí nhí: ``Em hông cố ý mà, ne\dots''

``Anh biết,'' cậu cười hiền, ``nhưng cảm ơn em, nhờ em mà anh có một kỷ niệm sinh nhật chắc chán không thể nào quên.''

Đám đông phía sau rộn ràng trở lại, không khí vui vẻ lại tràn ngập căn phòng. Những tiếng cười khúc khích, tiếng trò chuyện, tiếng gọi nhau vang rền khắp nơi như thể chưa từng có làn khói nào vừa che khuất cả phòng.

Rồi từ đâu đó, một giọng ai đó vang lên: ``Ê ê, đến tiết mục tiếp theo rồi đấy!''

``\textbf{Phải đó, phải đó!}'' một nhóm khác hưởng ứng.

Ai nấy lại rôm rả đứng dậy, rục rịch kéo nhau về phía góc phòng nơi đã chuẩn bị sẵn sàng cho phần tiếp theo, tiết mục tặng quà.

Suisei, Marine, Polka, Luna, Towa và Kanata là sáu người đã chuẩn bị sẵn những món quà đặc biệt cho Kuon, những món mà họ tin chắc rằng cậu sẽ thích.

Khi những chiếc hộp gói giấy đủ màu được bày ra trước mắt, Kuon lập tức ngờ ngợ. Mặt cậu thoáng ngạc nhiên, rồi hơi hoang mang: ``Đây là\dots\ quà của mấy đứa à?''

Nàng Ác ma cười khoái chí, giơ ngón tay cái như thể khoe công trạng cho cấp trên: ``Đúng rồi đó anh!''

Nàng Công chúa xứ Kẹo cũng líu lo chen vào, miệng chúm chím đầy tự hào: ``Bọn em đã mất nhìu thời dan lắm mới chọn được món đồ phù hợp với anh đó-nora\~{}!''

Kuon thở dài, tỏ vẻ lúng túng trước những hộp quà mà họ đã mang tới: ``Mấy đứa không nhất thiết phải-''

Nhưng ngay khi cậu còn chưa nói hết câu, nàng Thiên sứ lập tức bắt thóp, cất giọng như búa bổ đầu: ``Anh không cảm nhận được lòng thành của chúng em sao\dots?''

``Không phải thế, nhưng mà\dots'' cậu cố phân bua.

Chưa kịp nói hết, Suisei đã nhìn cậu bằng ánh mắt như sắp rưng rưng. Cô nàng phồng má, môi mím lại, giả vờ như sắp bật khóc, y như cách cô đã làm mỗi khi bị hội Thế hệ thứ Không ``cho leo cây'' trong những lần chơi nhóm.

``Kuon-senpai\dots\ không thích sao?''

Không khí trong căn phòng bắt đầu nghiêng về phía ``uy hiếp tình cảm'', nhưng chừng đó vẫn chưa đủ để khiến cậu hoàn toàn nao núng.

Flare và Mio liếc nhau rồi đồng loạt lên tiếng, như hai vị đại diện của ``phe nữ quyền nhằm bảo vệ cho quyền lợi của các thành viên''.

``Thôi nào, Kuon-senpai!'' nàng Dị giới bật cười, ``Ryuuhou-chan có nói với bọn em rằng anh không cần quà rồi\dots''

``\dots nhưng ít nhất anh cũng phải nhận quà mà mọi người tặng chứ!'' nàng Sói đen bồi thêm, nhướng mày như đang kiểm tra phản ứng từ cậu.

Kuon đưa tay ra hiệu từ chối một cách yếu ớt: ``Quà thì anh sẽ nhận thôi, nhưng vấn đề là-''

``Là gì?'' cả hội lập tức đồng thanh hỏi vặn, khiến cậu giật bắn mình.

``Là-'' Kuon ậm ừ, tìm từ, nhưng chưa kịp nói xong thì đám đông phía dưới đã đồng loạt lên tiếng.

``Là?''

``\textbf{\dots Là sao?}'' Mel, Mio và Haato gần như hét lên.

``\textbf{LÀ SAO!?}''

``\textbf{LÀ!?}''

Cậu Tổng rụt cổ lại, lùi về sau một bước, thở dài một hơi rồi lấy hết can đảm nói: ``\dots nhà không đủ chỗ.''

Cả căn phòng im bặt, cố tiêu hóa với lý do có phần hơi ngớ ngẩn này. Đặc biệt là từ sáu cô gái, những người đã chọn và tạng quà cho cậu.

Kuon nuốt nước bọt, giải thích như thể mình vừa thú tội: ``Nhà anh không đủ chỗ để chứa hết quà của mấy đứa tặng đâu\dots\ Tận ba mươi mấy đứa kia mà\dots''

Cả nhóm im lặng một giây, rồi đồng loạt phá lên cười. Suisei là người lên tiếng trước, vẻ mặt vừa hài hước vừa bất lực: ``Anh đúng là\dots\ Bọn em không tặng mấy quà to đâu!''

``\dots Hả?'' cậu ngơ ra.

``Chỉ là mấy món nhỏ nhỏ thôi mà,'' cô trấn an.

``Với cả bọn em cũng chỉ cử đại diện tặng quà thôi,'' Kanata bổ sung thêm. ``Chứ ai mà cũng tặn quà thì có người trùng nhau mất.''

Marine cũng chống hông, vỗ vai Kuon, giọng như trách yêu: ``Anh tưởng bọn em bê cả máy giặt với tủ lạnh tới để cúng cho anh à?''

Polka thì không nhịn được cười, búng tay: ``Có khi lần sau em tặng anh cái bồn tắm di động thật luôn đó!''

Kuon tròn mắt nhìn cả hai, rồi lùi lại một bước, mặt giãn ra, thở phào nhẹ nhõm: ``Thế mà mấy đứa dọa anh đến chết rồi\dots''

Cả nhóm phá lên cười một lần nữa, khi Kuon bắt đầu đưa tay tới những chiếc hộp gọn gàng trước mặt.

``Vậy để em bắt đầu nha-nora! Đây nè!'' Luna hớn hở kêu lên, hai tay giơ ra một chiếc hộp nhỏ hình vuông, to chừng nắm tay. Giọng cô vang lên trong trẻo như một đứa trẻ mẫu giáo lần đầu được đi phát quà sinh nhật.

Cậu cầm lấy chiếc hộp, thoáng ngơ ngác: ``Đây là gì? Trông nó giống như nhẫn đá kim cương gì đó\dots''

Cả Kaien và Ryujin cũng thoáng ngạc nhiên, cả hai người lớn thoáng nghĩ: ``\textit{Không lẽ con bé lại tặng đồ quý thật!?}''

Đáp lại những cái nhìn ngẩn tò te đó, nàng Công chúa xứ Kẹo vẫn giữ nụ cười ngây ngô, đáp lại: ``Anh cứ mử ra là bít lìn na!''

Kuon làm theo, nghĩ thầm: ``\textit{Chắc lại có món đồ đắt tiền gì cho mình đây mà\dots}''

Trong chiếc hộp đầu tiên là một chiếc hộp khác. Nhỏ hơn, màu xanh ngọc bích, có nơ cột hình tai thỏ.

Không chỉ cậu, mà cả Towa và Kanata - hai cô bạn thân thiết nhất của Luna trong cùng Thế hệ - cũng phải ngớ người, ánh mắt chứa đầy dấu hỏi: ``Con bé này lại bày trò gì nữa đây?''

Mọi ánh mắt đổ dồn về chiếc hộp nhỏ. Kuon chầm chậm mở ra, và ngay lập tức, tất cả bật thốt ngạc nhiên.

Bên trong là một bộ trang sức bạc nhỏ xinh, gồm một chiếc vòng cổ có mặt là viên đá pha lê trong suốt hình viên kẹo, và kèm theo đó là một chiếc vòng tay đồng bộ, cũng đính viên đá ngọt ngào như vậy. Cả hai món đều được chế tác tinh xảo, nhưng không phô trương. Chúng gợi cảm giác như một món quà thủ công do chính tay ai đó dành nhiều tâm tư chọn lựa.

Kuon sững người, còn Luna thì chỉ đứng đó, tay chắp sau lưng, ánh mắt long lanh như trẻ con tặng quà cho người bố của mình.

``C-Cái này\dots'' Kuon lắp bắp.

``Zòng cổ hình kẹo với zòng tay đó-nora!'' cô nàng xứ Kẹo cười toe toét. ``Em với Kanata-so với Towa-chan chọn á!''

``Không, ý anh là\dots\ tại sao lại là món quà này?'' cậu hỏi, giọng vẫn chưa hoàn hồn.

Một món quà có phần hơi đắt tiền so với những gì cậu nghĩ. Ừ thì, đúng là Himemori Luna là công chúa của một vùng đất, nên việc chi tiền cho một vật như thế cũng không khó nhằn gì.

Nhưng vấn đề là, tại sao? Tại sao cô nàng lại tặng một thứ như vậy tới cậu Tổng? Rõ ràng cậu được vị trí như thế nào mà được tặng như thế?

Ryujin cũng chen vào với vẻ ngạc nhiên thật sự: ``Đến chú cũng bất ngờ với lựa chon của cháu đấy.''

Kaien thì nhíu mày: ``Sao lại tặng món quà đặc biệt và trông đắt tiền thế cháu?''

Luna vẫn cười, lắc đầu nhè nhẹ như thể câu hỏi đó không hề đáng lo nghĩ: ``Zì anh í là quản lý của bọn cháu mà-nanora!''

``Nhưng nếu thế thì nghe vẫn cấn cấn thế nào ấy\dots'' Kuon phản ứng theo bản năng. Rõ ràng trong mắt cậu, cậu không xứng đáng được nhận một món quà như thế.

Luna nghiêng đầu, mỉm cười ngây thơ nhưng ánh mắt nghiêm túc một cách lạ thường: ``Vì anh là ngừu con chai di nhứt mà chúng em tin từng đó. Khum đúng xao ạ-nora?''

``\dots Hả!?'' Ryujin và Kuon đồng thanh, như thể ai đó vừa thả bom xuống giữa không trung.

Cậu không ngờ rằng một tài năng không tuổi, mới chỉ tập tễnh học nói, lại thốt ra những câu như vậy.

``N-Nhưng mà\dots\ anh tưởng mấy đứa đã từng tiếp xúc với các fan nam mà?'' cậu hỏi lại, dường như muốn phản biện.

Towa bước lên một bước, giọng bình tĩnh nhưng ấm áp: ``Em cũng đồng ý với cậu ấy. Bọn em cũng đã nói chuyện, thậm chí là bắt tay với nhiều người đàn ông ở ngoài rồi\dots''

``Thậm chí là làm vịc với nhiều người đàn ông khác nữ-nora,'' Luna bổ sung thêm. ``Luna bít chứ! Cơ mà\dots''

Kanata tiếp lời, ánh mắt khẽ cụp xuống rồi ngước lên nhìn cậu Tổng, người mà họ hết mực yêu quý: ``\dots chỉ khi ở gần anh, bọn em mới thấy thật sự thoải mái thôi.''

Mio cũng lên tiếng, nhẹ nhàng: ``Đúng đó\dots\ Dù cho anh nghiêm túc như em với Flare-chan như vậy\dots''

Flare mỉm cười dịu dàng, ánh mắt sâu lắng như người trưởng thành trong một gia đình đông đúc: ``\dots nhưng nếu thiếu anh, thì bọn em sẽ không thể gắn kết được như hôm nay đâu, Kuon-senpai.''

Cả căn phòng bỗng trở nên yên tĩnh, không phải vì ngượng ngùng, mà vì mọi người đều đang cảm nhận được sự chân thành trong từng câu nói. Kuon nhìn quanh, nhìn vào ánh mắt của từng người trong nhóm \hll\ ấy, nơi đầy những điều muốn nói nhưng không diễn tả thành lời.

Cậu không hề nhận ra được tầm quan trọng của mình đối với các cô gái. Theo bản thân cậu, dù có chức vụ rất cao, nhưng cậu tự nhận mình chỉ là một ``tên lính quèn,'' làm công ăn lương như những người khác, thậm chí còn không so bì được với hào quang của các cô gái nhà Tam Giác Xanh.

Nhưng cậu không ngờ rằng\dots\ vị trí của mình trong công ty cũng được chính các cô gái công nhận, thậm chí có thể được coi là tối quan trọng với họ.

Cậu ngước về phía gia đình mình, những người vẫn đang ngồi nhìn mọi chuyện diễn ra. Cậu cần một góc nhìn ngoài cuộc, mong chờ xem phản ứng của họ sẽ ra sao.

Điều mà cậu nhận được sau đó chỉ là những ánh mắt dịu dàng, nụ cười ý nhị từ Ryujin-Kaien và cái gật đầu như ngầm thừa nhận từ Ryuuhou. Cả ba người trong gia đình Hikawa cũng đồng tình với các cô gái, rằng: ``\textit{Con trai ta hẳn phải rất quan trọng với bọn chúng.}''

Cậu thở ra, bất lực nhưng cũng thấy nhẹ lòng. Dẫu trong tim vẫn còn cảm giác mình không xứng đáng, nhưng nếu điều đó khiến người khác hạnh phúc, thì có lẽ, chỉ lần này thôi\dots

``\dots Ờm, vậy thì\dots\ anh xin nhận.''

Lời vừa dứt, cả hội phía sau, đặc biệt là ba cô nàng Luna, Towa và Kanata, lập tức vỗ tay như pháo nổ, hò reo như thể ai đó vừa đạt huy chương vàng Thế vận hội.

``Phải thế chứ, Kuon-senpai!'' Polka năng nổ nói, trong khi hai cô bạn của mình là Polka và Nene cũng chỉ mỉm cười. Lamy thì khoanh tay, khẽ lắc đầu nhẹ, nhưng khóe môi cô nở nụ cười: ``Anh Tổng nhà ta đúng là được mọi người yêu quý thật.''

Towa và Kanata cùng bước lên. Hai cô gái không nói nhiều, chỉ trao cho cậu mỗi người một món: hai chiếc vòng tay nhỏ, một màu tím, một màu xanh dương, tương ứng với màu đặc trưng của họ. Chúng đơn giản, hơn, có phần không lấp lánh bằng, nhưng lại mang đến một sự tinh tế và nét dấu ấn riêng biệt.

Cậu đón nhận chúng bằng cả hai tay, lần này không một chút do dự, và nụ cười trên môi cậu nói lên tất cả.

Lại một tràng vỗ tay vang lên, như một điệp khúc quen thuộc cho những khoảnh khắc ấm lòng.

Tiếp sau bộ ba thiên thần vừa tặng vòng tay, ba người khác lại bước lên với nụ cười bí hiểm. Đứng giữa là Suisei, hai bên là Marine và Polka lắm trò. Mỗi người đều tay không, nhưng ánh mắt thì đầy tự tin như thể đang chuẩn bị tung ra tuyệt chiêu cuối cùng.

Nàng ``Sao chổi'' đưa cho cậu Tổng một gói quà lớn, rộng chừng hai gang tay, được gói bằng giấy bạc óng ánh xen kẽ sọc hồng và xanh lam.

Cậu cầm lên, hơi chao nhẹ một chút rồi nhíu mày: ``Ơ\dots\ Sao trông nó nặng vậy?''

``Anh cứ mở đi thì biết.''

Cậu lắc chiếc hộp nhẹ, nghe thấy những tiếng lọc xọc bên trong, tựa như một hộp kim loại chứa đầy bi sắt.

Kuon rướn mày, nghi ngờ nhìn về phía ba người: ``\dots Đừng nói với anh là ba đứa tặng anh cục đá đấy nhá?''

``Không có!'' cô nàng giả vờ giận, ``Nó còn tốt hơn cả cục đá nhiều!''

``Anh sẽ không ngạc nhiên nếu mấy đứa tặng cho anh một cục sắt đâu\dots'' cậu lẩm bẩm, cố nhịn cười khi mở gói quà.

Sau vài lớp giấy gói, món quà lộ diện: một chiếc hộp mang biểu tượng của máy chơi điện tử nhà Nin***do S*itch. Một chiếc máy màu trắng bạc, với hai tay cầm phối màu đỏ và xanh lá.

Cậu tròn mắt ngạc nhiên: ``Sao mấy đứa biết là anh muốn cái máy này?''

Suisei đáp: ``Nhờ công của Ryuuhou-chan đấy! Em ấy đã nói cho bọn em biết!''

Cậu liếc về phía cô em gái đang ngồi chơi điện tử. Ryuuhou ngẩng mặt lên ngay khi có người nhắc đến tên cô, và rồi đáp lại ông anh với cái nhún vai như muốn nói: ``Chẳng phải anh đã từng nói là muốn có một cái máy như vậy sao, Onii-san?''

Cậu hướng mắt về bộ ba, rồi gãi đầu gãi tai: ``Anh cảm ơn, nhưng\dots\ tại sao?''

Marine gật đầu, giải thích: ``Em thấy hầu hết mọi người trong \hll\ đều có S*itch để chơi game với nhau, mà anh thì lại không có.''

``\dots Giờ em mới để ý, đúng là thế thật. Nhưng vấn đề là\dots\ tại sao?''

Polka xen vào ngay, giọng líu lo: ``Bọn em cũng muốn anh tham gia chơi cùng đó!''

Suisei cười nghiêng đầu: ``Anh thì hay thích sự đơn giản, nên tụi em chọn cái phiên bản giới hạn này, vừa là bản OLED, lại không lòe loẹt nhưng vẫn nổi bật!''

``Chà\dots\ Ờm\dots\ Vậy\dots\ anh cảm ơn nhiều,'' Kuon lúng túng, nhưng ánh mắt không giấu nổi sự cảm động.

Cậu hướng về cô em gái một lần nữa, và cô bé giơ ngón tay cái lên. ``Thế là tốt rồi đó, Onii-san!''

Ngay khi cậu tưởng rằng món quà đã kết thúc, bỗng chốc Marine ngoài người ra phía sau, hô lớn: ``Gượm đã! Bọn em còn có đồ nữa!''

Cô nàng lôi ra một chiếc túi giấy lớn, nặng đến mức cả hai người phải vác được. ``Món chính vẫn còn chưa hết đâu!''

Cô đặt chiếc túi trước mặt Kuon. Không cần hỏi họ đó là gì, cậu lập tức mở ra, và ngay lập tức đập mặt cậu là gần chục băng trò chơi đủ thể loại, từ những trò chơi máu mặt từ nhà Nin, cho tới những trò chơi đa người chơi nổi tiếng, khiến cho tất cả mọi người đều choáng ngợp.

``C-Cái gì vậy kìa!?'' Korone thảng thốt. Theo sau là Miko, người cũng trố mắt ngạc nhiên: ``Có cả trò mà chị chưa bao giờ thấy kìa, ne!''

Lamy cũng ngạc nhiên với món quà mà nàng Thuyền trưởng và nàng Cáo xiếc tặng cho cậu: ``Hai chị mua nguyên cả thư viện cho máy ạ!?'' tới mức suýt làm rơi chai rượu quý của mình.

``Bọn em tặng kèm luôn cả bộ sưu tập game đấy!''  nàng Chủ xiếc cười khoái chí, ``Không sợ anh bị chán đâu!''

Cậu nhìn xuống đống băng đĩa trò chơi, rồi vừa nhặt băng vừa liệt kê: ``Tron đây có cả Spla*oon, Z**da, M*r*o\dots\ À, có cả Tetris 99 nữa này.''

Ngay khi cái tên của trò chơi xếp hình vừa thốt ra, bỗng dưng tất cả mọi người chững lại như bị đóng băng, những giọt mồ hôi chảy dài bên thái dương.

``T-T-Tetris!?'' Noel cảm thán. ``Cái trò chơi mà Sui-chan làm bà hoàng ấy ư!?''

``Cái trò chơi đã mần chúng ta ra bã đây mà\dots'' Fubuki xanh mặt, đôi tai cáo cùng cái đuôi của cô dựng đứng lên.

``Suisei-onee-san đúng là có cách chọn trò chơi dị thật đó\dots'' cô em gái của Kuon, dù không hiểu chuyện gì, nhưng cũng ngờ ngợ đoán được phản ứng của mọi người.

Suisei nghiêng đầu, tươi cười như thiên thần nhưng ánh mắt lại như của quỷ vương: ``Ừ đó. Nếu anh muốn sống sót qua sinh nhật này thì tốt nhất\dots\ nên luyện trước đi, Kuon-senpai\~{}''

Cả phòng cười rộ lên. Kuon thì vừa gãi đầu vừa thở dài, ánh mắt như thể đang viết di chúc trong đầu, mặc dù cậu biết rằng với trình độ hiện tại, đấu với cô nàng Ca sĩ chẳng khác nào lấy trứng chọi đá.

``Không biết là anh có chơi nổi với mấy đứa không nữa\dots\ nhưng thật lòng cảm ơn tấm lòng của tụi em.''

Và cũng như những lần trước, mọi người lại vỗ tay rầm rầm, nhưng lần này kèm theo cả tiếng cười không ngớt và những lời động viên từ mọi người, nhất là khi thấy mặt Kuon hơi tái lại lúc cầm hộp băng trò chơi Tetris 99.

Suisei có lẽ là người nở nụ cười mãn nguyện nhất.

``Chẳng biết sinh nhật năm nay liệu có yên ổn không đây\dots'' cậu Tổng tự nhủ, cười chát với tình cảnh hiện tại.

Sau khi những món quà mang tính công nghệ, giá trị và cả tinh thần lần lượt được trao, bầu không khí trong căn phòng dần yên bình trở lại. Ba người cuối cùng bước lên phía trước, không lộng lẫy, không tay xách nách mang, nhưng ánh mắt đều lấp lánh.

Sora, Rushia và Ayame.

Những người thuộc nhóm tặng quà lúc đi mua hàng trước đó cũng cảm thấy hơi ngạc nhiên khi họ bước lên.

``\hl{Ể? Sora-senpai cũng có gì tặng cho Kuon-senpai sao?}'' Aqua nói nhỏ với cô bạn.

``\hl{Chắc là thế. Tôi đâu có nhớ chị ấy cùng hội với Sui-chan hay Towa-chan đâu,}'' Shion đáp.

``\hl{Có khi họ có món quà gì đó đặc biệt để tặng họ chăng?}'' Botan chêm vào, mặt vẫn không tỏ biểu cảm gì.

Dưới những tiếng xì xầm của ba người, trên ``sân khấu'', cô nàng Thần tượng nhẹ giọng, khẽ cúi đầu, rồi nói: ``Cho bọn em xin lỗi anh vì chưa chuẩn bị quà gì hoành tráng cả\dots\ Bọn em chỉ làm vài thứ nhỏ thôi, trong cả buổi chiều chuẩn bị vừa nãy.''

Nàng Pháp sư thì rụt rè gật đầu: ``Nhưng mà\dots\ mỗi món tụi em làm đều mang theo một chút kỷ niệm với anh.''

Nàng Quỷ sừng thì cười ngượng, nhưng vẫn không giấu nổi tính lanh chanh vốn có của mình: ``Anh chắc là không ngại đâu nhỉ?''

Cậu Tổng nhìn cả ba món đồ tự làm, cười nhẹ: ``Không sao. Anh cũng muốn thử xem mấy đứa đã làm gì.''

Đầu tiên là Sora.

Cô lấy ra từ phía sau một chiếc hộp nhỏ, mở nắp, bên trong là một khung ảnh gỗ thủ công, được khắc thủ công khá khéo, viền bằng những miếng gỗ nhỏ ghép lại, dán keo trong suốt, bên trong là một tấm ảnh in lại từ điện thoại: đó là một bức ảnh chụp chính cô đang giả vờ mất trí nhớ mà quên mất cách đi đứng, cạnh cô là Miko và Roboco đang tranh cãi nhau.

\begin{figure}[H]
    \centering
    \includegraphics[width=15cm]{../../../../Image/Book?/3-8e.png}
    \caption{Chính xác hơn, đây là một cảnh lấy trực tiếp từ {\flaregothic ÉPISODE 37}, quyển số Ba.}
\end{figure}

Sora cười đầy tinh nghịch: ``Anh nhớ cái lần A-chan và chúng ta bàn nhau cách chơi khăm họ từ năm ngoái không?''

Kuon bật cười thành tiếng: ``Làm sao quên được. Chính chúng ta là những người hiến kế cho trò này chứ còn sao.''

Miko ngồi phía xa lập tức hét lớn: ``Pha đó làm dọa cho bọn em sợ chết khiếp rồi đó, ne!''

Roboco vỗ trán: ``Đến tận bây giờ mình vẫn không hiểu chuyện gì xảy ra luôn á\~{}''

Cô nàng cùng Thế hệ che miệng cười: ``Đó là lần đầu tiên mà mình cảm thấy muốn nghịch ngợm một chút\dots\ nên là mình khắc lại khoảnh khắc đó luôn.''

``Có bao nhiêu kỷ niệm hay hơn mà lại không lấy, ne\dots'' nàng Vu nữ phồng má lên.

``Thôi, thôi\~{} Pha đó dù gì cũng đáng nhớ mà, Miko-chi.''

Tiếp đến là Rushia.

Cô hơi lúng túng đưa ra một món đồ hình trái tim nhỏ nhắn, được làm bằng nhựa resin trong suốt, bên trong là một con búp bê nhỏ dễ thương mang dáng dấp hình cô được may bằng tay, phía dưới là dòng chữ tí hon được viết tiếng Anh khá vụng về: ``You're enough. (Với em, thế là quá đủ.)''

Cậu đọc dòng chữ trong đó, rồi nhìn về phía cô nàng đang lấp ló đằng sau món đồ. Cậu nghiêng đầu: ``Ru-chan, quà của em mang ý nghĩa gì đây?''

Cô nàng lí nhí đáp: ``Anh còn nhớ hồi anh làm tư vấn viên thay cho Pekora vì dám xúc phạm\dots\ núi của em\footnote{Xem lại phần {\abv Truyện phụ}, {\flaregothic ÉPISODE 111}, quyển số Chín.} không?''

``Này!'' nàng Thỏ lên tiếng. ``Tôi đã cố gắng hết sức để tư vấn cho bà rồi nhá, peko! Ca của bà sao mà khoai khoẳm vậy!''

Nhưng cô nàng không để tâm tới những gì Pekora nói, mà tiếp tục nhìn về phía cậu Tổng. Cô nàng với đôi má hơi ửng hồng, tiếp tục: ``Lần đó, anh có nói rằng\dots\ ê-tồ\dots\ `tốt gỗ hơn tốt nước sơn', đúng chứ?''

``Đúng là như thế thật,'' cậu đáp. ``Anh có nói rằng hãy cứ là chính mình, đừng để tâm tới những gì người ta đàm tiếu, đúng chứ?''

Rushia gật đầu nhẹ. ``Cũng chính sau hôm đó, em cảm kích lắm\dots\ Em đã thử làm một món gì đó để tặng, nên đã may con búp bê nhỏ này, nhưng mà\dots''

``\dots sợ là không đủ, nên cuối cùng đổ nhựa trong suốt lên sao?''

Cô cúi mặt, nhỏ giọng nói: ``Xin lỗi nếu như nó không đẹp-nanodesu\dots\ Em cũng không khéo tay lắm mà\dots''

Kuon đón lấy món quà như nâng báu vật: ``Anh thấy nó cũng đẹp mà. Anh thật sự vui vì em đã thực sự quý trọng lời khuyên đó.''

Towa huých Flare: ``Ai mà chê Rushia-senpai thì đúng là đồ vô cảm luôn á.''

Nàng Dị giới đáp lại: ``Đó cũng là một hình thức yêu quý của họ thôi mà. Cơ mà cũng có những người cực đoan lan tới thật\dots''

Cuối cùng là Ayame.

Nàng Oni tiến lên, tay cầm một món lạ lùng: một chiếc mô hình nhỏ, trông như một sân khấu thu nhỏ. Trên đó là hai nhân vật tí hon: một đang bơi trong hồ dung nham bằng biểu cảm tỉnh bơ - không ai khác chính là Miko\footnote{Liên hệ tới {\flaregothic ÉPISODE 81}, quyển số Sáu.}. Một cảnh khác là hình vẽ tí hon của Subaru đang bay lên cao, với chính cô đang cầm cây chùy vung vào chiếc ghế bành\footnote{Xem lại {\flaregothic ÉPISODE 104}, quyển số Tám.}.

Ayame cười khúc khích giải thích: ``Đây là bản diễn họa của hai lần anh giúp em quay video, nhớ không? Một là `cách chữa trầm cảm tháng Năm' với Subaru và Roboco-senpai, hai là `sưởi ấm bằng dung nham' với Miko-senpai. Em dựng lại cảnh y chang luôn á, yo!''

Nàng Vịt hoảng hốt khi nhìn thấy chính mình trong mô hình đó: ``Khoan đã, cái cảnh bay lên đó là tôi á!? Phản đối, phản đối kịch liệt!''

Nàng Vu nữ thì ôm đầu: ``Làm ơn đừng nhắc tới vụ dung nham đó nữa\dots\ Hôm đó chúng ta dọn mớ lộn xộn đó cũng đủ mệt người rồi đó, ne!''

Nàng Người máy cũng nhăn mặt: ``Tại ai mà chị suýt bị hỏng hả?''

A-chan ngồi bên cạnh bố mẹ của Kuon, vừa thở dài ngán ngẩm nhưng cũng vừa khúc khích: ``Nakiri-san đúng là có những ý tưởng điên rồ thật đấy.''

Nàng Quỷ sừng tủm tỉm: ``Em nhớ anh với Miharu-san bình luận đằng sau cũng hài hước lắm chứ! Xong rồi còn chỉnh ánh sáng, góc quay cho em, và voilà! Chúng ta có được những thước phim hoàn chỉnh.''

``Để rồi chúng ta phải quay lại chúng từ đầu,'' cậu Tổng châm chọc. ``Anh không nghĩ anh chịu được với lượng Nakirium toát ra từ người em tuôn ra như suối đâu.''

Tất cả mọi người trong căn phòng đều bất giác bật cười với lời bỡn cợt của cậu, trong khi nàng Oni thì cười khì.

Cô nàng lon ton dâng hai bộ mô hình của mình về phía chủ nhân bữa tiệc, nở một nụ cười tươi như nắng: ``Kuon-senpai, em muốn anh giữ lại một sân khấu nho nhỏ này.''

Kuon cầm mô hình, im lặng một lúc lâu, rồi nở nụ cười nhẹ nhàng: ``Anh không ngờ những việc nhỏ nhặt anh từng làm lại được các em trân trọng đến vậy\dots''

Cả căn phòng lại rộn lên một lần nữa, hòa cùng với những tiếng cười khúc khích và ánh mắt đầy ấm áp.

Ba món quà cuối cùng khép lại chuỗi tặng quà không thể nào quên trong sinh nhật của cậu. Chúng không lấp lánh như dây chuyền, không hoành tráng như máy chơi game, nhưng chứa đựng cả một bầu trời ký ức, gắn bó và yêu thương.

\ct

Sau khi món quà cuối cùng được trao tận tay Kuon, A-chan cất tiếng, vẫn với chất giọng đều đặn nhưng không giấu nổi sự hài lòng: ``Và sau đây, xin mời Hikawa Kuon-san có đôi lời phát biểu cảm nghĩ.''

Cánh tay cô vung về phía cậu một cách đầy trang trọng.

Ngay lập tức, cả căn phòng yên lặng đúng nửa giây trước khi mọi ánh nhìn dồn về nhân vật chính đang\dots\ há hốc miệng vì bị réo tên một lần nữa ngoài tính toán của cậu.

``L-Lại là em nữa ạ!?'' cậu Tổng thốt lên, gần như bật dậy ngay khi cậu vừa lui về gia đình.

Thấy cậu con trai vẫn chưa chịu lên, Ryujin nhướng mày, vỗ tay vào đùi một cái: ``\vie{Lên đi con. Làm người đàn ông thì phải biết đứng dậy nói chuyện đàng hoàng!}''

``\vie{Nói to, tròn vành rõ chứ nhé con!}'' Kaien cũng lên tiếng động viên.

Ryuuhou thậm chí còn lôi chiếc điện thoại ra, nhanh chóng tiếp lời, vẫy tay: ``\vie{Cố lên Onii-san! Em đang quay lại đây!}''

``\vie{Lại còn có cả trò quay phim nữa à\dots}'' cậu thở hắt, bất ngờ với hành động của cô em gái.

Tới đây, cậu khẽ ười chát, mắt đảo một vòng qua ba mươi mốt ánh mắt đầy tò mò trước mặt. ``Nhìn mọi người như thể đang làm chủ tòa cho mình vậy\dots''

Cậu hít một hơi thật sâu trong tiếng vỗ tay nửa khích lệ, nửa trêu ghẹo vang lên râm ran.

``Lên đi! Lên đi!'' những tiếng hô gọi từ mọi người vang lên, như thể đang thúc giục cậu lên nhanh hơn.

Bước ra giữa gian phòng, Kuon đứng thẳng lưng, không cầm lấy míc-rô, chỉ có giọng nói trầm ổn tự nhiên lan ra giữa khung cảnh ấm cúng.

Cậu hắng giọng một cái, rồi bắt đầu nói: ``Đầu tiên thì anh xin gửi lời cảm ơn tới tất cả mọi người\dots\ vì đã đến dự bữa tiệc sinh nhật này. Anh thật sự bất ngờ, hoàn toàn không đoán được gì cả.''

Một vài tiếng cười nho nhỏ bật lên. Kuon quay sang A-chan, nhẹ gật đầu: ``Cũng nhờ A-san là người đã lập kế hoạch tổ chức sinh nhật bất ngờ này ngay từ đầu, nên mọi chuyện đến tận bây giờ đã thuận buồm xuôi gió ạ.''

A-chan khẽ mỉm cười, gật đầu. Ở một góc, Kanata huých Miko, đùa: ``Công phu đến mức chị suýt ngã cầu thang luôn đó.''

Kuon đưa mắt về phía bố mẹ, tiếp tục bài diễn văn của mình, nói ra hết những gì tâm can trong lòng: ``Bố mẹ ạ\dots\ Con cảm ơn vì đã nuôi dạy con nên người. Mặc dù con vẫn còn ăn bám khá nhiều, và lười chảy thây nữa\dots''

Cả căn phòng bật cười, một tiếng cười không ồn ào nhưng đầy chân thành. Hiển nhiên cũng có một số tiếng há hốc mồm nhẹ khi cậu dám tự kiểm điểm chính mình một cách đầy thẳng thắn, nhưng với họ, dù gì cũng là một điểm tốt.

Ryujin khịt mũi, còn Kaien thì lắc đầu: ``Cậu ấy ít nhất cũng biết nhận ra là tốt rồi\dots''

Kuon cười nhạt, xoa gáy rồi quay lại phía ba mươi mốt người con gái ngồi đối diện mình.

``Anh cũng muốn cảm ơn tất cả mọi người trong \hll\ vì đã tổ chức một bữa tiệc sinh nhật hoành tráng và tỉ mỉ thế này. Anh không ngờ\dots\ mấy đứa ngốc nghếch và làm trò như các em mà cũng biết đa-zi-năng thiệt đấy.''

Lần này, bố mẹ cậu cười lớn nhất. Nhưng đám con gái thì cười gượng nhiều hơn.

Noel huých Flare: ``Vậy là anh ấy đang khen thật hay khen đểu vậy?''

``Kuon-senpai cứ phải châm chọc bọn này mới chịu hả, peko\dots'' Pekora càu nhàu, đôi tai thỏ của cô cụp xuống.

Kuon nhướng mày: ``Thì đúng mà! Phá hoại xong rồi đến mức cả anh lẫn A-san đều phải đền bù, sửa chữa đến phát ốm! Mắng xong rồi, mai lại đâu vào đấy, đúng không? Ngốc đến mức đó thì còn gì để cười nữa chứ!''

Một khoảng im lặng. Rồi tất cả lại cười trừ, vì\dots\ đúng là họ không thể phản bác.

``Quả nhiên chúng ta vẫn sẽ mãi như vậy à\dots'' Matsuri ngán ngẩm, tự trách bản thân.

``Nhưng nếu không như thế thì đâu còn là chúng ta, đúng chưa?'' Nene vỗ vai an ủi tiền bối mình.

Kuon cười nhẹ, nhưng giọng chuyển sang nghiêm túc hơn: ``Nói thế thôi. Không có những trò phá, không có những lần nghịch ngợm, không có tiếng ồn và những pha `hỗn chiến' kia\dots\ thì làm gì có cái gọi là \hll\ như hôm nay. Nếu không có những điều đó, liệu chúng ta có thể ngồi cạnh nhau thế này, trong một căn phòng tràn ngập tiếng cười và ánh đèn ấm áp không?''

Một vài người khẽ gật đầu. Một số khác thì nhìn nhau và bất giác mỉm cười.

``Và nếu không có những ngày tháng ấy, làm sao các em - những sinh vật vốn không thuộc về thế giới này - có thể sống hòa hợp, tìm được công việc, được khán giả đón nhận, và trên hết\dots\ tìm được ý nghĩa cho chính sự tồn tại của mình?''

Lần này, câu nói như một mũi tên thẳng vào lòng từng người, đặc biệt là từ các thành viên phi nhân. Không cần ai lên tiếng, mà chỉ ánh mắt bỗng trở nên long lanh của một vài cô gái là đủ.

Kuon cúi đầu một chút, rồi ngẩng lên: ``Vì thế\dots\ anh cảm ơn các em, vì đã cho anh một cơ hội được sống và làm việc cùng mọi người. Và cảm ơn bố mẹ, vì đã tạo ra một con người có thể giúp ích được điều gì đó cho thế giới này.''

Cậu ngẩng đầu, thở ra một hơi dài.
Đôi mắt ánh lên quyết tâm, giọng nói bỗng vang to hơn: ``\textbf{Tôi hứa rằng\dots\ trong năm nay, và những năm sau nữa, tôi sẽ nỗ lực hết mình để không phụ lòng tin của mọi người!
    Cả nhà thấy thế nào!?}''

``\textbf{Rất tuyệt!}'' một giọng ai đó hét lên.

``\textbf{Quá được luôn, Kuon-senpai!!}'' một tiếng khác vang lên từ góc phòng.

Và rồi, không ai bảo ai, những tràng pháo tay vang dội như chưa từng có.

Không chỉ là tiếng vỗ tay một cách đơn thuần, mà là tiếng trái tim cùng chung một nhịp đập, một câu trả lời chung mà tất cả các cô gái đều đang được mong chờ tới từ một người thủ lĩnh nhất ngôn, nghiêm túc và rộng lượng.

Cậu Tổng giơ nắm đấm lên, giật lấy cây míc-rô từ tay nàng Đeo kính, rồi hét lớn: ``\textbf{Vậy không để mọi người chờ lâu nữa. Chúng ta cùng bắt đầu bữa tiệc sinh nhật nào! Cả nhà sẵn sàng chưa!?}''

``\textbf{Sẵn sàng!!}'' ba mươi lăm người giọng cùng đồng thanh hô vang.

Căn phòng như bùng nổ trong một khung cảnh náo nhiệt và rộn ràng, một bữa tiệc không chỉ của bánh, hoa, hay quà tặng, mà còn là của tình cảm, của gia đình, và của một tập thể đầy vững mạnh.

Đây chắc chắn có thể sẽ là bữa tiệc sinh nhật đáng nhớ nhất trong cuộc đời của Hikawa Kuon, thậm chí còn ý nghĩa hơn so với lần sinh nhật ba ngày liên tiếp đó.

Và như Hikawa Kuon nói:

\begin{center}
    \uline{\eng{Let's get the party started.}}
\end{center}

\ct

Đồng hồ đã chỉ tới tám giờ kém mười phút.

Sau khoảng mười hai tiếng đồng hồ chạy đôn chạy đáo, chuẩn bị, di chuyển, đợi chờ, cười đùa, hồi hộp và bí mật\dots\ thì giờ phút mà ai nấy đều mong đợi cũng đã tới.

Các cô gái \hll, những người đã phải kìm nén cơn đói, hoặc lót dạ bằng mấy món vặt trên đường, hoặc thậm chí chẳng ăn gì từ ngày hôm qua như Mio, Matsuri và Sora, giờ đây đều ngồi quanh một mâm cỗ lớn, bày biện đủ nồi lẩu đang sôi ùng ục và chiếc chảo nướng đang bốc khói nghi ngút.

Kuon vừa đặt phần nguyên liệu cuối cùng xuống bàn vừa nói: ``Ai thích lẩu thì ăn lẩu, ai thích nướng thì ăn nướng nha!''

Chẳng mấy chốc, tất cả mọi người đã bắt đầu đụng đũa, nhưng không quên đồng thanh như tiếng sấm: ``\textbf{Itatakimasu!!/Mời cả nhà ăn cơm!!}''

Ngay tức khắc, tất cả mọi người bắt đầu thưởng thức bữa tối thịnh soạn trong ngày. Có người bắt đầu gắp lấy gắp để như sợ bị ai đó ăn phần, người thì tĩnh tâm lấy đồ ăn, tỏ vẻ không có gì xảy ra, người thì xung phong làm ``mâm trưởng'' nướng đồ và gắp nhúng đồ cho mọi người.

Matsuri nhanh tay gắp ngay một vỉ ba chỉ lợn thả lên chảo, đôi mắt cô sáng quắc lên như bị bỏ đói, trong khi Mio và Sora đồng thời đổ một đĩa rau và nấm vào nồi lẩu đang sôi.

``Mọi người ăn từ từ thôi, vẫn còn nhiều lắm!'' Choco lên tiếng nhắc nhở. Mel cũng cười: ``Tận từng này đồ nướng với lẩu chứ ít gì đâu!''

Shion cũng thêm vào: ``Cứ thoải mái đi nha mọi người! Hôm nay là ngày đặc biệt mà!'' trong khi tay vẫn đang luân hồi gắp đồ.

Cậu Tổng liền đi quanh từng bàn, tay cầm theo một thùng bơ to tổ bố và một rổ đầy những lọ sa tế - thứ mà nhà cậu đã mua dự trữ trước đó. ``Ai muốn ăn nướng ngon thì phết bơ lên! Còn ai ăn lẩu thì thử cái sa tế tôm này nha! Đảm bảo là ngon phê lưỡi luôn!''

``Cái gì? Bơ? Sa tế?'' Nene nghiêng đầu, thắc mắc với đồ tiếp tế từ cậu.

``Thì các cháu cứ ăn thử đi,'' Ryujin, người đang uống cốc bia với Lamy, cao hứng nói. ``Nguyên liệu bí mật của nhà Hikawa đấy!''

``Cụ thể là thế nào vậy?'' Aqua hỏi, cầm một lọ sa tế lên.

``À, cái này thì dễ mà\dots''

Cậu giải thích ngắn gọn cách dùng, rồi bắt đầu chia đều sa tế cho từng nồi lẩu. Thùng bơ thì chuyền tay nhau mỗi người lấy một thìa vừa đủ, thoa lên miếng thịt đang nướng hoặc vào chảo nướng đang xèo xèo.

Chỉ một lúc sau, mùi thơm đã lan khắp căn phòng. Hương vị bơ chảy ra từ thịt lợn, quyện với vị sa tế cay cay thơm nồng từ nồi lẩu, khiến bụng ai nấy đều réo lên.

Noel, người vốn thích các món đậm vị, thử nhúng một muỗng nhỏ sa tế vào phần nước lẩu đang sôi, rồi múc một muỗng chan lên tô cơm bò còn sót lại của mình.

Nàng Hiệp sĩ cho một miếng.

Ngay lập tức, một hương vị ran rát nhẹ bùng nổ trong miệng, hòa quyện với vị ngọt lịm của thịt bò và những hạt cơm trắng. Cô nàng trầm trồ, đôi mắt cô sáng như sao: ``Oa\dots\ vị của nó lan tỏa trong miệng này! Phải mua một ít đem về thôi!''

Bên phía thịt nướng, người ăn ngạc nhiên nhất chính là Lamy: ``T-Thịt lợn nướng ư? Từ trước tới giờ em toàn ăn thịt bò nướng\dots\ Thế mà cái này\dots\ mềm và béo quá\dots''

``Ryujin-san nói đúng thật, thịt ướp thế này ngon hơn hẳn!'' Watame cười, tay gắp thêm một miếng thịt từ chảo nướng.

``Cô Kaien-san cũng là người khá sáng tạo với nước lẩu mà!'' Kanata bổ sung, múc thêm một thìa nước vào bát.

``Còn chưa tính tới cái lớp bơ này nữa!'' Towa vừa nói vừa đưa miếng thịt lên miệng, mắt lim dim như đang ở thiên đường.

``Mấy đứa cứ ăn thoải mái đi, các cháu ăn được hết thì tốt quá,'' Ryujin cười hà hà.

``Các cháu không phải khách sáo đâu!'' vợ của người chồng nói, ``Cô đây còn sợ thiếu cơ.''

Thức ăn liên tục được thêm vào, hết thịt lợn, thịt gà, lòng gà, bò viên, rồi lại đến đậu phụ, rau muống, nấm kim châm, bắp cải thảo\dots\ Đặc biệt là rau, được chuẩn bị rất kỹ từ trước, cũng được mọi người săn đón.

Tuy nhiên, một người duy nhất vẫn khư khư né tránh rau bằng mọi giá. Hoshimachi Suisei.

Miko trề môi, dùng đũa gắp một nhúm rau muống cho vào bát của cô bạn: ``Sui-chan ăn rau đi nào, ne\~{}''

Lập tức, Suisei xua tay lia lịa như gặp rắn: ``Không, không, tuyệt đối không! Ứ thèm!''

``Ơ kìa chị!'' Haato cũng lên giọng châm chọc, ``Rau này người ta trồng hữu cơ á! Chẳng phải idol chúng ta phải có một chế độ ăn cân bằng sao?''

``Chị nghĩ chế độ ăn của chị là quá đủ cân bằng rồi, khỏi cần rau!'' Suisei vênh mặt. ``Kuon-senpai thậm chí còn không thích ăn rau kìa!''

``À, về chuyện đó thì\dots'' Sora lên tiếng, rồi ngoái nhìn về phía Kuon đang quây quần với Thế hệ thứ Tư.

Cả nhóm đồng loạt quay đầu lại. Và rồi\dots\ cảnh tượng trước mắt khiến Suisei như bị đóng băng.

Tại đó, nhân vật chính của bữa tiệc đang ngồi kế bên bố mình, tay gắp một miếng nấm đùi gà to tổ chảng từ nồi lẩu chan vào chén, rồi vừa thổi vừa ăn ngon lành như thể đó là thịt bò Kobe hảo hạng.

``Mmmm\dots\ mấy đứa đúng là biết cách chọn nấm thật\dots'' Kuon còn bình luận với Kanata và Towa, khiến cả bàn gật gù tán thành, còn Choco thì cười nhẹ.

Nàng Ca sĩ quay lại, mặt đơ như tượng, còn Miko với Haato thì cố gắng nén cười đến đỏ mặt.

``Cảm giác\dots\ bị phản bội\dots'' cô nàng rầu rĩ, hai tay chống má nhìn vào hư vô. ``Tư tưởng Kuon-senpai\dots\ hóa ra là ăn rau sao\dots?''

Dường như cô nàng cũng chẳng để tâm tới những sự việc trước mắt, khi ``đồng minh thân cận'' nhất của cô trong chiến dịch chống rau cuối cùng cũng đã ``đâm sau lưng'' với lòng thành của mình.

``Ne, ne\~{} Không ăn rau là không được lên sân khấu lần tới đâu nhaaa\~{}'' nàng Vu nữ đe dọa bằng giọng em bé nũng nịu.

``T-Tôi\dots\ Thôi được\dots'' nàng ``Sao chổi'' thở dài, rồi quay sang chén của mình đã đầy ắp rau từ khi nào. ``Kuon-senpai ăn được thì tôi cũng sẽ ăn được thôi.''

Cô hít sâu một hơi, rồi mạnh dạn gắp cả nắm rau muống bỏ vào nồi, miệng lẩm bẩm như tụng kinh: ``Vì tương lai của âm nhạc và sự nghiệp của mình!''

Cả bàn cười ồ, còn Haato thì giơ điện thoại lên nhưng bị Suisei tóm kịp thời, và cuộc chiến với rau vẫn tiếp tục, tuy lần này, Suisei có vẻ đã chịu thua.

Ở một chiến tuyến khác, thịt cũng là ``mặt trận'' được tất cả mọi người chú ý. Có hai ``chiến thần'' gây chú ý trong mảng này: Noel và Kuon.

Mỗi khi mọi người mới đặt thịt lên nướng, còn chưa kịp lật qua thì cậu Tổng đã nhanh tay gắp ngay, chấm nước sốt rồi ăn ngon lành. Nàng Hiệp sĩ thì không chịu kém cạnh: không những ``cày'' phần lẩu mà còn chủ động chiếm luôn cả vỉ nướng gần đó.

``Ê-ê! Hai người này là máy nghiền thịt hả!?'' Subaru trợn mắt, khi nhìn đĩa thịt quanh họ dần vơi.

``Giờ em mới biết anh ấy cũng là bản sao của Noel-senpai đó\dots'' Watame xanh mặt. ``Chỉ hy vọng hai người đó không làm hại gì đến em là được\dots''

``Có hai chiến tuyến, thì cần hai tướng lĩnh chứ!'' Kuon đáp tỉnh rụi, miệng vẫn nhai nhồm nhoàm.

Noel thì bĩu môi: ``Thịt là lẽ sống. Không ăn là phí. Phí là tội.''

``Hai anh chị đúng là kẻ sát thịt mà\dots'' Ryuuhou nói với giọng mỉa mai, khẽ bĩu môi.

Lời vừa dứt, cả nhóm cười rầm rầm, rồi tiếp tục ``chiến đấu'' cùng lẩu và nướng, cuộc chiến mà ai cũng mong sẽ kéo dài mãi.

``Ăn từ từ thôi con ơi,'' Kaien lên giọng nhắc nhở, ``còn để phần cho mọi người nữa chứ!''

Suisei, người vẫn đang vật lộn với bát rau trông cảnh tượng trước mặt, cũng chỉ thở dài tiếc nuối bản thân: ``Mấy người quên mất còn có máy nghiền thịt là tôi đây sao\dots?''

Không khí bữa tiệc vừa ấm cúng, vừa náo nhiệt, và gần như lúc nào cũng có tiếng cười. Những lời bàn tán rôm rả về những cuộc chuẩn bị buổi phát sóng sắp tới, những câu chuyện phiếm, và những kỷ niệm của họ với Kuon được vang ra khắp nơi, tạo nên một bữa tiệc sôi động.

Cuối cùng, sau khoảng một tiếng vui đùa với lẩu và nướng, toàn bộ số đồ ăn đã được chén sạch, có nồi lẩu thậm chí còn cạn cả nước, không còn lấy một giọt nào!

Bữa tối no nê cuối cùng cũng khép lại. Khi đồng hồ điểm tám giờ rưỡi, cả gian phòng vẫn còn vương mùi thơm còn vương lại từ những nồi lẩu, vỉ nướng vừa cạn dần. Đến lúc này, mọi người mới bắt đầu dọn dẹp chiến trường ẩm thực vừa diễn ra.

``Nene chưa từng ăn bữa nào đã như vậy luôn á!'' nàng ``Hải cẩu'' xoa bụng, ngồi phịch xuống ghế, mặt hớn hở như mèo con vừa được ăn no.

``Đầy bụng thiệt đó\dots'' Matsuri nằm dài ra sàn, ôm bụng thở dốc. ``Chị ăn hơn năm phần thịt luôn\dots''

``Ngon mà no quá\~{}'' Lamy vỗ tay vào má, mặt hồng lên vì no lẫn vì rượu sake nhẹ. ``Lần tới chắc em cũng sẽ phải học theo cách làm từ chú Ryujin-san và Kuon-senpai mới được\dots''

``Công nhận món sa tế gì đó ăn với cơm bò cũng được luôn à!'' Noel ngẫm nghĩ, tay vẫn chưa ngừng gắp dù mọi người đã lục tục đứng dậy dọn dẹp.

``Đây là tiệc sinh nhật mà Kuon-senpai xứng đáng có được đó!'' Kanata hô lên, rồi vác mấy nồi lẩu còn chút nước dùng lên trên mâm. ``Em mang cái này đi đổ ở tầng sáu nhé!

Từng người một bắt tay vào công cuộc dọn dẹp: người gom bát đũa, người xách nồi niêu, vỉ nướng. Một số thành viên thậm chí ``chịu chơi'' tới mức xách cả xoong vào nhà tắm tầng ba để rửa vì quá đông người chờ tới lượt ở bếp.

``Thịt lợn phết bơ ngon quá chừng luôn\dots'' Polka vẫn còn thòm thèm liếm ngón tay.

``Chị đâu có ngờ là bữa tối nay sẽ như một bữa tiệc như vậy đấy!'' Towa lẩm bẩm.

Mio đáp lại, tay lau bàn cùng với Watame và Flare: ``Cũng nhờ mấy đứa chuẩn bị kỳ công mà.''

Một số thành viên khác dọn ghế, dọn khăn trải bàn, nhường chỗ cho các tiết mục kế tiếp. Danh sách cho các hoạt động tiếp theo sau bữa ăn chính vẫn còn ít nhất bốn tiết mục nữa.

Kuon đưa mắt nhìn quanh một lượt, rồi bất chợt dừng lại về một góc phòng.

``Anh để ý thấy vẫn còn một thứ gì đó chưa được mở ra này,'' cậu nói, ánh mắt tỏ vẻ hơi tò mò. Trước mặt cậu là một tấm màn lớn màu tím che kín thứ gì đó cao gần sát trần.

Dường như vì cậu quá hăng say với bữa tiệc lẩu và nướng, hoặc vị trí của nó được ẩn với lớp sơn và giấy tím quá hoàn hảo, nên đến bây giờ cậu mới để ý tới sự tồn tại của tấm màn bí ẩn này.

``À, cái đó là để cho bất ngờ cuối cùng cho bữa tiệc sinh nhật này đấy,''Sora đáp, tay vẫn không ngừng quét chổi quanh khu nhà.

``Vẫn còn à?'' Kuon bật cười nhẹ, không khỏi ngạc nhiên vì độ chịu chơi của cả nhóm. ``Mấy đứa vẫn còn trò gì cho anh xem nữa?''

Ngay lúc đó, ba bóng người tiến lại gần: Shion, Pekora và Luna, mỗi người giữ một góc màn.

``Rồi anh sẽ thấy thôi, Kuon-senpai. Luna, Pekora, kéo nào!''

``Chuẩn bị nào! Hai, ba\dots\ peko!''

Luna vô tình kéo lệch nhịp, khiến màn nghiêng hẳn về một bên, suýt nữa đổ trùm lên người Shion.

``\textbf{Luna!!!}'' nàng Phù thủy hét lên, suýt nữa bị đồng đội phản bội mà kéo thẳng xuống đất.

Đáp lại tiếng hét như trời giáng vô tình thu hút sự chú ý của mọi người, nàng Công chúa xứ kẹo chỉ lè lưỡi xin lỗi tới líu lưỡi: ``Xin lũi chị mà, Shion-senpai\~{}''

Tấm màn được kéo xuống, để lộ ra một chiếc bánh sinh nhật cao tới sáu tầng, tầng dưới cùng to đến mức hai người ôm không xuể. Mỗi tầng lại chứa bốn đến sáu loại bánh khác nhau, đủ đến hai mươi mùi vị, từ trà xanh, sô-cô-la, vani, dâu, cam, chanh dây\dots\ cho tới cả hương sữa dừa, hạt dẻ, matcha sữa, phúc bồn tử\dots

Chiếc bánh được trang trí rực rỡ, với kẹo lấp lánh ánh kim, cookie hình idol cắm đều quanh viền tầng, si-rô dâu đỏ và bạc hà xanh đan xen nhau tạo thành từng đường sóng uốn lượn, và hoa quả tươi đủ màu sắc điểm xuyết trên đỉnh.

Nhưng đặc biệt nhất là một tấm giấy dài dán dọc theo mặt trước, được chia làm 31 ô vuông, mỗi ô là một lời chúc sinh nhật từ từng người trong gia đình \hll\ và của cô em gái Ryuuhou. Có nét chữ nghiêng nhẹ dịu dàng như của Sora, có chữ cứng cáp như của Towa, cũng có cả những dòng thơ vần vèo của Polka hoặc những hình vẽ đáng yêu của Aqua. Mỗi ô là một mảnh cá tính, và tất cả gộp lại thành một bản giao hưởng đầy màu sắc.

``Kuon-senpai thấy chiếc bánh đó như thế nào?'' Kanata húc nhẹ vào vai của cậu Tổng, người vẫn đang tỏ rõ sự ngạc nhiên với chiếc bánh khổng lồ đó.

``\dots Mấy đứa thực sự đã làm nên chiếc bánh này sao?''

``Thực ra, Ryuuhou-chan cũng tham gia làm chiếc bánh này đó,'' Choco tiếp lời.

Cậu hướng về cô em gái, người đang nở một nụ cười rất tươi: ``Em được chỉ định làm bếp trưởng cho món bánh kem này đấy, Onii-san!''

``Aku-tan với Shion-chan cũng là những người góp sức vào món bánh này mà-nanodesu,'' Rushia chỉ ra.

Hướng về cặp bạn thân vừa nhắc tên, cậu nhìn thấy những cặp má hơi đỏ vì ngượng nghịu.

``E-Em nghĩ rằng là một hầu gái, em nên làm một món ăn nào đó phục vụ bữa tiệc\dots'' nàng Hầu gái rụt rè. ``Thế nên em là người nghĩ ra ý tưởng này!''

``Đến cuối cùng tôi vẫn phải là người làm phần lớn mọi thứ đấy thôi,'' nàng Phù thủy buông lời, khẽ tặc lưỡi. ``Bà vẫn còn lơ đễnh lắm.''

``Đ-Đừng có nói vậy chứ, cái bà này!''

``Tôi nói có sai đâu?''

Nhìn thấy cặp bạn lại chí chóe nhau như mọi ngày, cậu bật cười khe khẽ. ``Thôi được rồi. Dù sao thì cũng cảm ơn hai đứa, cùng với Ryuuhou vì đã làm nên chiếc bánh này. Đến nằm mơ anh cũng không nghĩ là sẽ thấy một chiếc bánh khổng lồ như thế đâu.''

Nàng ``Hành tím'' chỉ e thẹn gật đầu, không biết phải nói gì, còn nàng ``Tỏi'' thì vờ quay ngoắt sang chỗ khác, dường như chưa chịu nổi lời khen đường đột này.

``Cái này mà đưa ra đấu giải bánh kem chắc quán quân luôn á\~{}'' Watame gật gù, trong lòng vẫn hơi tiếc nuối vì không tham gia làm bánh.

Cậu tiến tới gần chiếc bánh, ngó qua những tấm giấy chúc được cắm xung quanh chiếc bánh. Những lời chúc viết tay với đủ kiểu chữ, nét nghiêng, nét tròn, nét cứng cáp hay nguệch ngoạc, mỗi tấm lại mang một sắc thái riêng biệt:

% Được tạo bởi GPT 4.1 từ GitHub Copilot, ngoại trừ Rushia.
\begin{itemize}
    \item \textbf{Sora}: ``Chúc Kuon-senpai luôn vui vẻ, mạnh khỏe và mãi là người anh tuyệt vời của đại gia đình chúng ta!''
    \item \textbf{Suisei}: ``Năm nay nhớ luyện Tetris chăm chỉ nhé! Chúc sinh nhật vui vẻ và luôn tỏa sáng như một ngôi sao!''
    \item \textbf{Miko}: ``Chúc mừng sinh nhật ne, Kuon-senpai! Đừng quên chơi cùng bọn em nhiều hơn nha, ne!''
    \item \textbf{Roboco}: ``Cảm ơn anh đã luôn sửa chữa mọi thứ cho em! Sinh nhật vui vẻ nhé!''
    \item \textbf{Fubuki}: ``Chúc anh luôn giữ được nụ cười và sự bình tĩnh trước mọi trò nghịch ngợm của bọn em!''
    \item \textbf{Matsuri}: ``Sinh nhật vui vẻ! Nhớ giữ gìn sức khỏe để anh có thể tiếp tục chịu được độ tăng động của bọn em!''
    \item \textbf{Aki}: ``Chúc anh luôn có thật nhiều năng lượng và cảm hứng để dẫn dắt mọi người!''
    \item \textbf{Haato}: ``Chúc Kuon-senpai không già đi mà chỉ trưởng thành hơn thôi nha-chama!''
    \item \textbf{Aqua}: ``Cảm ơn anh đã luôn giúp đỡ mọi người! Sinh nhật vui vẻ!''
    \item \textbf{Shion}: ``Chúc anh luôn gặp may mắn và không bị dính phải phép thuật lỗi của em!''
    \item \textbf{Ayame}: ``Chúc anh mạnh mẽ như một Oni và luôn vui vẻ bên đại gia đình này, yo!''
    \item \textbf{Choco}: ``Chúc anh sức khỏe dồi dào, ăn ngon ngủ ngon và không bị stress nhé!''
    \item \textbf{Subaru}: ``Chúc Kuon-senpai luôn là thủ lĩnh vững vàng của cả nhóm! Sinh nhật vui vẻ!''
    \item \textbf{Mio}: ``Cảm ơn anh đã luôn lắng nghe và dẫn dắt mọi người! Chúc sinh nhật thật nhiều niềm vui!''
    \item \textbf{Okayu}: ``Chúc anh luôn có những bữa ăn ngon và thật nhiều khoảnh khắc ấm áp, meo\~{}''
    \item \textbf{Korone}: ``Chúc Kuon-senpai luôn vui vẻ và không bị chó cắn nha, gâu gâu!''
    \item \textbf{Pekora}: ``Chúc anh luôn mạnh mẽ, không bị bẫy cà rốt của em và mãi là người quản lý tuyệt vời, peko!''
    \item \textbf{Rushia}: ``Kuon-senpai, cảm ơn anh đã trở thành chỗ dựa tinh thần cho em\dots\ Iu anh nhiều lắm-nanodesu!''
    \item \textbf{Flare}: ``Chúc anh luôn vững vàng, dịu dàng và là chỗ dựa cho tất cả mọi người!''
    \item \textbf{Noel}: ``Chúc anh luôn khỏe mạnh để còn ăn thịt cùng em thật nhiều nhé!''
    \item \textbf{Marine}: ``Chúc Kuon-senpai luôn trẻ trung, đẹp trai và không bị dính bẫy của hải tặc!''
    \item \textbf{Kanata}: ``Chúc anh luôn bay cao như thiên sứ và không bao giờ gục ngã!''
    \item \textbf{Watame}: ``Chúc anh luôn dịu dàng, ấm áp như cừu bông và không bị stress vì công việc!''
    \item \textbf{Towa}: ``Chúc anh luôn mạnh mẽ, quyết đoán và không bị ác ma nào quấy rối!''
    \item \textbf{Luna}: ``Chúc anh luôn ngọt ngào như kẹo và có thật nhiều niềm vui, nanora!''
    \item \textbf{Lamy}: ``Chúc anh luôn bình an, hạnh phúc và được mọi người yêu quý!''
    \item \textbf{Nene}: ``Chúc anh luôn tràn đầy năng lượng và không bao giờ cảm thấy cô đơn!''
    \item \textbf{Botan}: ``Chúc anh luôn vững vàng như sư tử trắng và thành công trên mọi mặt trận!''
    \item \textbf{Polka}: ``Chúc anh luôn vui vẻ, rực rỡ như rạp xiếc và không bao giờ hết bất ngờ!''
    \item \textbf{Ryuuhou}: ``Chúc Onii-san luôn hạnh phúc, bớt lười và mãi là anh trai tuyệt vời nhất!''
    \item \textbf{Ryujin}: ``Chúc con trai luôn mạnh khỏe, vững vàng và hạnh phúc bên gia đình!''
    \item \textbf{Kaien}: ``Chúc con luôn thành công, giữ vững bản lĩnh và sống trọn vẹn từng ngày!''
    \item \textbf{A-chan}: ``Chúc Hikawa-san mãi là người đồng nghiệp tuyệt vời, luôn được yêu thương và trân trọng!''
\end{itemize}

Cậu đọc từng lời chúc, cảm nhận được sự quan tâm, tình cảm và cả những nét cá tính riêng của từng người gửi gắm. Một nụ cười nhẹ nở trên môi, và trong lòng cậu dâng lên một cảm giác ấm áp chưa từng có.

Tuy nhiên, khi mọi người chuẩn bị chia bánh\dots

``Giờ ta mới để ý, ta quên mất phải mua thêm đồ đựng rồi\dots'' Okayu lên tiếng, tay chống cằm nhìn về phía chiếc bánh.

``Ừ ha! Giờ tôi mới để ý!'' Korone thảng thốt đáp. ``Nhóm mình ai là người phân công mua đĩa với dĩa nhỉ?''

``Chị không nhớ là chị có phân ai đó mua đâu,'' Sora ngẫm lại cuộc mua sắm hồi trưa hôm đó, rồi lắc đầu nhẹ. ``Chị tưởng các em đã tính đến rồi chứ.''

``Kuon-senpai, nhà anh còn bao nhiêu đĩa vậy, peko?'' Pekora quay sang hỏi cậu Tổng.

Kuon gãi đầu, bối rối: ``Anh nghĩ là không đủ cho toàn bộ ba mươi lăm người chúng ta đâu\dots\ Có khi chúng ta sẽ phải ăn theo kiểu lần lượt đấy\dots''

``Biết ngay là kiểu gì cũng thiếu mà\dots'' Fubuki thở dài, đôi tai cáo cụp xuống như biểu hiện nỗi buồn chân thật nhất của giống loài ấy. ``Xin lỗi Kuon-senpai, bọn em chưa tính đến điều đó\dots''

``Bây giờ chạy ra mua liệu còn kịp không nhỉ\dots?'' cậu Tổng lấy điện thoại ra, ánh mắt liếc về góc màn hình.

Chín giờ kém mười lăm. Đồng nghĩa với việc cửa hàng bách hóa gần nhà cậu đã đóng cửa.

``Có khi còn chẳng kịp ấy\dots'' Subaru nhăn trán. ``Siêu thị gần nhất chắc là vẫn còn mờ đấy chứ?''

``Ờm, không. Hôm nay là thứ bảy, nên chắc giờ này họ chuẩn bị đóng cửa rồi.'' Cậu quay về phía cô nàng Phù thủy nhỏ: ``Shion-chan, em có thể-''

``Nói trước, em cũng đã hết phép rồi đấy,'' Shion vội cắt ngang lời cậu, như càng đổ thêm dầu vào lửa. ``Cả ngày hôm nay em đã dùng hết cho việc làm bánh với sửa chữa rồi.''

Không khí đột ngột chững lại, trong khi chiếc bánh sinh nhật vẫn đứng đó, lộng lẫy và thách thức, như một nữ hoàng chưa tìm được ngai vàng của mình.

Cho đến khi\dots

``Hửm? Sao mọi người nhìn mình vậy?''

\dots tất cả mọi người đồng loạt quay đầu về phía Aki, người vừa thong thả từ tầng sáu xuống sau khi rửa bát và mới yên vị bên cạnh Roboco. Cô nàng vẫn còn đang lau tay bằng khăn bếp, gương mặt thoáng nét ngơ ngác như chưa hiểu tình hình.

Cô liếc quanh nhìn mọi người, rồi chạm ngay vào ánh mắt với sự thỉnh cầu bất đắc dĩ của Kuon.

``Xin lỗi Aki-chan nha, phải nhờ bọn em rồi\dots'' cậu thở dài, tay chống hông như thể đang phải gọi một người đi nghĩa vụ khẩn cấp.

``Có vẻ như bọn em lại phải nhờ đến cái `cơ chế' đó của chị nữa rồi\dots'' Subaru lắc đầu, nét mặt thể hiện rõ sự\dots\ cam chịu trước thế giới này.

``Ể?'' nàng Dị giới nghiêng đầu, khuôn mặt ngây thơ pha chút lo lắng. ``Cơ chế nào cơ\dots?''

\ct

Và thế là, trong ánh đèn dịu nhẹ, giữa gian phòng khách tạm thời biến thành phòng y tế dã chiến, cô nàng Dị giới giờ đang nằm ngửa trên lòng của nàng Người máy, hai chân dang rộng ra, dưới ``trung tâm xuất phát'' là một tấm khăn vải lót sẵn. Xung quanh là Subaru, Shion, và Kuon, mặt ai nấy căng như dây đàn, giống như đang đỡ đẻ thời chiến.

Roboco thì giữ chặt vai cô nàng tóc vàng như hộ sinh chuyên nghiệp, Subaru đứng cầm khăn và quạt tay cho ``bà đẻ'', còn Kuon thì\dots\ đanggồng tay như thể sắp bắt được món gì quý giá.

``\vie{Cảnh này\dots} Ryujin nhíu mày, mồ hôi chảy ròng ròng. Ông quay sang vợ, hỏi: ``\vie{Em\dots\ em có chắc đây là idol không?}''

Kaien thở dài: ``\vie{Chắc chứ. Em từng thấy cảnh này một lần rồi mà. Hồi Tết năm ngoái, nó rặn ra nguyên cái chậu nhựa\footnote{Xem lại {\sen Partie 2}, quyển T.} đấy.}

``\vie{Hả!?}'' ông thảng thốt. ``\vie{Rặn ra cái gì cơ!?}''

``\vie{Cái chậu giặt quần áo. Loại chậu I(no)mata chín lít rưỡi. Mới nguyên tem, lại còn bền nữa.''}

``M-Mọi người đang nói gì vậy\dots?'' Aki bất ngờ kêu lên, rõ ràng đã nghe thấy phần nói chuyện, nhưng giả vờ không để tâm. ``Đừng làm mình mất tập trung-''

Bất chợt, cô nàng rên lên. ``Aaa\dots\ \textbf{Aaa!}'' Toàn thân cô gồng cứng lại. Ngay lập tức, Ryuuhou - người vừa tò mò đến xem thử - lập tức lảng đi chỗ khác, bịt tai mình lại: ``Thôi, thôi nhá!''

Một tia sáng lóe lên từ giữa hai đùi cô, và *bộp!*, một cái dĩa bằng sứ trắng xuất hiện trên khăn.

``\textbf{Rồi, được 4 cái đĩa với cái dĩa rồi!}'' Subaru hét lên, giơ tay ăn mừng như một huấn luyện viên đang reo hò khi đội tuyển chiến thắng vòng loại giải đấu thể thao điện tử.

``Cố gắng lên nào, Aki-senpai!'' Shion tiếp tục cổ vũ.

``\textbf{Aaa-! A\dots\ aaaaah!}'' nàng Tóc bay tiếp tục quá trình ``sinh sản'' trong tiếng hô hào của mọi người.

Từ phía xa, Rushia và Flare, người chắp tay cầu nguyện, người nhăn mặt vì không hiểu chuyện gì đang xảy ra.

``Liệu Aki-senpai có ổn không đây\dots?'' nàng Pháp sư lẩm bẩm cầu nguyện.

``Cái quái gì vậy trời\dots'' nàng Dị giới toát mồ hôi hột vì vừa đang nín thở theo dõi, vừa khó hiểu trước những gì đang diễn ra trước mắt.

Trong những tiếng reo hò của nàng Tomboy và nàng Phù thủy và tiếng rên la đau đớn của Aki, Ryujin tiến lại gần con trai, vẫn chưa hoàn hồn: ``\vie{Cái gì vậy con\dots?}''

``\vie{À\dots\ đây là cách mà Aki-chan tạo ra đồ vật đấy ạ\dots}'' Kuon đáp lại, không biết phải giải thích cho bố cậu như thế nào.

``Sao trông giống\dots\ đi đỡ đẻ thế?''

``Con cũng nghĩ y chang bố khi lần đầu chuyện này\footnote{Xem lại \flaregothic{ÉPISODE 10}, quyển số Một.} xảy ra ạ. Đây là lần thứ ba rồi.''

``\vie{Lần ba!? Hai lần kia là thế nào?}''

``\vie{Một lần là\dots\ dòng nước đen ngòm. Một lần là\dots}'' cậu chỉ tay về phía Subaru, người đang đông viên qua chiếc loa cầm tay luôn giắt trên hông mình, ``\vie{\dots là con quạc quạc kia đấy ạ.}''

``Giờ tôi mới biết là Aki-chan cũng đã từng trổ tài làm bài này một lần nữa rồi đấy,'' Haato chống hông, nhìn đồng đội cô với vẻ mặt e ngại.

Tới đây, Ryujin không nói nên lời. Cuối cùng, ông chỉ thốt lên: ``\vie{Sao lại có cái cách tạo củ chuối thế\dots}''

``\vie{Con biết là nói thế bố sẽ không tin, nên con đã mang cả tờ hướng dẫn tới rồi đây.}''

``\vie{Có cả tờ hướng dẫn à?}''

``\vie{Vâng.}'' Cậu rút ra một tờ giấy gấp tám từ túi áo, giở ra và chỉ vào một phần trên bản phác thảo: ``Đây này, khu vực `trọng yếu' ở bên trong đùi là nơi có thể sản xuất \emph{mọi thứ}, ít nhất là những gì Fubu-chan nói vậy.''

\img[15cm]{1-6b}

``\vie{Thật là\dots\ dị hợm. Bố ngạc nhiên là thậm chí còn có một người như vậy đấy\dots}'' ông chốt lại.

``\vie{Vì \hll\ là như thế đấy ạ,}'' cậu quản lý trẻ lặng lẽ kết luận.

Ngay lúc đó, Aki gào lên một tiếng cuối cùng như vượt cạn thành công và một loạt đĩa, dĩa đủ cỡ xếp chồng lên nhau trên khăn, chia đủ cho ba mươi lăm người. Mọi người vỗ tay như mừng sinh nhật của một người mới chào đời.

``Cảm ơn chị nhiều, Aki-senpai!'' cả nhóm vây quanh dìu cô đứng dậy, để cô thở hổn hển, mắt mờ lệ, như vừa trải qua một hành trình mang nặng đẻ đau của nền văn minh thần tượng đây kỳ quặc.

``\vie{\dots Nó có dùng được không vậy?}'' Ryujin tiếp tục nhìn đống đĩa và dĩa mới toanh, không một tì vết ấy.

``\vie{Anh có tin là nó có thể dùng tốt không?}'' Kaien lên tiếng.

``\vie{Cái chậu đó thì nói làm gì!}'' người chồng hơi cao giọng, vặn lại vợ.

``\vie{Không tin thì thôi.}'' Bà sau đó bảo cậu con trai chạy lên lấy một chiếc đĩa và dĩa ở trên bếp tầng sáu, mặc dù Ryujin từ chối.

Và thế là, họ đã có đủ đồ đựng cho chiếc bánh khổng lồ này.

Có đĩa, có dĩa rồi, mọi người gần như đã sẵn sàng thưởng thức chiếc bánh được cứu vớt từ sự thiếu thốn lúc nãy. Nhưng rồi Subaru khựng lại giữa không trung khi chuẩn bị cắt bánh.

``\dots Mà, gượm đã, dao để ăn cắt bánh đâu?''

Cả đám khựng lại. Một lần nữa, không một ai có thể ngờ tới trường hợp này.

``\dots Ừ nhỉ,'' Aqua nhíu mày. ``Mấy con dao hình như đã mang lên bếp hết rồi.''

``Không ai mang dao xuống sao?'' Roboco hỏi, nhìn quanh.

``Lúc nãy bọn em lo rửa bát, ai để ý đâu\dots'' Fubuki gãi đầu.

``Đáng ra bà phải nhắc từ trước chứ!'' Mio nhíu mày. ``Biết cắt bánh rồi thì ít nhất cũng phải mang dao xuống chứ!''

``Giờ quay lại bếp lấy thì lại mất công quá\dots'' Haato thở dài. ``Ai lên giờ này nữa\dots''

Không cần ai nói ra, tất cả đồng loạt nhìn về phía Aki, một lần nữa, người vẫn còn đang thở dốc, tóc rối bời, ôm một cái khăn mặt và quấn chiếc chăn mỏng quanh người như vừa thoát khỏi một nghi lễ ``đẻ đau'' cổ đại.

``L-Lần này lại gì nữa\dots?'' cô nàng hoang mang khi thấy ánh nhìn tập thể nhắm thẳng vào mình.

``Bọn tôi lại phải nhờ bà lần nữa rồi, AkiRose\~{}'' Matsuri nhảy tới, tay chống hông, cười toe.

Lập tức, nàng Tóc bay kêu lên, lùi sát vào ghế, mắt ngân ngấn nước: ``Không! Không! Không phải lần nữa chứ! Mình không chịu nổi nữa đâu!''

Ngay khi các cô gái định sồn tới để thúc ép Aki một lần nữa, Kuon liền bước tới, giơ tay ngăn đám đông đang bắt đầu tụ tập như định làm nghi lễ lần hai: ``Thôi, không cần đâu. Để anh chạy lên bếp con dao xuống.''

Sora khoanh tay, nhíu mày nhìn mọi người: ``Đừng có trêu Aki-chan nữa mọi người. Em ấy đã không muốn rồi thì đừng có ép chứ.''

``\dots Xin lỗi mà, Aki-senpai,'' Fubuki lè lưỡi.

``Hehe, tôi chỉ đùa tí thôi\~{}'' Matsuri cười trừ, gãi má.

``Em\dots\ em cũng không nghĩ là chị khóc thật\dots'' Shion nhỏ giọng nói, mắt nhìn sang hướng khác.

Aki không đáp, chỉ ôm gối ngồi trên ghế bành, gò má phồng lên như cái bánh bao hấp, mắt liếc ngang liếc dọc với vẻ dỗi hờn. Roboco vỗ nhẹ lưng cô an ủi, trong khi Subaru mang tới một cốc nước ấm, như một y tá tận tâm chăm sóc sản phụ.

\ct

Sau một hồi dỗi nhẹ và được dỗ ngọt bằng đủ loại nước uống và lời xin lỗi, Aki cũng chịu bỏ gối xuống và tập hợp với mọi người. Lúc này, Kuon từ tầng sáu đã trở lại với một con dao lớn trong tay, loại chuyên dùng để cắt bánh cưới, trông còn sáng loáng như mới mài.

``Được rồi, ai là người sẽ được vinh dự để cắt chiếc bánh khổng lồ này đây?'' cậu đưa con dao lên cao, hỏi mọi người.

``Em, em! Em muốn làm!'' Suisei giơ tay lên, đôi mắt sáng rực như học sinh tiểu học giành quyền trả bài. ``Em muốn làm việc này từ lâu rồi!''

Cậu Tổng đưa con dao cho cô, nhoẻn miệng trêu: ``Đừng dùng nó để đi xiên người ta đấy.''

Đáp lại cậu chỉ là một cái phồng má nhẹ. Cô xì khói, hơi gằn: ``Anh nói như thể em là sát nhân ấy, Kuon-senpai!''

Cô cầm chiếc dao, đi về phía chiếc bánh cao tới gần chạm trần kia. Sau một hồi ngó quanh để tìm cách cắt, cô gãi đầu gãi tai, rồi quay sang hỏi mọi người: ``Cắt thế nào đây mọi người?''

Choco chỉ tay lên tầng cao nhất, giọng vẫn giữ vẻ từ tốn như đang dạy trong lớp học: ``Cắt từ trên xuống là dễ nhất đó, xong rồi ta cứ từ từ cắt dần xuống thôi. Chúng ta đã chia sẵn thành các phần từ các vị rồi đấy.''

``Ừ nhỉ.'' Suisei gật gù, quay qua phó thác cho một người: ``Noel-chan, giúp một tay được không?''

``Cứ để đó cho Danchou!'' nàng Hiệp sĩ chạy lên trên tầng, rồi nhanh chóng vác cả một cái thang gấp xuống từ tầng sáu, kê ngay sát chiếc bánh, giữ chân thang cẩn thận.

Nàng Ca sĩ leo lên, hai tay cầm dao, mắt chăm chú nhìn chiếc bánh như đang giải toán. Bốn vị tầng trên cùng (bánh quy, caramel, đậu đỏ, cà phê) đang đợi phán quyết từ thiên hà nữ vương.

Dưới đất thì bắt đầu náo loạn, người nào người nấy đếu đang vọng lên những tiếng nói xin phần bánh.

``Cho en phần caramel nha!'' Fubuki hét lên, vẫy tay.

``Không! Em gọi trước rồi mà!'' Lamy chen vào.

``Bánh quy! Ai lấy bánh quy của tôi là không yên đâu đó!'' Nene giơ nắm đấm cảnh cáo.

``Choco-sensei\~{} chia phần em caramel với cà phê được không ạ?'' Mel nhỏ giọng năn nỉ.

``Được chứ, Mel-sama,'' Choco cười, vuốt tóc cô. ``Vẫn còn nhiều vị khác mà, mọi người cứ bình tĩnh thử nhé!''

``Tôi xí đậu đỏ! Tôi thích đậu đỏ nhất á, ne!'' Miko giơ hai tay, làm dấu chữ thập.

``Em tưởng đậu đỏ để ăn sáng mà?'' Aqua ngơ ngác.

Và, cứ như vậy. Hàng loạt những lời gọi xin phần cứ thế tuôn ra như suối, trong khi Suisei ỏ trên vẫn đang căn chỉnh hướng dao để cắt, nhưng vẫn không quên vọng lại: ``Được rồi, bình tĩnh nào! Ai cũng sẽ có phần!''

Nhưng rồi, bất chợt\dots *CẠCH!*

Chiếc thang bỗng chao nhẹ. Một bên chân đã bị tuột vít.

``Khoan\dots!''

Nhưng, mọi thứ diễn ra sau đó xảy ra quá nhanh tới mức không để bất cứ ai phản ứng. Chiếc thang dần đổ nhào xuống\dots

\textbf{*RẦM!*}

Suisei không kịp giữ thăng bằng, rơi thẳng xuống như thiên thạch rớt khỏi quỹ đạo, rồi đáp cái *phịch!* xuống tầng lớn nhất của chiếc bánh, ngay góc dâu tây đỏ rực.

*ỤP!*

Cả phần lưng của Suisei bị lem nhem phần kem hồng, một phần nhỏ của chiếc bánh bị móp một cách thảm thương.

``Trời ơi, bẩn mất bộ cánh của mình rồi\dots'' Suisei thốt lên tỏ vẻ hơi tiếc của. ``Mất công giặt giũ hôm qua\dots''

Nhưng, sự hỗn loạn vẫn chưa kết thúc.

Ngay khi nàng Ca sĩ bị ngã, con dao trên tay cô cũng bị văng ra khỏi tay, lộn vài vòng trên không rồi\dots

*KENG!*

\dots cắm thẳng ngay ở dưới sàn, gần chân phải của Watame.

``\textbf{Beeee!!}'' Watame kêu lên thất thanh, mắt trợn trắng, người run cầm cập ngay khi cô suýt bị con dao đâm vào, trông như sắp ngất tới nơi.

``W-Watame-senpai!?'' Polka nhào tới, đỡ cô nàng Cừu đang lẩy bẩy kia đang chuẩn bị chạm đất: ``Chị ổn chứ!?''

``C-Chị không sao\dots'' cô nàng run bần bật, tay bám víu vào Botan. ``Suýt tí nữa thì tè ra quần rồi\dots''

``Em thấy chị đang xả lũ luôn rồi ấy chứ,'' nàng Sư tử trắng khẽ bật cười đáp lại.

``Thế á!?'' nàng Cừu hốt hoảng, vội đứng dậy và quay sau lưng lật váy mình lên. Quần lót vàng của cô đã bắt đầu ướt đôi chút, thậm chí đã bắt đầu rỉ nước.

Cô nàng chậm rãi quay lại, hai tay sà xuống vội nắm lấy phần hạ bộ của mình, rồi lí nhí nói với mọi người: ``C-Chị xin phép\dots'' và ngay lập tức lao vào nhà tắm.

Một số thành viên, bao gồm cả Kuon và Ryuuhou, đều nghe lỏm được đoạn hội thoại này, cũng đã bụm miệng cười trước tai nạn khó đỡ nhưng nguy hiểm vừa rồi.

Botan nhìn tiền bối cô hớt hải chạy vào toa-lét, chống hông cười tươi nói mỉa: ``Em tưởng là Suisei-senpai không đi xiên người mà?''

``Cũng may là không có ai bị thương đấy,'' cô em gái của cậu thở phào. ``Cái thang đó dù gì cũng cũ rồi mà.''

``Thiệt hại duy nhất ta có là quần chíp của nàng Cừu thôi,'' cậu giở giọng đùa, ``nếu không vì ai đó đấy.''

``Cái đó hoàn toàn là tai nạn mà!'' Suisei minh oan bản thân.

``Anh biết, anh biết chứ. Cơ mà\dots'' cậu nhìn lại phần bánh dâu đã vô tình bị phá hỏng kia, rồi chưng hửng hỏi: ``Mấy đứa định xử lý thế nào đây?''

Korone và Okayu vội nhào vào, lấy kem dự phòng và thìa trét lại phần móp, rồi tuyên bố như một vị giám khảo: ``Ổn rồi! Bánh móp vẫn ăn được! Kuon-senpai ăn phần này nha!''

Mio, nhìn hiện trường, chỉ thở dài rồi tuyên bố như một vị giám khảo: ``Có lẽ là vậy thôi, Kuon-senpai. Nếu anh không ăn thì để em ăn cũng được.''

ất cả đều quay sang nhìn cậu trai duy nhất trong phòng. Kuon cười trừ, nhìn phần bánh bị sụp, hơi nghiêng sang một bên như một công trình đang thi công dở: ``Anh đoán là chẳng ai muốn ăn phần này ngoài anh cả.''

``Vậy thì Kuon-senpai ăn phần đó nha!'' Kanata tuyên bố, rồi chắp tay trước mặt cậu: ``Cảm ơn vì đã giúp!''

``Ai bảo là nam chính thì phải gánh hết!'' Fubuki vỗ vai cậu, không hề giấu nụ cười tinh quái.

``Chuẩn cơm mẹ nấu luôn.''

\ct

Chiếc thang gấp đã gãy một bên chân, nằm chỏng chơ giữa sàn như một chiến binh già hết thời. Mọi người nhìn nó, rồi quay sang nhìn chiếc bánh cao chót vót vẫn còn năm tầng chờ đợi.

``Giờ ai leo lên cắt tiếp đây?'' Sora hỏi, ánh mắt quét một vòng, nhưng không ai dám giơ tay sau cú ngã vừa rồi.

Suisei rũ áo, phủi kem trên mặt và khẽ lắc đầu: ``Tôi xin kiếu. Một lần rơi từ trên trời xuống là đủ rồi\dots''

``Vậy để mình phụ giúp nha\~{}'' Roboco bước ra, tay chỉ vào lưng mình. ``Sui-chan, trèo lên đi. Tôi có thể bay lên và giúp bà được đấy.''

``Ể\dots? Liệu có ổn không vậy?''

``Đảm bảo là ổn luôn! Dù gì thì tôi vốn dĩ là rô-bốt có thể bê vác được chục cân đồ mà.''

``Hy vọng là lại không làm gì đón đại PON là được\dots'' Kuon lầm bầm, tỏ vẻ không tin tưởng lắm tới cô nàng Người máy.

Dĩ nhiên, câu nói đó khiến cho cô không vui. ``Kuon-senpai, mấy cái câu PON của anh bắt đầu nhàm rồi đó\~{}!''

``Vậy thì để em với Ogayu đỡ bánh nha!'' Korone hớn hở chen chân.

``Đồng quan điểm đó, Koro-san\~{}''

Một lần nữa, nàng Ca sĩ lại leo lên đỉnh, lần này là nhờ lưng của đồng đội Rô-bốt làm bệ phóng. Bên dưới nàng Chó và nàng Mèo huẩn bị đĩa, dĩa và khay bánh. Cặp đôi bạn thân này phân công nhau rõ ràng: Korone nhận bánh từ Suisei, đưa cho Okayu, rồi Okayu trao tận tay các thành viên đang ngồi thành vòng tròn quanh bàn.

Bắt đầu từ mẻ trên tầng lớn nhất, gồm bốn vị: bánh quy, caramel, đậu đỏ, và cà phê.

Ayame ăn phải những mẩu bánh quy, cười tủm tỉm. ``Cứng cứng, giòn giòn, y như tâm trạng khi ngủ dậy bị gọi livestream.''

Ryuuhou nhíu mày: ``Đó lại là cách\dots\ ví von lạ nhất em từng thấy đấy.''

Kuon cắn miếng cà phê, trầm ngâm: ``Đắng vừa, thơm. Vị cà phê mạnh như bầu không khí buổi sáng.''

Watame - người vừa quay lại sau vụ ``tràn đê'' ban nãy - thử caramel rồi mỉm cười nhỏ nhẹ: ``Thơm thật đấy\dots\ lại còn ngọt nữa chứ. Làm em liên tưởng đến mấy viên kẹo bọn trẻ con hay ăn.''

Xuống đến tầng tiếp theo, bao gồm hạt dẻ, dừa và đào.

Lamy ăn hạt dẻ, mỉm cười mơ màng: ``Tuyệt vời. Vị như mùa thu trên núi.''

Kanata thử miếng dừa rồi nhăn mặt, nhằn phải mấy sợi dừa còn lấp ló bên trong, rồi giương ra bộ mặt khó tả: ``Ơ\dots\ sao giống ăn bánh xà phòng vậy ta?''

``Em nói xấu bánh của bọn chị hả?'' Suisei từ trên lưng Roboco la xuống, giơ con dao lên cao.

``Không! Không! Chỉ là em chưa quen thôi mà!'' Kanata giơ hai tay xin tha.

Tiếp đến là tầng bao gồm ba vị: dưa gang, sơ-ri và phô mai.

Marine ăn sơ-ri, rồi rít lên vì sung sướng: ``Trời ơi\~{} ngọt như nụ hôn đầu luôn á\~{}!'' Cô chu miệng về phía Rushia, người đang nếm thử một vị khác. ``Rushia, cho tôi hun đi\~{}''

``Kinh quá!'' nàng Pháp sư hốt hoảng, hướng miếng bánh đang ăn dở lên cao. ``Tránh xa tôi ra!''

Cặp đôi Chó-Mèo cũng có phần riêng cho mình. Nàng Mèo bất chợt hỏi đùa: ``Ể\~{}? Marine-Senchou hôn từ khi nào vậy?''

``Chuyện của người lớn! Vô duyên!'' nàng Thuyền trưởng phát cáu.

``Người ta chỉ hỏi thôi mà\~{}''

Một lần nữa, cậu Tổng nếm thử vị phô mai, rồi khẽ gật gù: ``Phô mai làm béo, mịn màng\dots\ mỗi tội hơi ngấy. Mấy đứa hình như cho hơi nhiều phô mai nhỉ?''

``Em nghĩ là thế\dots?'' Aqua nghiêng đầu, rồi xúc một miếng ăn thử. Sau một hồi nghiền ngẫm, toàn thân cô khẽ run lên, dường như khẳng định với suy nghĩ của mình.

Kế đến tầng dưới là các vị bạc hà, mâm xôi và nho.

Sora thử miếng bánh mâm xôi, rồi khẽ giật mình vì vị đăng đắng. ``Hơi gắt, nhưng thơm. Đúng kiểu vị của những bài ballad buồn.''

``Sora-senpai thực sự cảm nhận được như thế ư?'' Towa liếc về phía cô nàng đang rơi nước mắt kia.

Korone thì chớp lấy phần bánh vị bạc hà. ``Úi chà! Hơi cay cay, nhưng thơm mồm phết!''

Ayame thì chọn bánh nho. ``Hmm\dots\ nhẹ, dễ ăn. Nhưng hơi giống món bữa sáng hay ăn\dots\ Thêm một chút sữa nữa là vừa ý, yo.''

Tiến tới tầng áp chót là một bữa tiệc bốn vị: cam, dứa, kiwi và táo.

Haato ngậm miếng kiwi, rồi la lớn tới mức suýt phun bánh ra ngoài: ``Ặc! Chua dữ thần luôn á!''

Marine ngược lại thì rất mê mẩn với vị dứa: ``Ngon ghê ha~{} Vị ngọt dịu, lại còn dậy mùi nữa. Có cảm giác như ăn dứa giữa biển Nam Thái Bình Dương vậy\~{}''

``Tôi cũng không ngờ tới vị cam cũng khá hợp lý cho món bánh kem đấy,'' Botan từ tốn ăn miếng bánh của loại trái cây này.

``Cơ mà vị táo này hơi chua đấy,'' Flare xắt ra một miếng kem và nếm thử. ``Thêm một chút đường nữa là ổn.''

Cuối cùng, tầng dưới cùng là đại hội bốn hương vị cổ điển và phổ biến nhất: dâu, sô-cô-la, trà xanh, vani, ngô.

Kuon nhận một miếng dâu - cũng là miếng mà Suisei phá hỏng trước đó - và một miếng trà xanh. Cậu nhắm mắt nhai, rồi gật đầu nhẹ, chậm rãi thưởng thức: ``Dâu vẫn ngon nhất. Vị trà xanh khá đậm, giống kiểu matcha của Kyoto.''

Botan chọn sô-cô-la, rồi khẽ rướn mày với hương vị mà cô cảm nhận được: ``Hm. Đậm vị thật. Nhưng\dots\ hơi ngọt hơn tui tưởng.''

Lamy ăn miếng bánh ngô, mắt sáng lên: ``Ngon lắm đó! Bánh ngô mềm như tan ra, thơm như bắp nướng ở Hokkaido!''

Watame thử miếng vani. Cô nheo mắt: ``Hơi đơn điệu\dots\ nhưng cũng dễ chịu. Giống như vị của một buổi trưa mùa hè.'' Mặt cô giãn ra, vui vẻ thưởng thức với hương vị này.

Trong bữa tiệc hương vị bánh đó, Miko lại\dots\ không ăn ngay.

Cô len lén nhìn phần bánh dâu đã bị móp ban nãy, phần mà Suisei đã ngã tại đó. Kem đã được trét lại nhưng vẫn nhìn ra dấu vết của cú tiếp đất.

``Miko muốn ăn phần đó,'' cô chỉ tay về phần còn lại của miếng móp kia.

Lập tức, nàng Ca sĩ bất chợt đỏ mặt. Cô lắp bắp, định ngăn cô bạn của mình: ``N-Nhưng, bà biết đó là phần bị\dots''

``Chính vì đó là phần bà bị rơi xuống đó nên nó còn có thêm cả hương vị của Sui-chan nữa đó, ne!'' nàng Vu nữ lắc lắc người, giọng ngân như cầu nguyện.

``C-Cái gì cơ!?''

``Như thể đây Kami-sama đã ban phước cho tôi vậy, ne\~{} Có Suisei rơi trúng rồi mà còn không thành vật phẩm thánh thần rồi thì còn là gì nữa, ne\~{}''

Mọi người ồ lên, cười nghiêng ngả. Suisei mặt đỏ như trái cà chua, còn Miko thì đã cắn một miếng bánh được ``ban tặng'', mắt long lanh như đang ăn món linh thiêng nhất vũ trụ.

\ct

Chín giờ năm phút.

Trong lúc những lát bánh tiếp tục được phục vụ cho mọi người, một nhóm các thành viên khác (những người chưa động vào bánh) đã nhanh chóng bắt tay chuẩn bị cho hoạt động tiếp theo: hát karaoke.

Theo kế hoạch, họ muốn có một buổi hát thi đấu và dự đoán kéo dài lên đến một giờ, nhưng vì thời gian hiện tại cộng với hoàn cảnh không cho phép, họ đã cắt giảm xuống bốn mươi phút và hát theo dạng karaoke gia đình.

Chiếc ti-vi cỡ lớn cùng bộ dàn âm thanh cồng kềnh được Noel vác ``như không'' từ phòng của Kuon xuống phòng khách, đi kèm là một bộ loa sân khấu và bộ âm-li (amplifier) mà chỉ nhìn thôi cũng biết không dành cho người tay mơ.

Trong khi Mio và A-chan xử lý phần lắp đặt, Ryujin cùng Kaien, hai vợ chồng yêu thích âm nhạc của nhà Hikawa, bắt đầu chỉnh các thiết lập âm thanh, kết nối dàn mixer, micro không dây và hệ thống ánh sáng mini mà Lamy đã mang đến từ trước.

``\vie{Muốn chơi karaoke hay,}'' người chồng vừa nói vừa quay các núm chỉnh EQ, ``\vie{phải cắt dải trung đúng chỗ và đẩy bass nhẹ thôi. Không là giọng ai cũng thành tiếng vịt.}''

``\vie{Tăng hiệu ứng phòng vừa thôi,}'' người vợ tiếp lời, ``\vie{vang quá là thành hát trong nhà tắm đấy.}''

Để thử âm thanh, họ mở lần lượt các bài nhạc kinh điển từ thập niên 80, 90 và đầu 2000, từ những bài hát nước ngoài nổi tiếng, cho đến những ca khúc có tiếng tăm trong xứ Etsunan và Nhật Bản vào thời điểm đó.

Mỗi lần điệp khúc vang lên, một vài thành viên ngồi gần đó lại khẽ lắc lư, ngân nga theo những câu hát đã gắn bó với thời tho ấu của họ.

``Cô chú đúng là có gu nghệ thuật thật đấy,'' Sora vừa đung đưa trước giai điệu vừa tấm tắc khen.

``Cô ngạc nhiên là cháu cũng có sở thích âm nhạc này đấy,'' Kaien mỉm cười đáp lại. ``Cháu thích nghe những bài hát từ những thập kỷ trước sao?''

``Dạ! Cháu với AZKi-chan\footnote{Sẽ được giới thiệu ỏ quyển số Mười Một.}'' thích nghe nhạc mà. Nhiều khi cậu ấy hát những bài này còn nhiệt tình hơn cả những bài idol pop đấy chứ!''

Vợ chồng nhà Hikawa khẽ bật cười, tỏ vẻ hứng thú với câu chuyện này. Ngay khi họ đang chuẩn bị nói tiếp, một giọng nói từ một người đã đứng cạnh từ nãy tới giờ cất lên.

Từ Miko-chi. ``Miko cũng muốn thử! Cô chú cho Miko thử âm thanh với, ne\~{}'' cô nàng hớn hở, đôi mắt sáng như đèn pha ô tô.

``Được thôi,'' người thợ cơ khí hiền từ đáp, đưa chiếc điều khiển (mà thực chất chỉ là một ứng dụng xem video trên điện thoại) của ông cho nàng Vu nữ.

Ngay khi cô nàng giành được quyền điều khiển, Miko lập tức bấm bừa một thước phim ngay trên màn hình. Một giây sau\dots

\href{https://www.youtube.com/watch?v=I7_tBbluNZ4&t=3270s}{``\textbf{Ái cha, ái cha, ái cha, ái cha!!! AAAAAAaaaaaAAAAaa!! Cứu tôi với!!}''}

Một đoạn gào thét vang trời bất ngờ phát ra từ dàn loa chính, khiến mọi người giật mình, vài người còn suýt đánh rơi cả đĩa bánh.

``C-Cái gì vậy?'' Ryuuhou và Lamy đồng thanh, ngước nhìn về phía cô nàng Vu nữ, người cũng vừa giật mình khi nghe thấy tiếng hét thất thanh đó.

``Nóng\dots\ nóng\dots\ nóng quá!'' giọng nói đậm chất nhịu của nàng Vu nữ tiếp tục vang lên, kế sau đó là tiếng ngựa hí.

Tiếng gào ấy không lẫn vào đâu được. Không gì khác, đó chính là tiếng hét ``huyền thoại'' của Miko khi cô cho dung nham vào lu rồi tự thiêu chính mình trong M*necr*ft, đoạn clip đã trở thành trò đùa sống trong cộng đồng người hâm mộ từ năm nào.

Miko chết sững. Mặt cô đỏ ửng trong khi một số thành viên bật cười nghiêng ngả, đặc biệt là Suisei và Pekora.

``Ôi trời\dots'' nàng Ca sĩ nhoẻn cười, ``tôi còn tưởng Miko-chi là người điều khiển mic, ai ngờ lại là dung nham điều khiển bà cơ chứ!''

Nàng Thỏ thì cười nghiêng ngả đến bật ngửa ra sau: ``Em không thể nghĩ là chị PON đến mức đó đấy, Miko-chi ạ!''

``Không phải là hai người chứ!!'' Miko mặt đỏ phừng, trông như sắp phóng bộ điều khiển của họ tới nơi, nhưng Suisei chỉ nhún vai, vờ như vô tội.

``Cái đó gọi là quả báo đấy, peko. Sau cái vụ chị cười đùa với cái cú ngã của Sui-chan ban nãy mà,'' Pekora lau nước mắt, rồi quay sang Suisei đập tay.

Miko chu môi, quay mặt giận dỗi, nhưng không giấu được cái vành tai vẫn đỏ bừng vì xấu hổ. Cô đành lùi ra ghế, tay ôm gối, lầm bầm không rõ lời, để mặc cho tiếng hét và tiếng nói tiếp tục vang lên rộn rã khắp căn phòng.

\ct

Dàn loa đã được chỉnh xong. Sau một thoáng lặng ngắt và điều chỉnh ngắn ngủi, chương trình karaoke - hoạt động chính đầu tiên của bữa tiệc sinh nhật bắt đầu.

Sora là người mở màn, không ai khác xứng đáng hơn. Giọng hát ngọt ngào, dịu dàng như gió xuân của cô khiến cả phòng khách đột nhiên trầm lắng hẳn đi. Mọi người tự giác đặt bánh và cốc nước xuống, ngẩng đầu lắng nghe. Đôi mắt Sora khẽ nhắm lại khi ngân dài nốt cao, và khi kết thúc, tiếng vỗ tay vang lên đều khắp cả gian nhà.

Tiếp theo là Nene, đầy năng lượng và rực rỡ như thường lệ. Cô chạy quanh căn phòng khi hát, nhảy tưng tưng trên ghế, thậm chí còn lôi cả Aki đang ngồi gối ôm ra nhảy theo. Không ai hiểu sao cô vẫn hát được dù nhịp thở rõ ràng như người vừa chạy ba vòng sân vận động. Nhưng cô vẫn hát tốt, thậm chí là rất tốt.

Kanata thì chọn bài hát chậm hơn, lặng lẽ đứng một góc, mic cầm bằng hai tay. Giọng của cô vang lên khàn khàn, nhưng lại có sức nặng đến kỳ lạ. Mỗi câu hát của cô như chạm vào ngực người nghe, tạo thành một cảm xúc lắng đọng khó gọi tên. Không ai nói gì khi cô hát xong, chỉ vỗ tay thật nhẹ, như sợ phá vỡ dư âm đang đọng lại trong không gian.

Và cứ thế, một vài thành viên khác cũng xung phong lên trình diễn bài hát sở trường của mình.

Giờ đã gần chín giờ năm mươi.

Ryujin nhìn đồng hồ, rồi quay sang các cô gái thông báo: ``Chúng ta có lẽ sẽ kết thúc sớm tiết mục này thôi, không lại ảnh hưởng đến các nhà xung quanh. Ai muốn là người hát cuối cùng nào?''

Như thể đã đồng thuận nhau từ trước, cả phòng đồng loạt hướng ánh mắt về một người: nhân vật chính của bữa tiệc sinh nhật này.

\begin{dialogue}
    \speak{Polka} Kuon-senpai!!
    \speak{Rushia} Phải là anh ấy chứ còn ai nữa-nanodesu!
    \speak{Towa} Sinh nhật mà không hát một bài thì không được!
    \speak{Haato} Đừng có mà trốn, Kuon-senpai!
\end{dialogue}

Kuon nhắm mắt, thở dài. Đúng như cậu dự đoán, kiểu gì cũng sẽ bị họ réo tên. Cậu khoát tay từ chối, cười trong sự khổ sở: ``Thôi\dots\ Để người khác hát đi. Anh thì hát dở lắm. Bây giờ mà cất giọng là bao nhiêu không khí xúc động bồi hồi là đi tong hết.''

Tiếng cười rộ lên vài chỗ, nhưng ngay lúc ấy, một âm thanh khác đột ngột vang lên từ dàn loa, bản nhạc nền huyền thoại trong một bộ phim Titanic, đi kèm là tiếng hai người đàn ông hát\dots\ cũng không đến nỗi tệ.

``\eng{Every night in my dreams\~{} I see you\~{} I feeeeeel you\~{}}''

Một giọng hát cực trầm trong hai giọng ca đó, các cô gái luận ra ngay đó là Ryujin, người cũng đang mỉm cười nghe lại.

Giọng hát còn lại, có phần cao hơn, nhưng dường như lại tụt hơi khá nhanh, đó lại là giọng ca ``dở tệ'' của chính cậu Tổng.

Mọi ánh mắt đổ dồn về phía Kaien, người đang cầm điện thoại với gương mặt thản nhiên như không, bật đoạn ghi âm từ lần hai bố con Hikawa cùng hát hò trong một buổi tối rảnh rỗi.

``Ai bảo là con không hát được chứ?'' bà bình tĩnh nói, tay vẫn không dừng phát đoạn tiếp theo.

Kuon nhăn mặt, thở dài ngán ngẩm: ``Mẹ ơi, con đã không muốn hát rồi lại còn thúc người ta làm gì chứ\dots''

Tiếng cười lại bùng lên, kèm theo đó là những lời nói châm chọc. Kuon lập tức quay mặt đi, tay giơ lên đầu như muốn đầu hàng, đỏ mặt thấy rõ.

``Kuon-senpai thấy ngại nên là không hát chứ gì\~{}'' Marine cười khà khà, uống một ngụm nước ép.

``Anh không cần phải thấy xấu hổ đâu, Kuon-senpai. Tự tin lên chứ!'' Fubuki động viên, khẽ đẩy cậu lên trên ``sân khấu''.

Cuối cùng, Kuon cũng đành phải chịu khuất phục. Cậu tặc lưỡi, chống hông: ``\dots\ Thôi được. Hát một bài. Một bài thôi. Nhưng ai cười thì anh dỗi, đếch thèm chơi nữa đâu đấy.''

Và thế là, giữa tiếng cổ vũ rôm rả, Kuon cầm mic đứng dậy. Nhạc vang lên, và cậu bắt đầu hát, dù không hay như các idol, nhưng đủ ấm và đầy cảm xúc, khiến mọi người bất giác im lặng và lắng nghe.

Khi nốt cuối cùng ngân lên, tiếng vỗ tay vang dội khắp phòng, không chỉ vì giọng hát, mà còn vì người đang hát cũng là nhân vật chính của đêm nay.

\ct

Chín giờ năm mươi lăm phút.

Hoạt động tiếp theo trong chuỗi bốn hoạt động mừng sinh nhật Kuon là giải đấu trò chơi điện tử.

Ngay khi Kuon vừa dứt bài hát, chiếc vô tuyến đã được chuyển giao sang nhiệm vụ tiếp theo: từ việc phát nhạc cho karaoke sang việc phát màn hình cho trò chơi trên máy S*itch.

Các dây kết nối đã được lắp từ lúc trước, và nhóm GAMERS, vốn đã thủ sẵn máy Switch trong ba lô, nhanh chóng thiết lập một phòng chơi cho trò M*r*o Kart.

``\textbf{Và bây giờ là Giải thi đấu M*r*o Kart mừng sinh nhật Hikawa Kuon-senpai!!}'' Sora cầm chiếc míc-rô vừa đặt xuống kia, hào hứng thông báo với mọi người.

``Sinh nhật Kuon-senpai thì phải có hoạt động gì\dots\ kịch tính chút mới được,'' Fubuki cười ranh mãnh, vừa mở túi lấy tay cầm bộ điều khiển của riêng mình.

``Anh có thể tận dụng luôn chiếc S*itch vừa được Sui-chan tạng ban nãy đấy,'' Mio nói thêm, mắt liếc sang hộp quà đã được mở ban nãy.

``Ừ, sao?'' Kuon đáp, mắt vẫn nhìn vào màn hình chiếc máy để hoàn tất thiết lập.

``Giờ là lúc để xem trình độ lái xe của chủ nhân bữa tiệc này tới đâu!''  Korone thêm vào, trong lúc Okayu cặm cụi thiết lập kết nối J*y-C*n.

Đám đông tụ lại quanh khu vực chiếc ti-vi, nơi chiếc máy mới của Kuon đã được cắm vào. Cậu đành cười ngại ngồi xuống, nhận cặp J*y-C*n đang được kiểm thử từ Suisei, người vừa nháy mắt: ``Lái cho cẩn thận nhé, quà của em đấy.''

``Nếu mà anh lao xuống vực thì đó là lỗi tại em đấy,'' cậu Tổng đáp lại nửa đùa nửa thật.

Vòng đầu tiên của giải đấu bắt đầu, một lượt đua free-for-all (tự do) giữa Kuon và bốn thành viên GAMERS. Mỗi người chọn nhân vật mình thích: Fubuki chọn F(o)x B(o)y, Mio chọn I(nkl)ing G(ir)l, Korone trung thành với Sh(y) Gu(y)*, Okayu chọn C*t P(ea)ch.

Về phía Kuon, sau một thoáng ngập ngừng, đảo qua đảo lại qua các nhân vật, cậu quyết định chọn L**gi.

Nàng Khuyển liếc sang cậu, tỏ vẻ hơi ngạc nhiên với lựa chọn của cậu: ``\dots Anh chọn Luigi thật à?''

``Ừ. Anh thấy trong mắt ai cũng như nhau cả mà.Có gì sai sao?''

``Bọn em thường không thấy mấy ai chọn nhân vật đó thôi,'' nàng Miêu chiêm vào.

``\dots M*i*o hay L**gi trong cái trò này anh thấy ai cũng như ai thôi mà. Có sức mạnh nào đâu.''

Đường đua đầu tiên được chọn là một đường đua kinh điển: M*r*o Circuit. Đây là một đường đua tọa lạc trên ngọn đồi xanh, với một vài lâu đài thời Trung Cổ lênh bênh xung quanh.

Một lia máy quay ghi lại một góc của đường đua, với các nhân vật do các cô gái chọn lần lượt tiến về vạch xuất phát.

Đồng hồ bắt đầu đếm ngược.

``Ba, hai, một\dots\ \textbf{Go!}''

Ngay từ những giây đầu tiên, Korone đã chiếm thế thượng phong, bẻ lái chính xác ngay góc cua đầu tiên như thể đây là sân nhà của mình.  Cô nàng hét lớn, cười cành cạch: ``\textbf{Ha ha ha! Đường đua này là của chị!!}''

Okayu thì lặng lẽ bám theo, nước đi mượt mà không kém, đôi mắt nheo lại tập trung. ``Tôi sẽ đuổi kịp bà thôi, Koro-san\~{}''

Mio và Fubuki phía sau cũng không hề buông lỏng, mỗi người một chiến thuật: một người giữ tốc độ ổn định, người còn lại làm chuyên gia ``phá đám'' bằng cách sử dụng vật phẩm.

Kuon thì\dots\ hoàn toàn ngược lại. Trong con mắt của họ, cậu không khác nào một tay chơi mới vào nghề.

``Ơ này\dots\ cái bản đồ đâu!? Không có mini-map sao!?'' cậu la lên giữa lúc tay vẫn cố xoay cần lái theo phản xạ. ``Sao đoạn này cua gắt thế!?''

Một khúc cua trái bất ngờ khiến nhân vật của cậu điều khiển đâm sầm vào bức tường đá.

Xui xẻo hơn, ngay khi cậu vừa lấy lại tốc độ thì một trái chuối của Fubuki đã khiến xe trượt bánh quay vòng như con quay.

``N-Này, Fubu-chan! Đây là lần đầu anh chơi cái trò này đấy!'' cậu Tổng kêu lên trong khi không nhịn được cười.

Suisei ở phía sau chỉ nhoẻn miệng đáp: ``Kuon-senpai đúng là hơi đen rồi đấy. Đối đầu với GAMERS thì làm sao chơi lại được với họ!''

``Ừ. Đáng ra anh phải nghĩ đến điều đó.''

Cậu biết. Cậu biết mình đang phải đối đầu với một bộ tứ mà từng người đều có kinh nghiệm hàng trăm giờ trong thế giới trò chơi điện tử. Nhưng điều đó không làm cậu nản.

Vì dù gì, cậu vẫn còn một vũ khí bí mật nữa nắm trong lòng bàn tay. Nhưng trước mắt, cậu sẽ phải làm quen với cơ chế của trò chơi này trước.

Dù liên tục về cuối, Kuon vẫn cố giữ tay lái ổn định, tận dụng các boost pad (tấm tăng tốc độ) để đuổi kịp, thi thoảng còn bắn được vài quả shell (mai rùa) để trả đũa nhẹ nhàng. Một lần, cậu bứt lên được vị trí thứ ba, nhưng ngay lập tức bị Okayu dùng nấm tăng tốc vượt mặt như chớp.

Cuối cùng, vòng đua khép lại với Korone về nhất, Okayu thứ nhì, Mio ba, Fubuki tư, và Kuon, như đã đoán định trước, về vị trí thứ năm, trong giải đua chỉ có năm người.

``Chà, không tệ đâu Kuon-senpai!'' nàng Cáo trắng vừa nói vừa vỗ nhẹ lên vai cậu.

``Ừm, thua nhưng mà vui,'' Kuon cười xòa, đặt tay cầm xuống. ``Thật ra cũng dễ hiểu. Không map, lần đầu chơi, mà bị ném vào chuồng cọp thế này thì\dots\ sống sót được cũng là kỳ tích rồi.''

``Vẫn còn vòng hai chứ, đúng không mọi người?'' nàng Khuyển nháy mắt, tay xoay chiếc tay cầm như vũ khí chiến thắng.

``Mấy đứa định làm mấy vòng vậy?''

``Ba vòng!'' nàng Miêu chậm rãi đáp lại. ``Nhanh lên nào, chuẩn bị bắt đầu rồi đấy!''

Cậu Tổng cầm chiếc J*y-C*n lên, sẵn sàng cho vòng thi đấu thứ hai.

\ct

Bản đồ thứ hai được chọn là một đường đua uốn lượn quanh khu phố đêm kiểu Nhật, với đèn lồng đỏ rực, biển hiệu neon nhấp nháy, và những khúc cua hẹp ẩn sau dãy nhà gỗ cổ kính. Không khí trên màn hình trái ngược hẳn với bầu không khí náo nhiệt bên ngoài phòng khách, nơi mọi người vẫn đang cổ vũ, hò reo theo từng cú tăng tốc.

Tất cả mọi người vẫn giữ nhân vật như cũ, vì ``giờ mà đổi thì có khi còn loạn thêm nữa.'' Họ nắm chặt tay cầm, mắt chăm chăm vào màn hình, dặn mình không được để bị đối phương vượt mặt thêm lần nào nữa.

``\textbf{Go!}''

Ngay từ những đoạn đường đầu tiên, Okayu đã lên vị trí dẫn đầu, lợi dụng đoạn đường thẳng để xài nấm tăng tốc đúng lúc. ``Đã bảo là tôi sẽ cho bà hít khói mừ\~{}''

Korone lại dùng lối đánh ``cá tính'', bằng việc cố tình chạy gần xe Kuon để tung chuối, khiến cậu trượt dài sang bên mép cầu.

``Tròi ạ, tha cho anh đi mấy đứ!'' cậu rên rỉ, nhưng không ngừng cười khi thấy nhân vật mình quay như chong chóng.

``Chỉ là trò chơi thôi mà, Kuon-senpai\~{}'' nàng Khuyển vừa đáp lại vừa\dots\ thả thêm trái chuối nữa. Cậu chỉ còn biết lắc đầu mà cam chịu.

Tuy vậy, cậu đã thích nghi nhanh hơn hẳn so với vòng đầu tiên. Dù bản đồ không có chỉ dẫn, cậu đã dần nhớ được thứ tự khúc cua, bắt đầu bo lái gọn gàng hơn và dùng vật phẩm hợp lý hơn. Một lần, khi bị Fubuki vượt lên, Kuon đã khéo léo bắn quả shell đỏ khiến cô bị khựng lại đúng đoạn dốc.

``Này!'' nàng Cáo trắng rít lên, trong khi giấu khỏi sự ngỡ ngàng khi bị bắt bài đúng thời điểm.

``Ồ, senpai đúng là học nhanh thật đấy,'' Mio  Mio buột miệng khen khi bị Kuon vượt mặt trong thoáng chốc.

``Cũng bình thường thôi mà em,'' cậu đáp lại, mặt vẫn tập trung vào màn hình.

Ở vòng cuối cùng, Kuon và Okayu bám nhau sát nút. Cậu vừa bấm drift (lết bánh), vừa nghiêng đầu theo từng đoạn cua như thể chính mình đang cầm vô-lăng thật.

Đích đã ở ngay trước mặt. Cậu nghĩ thầm: ``\textit{Chắc vòng đấu này sẽ có cải thiện được phần nào nhỉ?}''

Nhưng, chỉ tiếc là\dots\ một quả shell xanh từ nàng Sói đen bay tới làm phá tan mọi hy vọng về nhất của cậu. Xe của Kuon lăn vài vòng rồi khựng lại ngay khúc cua cuối.

Cậu gục mặt xuống, bụm miệng cười, như thể đã đoán định được tình huống này. ``\textit{Tất nhiên là làm gì có chuyện ngon ăn như thế.}''

``Cảm ơn nhé, Mio\~{}'' nàng Cáo trắng giơ ngón cái lên, khuôn mặt tỏ vẻ thỏa mãn vì đã được trả thù.

Nhưng dù sao, Kuon cũng đã có điều mình muốn: cải thiện thứ tự xếp hạng, với việc cậu về thứ ba chung cuộc. Hai người đứng trước cậu là Okayu và Korone, và Mio cùng Fubuki về sau cậu.

``Khá phết đó\~{}, không tệ, không tệ\~{}'' Fubuki cười, chìa tay đập nhẹ vào tay cậu.

``Quả vừa rồi nếu không bị dính đạn thì chắc là về nhất rồi đấy,'' cậu bắt tay lại với đối phương.

``Thì ai bảo chạy sát đuôi Okayu làm chi?'' Mio cười hiền, nhưng giọng của cô trông không vẻ gì vô tội cả.

``Thôi được, hai vòng rồi, giờ anh bắt đầu quen với cơ chế này rồi đây\dots\ Anh nghĩ ta nên chỉ thêm một vòng nữa thôi, để còn kịp cho các hoạt động khác,'' cậu vươn vai, quay sang nhìn chiếc màn hình đang chuẩn bị cho vòng đấu cuối cùng.

``Mà này, đừng nói là map tiếp theo là\dots\ Rain(bo)w R(oa)d nhá?'' cậu nhìn thấy một trong những đường đua mà cậu đa từng thấy trước đó.

Korone cười nhăn răng, ánh mắt long lanh như trẻ con: ``Hồi hộp không Kuon-senpai?''

``\dots Có chứ. Chiến luôn.''

Mọi người lại cười ồ. Không khí dần chuyển sang cao trào, chuẩn bị cho vòng đua cuối cùng, nơi không chỉ cần tốc độ, mà cả\dots\ niềm tin vào mặt đường cũng không bị rơi mất.

``Cơ mà anh vẫn cần mấy đứa giải thích kỹ hơn cho anh về cơ chế trò chơi này.''

``Rất sẵn lòng!'' cả bốn cô gái GAMERS đồng thanh.

\ct

Âm thanh vui nhộn đặc trưng của M*r*o K**t vang lên một lần nữa khi màn hình lựa chọn nhân vật và xe hiện ra. Các cô gái nhóm GAMERS vừa thao tác vừa giải thích rôm rả, lần lượt chỉ cho Kuon cách xây dựng một chiếc xe chuẩn chỉnh nhất.

``Nhân vật thì anh thấy đấy, được chia ba loại: Small (nhỏ), Medium (vừa) và Large (lớn),'' Mio vừa điều khiển tay cầm, vừa giảng giải. ``Cỡ nhỏ thì nhẹ, dễ điều khiển, tăng tốc nhanh, nhưng dễ bị hất văng. Cỡ lớn thì ngược lại, khó ôm cua hơn, nhưng mạnh và `tông' được tụi nhẹ cân.''

``Xe và diều cũng vậy luôn đó,'' Fubuki gật đầu, nhấn mạnh thêm. ``Tùy loại sẽ thay đổi tốc độ, khả năng tăng tốc, khả năng drift, bám đường và cả chỉ số khi bay.''

``Nói cách khác, xe to thì dễ tông nhưng sẽ chạy chậm và ngược lại, anh cứ hiểu như thế đó!'' Korone chốt lại.

Okayu nhẹ nhàng liếc sang màn hình của Kuon: ``Nhớ xem cả bánh xe nữa nhé. Nhìn nhỏ xíu chứ ảnh hưởng nhiều lắm đó\~{}''

Kuon gật đầu lắng nghe, mắt dán vào bảng chỉ số chi tiết như một người đang phân tích dữ liệu chứ không phải chơi game. Các chỉ số như Speed (tốc độ), Acceleration (tăng tốc), Handling (độ phản hồi), Weight (cân nặng), Traction (độ bám), và Mini-Turbo (gần giống nitro) lần lượt hiện lên trước mắt cậu như những biến số trong một bài toán tối ưu.

``\textit{Nói chung là được cái lọ thì phải mất cái chai.}''

Cậu không thay đổi nhân vật, vẫn chọn L**gi làm người bạn đồng hành. Nhưng ở phần xe, cậu chọn một khung xe loại ``Standard Bike'', bánh ``R(olle)r'' giúp tăng khả năng tăng tốc và ôm cua, và một chiếc dù ``Cl(ou)d G(li)der'' để giữ sự ổn định trên không.

Không ai trong nhóm biết rằng, đằng sau nụ cười hiền lành và vẻ ngoài thư sinh kia là một ``huyền thoại ngầm'' từng chinh chiến qua hàng trăm giờ A(sph)alt 8 (A8) và A(sph)alt 9\footnote{Nay là \textit{Asphalt Legends Unite} sau ngày 17/7/2024.} (A9) trên điện thoại và trên máy tính. Kuon đã từng tính toán từng đoạn drift (trượt bánh), từng pha knockdown (tông phá xe), từng góc cua tử thần trên những cung đường đầy bụi và nitro.

Chỉ có Ryujin và Ryuuhou - những người gắn bó với cậu lâu hơn cả - biết rõ điều đó. Cô em gái của cậu đã từng thua không dưới năm lần trên các đường đua ở Tokyo, Nevada và Iceland trong A8, và đến tận bây giờ cô vẫn còn ấm ức.

Và lần này, tinh thần chiến binh đua xe trong cậu bỗng chốc trỗi dậy ngay tại thời khắc này. Ánh mắt Kuon sắc hơn thường lệ, sống lưng thẳng tắp, tay nắm chắc J*y-C*n.

Trông thấy ``ngọn lửa nhiệt huyết'' bùng lên từ cậu, Korone nhướng mày, hơi ngả đầu: ``Kuon-senpai\dots\ Anh trông có vẻ nghiêm túc lắm đây.''

``Coi bộ bắt đầu hiểu cơ chế rồi đó nha\~{}'' Fubuki cười khúc khích, nhưng ánh mắt có phần thận trọng.

``Lần này bọn em sẽ không nương tay với anh đâu đấy!'' Mio cười mỉm, nhưng rồi cũng phải hơi rén.

Còn Kuon? Cậu chỉ khẽ mỉm cười.

``Không hẳn là hiểu,'' cậu đáp, giọng nhẹ như gió thoảng. ``Nhưng\dots\ thử xem sao.''

Người bố của gia đình Hikawa dường như là người duy nhất nhận ra: nụ cười đó y hệt như cái ngày cậu bé Hikawa Kuon lần đầu vượt mặt ông trong cuộc đua trên đảo Barcelona mà không bằng sức mạnh, mà bằng tính toán tỉ mỉ và một trái tim nóng bỏng.

``\vie{Coi bộ thằng Kuon bắt đầu nóng lên rồi đấy,}'' ông buông ra một giọng nghiêm túc.

Các nhân vật đã chọn. Xe đã lên khung. Đường đua tiếp theo đang chờ đợi họ.

``\eng{Ta cùng bắt đầu vòng cuối nào.}''

\ct

Rainbow Road. Tên bản đồ hiện lên rực rỡ trên màn hình với những dải màu cầu vồng lấp lánh, vắt mình qua không gian vũ trụ. Dưới nền nhạc điện tử rộn ràng, các tay đua đã vào vị trí. Những bánh xe rít lên từng tiếng sốt ruột.

``\textit{Rainbow Road à\dots}'' cậu khẽ nhíu mày. Dù là lần đầu chạm mặt, cậu đã nghe danh về tỷ lệ rơi xuống vực cao đến vô lý của bản đồ này. Cũng vì lý do ấy, đây là một đường đua mà dân chuyên nghiệp thì hào hứng, còn lính mới thì chỉ biết toát mồ hôi lạnh.

Bản nhạc ``Animals'' của Martin Garrix bắt đầu phát lên trong tâm trí của Kuon. Một bài hát mang phong cách điện tử mà cậu đã nghe rất nhiều lần hồi cậu chơi A8. Cậu hy vọng rằng bài hát ``tủ'' sẽ giúp cậu gặp may.

Lượt đua thứ nhất bắt đầu. Các tay đua GAMERS nhanh chóng bung tốc như thường lệ. Mio chọn một nhân vật nhẹ cân, xe nhẹ, nhưng tốc độ tăng vọt; Korone thì vẫn trung thành với xe nặng, mạnh về va chạm; Okayu dùng xe có độ bám đường cao; Fubuki lại đi theo hướng cân bằng, dễ drift. Họ đã quen thuộc với từng khúc cua, từng đoạn trượt, từng cú phi thân trên bản đồ không có lề đường này.

Nhưng Kuon thì không vội. Ngay khi có hiệu lệnh, cậu không tăng tốc ngay, mà đeo bám các đối thủ khác. Cậu chạy ở tốc độ trung bình, như một học sinh đang dò đề kiểm tra.

Từng lần văng bánh, từng cú bay của đối thủ, thậm chí từng loại diều mà họ dùng, tất cả được ghi vào trí nhớ của cậu trước đó như một chuỗi mã hóa hoàn chỉnh.

Cậu nhìn một cách lặng lẽ, não phân tích đối thủ rất cặn kẽ, rồi tính toán lại trong đầu.

``\textit{Fubu-chan nhẹ hơn mình. Mio-chan là nhẹ nhất. Korone-chan thì nặng hơn, không chơi kiểu hạ gục được. Okayu thì một chín một mười với mình.}''

``\textit{Glider của Korone-chan rộng, có thể kéo dài thời gian trên không, coi như là bù cho độ tăng tốc chậm. Fubu-chan dùng loại tăng tốc nhưng thiếu ổn định\dots}''

``\textit{Nếu mình có thể dùng Shell đúng thời điểm\dots\ mình có thể ép được xe của họ.}''

Nói cách khác, Kuon không khác gì một chiếc máy tính sống đang điều khiển Luigi, một trong những nhân vật bị dè bỉu nhất trong M*i*o Kart.

Trong lúc đó, cả bốn thành viên của GAMERS vẫn nô đùa triệt hạ lẫn nhau, hoàn toàn quên đi sự hiện diện của một con quỷ đội lốt ở đằng sau.

Để rồi, ngay khi lượt hai tới, cơn ác mộng dành cho bộ tứ GAMERS mới thực sự bắt đầu.

Ngay khi Kuon vượt qua vạch vôi lần thứ hai, cậu bất ngờ tăng tốc. Một cú drift gắt khiến chiếc xe nghiêng về một bên, rồi bung ra một luồng mini turbo màu lam.

``Ể? Kuon-senpai tăng tốc rồi kìa!'' Mio khẽ giật mình khi thấy chiếc xe của cô bị Kuon bất ngờ vượt mặt.

Cậu nhặt được Red Shell, mắt lia tới mục tiêu trước mặt đầu tiên.

``Đỡ lấy!''

Cậu quẳng chiếc mai rùa thẳng vào bánh xe của Okayu. Bánh trước lệch góc. Xe lật nhẹ sang phải và\dots

``N-Này!!'' Okayu hét lên. Xe của cô trượt khỏi rìa cầu vồng, rơi xuống vực trong ánh nhìn bất lực.

``Ể! Kuon-senpai chơi không đẹp nha!'' nàng Miêu thốt lên, cười nửa giận nửa thương.

``Cái đó gọi là chiến thuật em ạ,'' cậu đáp lại bằng giọng nhẹ tênh.

``\textbf{Tôi sẽ phục thù cho bà, Ogayu!!}'' nàng Khuyển thốt lên đầy nhiệt huyết.

Cả khán phòng như nín thở sau cú hạ gục hai đối thủ của cậu con trai.

``Trận này căng hơn tôi tưởng nhiều peko\dots'' Pekora ngồi chồm hổm trên sàn, ôm gối, mắt không rời màn hình. ``Rõ ràng từ nãy anh ấy đang chơi giả nai peko!''

Shion thì ngồi vắt chân, nghiêng đầu: ``Ô hô? Làm ra vẻ vô hại mà lại tính toán kiểu đó\dots\ đúng kiểu nhân vật bí ẩn trong anime đấy.''

Đến một đoạn cua khá hẹp. Nàng Khuyển đang ở ngay sau cậu Tổng. Cô lao tới với vận tốc cao, định tông vào xe cậu ra khỏi đường đua\dots

``Nhận lấy này!!''

\dots nhưng rồi, như đã bắt trọn được đối thủ từ trước, cậu nhanh chóng nhảy diều đúng lúc bẻ cua, khiến cho cú va chạm của cô đi chệch hướng.

``Ể\dots\ Ể!?''

Nàng Khuyển hoảng hồn, lạc hoàn toàn tay lái, rồi nghiễm nhiên rơi thẳng xuống vực không gian.

``Cái gì vậy!?''

Để rồi, Kuon đã thành công cướp ngôi ở vị trí thứ hai từ nàng Chó đen đủi kia.

``Ôi giồi ôi\~{} nhìn cách cậu ấy xử lý khúc cua kìa\dots'' Lamy khẽ gật đầu, tay vẫn cầm ly nước nhưng không dám uống. ``Cái đó\dots\ không phải lần đầu chơi đâu nhỉ?''

``Uưưh\dots\ Kuonsa-senpai bắn ngừi ta dơ quá naaaaaa\~{}!'' Luna ôm đầu lắc lư khi thấy Korone bị hất khỏi đường đua.

``Anh ấy bắt đầu điều tiết đường đua rồi\dots'' Lamy nói, tay ôm vai lo lắng.

``Ừ thì, nếu game này là chiến trường\dots\ thì Kuon đang làm chủ bản đồ đấy,'' Ayame gật đầu. ``Anh ấy có vẻ không phải là tay dễ chơi đâu\dots''

Tiến tới lượt cuối cùng. Khi này, chỉ còn lại hai tay đua là ``Fubuking'' và Kuon. Hai người họ kì kèo nhau, ăn miếng trả miếng lẫn nhau, không ai chịu thua ai.

Fubuki vượt lên trước từ đầu và giữ thế áp đảo. Cô nàng chạy băng băng trên đường, không thèm đếm xỉa tới những đối thủ khác.

Ở bên dưới, những người chứng kiến trận đấu hò hét cổ vũ cho cô nàng Cáo trắng.

``Go! Fubuki-senpai ganbare\~{}!'' Luna vừa hò vừa vẫy hai tay.

``Đừng để anh ấy bứt lên nha! Hạ cậu ta đi!'' Shion nheo mắt trêu chọc.

``Được thôi! Fubuki ta đây sẽ giành chiến thắng!'' nàng Cáo trắng hạ quyết tâm, tiếp tục tăng tốc độ lên mức cao hơn.

Hẳn nhiên cậu Tổng hoàn toàn bỏ ngoài tai những lời như vậy mà tiếp tục tập trung vào cuộc đua.

Vì trong tay cậu vẫn còn một lá bài tẩy nữa. Một lá bài tẩy chí mạng.

Ngay khúc cua cuối, nơi đường đua chỉ còn rộng bằng một thân xe, cậu kích hoạt Mushroom. Bánh xe cháy lên tia lửa.

``Cho anh xí cái ngôi nhất nha.''

Một cú nghiêng nhẹ về phải. Xe Luigi va sườn vào Fubuki bằng tổng trọng lượng của nhân vật, xe và diều cộng lại.

*PANG!*

Chiếc xe cáo trượt bánh, lạc tay lái, và rồi\dots

``\textbf{Khôôôôôôôônggggggggggggg!!}''

Nàng Cáo tuyết đã bị Kuon hất khỏi đường đua. Điều đắng cay hơn là, cô nàng bị loại khỏi cuộc chơi ngay trước vạch đích.

Và dĩ nhiên, chiếc xe Luigi, một trong những người bị cộng đồng hắt hủi nhất, lại là xe giành chiến thắng trước sự ngỡngỡ ngàng của tất cả mọi người.

Fubuki chưa kịp trở lại đường đua thì đã nghe tiếng reo hò từ sau lưng.

``Tệ quá\~{} anh ấy trông ngầu đét chưa kìa\dots'' Ayame đập tay lên trán, mặt đỏ rần vì cười.

``Quả đúng là nam chính giấu nghề có khác peko\dots'' Pekora gật gù.

``Kuon-senpai bây giờ khác hẳn hồi nãy rồi à\dots'' Lamy thì cười nhẹ, mắt vẫn không rời màn hình.

Cùng lúc này, cô nàng vừa bị Kuon trở thành nạn nhân ban nãy quay ngoắt về phía cậu Tổng, phồng má đến mức tưởng sắp nổ: ``\textbf{Kuon-senpaiiiii!! Không công bằng!!}''

Cậu chỉ đáp lại bằng một cái nhún vai: ``Ờ thì\dots\ cái trò này có bao giờ công bằng ngay từ đầu đâu. Chẳng phải em nói là chơi để cho vui à?''

``Anh đã cướp mất chiến thắng của em rồi đó!!'' cô nàng giậm chân rầm rầm, đôi mắt trợn tròn như thể sắp rơi nước mắt đến nơi.

``Thì kệ xác em thôi. Anh cũng đếch quan tâm.''

Tiếng ``\textbf{HẢAAAA?!?}'' của Fubuki vang lên đến mức cô chó Mi ở trên tầng năm đang nằm cũng giật mình bật dậy.

Cả phòng khách như vỡ tung bởi tiếng cười vì phản ứng không ai ngờ tới này. Ryujin thì vỗ đùi cười rũ rượi: ``Thằng bé mà đã chơi đến đua xe là xác định\dots'', còn Ryuuhou ở sau lưng giơ ngón cái: ``Biết ngay mà! Onii-san đã lật kèo là y như trong A***alt luôn!''

Fubuki đùng đùng tức giận, hét lớn và chỉ thẳng vào mặt Kuon: ``\textbf{Hikawa Kuooooon!!! Em nghỉ chơi với anh một tuần!! Một tuần luôn đấy!!}'' rồi gập người xuống ôm chân, mặt gục vào đầu gối, trông chẳng khác gì đứa trẻ bị cướp mất cây kẹo yêu thích.

``Thôi nào, đừng trẻ con nữa. Cái dáng đó làm anh nhớ đến hồi em bị anh với Towa-chan cướp mất cái đồ chơi orc để làm lễ\footnote{Xem lại {\flaregothic ÉPISODE 98}, quyển số Tám.} đấy.''

``Oi oi\dots\ Tới khúc `hít-le' rồi đấy, peko!'' nàng Thỏ vừa nói vừa cười không thở nổi, suýt lăn xuống khỏi ghế.

``Chết cười luôn trời ơi\dots\ Fubuki-chan lần đầu bị chơi khăm hả?'' nàng Phù thủy chống cằm, cười chảy nước mắt.

``Lại còn là một vố đau nữa chứ!'' nàng Hầu gái chêm thêm.

``Uwaa\~{} Fubuki-chan mặt đỏ như cà chua kìa\~{}'' nàng Yêu tinh tuyết chỉ vào màn hình, tay vẫn run vì cười.

``Người chiến thắng thành nhân vật phản diện luôn kìa!!'' nàng Oni hét to, quăng luôn gối ra sau.

``Uuuuuu\~{} Kuonsa-senpai đáng gét quá điiiiii\~{}!!'' nàng Công chúa xứ Kẹo lăn lộn trên sàn, vừa cười vừa giả vờ gào thét giận dữ theo phong cách em bé.

Trong lúc đó, Mio bước lại gần cô bạn thân, nhẹ nhàng vỗ lưng cô nàng: ``Thôi nào, Fubuki. Chỉ là trò chơi thôi mà. Vui lên đi!''

``\textbf{Không vui vẻ gì hết!} Tôi đã chạm tay vào chiến thắng rồi mà!!''

``Ừ thì, vẫn có lần sau mà,'' Mio mỉm cười dỗ ngọt như mẹ dỗ con. ``Bà quên mất rằng đây chỉ là giao hữu thôi à?'

``Bọn em sẽ để chị thắng lần tới mừ\~{}'' Korone chiêm vào, xoa lưng cho đàn chị.

``Nhưng không phải chị thắng mà là anh ấy kia kìaaaa!!'' nàng Cáo tuyết hậm hực chỉ thẳng vào Kuon vẫn đang cười tỉnh bơ bên cạnh, như thể không biết trời đất gì.

``Thôi nào, ngầu một hồi rồi lặng lẽ ăn vạ đi, giữ chút hình tượng chứ\dots'' Okayu chống nạnh, vừa thở dài vừa cười rúc rích.

Sau một hồi châm chọc nàng Cáo trắng và cậu Tổng và những tràng cười đến ná thở, mọi người cũng dần lấy lại bình tĩnh.

Sora bước lên giữa phòng, tuyên bố tổng kết giải đấu: ``Được rồi, tổng kết. Giải `M*r*o K**t Friendship Cup' mừng sinh nhật Hikawa Kuon của chúng ta xin được kết thúc tại đây! Người thắng chung cuộc là\dots\ Hikawa Kuon-senpai!''

Cả phòng đồng thanh nhao nhao: ``\textbf{EHHHHHHHHHHHHHHHHHHHH!!?!?!??}''

``Anh tưởng là chúng ta tính theo kiểu lấy thứ hạng ba vòng chia trung bình hay là tổng điểm của ba vòng?'' Kuon gãi đầu gãi tai, tỏ vẻ thắc mắc với chiến thắng có phần quá dễ dàng và phi lý này.

``Em cũng thấy vô lý hết sức\dots\ Một lần anh về chót, một lần về ba, một lần về nhất, chắc gì được cao nhất\dots'' Mio thở dài, đồng quan điểm với cậu.

``Giải gì nghe kỳ cục vậy\dots'' nàng Miêu lắc đầu, khươ tay lên như thể không hiểu chuyện gì.

``Bây giờ mấy đứa mới để ý cái tên à?'' cậu nhíu mày, nhìn về phía nàng Thần tượng đang làm người dẫn.

``Cái tên đấy không phải đến để giao lưu đâu, đến để hủy diệt đó chứ,'' Korone vắt chân, vừa uống nước vừa nói như thể không liên quan gì.

``Ờ thì\dots\ trò này đâu có công bằng từ đầu mà?'' cậu Tổng lại buông câu cũ với nụ cười vô hại như thiên thần.

Để rồi đáp lại cậu là một tiếng ``\textbf{Im đi!}'' từ Fubuki, và một chiếc gối ném thẳng về cậu, khiến cả phòng nổ tung một lần nữa trong tiếng cười.

\ct

% Đoạn này ở phiên bản gốc sẽ nằm ở Partie 6 (C-3F), nhưng để giữ nhịp, phần này được đẩy lên đây.
Đồng hồ đã chỉ mười giờ ba mươi phút. Tất nhiên, các hoạt động mừng sinh nhật vẫn chưa kết thúc, mà dường như mọi thứ càng lúc càng náo nhiệt hơn.

``Giờ thì, quý vị và các bạn đang có mặt trong lễ sinh nhật đặc biệt hôm nay\dots'' nàng Cáo sa mạc hô to, đứng trên bục cao (vốn chỉ là mấy hộp bìa mà họ tìm thấy dưới kho) với ánh đèn chiếu rọi mà Lamy đã đem hồi nãy, ``Polka-chan xin tuyên bố tiết mục ảo thuật đầu tiên. Đó chính là\dots''

Cô giơ hai tay lên trời, ánh sáng đèn chiếu lóe lên từ phía sau như sân khấu Broadway: ``\textbf{The Mirage of Six Fates!!} (Ảo ảnh sáu lá định mệnh)''

Cả khán phòng vang lên tiếng pháo tay (chắc là do Polka tự bấm nút hiệu ứng từ xa), trong lúc cô lấy ra bộ bài Tây mới toanh mà mình đã mua khi cả nhóm đi sắm đồ chuẩn bị cho buổi tiệc sáng nay.

``Đúng là em ấy có mua bộ bài Tây đó thật,'' Marine ngồi khoanh tay nói với Suisei, ``lúc mà chúng ta mua quà tặng cho Kuon-senpai ấy?''

``Lúc đó chị còn tưởng em ấy mua về để chơi Tiềm Ô Quân kìa\dots'' nàng Ca sĩ ngước mắt lên bộ bài, đưa hai tay sau gáy.

``Sau cái vụ đó thì ta đã bị Kuon-senpai với A-chan cấm cửa rồi đấy\dots'' nàng Thuyền trưởng thở dài, trước khi tập trung vào tiết mục trước mặt đang chuẩn bị bắt đầu.

Trên bục ``sân khấu'', nàng Cáo xiếc vẫn đang luyên thuyên với bộ bàu của mình: ``Đây không chỉ là bộ bài bình thường đâu nhé\~{} Nó đã được Polka đây chế thêm chút `pol-pol energy' rồi!''

Sau màn tung hứng bộ bài - mà cô làm khá vụng về làm suýt rơi vài lần, khiến cho khán giả xôn xao e ngại - Polka đưa bộ bài cho Ayame và mời gọi: ``Đây, lấy một con bài bất kỳ đi nào, Ayame-tan.''

Nàng Quỷ Sừng nhẹ nhàng bốc ra một lá: bốn cơ. ``Đây nè\~{}''

Kế đến, Botan bốc ra một lá át (A) tép. Nene chọn được lá tám rô, Pekora là K bích, còn Korone lôi ra mười cơ.

Đến lượt Kuon, cậu nhìn thoáng xấp bài được xòe ra một chút rồi rút ra bảy tép. Một con số khá\dots\ trung tính, như chính vẻ mặt của cậu hiện tại. ``Anh là lá này.''

``Bây giờ mọi người hãy đặt lại lá mình vừa chọn vào trong bộ bài này nha!'' Polka tiếp tục hướng dẫn, và sáu người làm theo, mỗi người đặt lá bài của mình vào chỗ bất kỳ trong bộ bài. Người đặt ở đầu, người lại đặt lá chính giữa, người lại để lá ở cuối.

Sau đó, cô chia bộ bài ra làm hai nửa: ``Một nửa cho Ayame-senpai, một nửa cho Botan! Tráo giúp mình đi!''

``Được thôi\~{}'' nàng Oni nhận lấy nửa trên của bộ bài.

``Tráo kiểu gì cũng được ha?'' nàng Sư tử trắng tiếp lời, cầm lấy nửa còn lại.

``\textit{Dám chắc là họ sẽ tráo bài theo phong cách không bình thường cho mà xem,}'' Kuon thầm nghĩ, miệng cười khoáy trên môi.

Và đúng như những gì cậu dự đoán, Ayame bắt đầu vừa xáo bài vừa hát nhỏ giai điệu ``Docchi Docchi'', động tác nhẹ nhàng như đang múa quạt, thậm chí còn nhấp nhổm trong lúc trộn bài.

Botan thì\dots\ xáo theo cách ``hỏi chấm hơn''; cô thẳng tay ném bộ bài xuống sàn như đang ném lựu đạn trong các trò bắn súng, rồi nhăt lên xếp bừa và lặp lại như thế, khiến cho những người ngồi ở hàng đầu hoảng hồn.

``\hl{Shishiron-san đúng là có cách xáo bài dị thật,}'' A-chan, người đang ngồi ở hàng ghế cuối, cũng phải khẽ rùng mình mỗi lần tiếng *PẸP!* vang lên.

``\hl{Em cũng không ngờ tới luôn đó\dots}'' Sora cũng chỉ cười trừ. ``\hl{Thà rằng như Ayame-chan có khi còn trông vui vẻ hơn kìa\dots}''

Polka hài lòng gật gù: ``Tiếp theo, chuyển cho Nene-chi và Peko-senpai nào!''

Nàng ``Hải cẩu'' đỡ lấy, gồng tay, hét lớn: ``\textbf{Sức mạnh của mùa xuân!!}'' rồi cầm bộ bài xáo như thể đang ép rau má. Từ phía dưới, Miko giơ ngón cái lên, gật gù vẻ hài lòng.

Pekora thì vừa cười ``PEKO\uparrow\ PEKO\downarrow'' vừa vung tay xáo bài loạn xạ, trông không khác gì một con thỏ lên cơn động dục.

``Còn giờ thì đưa cho Korone và Kuon nữa\~{}'' Polka tiếp tục.

Nàng Khuyển chộp lấy bộ bài từ Nene, cười tít mắt: ``Xem đây\~{}!!'' và trộn những lá bài trong khi đang thực hiện những động tác trông như nhảy dây.

Cậu Tổng thì nhìn bộ bài, rồi nhìn lại những con người vừa xáo bài một cách không bình thường kia, cậu cũng chỉ tặc lưỡi thở dài, như thể phải làm một cái gì đó điên rồ.

Nói là làm, cậu cũng bắt đầu\dots\ tráo bài theo một cỗ máy đếm nhịp vô hình - cả theo nhịp và cách chuyển động của nó - làm cho một vài thành viên khán giả ngồi dưới cũng không nhịn được cười. Botan buột miệng: ``Trông anh ấy làm như đang lập trìn thuật toán vậy.''

Cuối cùng, cả hai người quẳng hai bộ bài đã được sắp lộn xộn lại về cho nàng Chủ xiếc. Polka bắt lấy gọn ghẽ, nâng nó lên trước mặt khán giả như thể đang cầm một thánh vật.

``Bây giờ tôi đã có bộ bài trên tay rồi, tôi sẽ-''

``\dots Đoán xem lá bài đó là của ai chứ gì?'' Mio chống cằm, vô tình thốt ra.

Không khí chợt tĩnh lặng như bị rút điện. Mọi ánh nhìn đổ dồn về phía nàng Sói đen, người vừa lỡ miệng ``cướp cò.''

``M-Mio! Đừng có làm người ta mất hứng chứ!'' Fubuki ngồi bên cạnh nhắc nhở, đù cô vẫn còn hơi ấm ức sau cú vuột mất chiến thắng từ hoạt động trước.

``X-Xin lỗi\dots\ Cái chất `thanh niên nghiêm túc' của tôi lại tuôn ra\dots'' Mio vội bịt miệng lại, cụp tai xuống nhìn về phía Polka.

Nàng Cáo sa mạc sững lại sau khi bị khán giả ``hớ''. Chiếc tai cáo phải của cô rũ xuống, miệng mếu máo: ``Gượm đã\dots\ chị biết rồi thì ai còn xem nữa!''

``Miosha đúng là vô duyên quá\dots'' Subaru ngồi đằng sau, vỗ vào lưng của ``người mẹ.''

``Đừng có phá hỏng bầu không khí chứ cháu,'' Kaien cười trừ. ``Kể cả biết rồi thì cũng phải im lặng, cháu nói thế có khác gì phá hỏng bầu không khí đâu.''

Nhận thấy rằng càng lúc càng có nhiều người la ó với hành động bất ngờ của cô, nàng Sói đen vội đứng dậy, hoảng hốt vẫy tay lia lịa: ``C-Chị xin lỗi mà! D-Diễn tiếp đi em!''

``Đó, Mio-chan xin lỗi rồi kìa,'' Kuon bật cười nhẹ. ``Pol-pol có tha cho em ấy không kìa?''

Cả khán giả cũng cười ồ lên, trong đó Luna vừa vỗ tay vừa la: ``Po-ka-chi\~{} cố lênnn\~{}!!''

Pekora thì ngồi sau hét lớn: ``\textbf{Đừng để bị phá nữa peko! Giành lại sân khấu đi peko!!}''

Sau một hồi động viên từ mọi người, cuối cùng nàng Cáo sa mạc mới có thể tiếp tục màn trình diễn của mình, mặc dù cảm thấy hơi mất hứng sau pha vừa rồi.

Polka rướn người về phía trước, giơ cao một chiếc băng bịt mắt màu đỏ lấp lánh như đồ diễn xiếc chuyên nghiệp: ``Giờ thì\dots\ tôi sẽ đeo băng bịt mắt này nha\~{}!!''

Một số tiếng xì xào, tiếng hú hét ``Wooo\~{}'' của những thành viên có trái tim ngây thơ và cả ``Oaa\~{}'' đã chuếnh choáng men rượu (chủ yếu là từ Aqua và Lamy) nổi lên từ phía dưới khán giả.

Sau khi bịt mắt xong, cô quay vài vòng để thêm phần kịch tính, rồi nhẹ nhàng rút ra một lá bài từ trong bộ. Tay cô khẽ run lên vì hồi hộp,hoặc vì muốn làm người ta tưởng là hồi hộp.

Polka bỏ chiếc bịt mắt ra, mắt đảo qua lá bài cô vừa rút, rồi mỉm cười một cách tự tin, lớn giọng hỏi: ``Ayame-senpai, đây có phải là lá của chị không!?''

Một lá bốn cơ.

Ayame nhấc tay reo lên ngay tức thì, đôi mắt cô rạng rỡ: ``Đúng rồi đó, Chủ xiếc-chan oi!''

Một tràng vỗ tay giòn giã vang lên phía dưới. Đến bản thân Polka cũng không tin được rằng cô nàng đã đoán đúng, mặc dù cô nàng bản thân là một chủ xiếc.

Tiếp theo là Botan. Polka lại rút đúng lá át tép.

Nen là đúng lá tám rô.

Pekora, K bích, không sai vào đâu được.

Korone, chính là mười cơ.

Không ai nói thành lời, tất cả chỉ có thể nhìn nhau rồi vỗ tay lần nữa. Cứ mỗi lá bài nàng Cáo sa mạc đoán đúng, những tràng pháo tay lại một lúc lớn hơn đi kèm với những tiếng hò reo được\dots\ vang lên khe khẽ, vì bây giờ cũng không còn sớm sủa gì.

Haato cười khúc khích: ``Quả đúng là Chủ xiếc nhỉ. Cái gì cũng biết!''

Pekora nghiêng đầu, nhìn khuôn mặt nửa hớn hở, nửa ngạc nhiên (mà cô không biết rằng đó là cảm xúc thật của Chủ xiếc) đáp: ``Không biết em ấy có bàn tay ma thuật gì nữa, peko\dots''

Nene nhí nhảnh nói: ``Đoán đúng năm lá liên tiếp thế này thì đúng là giỏi thật!''

Mio chỉ cười nhẹ, nhìn về Polka với đôi mắt thán phục: ``Để xem em ấy có đoán đúng được lá của Kuon-senpai không.''

Fubuki thì tỏ vẻ như đã đoán định từ trước: ``Dám chắc là `đồng chí' của chúng ta đoán đúng thôi.''

Polka mỉm cười, cảm giác hưng phấn lan tỏa trên khuôn mặt cô như một đứa trẻ sắp được tặng thêm bánh kem. Cô đưa tay vào rút lá cuối cùng với vẻ chắc chắn tuyệt đối.

``Đây có phải lá của anh không, Kuon-senpai?'' cô hỏi, giơ ra một lá bài úp mặt về phía mình, nhưng quay về phía Kuon.

Một lá bảy tép. Đúng lá bài mà cậu rút ra.

Nhưng thay vì cậu nói lá bài đó đúng hay sai, cậu chỉ nhướng mày, rồi nở một nụ cười đầy ẩn ý: ``Hmm\dots\ Có đúng không nhỉ?''

Im lặng.

Toàn bộ mọi người trong căn phòng bỗng chốc như ngừng thở.

Không phải vì họ hồi hộp, mà vì không tin nổi là cậu Tổng lại\dots\ chơi đòn tâm lý ngược.

Polka chớp mắt vài lần, hoàn toàn bị động với câu hỏi này: ``Ể\dots? Ý anh là sao?''

Kuon ngả nhẹ người ra sau, mắt vẫn dán vào lá bài: ``Em thử nói anh nghe xem\dots\ anh đã bốc lá gì?''

Từ bên dưới, Watame chợt lên tiếng: ``Này, Kuon-senpai, chẳng phải điều này là hơi quá sao?''

Cậu khẽ lắc đầu, nhẹ nhàng đáp lại: ``Không quá đâu. Tại anh thấy em ấy đoán đúng liên tiếp quá nên anh sinh nghi thôi.''

Cậu ngoảnh đầu về phía nàng Chủ xiếc đang toát mồ hôi, tiếp tục gây áp lực: ``Vậy thì Polka-chan? Em trả lời được không?''

Nàng Cáo sa mạc lắp bắp, nhưng vẫn cố nở một nụ cười: ``Ê-tồ\dots\ Một con\dots\ tép ạ?''

``Đúng.''

Khán giả ở dưới bắt đầu thì thầm, bắt đầu có cùng chung suy nghĩ với Kuon.

``Kuon-senpai nói cũng có lý,'' Shion lên tiếng suy luận, ``Đoán trúng được tận năm lần liên tiếp như vậy rất là khó luôn.''

``Chưa kể rằng em ấy còn có bài gì giấu chúng ta nữa chứ\dots'' Matsuri nói góp.

Trên ``sân khấu'', Kuon tiếp tục hỏi, vẫn với giọng bình bình: ``Số mấy?''

Cô nàng Chủ xiếc cắn môi, ánh mắt đảo loạn như chong chóng rồi lí nhí cất giọng đoán bừa: ``S-Số bảy\dots ạ?''

Kuon im lặng. Không nói một lấy một câu, cũng không có một cử chỉ nào để tuyên bố xem câu trả lời là đúng hay không.

``K-Kuon-senpai?'' Polka run rẩy, tay vẫn cố cầm vững lá bài.

Aki cảm nhận thấy bầu không khí bỗng nặng trĩu, ho he lên tiếng: ``Polka-chan trông có vẻ mất tự tin ghê nhỉ\dots''

Roboco ngước nhìn cô nàng, đôi mắt tỏ vẻ lo lắng: ``Cố lên nào, Polka-chan!''

``Chắc là cậu ta bắt đầu nổi da gà rồi đấy\dots'' Botan nhìn bản mặt hoảng hốt của người bạn thân, rồi nhíu mày cười. ``Anh ấy thích tăng độ khó cho game cơ đấy.''

Ba mươi giây trôi qua. Vẫn không có một chỉ thị nào từ Kuon. Thứ duy nhất mà Polka nhận chỉ là ánh mắt trông như hình viên đạn và vẻ mặt lạnh tanh.

``Kuon-senpai! Làm ơn nói gì đi chứ!'' nàng Cáo xiếc rít lên, toàn thân cô run lên như muốn hối thúc cậu thật nhanh.

Im bặt.

Một giây. Hai giây. Ba giây.

\dots

``\textbf{Chính xác!!!}'' Kuon bất chợt hô lớn, vang lên đến cả căn nhà, như thể cậu vừa đậu cấp ba.

Những tràng vỗ tay bùng nổ, những lời có cánh được đưa tới Polka, người vẫn còn đang đơ nghệch, không hiểu chuyện gì vừa xảy ra.

Thậm chí, nàng Song tính còn bật ngửa ngã ra sau tiếng hô trời giáng kia.

``Giỏi quá, Polka-senpai-nora!!\~{}'' Luna vỗ tay bôm bốp, cười tít mắt sau màn trình diễn vừa rồi.

``Phải công nhận là Polka-sama siêu thật\dots'' Choco mỉm cười, cầm lấy cốc nước ép của mình và uống một ngụm sau tình huống áp lực kia.

Về phía Polka, ngay khi cô biết được cô đã đoán đúng và được Kuon ``giải thoát,'' cô gục xuống sàn, tai phải của cô dựng đứng, rồi hét tới lạc cả giọng: ``\textbf{Anh dọa em một vố rồi đó, mồ\~{}\~{}! Kuon-senpai đúng là đồ xấu tính!}''

Kuon chỉ cười nhẹ đáp: ``Áp lực mỗi tí mà đã đằng đặc thế kia\dots\ Nhưng mà phải công nhận, Pol-pol giỏi thật đấy. Sáu lá bài liền, bịt mắt mà vẫn đoán đúng. Tập luyện nhiều rồi đấy hỏ?''

Nghe thấy thế, Polka nhảy dựng lên, búng tay cái tách: ``Chuyện nhỏ!! Em mà lị! Em còn nhiều trò hay nữa cơ!''

Nàng Sư tử trắng vừa vỗ tay vừa lẩm bẩm: ``Kuon-senpai chơi đùa với người ta y như người lớn trong phim hoạt hình vậy.''

Nàng Sói đen thở hắt ra, ôm lấy lồng ngực: ``Chị suýt tưởng Polka ngất thật\dots''

Nene thì ôm chầm lấy cô bạn cùng Thế hệ của mình: ``Làm thêm trò nữa đi\~{}!''

Và đúng như lời tuyên bố, nàng ``Chủ xiếc đại tài'' ngay sau đó tiếp tục dẫn dắt những tiết mục tiếp theo, vừa hài hước, vừa kỳ quặc, và tất nhiên, vô cùng đậm chất nàng Cáo xiếc.

\ct

Polka tiếp tục bước lên sân khấu lần nữa, kéo lê một cái hộp gỗ to như hòm kho báu (mà thực chất là cô ``chôm'' được ỏ quán cà phê bên cạnh), phía trên dán đầy các hình dán màu mè và mấy dòng chữ viết tay kiểu như ``Cấm mở khi chưa niệm chú!'' và ``Quà sinh nhật cực sốc cho vị khách đặc biệt!''.

Cô nàng vỗ tay gây sự chú ý của khán giả, những người vẫn còn hơi choáng sau màn ``đấu trí'' bất ngờ giữa nàng Cáo xiếc và cậu Tổng. ``Mọi người ơi, mọi người ơi!! Bây giờ chúng ta sẽ đến với tiết mục tiếp theo trong màn trình diễn của Ảo thuật gia Omaru Polka nha! Tôi gọi nó là `The Mysterious Gift!' (Hộp quà bí mật)'

Tiếng vỗ tay và vài tiếng huýt sáo vang lên, trong đó có tiếng ``Peko\~{}'' từ khán giả hàng đầu, mặc dù thế họ vẫn tự bảo ban nhau khẽ tiếng xuống.

``Không biết lần này Pol-pol của chúng ta sẽ có màn trình diễn gì tiếp theo đây!'' Aki hào hứng khi nhìn cô nàng Chủ xiếc.

``À, cái hòm bí ẩn đây mà,'' Rushia cười. ``Một trong những trò ảo thuật kinh điển nhưng không bao giờ chán.''

``Dám chắc lại là mấy con chim bồ câu thôi ấy mà,'' Botan đáp lại, mặt cố tỏ vẻ lạnh lùng, nhưng trong thâm tâm, cô cũng hoàn toàn đồng ý với nàng Pháp sư.

Trong khi đó, Polka chống nạnh, chỉ vào cái hộp, cố gắng tạo ra một bầu không khí đầy bí ẩn: ``Trong chiếc hộp này là một sinh vật huyền bí đến từ thế giới song song! Nói chính xác hơn\dots\ \textbf{là một con thú đến từ tương lai!}'' cô nàng cao giọng ở nửa sau câu. ``Món quà sinh nhật đặc biệt dành cho anh đó, Kuon-senpai!''

Lập tức, tất cả mọi người đồng loạt quay về phía Kuon. Cậu rướn mày, tỏ vẻ nghi ngờ: ``Em nói cái quái gì vậy? Nhà anh vốn dĩ đã có chó rồi mà, em nhớ chứ? Con Mi nhà anh đấy?''

Một số tiếng cười rúc rích vang lên khi các cô gái còn nhớ đến con chó Bắc Kinh lai Nhật lông trắng lúc nào cũng nằm ườn bên cửa phòng của Kuon.

``Nói mới nhớ, ta vẫn phải tạ ơn em ấy một lần đấy,'' Korone nói.

``Ừ. Nếu không có Mi-chan đuổi gián đi thì có lẽ chúng ta đã ngất ngay tại trên tủ quần áo rồi,'' Okayu đáp lại, giọng thản nhiên.

Hẳn nhiên, nàng Cáo sa mạc tiếp tục dùng chiêu bài thuyết phục. Cô một mực khẳng định: ``Con này khác nha! Không phải là chó bình thường đâu, mà là quái thú xuyên không luôn! Nhìn vô là khiến cho anh bất ngờ tới bật ngửa luôn á!''

``Nghe phét lác vãi,'' Kuon bật cười. ``Nhưng mà thôi được rồi, để xem em mang con gì tới nào.''

Một vài thành viên dưới đó cười ồ lên với lời nói có phần khá mạnh bạo từ nàng Chủ xiếc. Fubuki cười nhẹ, còn Mio thì không giấu được cái nhăn mặt: ``\textit{Làm ơn đừng mang con gì quá đáng là được\dots}''.

``Và dĩ nhiên, tiết mục này vẫn cần có người hỗ trợ!'' Polka quay ra, rồi chỉ về phía nàng Khuyển đang trò chuyện vui vẻ với cô bạn Miêu của mình, nổi giọng hứng thú: ``Korone-senpai, em lại phải nhờ chị lần nữa rồi.''

``Được thôi\~{}'' nàng Khuyển nâu vui vẻ đáp lại, chạy tung tăng lên trên bục. Nhưng chưa kịp để người dẫn nói gì thêm, cô mon men lại gần tới chiếc hộp.

``A-lê-hấp! Để chị xem nào!!''

``T-Từ từ đã nào, Korone-senpai!''

Cô mở nắp hộp ra. Cả khán phòng dán mắt nhìn về khuôn mặt rạng rỡ của nàng Chó, để rồi chứng kiến cô nàng nghiêng đầu, hồn nhiên hỏi: ``Ể? Chị có thấy cái gì trong này đâu Polka-chan? Trống trơn đây nè\~{}''

Polka bật cười, đóng chiếc hộp xuống rồi bảo: ``Hấp tấp như thế không được đâu! Chị phải niệm câu thần chú mới có tác dụng chứ!''

``À, ra thế!'' nàng Khuyển reo lên, tự cốc đầu vào mình lè lưỡi: ``Chị quên mất tiêu!''

``Chẳng trách chú thấy có gì đó sai sai,'' Ryujin lên tiếng. ``Đến cả những kẻ nghiệp dư cũng biết đến điều đó mà.''

Polka ấn chặt chiếc nắp hộp, một tay chống nạnh, tiếp tục hướng dẫn tiền bối của mình: ``Bây giờ, chị phải nói `Pol-kotsu Magic!' thật là lớn! Chị chỉ cần nói một lần thôi!''

``À, được thôi!''

Không một chút nghi ngờ gì, nàng Chó nâu làm theo ngay lập tức: ``\textbf{Pol-kotsu Magic!}'' thậm chí để làm cho người ta tin tưởng hơn, cô còn múa may tay và ``phả ma thuật'' vào thẳng chiếc hộp, cùng lúc đó, Polka thổi một cái kèn đồ chơi nhỏ (cái đó cô cũng mua về hồi sáng) và đập vào cạnh hộp.

Shion nhìn hai cô nàng cùng giuộc đang làm những trò ngớ ngẩn, buột miệng trêu: ``Hai người đang gọi hồn về đấy à?''

``Câu đó phải là để em nói mới đúng đó-nanodesu,'' Rushia lên tiếng, khẽ cau có nhẹ vì bị cướp mất lời.

``Rồi đó! Chị mở lại chiếc hộp đi!''

Korone mở chiếc hộp gỗ một lần nữa, cô nhìn thấy\dots

``Polka-chan! Trong này còn có cả một nàng chó nữa nè\~{}!''

\dots một cô gái có cặp tai chó một bên màu nâu, một bên màu trắng đang nằm cuộn mình bên trong với dáng mèo ngủ. Mái tóc trắng dài mượt cùng với phần đỉnh đầu có điểm chút màu nâu, đôi tai dựng khẽ giật giật cùng chiếc đuôi dài lông xù. Cô nàng người thú (kemonomimi) này diện một bộ áo hai dây màu lam cùng chiếc váy trắng, cùng tông màu với bộ tóc của mình.

``Ờm\dots\ Pol-pol, em kiếm em ấy ở đâu ra vậy?'' Kuon đứng dậy, nhìn thấy sự xuất hiện của nàng chó xa lạ kia.

``Thì em nói rồi đó, cô ấy là một con thú đến từ tương lai!'' nàng Chủ xiếc đáp lại.

Cậu Tổng định mở miệng ra tiếp tục hỏi, thì cô nàng chó xa lạ kia từ từ mở mắt, có lẽ vì tiếng động và tiếng hô lớn của nàng Cáo sa mạc.

``Ồ, em làm cô ấy dậy rồi kìa.''

Với vẻ mặt vẫn còn ngái ngủ, cô vươn vai lên trên cao, ngáp một hơi rất lớn, như vừa mới ngủ dậy.

``Đến giờ ăn trưa rồi à\dots'' cô nàng tóc trắng cất giọng, một giọng nói trầm ấm và có phần nhẹ nhàng hơn. Không ít thành viên trong nhóm bày tỏ sự ngạc nhiên ngay cả khi họ đã quen với những người thú nói tiếng người, kể cả là Kuon.

``À-rế\dots? Mọi người làm gì ở đây vậy\dots? Đông vui ghê\~{}'' cô ngước nhìn lên mọi người, đầu nghiêng hẳn sang phải, đôi tai hơi cụp xuống. Cô nhìn quanh căn phòng, rồi chợt dừng lại về phía gương mặt quen thuộc, người vẫn đang ngơ ngác dưới sân khấu.

Cô chớp mắt một hồi lâu, rồi mỉm cười cất giọng nói trong trẻo đó: ``Cậu chủ!''

Bầu không khí lặng ngắt. Những tiếng xì xào bàn tán bắt đầu nhen nhóm ở dưới, cùng với đó là những biểu cảm ngạc nhiên khi bỗng chốc một cô gái tai chó mà họ chưa từng gặp bao giờ lại gọi Kuon với danh ``cậu chủ.''

Người ngạc nhiên nhất, không ai khác ngoài Kuon, kẻ vẫn đang ngồi ở hàng đầu với đôi mắt mở to, hoàn toàn chưa thể tin vào những gì mình nghe được.

Cậu đứng bật dậy, chỉ tay vào cô gái đang nằm kia: ``Ch-Chờ đã\dots\ Em là ai vậy!?''

Cô nàng tóc trắng với phần mái nâu lồm cồm bò ra khỏi hộp, phủi bụi chiếc váy trắng của mình. Cô bước tới gần cậu, nụ cười mỉm của cô càng rõ ràng hơn, như vừa gặp lại người thân sao bao nhiêu ngày xa cách.

``Cậu chủ\dots\ Sao trông tóc anh bù xù vậy? Em nhớ là tuần trước anh đã cắt rồi mà nhỉ?''

Nghe những lời bộc bạch có phần hơi ngây thơ của cô, cậu con trai càng chết lặng hơn. Cậu muốn thốt lên: ``\textit{Em đang nói cái gì vậy? Cắt tóc gì?}''

Cô nàng tai chó này do thám toàn thân cậu, từ đầu đến chân, rồi tiếp tục hỏi với giọng nhẹ nhàng: ``Trông anh cũng hơi béo hơn á\~{} Hôm qua ăn nhiều bánh kem quá ạ?''

``\dots Hả?''

Cô gái cười khúc khích, rồi bất ngờ tiến tới nắm lấy tay cậu, ánh mắt tím long lanh trìu mến: ``Không sao đâu, Cậu chủ\~{} Dù anh có như thế nào thì anh vẫn là Cậu chủ của em mà. Chúng ta đã sống với nhau mười bảy năm rồi mà nhỉ\~{}''

Cậu Tổng nuốt khan, càng nghe càng không hiểu những gì mà cô nàng đang thốt ra.

Đáp lại với những sự ngỡ ngàng hiện rõ trên khuôn mặt tới không nói nên lời của cậu, cô nàng chỉ bật cười khe khẽ. ``Mặt anh trông ngố tàu quá!''

Ngay trong bầu không khí trầm lặng này, một giọng nói nhắng nhít cất lên phá hỏng bầu không khí này: ``Oa\~{} Tai này là tai chó thật này\~{} Em là giống gì vậy? Beagle hả? Hay là Golden Retriever? Trông em nhỏ con thế này chắc là giống chó cảnh nhỉ?''

Cô gái bí ẩn chớp mắt nhìn Korone, người đang đùa nghịch với đôi tai chó dựng và mái tóc trắng dài qua vai của mình đang để xõa. ``A! Korone-san! Chị hết dỗi rồi ạ?''

``Ể, Ể? Mà, em đang nói gì vậy?'' nàng Khuyển nâu khẽ nghiêng đầu.

``Lần trước chị thách đấu với em trò T*tr*s 99, và rồi chị đã cáu đến mức dỗi nửa ngày đó\~{}''

``Chị thực sự không hiểu em đang nói gì cả\dots\ Mà, em là ai vậy?''

Cô nàng tai chó trắng khẽ phồng má: ``Chị đừng có trêu như thế chứ! Em là Miyako đây mà, sao chị chóng quên thế? Em cũng từng là người cứu chị với Okayu-san khỏi đám gián hồi Tết nguyên đán hai năm trước đấy!''

``Nói cách khác, em chính là Mi-chan hồi đó sao?'' nàng Miêu cất giọng hỏi, không giấu nổi sự thích thú.

``Vâng!'' cô nàng tự xưng Miyako kia đáp với đầy sự hớn hở. ``Giờ em là Miyako rồi đó!''

Cả nhà Hikawa chấn động, sửng sốt nhất lại chính là cô em gái của Kuon, người mà đã chủ động bắt con chó Mi từ cửa hàng thú cưng về từ nhiều năm trước.

Ryuuhou đứng bật dậy, mặt mày xanh xao như tàu lá chuối: ``K-Khoan\dots\ Mi? Em là Mi-chan đó hả!? Cái con chó mà em bế về hồi Onii-san mười tuổi á hả!?''

Ryujin chỉ có thể nói đúng một chữ: ``Mi\dots?'' Đến bản thân ông cũng không tin nổi những gì xảy ra trước mắt.

Riêng Kuon, cậu hoàn toàn đứng chôn chân. Mặt đỏ lên, tay vẫn đang bị nắm, nhưng đầu óc như bị đóng băng toàn phần vì quá tải.

``Em nói thật hả\dots? Em là con Mi\dots\ mà nhà anh nuôi từ hồi đó á?''

Miyako mỉm cười, gật đầu như thể đó là điều hiển nhiên: ``Ừm! Mà cũng lâu lắm rồi còn gì. Em từ\dots\ khoảng hai năm sau thời điểm \emph{ấy} thì được Koyori-sama\footnote{Sẽ được giới thiệu ở quyển số Mười Ba.} cho em một cuộc đời mới rồi á\~{}''

``Ờm\dots\ Hả? Koyori? Em đang nói tới--'' cậu Tổng thảng thốt, nhưng rồi Polka bật cười cắt ngang lời trong lúc vẫn đang nấp đằng sau nộp: ``Không thể ngờ rằng đó là một bất ngờ cực lớn đấy, thưa quý vị!''

``Cái này gọi là dịch chuyển xuyên không-thời gian luôn á\~{}'' nàng Miêu đứng ở dưới gọi vọng lên. Ngay sau đó, khán giả đồng loạt cười rộ lên vì câu nói vừa đúng tình huống vừa đầy bất ngờ.

Nàng Cáo xiếc cười mỉm về phía Miyako, đồng thời vỗ tay tiếp tục gây chú ý: ``Rồi, rồi! Ta phải đưa Miyako-chan về chứ không là rối loạn dòng thời gian này mất!''

Nói rồi, cô lấy ra một bộ điều khiển (cũng là đồ cô xin được ở cửa hàng điện tử cạnh nhà) và bấm một cái nút, trên chiếc hộp lập tức bốc ra một làn khói lấp lánh.

Cô gái tai thú bị cuốn vào làn sáng mờ nhạt, và trước khi biến mất, còn kịp thì thầm với Kuon: ``Nhớ giữ ấm cho em như hồi xưa nha, cậu chủ\dots\ Em iu anh lắm\~{}''

*BỤP!*

Và rồi, Miyako tan biến trong làn sương khói đen mờ ảo.

Kuon vẫn đứng đơ như tượng sau sự xuất hiện của cô nàng tai chó kia, chỉ có thể thốt ra một câu: ``\dots Pol-pol đúng là làm anh bất ngờ tới ngã ngửa thật.''

Korone thì vừa tỏ vẻ thất vọng khi đồng loại biến mất, nhưng rồi tỏ vẻ phấn khích sau cuộc tiếp xúc ngắn ngủi kia: ``Ể\~{} Vừa mới có `đồng chí' mới mà biến mất rồi! Nhưng mà dễ thương ghê\~{}''

``Rồi chị sẽ gặp lại đồng đội của mình sớm thôi!'' nàng Cáo xiếc đáp lại. ``Cảm ơn Korone-senpai vì đã hỗ trợ em cho màn trình diễn này!''

Những tràng pháo tay rộn rã được trao tới nàng Cáo sa mạc. Luna, như mọi lần, lại là người vỗ lớn nhất: ``Tuỵt vừi quá-nanora! Một màn hô bín lun á!''

Matsuri thì nở một nụ cười mãn nguyện: ``Chị muốn được nựng em ấy ghê\~{}''

Kuon sau một hồi tiêu hóa tất cả những gì vừa diễn ra, cũng chỉ có thể cười trong bất lực trước màn trình diến của cô, nhưng đồng thời không thể không đặt câu hỏi về sự tồn tại thật sự của Miyako và sự liên quan tới con chó cảnh Mi nhà cậu.

``Mà\dots\ thôi kệ. Có khi đó lại là một ảo ảnh đó thì sao\dots'' cậu nhún vai, trước khi tạm thời chấp nhận với thực tại và tiếp tục theo dõi buổi trình diễn.

\ct

Chiếc hộp gỗ quay trở lại lần nữa, lần này được dán thêm giấy màu hồng chói lọi, đi cùng với một tờ nhắn màu hồng đề: ``Cảnh báo: không được mở trong khi đang ăn!''

``Chiếc hộp bí mật của chúng ta vẫn còn chứa bí mật nữa! Và tôi hứa, không có ai bay từ tương lai tới hết!'' nàng Cáo sa mạc nói trong khi khiêng chiếc hộp lên. ``Đây sẽ là một bài kiểm tra thách thức thần kinh cực đại, ai dám mở chiếc hộp bí ẩn này sẽ nhận được một món quà độc quyền không có ngoài thị trường!''

Nghe những lời nói đao to búa lớn của Polka, không ít thành viên ở dưới bắt đầu tỏ vẻ rất cảnh giác.

``Tôi thấy có mùi nghi ngờ rồi đấy\dots'' Flare nhíu mày nghi ngại.

``Tôi dám cá là lại một thứ gì đó mờ ám cho mà xem,'' Kanata căng mặt ra.

Khi chiếc hộp được đặt ở trên bục, thay vì gọi từ dưới khán đài lên như lần đầu, Polka chủ động đi xuống.

``Miko-senpai! Chị là người vinh dự làm người hỗ trợ cho em cho màn này!'' nàng Cáo sa mạc ngoắc tay gọi nàng Vu nữ.

Nghe thấy vậy, Miko cười khà khà, trông như vừa được trúng số: ``He he he\dots\ Để đó cho Vu nữ Ưu tú Miko-chan này lo!'' Thậm chí, cô nàng còn không thèm giấu vẻ mặt dâm đãng của một ông chú thèm khát có bạn gái tới cùng cực.

Và rồi, cô không hề ngần ngại mà gọi lớn: ``Kuon-senpai! Mau lên phụ em cái coi, ne\~{}''

Cậu Tổng, người vừa ngồi xuống chưa được bao lâu, đã phải thở dài ngán ngẩm vì bị réo tên lần nữa: ``Sao mấy đứa cứ thích gọi anh lên vậy?''

Đáp lại cậu chỉ là những tiếng cười từ khán giả, dù lần này họ tỏ vẻ ngượng ngùng hơn hẳn. Không còn cách nào khác, cậu lại phải lôi cái xác của mình lên ``sân khấu''.

``Được rồi, giống như lần vừa rồi nhá!'' Polka hô lớn. ``Lần này tôi cũng sẽ là người hô luôn!''

Cả ba cùng nhau đồng thanh, với một giọng tràn đầy năng lượng, một giọng tràn đầy sự ham muốn, và một giọng như bị gượng ép: ``\textbf{Pol-kotsu\dots\ Magic!}''

Chiếc hộp rung lắc nhẹ, không hề đứng yên như lần trước, khiến cho những người ngồi hàng đầu cũng phải khẽ giật mình.

``Ể? Có cái gì trong kia kìa\dots'' Towa chỉ vào chiếc hộp.

``Liệu có khi nào em ấy triệu hẳn ra một con gì đó không?'' Sora lên tiếng.

Polka vẫn tiếp tục màn trình diễn của mình, bỏ ngoài tai những lời bàn tán dưới sân khấu. ``Được rồi, Miko-chi, Kuon-senpai! Mở chiếc hộp ra nào!''

Chiếc hộp gỗ được mở tung ra, và thứ bên trong chiếc hộp\dots

\dots làm cho hai người mở hộp có hai thái cực khác nhau.

Một đống đồ lót phụ nữ sặc sỡ bật tung ra, rơi lả tả khắp xung quanh bục như mưa xuân. Ren, nơ, hoa văn đủ màu, có cả những cái còn gấp gọn, và một vài cái thì không!

Miko sáng bừng mắt, gần như phát sáng toàn thân khi phát hiện ra một ``kho báu'' đầy ý nghĩa cho một tín đồ trò chơi người lớn.

``K-Khoan, kia chẳng phải là\dots?'' Noel chỉ về phía vật mà nàng Vu nữ Ưu tú tự xưng đang cầm trên tay.

``Ooohhh\~{}!! Quần lót của Fubuki-senpai nè!! Cái này là của Okayu-senpai hình dâu tây! Còn cái có tai thỏ này là của Noel-senpai nè!!'' cô nàng cầm lấy ba chiếc quần lót từ trong hòm ra, vung vẩy trên không trung.

Nhìn thấy đồ nội y của mình đang được ``bêu rếu'' ngay tại bàn dân thiên hạ, Okayu, Fubuki, Noel lập tức đỏ mặt quay đi.

Fubuki hét lớn, mắt há hốc vì bất bình: ``P-Polkaaaaaa!!! Em lấy cái này lúc nào vậy hả!?''

Noel quay qua chỗ Flare, úp thẳng vào lòng của cô bạn vì ngượng: ``Em thề là không biết tại sao nó lại nằm trong đó!!''

Thậm chí là Okayu, vốn dĩ rất bình tĩnh và điềm đạm, cũng đỏ rực mặt, cố với lấy chiếc quần của mình, la lớn: ``Ơ\dots\ cái đó là hàng giới hạn! Polka-chan!! Trả lại cho chị!''

Đối nghịch hoàn toàn với sự hớn hở và điên cuồng của Miko, Kuon lại lặng lẽ quay đi, dần dần lùi về cuối phòng, cố giữ bản mặt bình tĩnh nhất có thể, không nói một câu gì ngoài lặp đi lặp lại các câu: ``Anh vô can nhá!'' và ``Đừng có lôi anh vào!''

Thậm chí là cả Kaien, Ryuuhou và A-chan, những người đang ngồi ngoài cuộc (một phần) cũng phải đỏ mặt vì thứ mà nàng Vu nữ đang cầm trên tay.

Với một cô nàng có bộ não ngập tràn eroge, cô ấy thậm chí còn nâng độ quái dị của mình lên một nấc thang mới, tới mức có thể gọi là ``phá hủy hoàn toàn hình tượng''.

``\textbf{Ê, xem này!} Chiếc áo này là của ai đây ne!?'' cô nàng cầm chiếc áo lót có sợi chỉ mảnh mang hoa văn báo lên, rồi\dots\ ngửi thẳng vào đó, hít lấy hít để như không có ngày mai.

Hành động quái dị đó đương nhiên khiến cho một vài thành viên la ó: có người bất mãn, có người ngượng ngùng, thậm chí có người còn muốn ăn thua đủ với kẻ bày ra trò này.

``\textbf{Miko!!! Đừng!!!}'' Kanata hét lên, bắt đầu rơi nước mắt vì tức giận.

``Oe\~{} Miko-chi đừng có phơi ra cái quần hình Daifuku chứ\~{}'' Lamy thì che mặt hoàn toàn, đôi tai nhọn hoắt của cô đỏ ửng lên.

Hết Miko, họ bắt đầu chuyển sang công kích Polka, người dường như vẫn đang cười cành cạch trước cảnh tượng này.

``Này, Polka! Mau tìm cách ngăn chị ấy đi chứ!'' Nene thảng thốt, định sồn sồn lên trên nhưng bị Botan cản.

``Polka-chan đúng là đồ dâm dê\~{}!'' Rushia rền rĩ.

Nhưng rồi, trong những tiếng ồn hỗn loạn kia, một thứ mà Miko bới được khiến cho tất cả mọi người chú ý.

Một chiếc quần lót màu cam chấm bi, với chiếc nơ hình đầu hề rõ ràng. ``Cái này hình như là của Polka-chan đúng không ne?''

Toàn bộ căn phòng im phăng phắc, hàng chục cặp mắt đổ dồn về chiếc nội y kia.

Lập tức, nàng Cáo xiếc đang cười bỗng chốc đứng hình. Sau câu nói nhẹ tênh nhưng chí mạng đó, cô nàng bắt đầu đổ mồ hôi lạnh.

Cô lắp bắp: ``C-Cái đó\dots\ h-hình như là\dots\ đồ đạo cụ á! Chắc là của\dots\ của\dots\ hình thú bông Polka-chan\dots?''

Kuon, người vẫn không ho he một lời ngay khi sự cố xảy ra, đáp gọn lỏn: ``Nó chính là cái quần mà em đang mặc đó, Pol-pol ạ.''

``Ể!?''

Polka vội lật váy lên kiểm tra. Cô nàng\dots\ đang tô hô trước mặt mọi người, khiến cho hàng chục con người sững sờ. Cái ``chiến lợi phẩm'' mà Miko đang cầm kia chính là thứ mà nàng Cáo sa mạc đang mặc!

``M-Miko-senpai!! Trả nó lại cho em!!'' cô nhào tới giật lại cái quần chấm bi từ tay cô nàng.

Nhưng nàng Vu nữ chỉ nhoẻn miệng cười: ``Ể? Polka-chan chả phải vừa bảo đây là phần thưởng dành cho người can đảm nhất sao, ne?''

Sau câu nói châm chọc đó, nàng Cáo xiếc ức tới mức không cãi lại được vì quả thật bản thân cô đã nói vậy, trong khi khán giả lại có thêm một tràng cười nữa vì ``quả báo tới nhãn lồng.''

``Polka-chan đúng là\dots'' Haato thở dài, không tỏ vẻ tiếc thương gì tới nàng Cáo xiếc: ``Đúng là quả báo nhãn lồng mà.''

``Bị như thế đúng là quá xứng đáng chung thuyền với Marine luôn-nanodesu!'' Rushia nói với sự đanh thép, tới mức cô bạn Senchou ngồi cạnh cũng điếng người: ``Sao bà lại lôi tôi vào!?''

Đến cuối cùng, Polka giật lại chiếc quần, mặc vào một cách vội vàng rồi tiếp tục bị mọi người châm chọc. Miko thì dẫu nhận về một vài ánh mắt ghê tởm, cô nàng vẫn không mảy may quan tâm mà vẫy tay chào tạm biệt như một thần tượng vừa biểu diễn thành công. Còn vè chiếc hộp, Mio đã đem đi niêm phong để không cho bất kỳ kẻ tội lỗi nào bén mảng nữa.

Ngoại trừ\dots\ một ai đó đang cố gắng tránh né chuyện này. Các bạn cũng đoán được đó là ai rồi.

\ct

Sân khấu trở nên yên tĩnh, ánh đèn dịu lại. Polka bước ra một lần nữa, vẻ mặt nghiêm túc bất thường. Màn trình diễn thứ ba, và có lẽ là cuối cùng của nàng Cáo xiếc sẽ được bắt đầu ngay sau đây.

Cô nàng bước lên bục, hắng giọng một tiếng, rồi tiếp tục dẫn với phong cách chuyên nghiệp nhất có thể: ``E hèm\dots\ Sau màn trình diễn khiến mọi người nóng mặt vừa rồi, mình muốn lấy lại lòng tin từ mọi người bằng một tiết mục đỉnh cao khác: `\textbf{The Hypnotize Coin!} (Đồng xu thôi miên)''

``Chắc là lại trò rẻ tiền gì đó rồi đó-nanodesu,'' nàng Rushia chẹp miệng.

``Bài đó hình như là của tôi thì phải?'' Ayame lên tiếng, liền nhận ra ngay khi cái tên cho màn trình diễn này được đưa ra.

``Đúng là dễ đoán thật đấy,'' Matsuri đáp, rồi quay sang hỏi nàng Chủ xiếc: ``Này Polka-chan, em học chiêu đó từ đâu vậy?''

Nàng Cáo xiếc chống hông đầy hãnh diện: ``Bí kíp này em luyện được từ một bậc tiền bối đã thuần hóa con sư tử trong đoàn xiếc ngày trước!''

Một số thành viên vỗ tay bôm bốp, thậm chí có người còn khen khả năng của Polka với một tài tuy đơn giản nhưng lại thú vị này.

\dots Đó là cho đến khi cô nàng bất chợt rụt rè hơn hẳn, khúm núm gãi đầu: ``\dots Cơ mà con sư tử kia đã cắn đứt sợi dây xích rồi đuổi theo ông ấy chạy đến trối chết\dots''

Tới đây, một số thành viên trong khán giả cũng chỉ cười trừ, một số còn bi quan tới mức đã đoán định được kết quả.

``Biết ngay là kiểu gì cũng thất bại mà-nanodesu,'' nàng Pháp sư chu môi, tỏ vẻ tự tin với dự đoán của mình.

``Cứ để xem em ấy sẽ làm gì, đừng có bình phẩm vội chứ, mọi người,'' Sora lên tiếng trấn an. Dù thế, có một cơ số người không hoàn toàn tin tưởng lắm, vì đến bản thân nàng Thần tượng cũng không chắc chắn với lời nói của mình.

Polka rút ra một đồng xu năm yên được buộc bằng một sợi chỉ, giơ lên cao trước mọi người: ``Người vinh dự được mời lên đây sẽ là người thử nghiệm! Nào, mời Kuon-senpai lên sân khấu!''

Cậu Tổng nghe thấy tên mình lần nữa cũng chỉ biết thở dài cam chịu. Cậu biết chắc là cậu sẽ lại trỏ thành mục tiêu lần nữa nên không quá bất ngờ.

``\textit{Cái lờ gì cũng là anh vậy\dots}'' cậu lẩm bẩm, nhẫn nại lên bục một lần nữa. Rút kinh nghiệm lần trước, cậu tỏ vẻ cảnh giác hơn.

Bây giờ, cậu đang đứng trước mặt nàng Chủ xiếc, người đang giơ đồng xu, đứng ở tư thế chuẩn bị.

``Bây giờ Kuon-senpai làm theo em! Thả lỏng người lại! Hít vào, thở ra!''

Cậu Tổng không ngần ngại mà làm theo, dù lần này cẩn trọng hơn.

``Bây giờ anh hãy nhìn vào đồng xu này, nhìn thật kỹ, đừng chớp mắt.''

``\dots Được rồi. Anh cũng hiểu được sơ sơ em định làm gì rồi đấy.''

Dứt lời, Polka bắt đầu đung đưa đồng tiền xu qua đi qua lại một cách chầm chậm.

``\textit{Biết ngay.}''

Cô nàng nói với giọng nghiêm nghị: ``Tôi là thỏ, tôi là thỏ\dots\ Nói theo em nào!'' trong khi tay không ngừng đưa qua đưa lại đồng xu.

Cậu Tổng dù thấy rằng điều cô đang làm khá lố bịch và có lẽ sẽ không có tác dụng, nhưng cậu cũng đành làm theo. Cậu đưa con mắt của mình theo đồng xu, lặp lại: ``Tôi là thỏ, tôi là thỏ, tôi là thỏ\dots''

Đồng xu tiếp tục lăc lư, nàng Chủ xiếc nhìn chằm chàm về phía cậu như đang cố điều khiển tâm trí vốn dĩ khó đoán này. Cả hai người họ tiếp tục lặp lại những câu nói tưởng chừng vô thưởng vô phạt đó.

Ba mươi giây trôi qua.

Kuon vẫn chưa có dấu hiệu gì của việc bị chiếc đồng xu kia điều khiển. Cậu vẫn nói với giọng bình thường, không hề chậm lại một nhịp. Thậm chí để chắc ăn, cậu đã nhắm mắt lại để tự ``đưa vào ảo giác của bản thân,'' nhưng có vẻ không có tác dụng.

``Tôi là thỏ, tôi là thỏ, tôi là thỏ\dots''

Một vài người theo thiên hướng ``ăn xổi'' đã bắt đầu sốt ruột. Đặc biệt là Matsuri, cô nàng vẫn đang dõi theo hai người họ: ``Lâu quá! Đến bao giờ mới xong đây!''

``Thì nó đã không hiệu quả ngay từ đầu rồi mà,'' Rushia thở dài. ``Mà khoan, từ từ đã\dots''

``Có chuyện gì xảy ra rồi kìa!'' Watame chỉ về phía một người đang mất dần sự tỉnh táo.

Không, không phải Hikawa Kuon. Mà chính là người đang cố gắng thôi miên cậu, Omaru Polka.

Chính cô nàng Cáo xiếc không biết vì là tay mơ hay do Kuon quá mạnh mà cô lập tức bị đánh ngược lại.

Cô đứng im như tượng. Mắt long lanh. Tai gật gù. Miệng lí nhí lặp lại: ``Pol\dots\ Pol\dots\ Polka là thỏ\dots\ Pol\dots''

Để rồi khi cậu Tổng mở mắt ra, cậu đã thấy nàng Cáo sa mạc đang nhảy tưng tưng y như một con thỏ thật.

Khán giả chứng kiến cảnh người thôi miên bị bật ngược lại đã có những phản ứng khác nhau.

Aki thì hốt hoảng khi thấy người dẫn đã bị chính đồng xu dụ khị: ``Polka-chan bị thôi miên rồi kìa!''

Botan thì thở dài: ``Tôi cũng chẳng ngạc nhiên lắm với kết quả này.''

Riêng Kuon thì ngơ ngác nhìn Polka đang hành xử như một con thỏ, thậm chí còn giống thỏ hon cả Pekora: ``Ủa\dots? Rồi ai mới là người bị thôi miên đây?''

Mà nhắc đến Pekora, ngay lúc đó, một tiếng hét quen thuộc vang lên từ phía khán đài: ``\textbf{Nè!! Đó là của cà rốt của chị mà!!}''

Nhân lúc không một ai để ý, Polka đã nhẫn tâm cướp củ cà rốt gắn trên lọn tóc của Pekora rồi nhảy tưng tưng đi mất trong sự ngỡ ngàng của chính bản gốc.

``\textbf{Polka!! Mau đền củ cà rốt cho chị mau lên!!!}'' nàng Thỏ tức giận, đứng bật dậy đuổi theo. ``\textbf{Với cả, cái đó là thương hiệu của chị mà!!!}''

Nhưng Polka vẫn dường như không nghe thấy gì. Cô tiếp tục lẩm bẩm trong vô thức: ``Polka là thỏ\dots\ Polka là thủ lĩnh của loài thỏ\dots\ Không phải\dots\ Peko\dots''

``\textbf{Đồ khốn nạn!!!}''

Thế là, Pekora đuổi theo Polka vòng vòng sân khấu, gào thét khản cả cổ, thậm chí còn xuống cả khán đài, xuống hẳn cả tầng một, chạy một vòng quanh căn nhà và quay trở lại tầng ba, còn Polka thì vẫn cứ vừa chạy vừa kêu ``Pyon! Pyon!'' như thỏ trốn sói.

Một số thành viên, bao gồm cả Kuon và Ryuuhou đều không thể nhịn được cười trước tai nạn khó đỡ này của nàng Cáo xiếc: vừa bị phản tác dụng, lại vừa vô tình cướp đi ``danh tiếng'' của nàng Thỏ khét tiếng.

Shion cười đến lăn lộn đất, suýt đạp đổ cốc nước: ``Ôi trời ôi! Hai người này diễn tuồng còn hay hơn cả bọn mình luôn!''

Luna thì cười ôm bụng tới mức chảy nước mắt: ``Polka-chan với Peko-senpai chơi dượt nhau chông zui quớ-nora!!''

A-chan thì chỉ bất lực với hai người đang đuổi bắt đó: ``Omaru-san đúng là không thể nào nghiêm túc quá ba phút\dots'' để rồi Kuon đáp mức: ``Cái máu xiếc của em ấy đã ăn sâu tới mức đó rồi, bỏ làm sao được.''

Đến cuối cùng, đích thân vị khách được mời lên lại chính là người hóa giải, đưa Polka thoát khỏi thuật thôi miên đó.

*BỐP!*

Cậu chỉ cần vỗ tay một cái vang dội ngay trước mặt cô và đôi mắt tím lẫn hình dán các quân bài đã trở lại bình thường.

Nàng Cáo xiếc giật mình, ngờ nghệch nhìn mọi thứ xung quanh, đôi tai dựng đứng lên vì bất ngờ. Cô lắc đầu vài cái, rồi ngơ ngác hỏi: ``Ể\dots\ Em vừa làm gì vậy? Sao mọi người lại nhìn em như nhìn sinh vật lạ thế này?''

Cậu Tổng lẩm bẩm, cười chát với thuật của nàng Cáo xiếc: ``Thuật thôi miên gì mà còn dễ phá hơn game flash mười năm trước thế.''

Cả phòng phá lên cười, còn Pekora thì vẫn ôm củ cà rốt vừa giành lại được, mặt hậm hực: ``Lần sau đừng có mà cướp thương hiệu của chị nữa, peko!''

Polka gãi đầu, cười trừ: ``Xin lỗi mà\~{} Em cũng không hiểu sao lại thành ra như vậy nữa\dots''

Kuon vỗ vai cô, mỉm cười: ``Thôi thì, màn trình diễn của em coi như thành công rồi đấy. Cảm ơn Polka-chan vì đã mang lại tiếng cười cho mọi người nhé.''

``\textit{Thế này thì lấy lại danh dự kiểu gì đây\dots}'' cô nàng lẩm bẩm, quay về vị trí ngồi của mình trong những tiếng cười vang của mọi người.

Tiếng vỗ tay lại vang lên, khép lại chuỗi tiết mục ảo thuật đầy bất ngờ và đậm chất ``xiếc'' của nàng Cáo sa mạc, đòng thời sẽ tiến đến hoạt động cuối cùng của bữa tiệc mừng sinh nhật Kuon.

\ct

Mười một giờ đêm. Chỉ còn một tiếng nữa, sinh nhật thứ hai mươi lăm của Kuon sẽ qua đi.

Trước mặt ba mươi lăm người là một đống dưa hấu căng mọng đã được chất gọn thành núi. ``Chúng ta sẽ tận dụng chỗ dưa này!'' cô nàng chống nạnh, trưng ra chỗ dưa hấu còn sót lại.

Lập tức, một làn sóng ``\textbf{Ohhh!}'' vang lên, nhưng rồi chính họ đã tự bịt miệng vì đã gần tới nửa đêm.

Subaru nhíu mày, đổ mồ hôi hột: ``Vậy là vẫn nhất quyết theo ý tưởng đập dưa này à\dots''

``Thì đập xong rồi ăn thôi\~{} Truyền thống mùa hè của Nhật mà, Subaru-chan!''

Choco ngước nhìn đống dưa có, tỏ vẻ nghi ngại: ``Ngẫm nghĩ lại chúng ta mua hơi nhiều thật\dots\ Thôi thì đập bớt mấy quả để cho vui chắc cũng không sao.''

Kuon cũng tiếp lời, đồng quan điểm với nàng Y tá: ``Anh cũng không thích ý tưởng đập dưa lắm đâu\dots\ Nhưng mà hôm nay sinh nhật, nên\dots\ cứ thoải mái đi. Với cả, mùa hè mà không có dưa hấu thì chẳng khác gì ăn phở không có quẩy.''

``Onii-san so sánh kỳ quặc thật đấy\dots'' Ryuuhou toát mồ hôi hột.

Sora vỗ tay, vội dập tắt những hoài nghi của mọi người: ``Rồi rồi, nếu mà Fubuki-chan muốn vậy và mọi người không có ý kiến gì, mọi người phụ nhau dọn chỗ này ra, để dựng sân chơi nhé!''

Okayu và Korone nhanh chóng bê từng bao cát chữa cháy (lấy từ kho xây dựng của nhà Hikawa, được Kaien bật đèn xanh) đã chuẩn bị sẵn ra, rồi đổ đều lên khoảng sân giữa để tạo cảm giác như bãi biển.

``Rải cát xung quanh nào, Ogayu!!''

``Đồng ý\~{}''

Subaru đi ngang qua nhìn cặp bạn thân đang chạy vòng quanh đổ những xô cát, càm ràm: ``Mọi người định dựng lại bãi biển thật đấy à\dots''

``Chẳng phải em nói hồi sáng là trò đập dưa phải ở trên bãi biển mới đúng kiểu mà, đúng chứ?'' Fubuki cười gian.

``\dots Đúng là thế thật. Nhưng mà chúng ta sẽ dọn dẹp thế nào đây!?'' nàng Vịt hỏi với ánh mắt lo ngại.

``Hốt cái máy hút bụi trên nhà là xong chứ gì\~{}''

Subaru chỉ có thể thở dài với câu trả lời tỉnh bơ đó của nàng Cáo trắng.

Một hồi rải cát trôi qua, và hoạt động cuối cùng trong chuỗi sắp sửa bắt đầu. Fubuki cầm sẵn bốn ba chiếc chày gỗ đánh bóng (cô đã mua sẵn thừa lúc không ai để ý trên đường tới quầy thanh toán), đặt trước mọi người.

Cô nàng hắng giọng, bắt đầu công bố luật chơi: ``Chắc mọi người cũng biết trò chơi này chơi với nhau rồi nhỉ?''

Tất cả mọi người cùng đồng thanh: ``Đúng thế!'' / ``Ưm!'' / ``Tất nhiên!'', đồng thời gật đầu.

``Ta sẽ bịt mắt, quay ba vòng, rồi dùng chày đập dưa! Trúng quả dưa nào thì được thưởng luôn quả đó!'' Haato phổ biến, ``Mấy cái luật này chắc ai cũng quen rồi nhỉ?''

``Với cả, nếu đập trúng ai thì người đó chịu nha\~{}'' Fubuki nói với giọng sắc lẹm, mắt lia về cậu Tổng, người vừa mới bị chơi một vố đau sau hoạt động đua xe ban nãy.

Kuon thở dài, lường trước được mọi thứ: ``\textit{Đây là cái cớ để phục thù mình đây mà\dots}'' Nhưng dẫu vậy, cậu vẫn bước vào trò chơi với tâm thế thoải mái nhất.

\ct

Lượt chơi đầu tiên bắt đầu. Nàng Cáo trắng lập tức bắt đầu chỉ điểm tới ba người ứng với ba đại diện: Noel cơ bắp, Rushia thất thường và Ayame quỷ kiếm sĩ.

Nàng Hiệp sĩ là người ra sân đầu tiên. Cô nàng mang cho mình giáp sắt này tràn đầy khí thế, buộc khăn che mắt xong thì giơ cao cây gậy như chuẩn bị xuất quân. ``Mình sẽ cho mấy quả dưa biết thế nào là cơ bắp chân chính!''

Cô chống chiếc chày xuống, tự xoay mình đúng ba vòng rồi cầm chiếc chày lên trên cao.

Nhưng chỉ sau đúng hai bước, cô nàng bắt đầu có dấu hiệu chập choạng, lệch dần sang bên trái.

Mọi người nín thở dõi theo. Có người còn không tiếc lời cổ vũ cho Noel, thậm chí có người còn cố gắng chỉ dẫn cho cô nàng về đúng vị trí của quả dưa hấu.

\begin{dialogue}
    \speak{Haato} Noel-chan! Hướng này, hướng này!
    \speak{Matsuri} Đi lệch đường rồi kìa! Quay lại!
\end{dialogue}

Trong những tiếng hò reo cổ vũ, nàng Hiệp sĩ tỏ vẻ hoang mang: ``Ể\dots\ Hướng nào đây? Hướng nào?''

``Đằng trước bà đó! Đó hướng đó- không! Xoay người lại!'' Marine có lẽ là người chỉ đường tận tình nhất, dù rằng Noel không tiếp nhận lắm.

``Đằng trước à!? Được rồi!''

Và thế là, cô vung chiếc chày với tất cả sức lực, không thèm dò lại để xác nhận.

``\textbf{Hây da!}''

*RĂC!*

Một tiếng nghe như gỗ vỡ vang lên khô khốc khiến cho mọi người giật mình.

``Đập trúng chưa vậy mọi\dots\ người?''

Sora cất giọng, hơi run: ``Noel-chan, đó không phải là quả dưa đâu, mà là\dots''

Nàng Hiệp sĩ tháo bịt mắt ra, hớn hở nhìn thứ mà cô đập trước mặt, để rồi nhíu mày khó hiểu. Trước mặt cô chẳng có quả dưa hấu nào cả, mà đó lại là\dots\ một cái bàn làm việc gỗ từ công ty trước để lại, giờ đây đã bị đập bể làm đôi và đổ sập xuống đất, để lại đám bụi và mảnh vụn.

``Noel-chan hụt rồi nha\~{}'' Okayu ngân nga. ``Cũng may đó không phải là ai đó đấy.''

Kuon ngẩn người đứng nhìn chiếc bàn đó, thở dài: ``Đằng nào bọn anh cũng định vứt cái bàn đó rồi. Không sao đâu.''

Dĩ nhiên, quả đó coi như là không đập trúng. Noel tỏ vẻ hơi tiếc nuối về vị trí, nhưng cũng cảm thấy may mắn khi không ai thương vong.

\ct

Lượt chơi thứ hai cho Rushia lại không suôn sẻ như vậy. Ngay khi nhìn thấy bản mặt nghiêm túc cực độ, tất cả mọi người bắt đầu thấy hơi lo.

``Cái bản mặt của Rushia trông đáng sợ quá\dots'' Marine toát mồ hôi lạnh, cảm nhận được sát khí bừng bừng trong người.

Cô nàng buộc chiếc khăn lại, lẩm bẩm: ``Lũ dưa kia đừng hòng chạy thoát\dots\ Ta sẽ đập nát các ngươi!''

Cô quay ba vòng một cách khá gọn gàng và nhanh gọn, rồi, không thèm nghe những lời chỉ dẫn xung quanh, lao thẳng về phía trước.

``\textbf{Aryaaaaaaaaaa!!!!!}''

Polka, người đang ngồi bệt ở ngay trước quả dưa hấu, chưa kịp hiểu rõ ngọn ngành mọi chuyện trước mắt, chỉ có thể lớ ngớ la lên: ``Ể!? R-Rushia-senpai!! \textbf{N-Này!!!}''

*ĐOẠNG!*

Một cú đập xuống chuẩn xác. ``Thế nào? Trúng rồi chứ?''

Nàng Pháp sư bỏ bị mắt xuống, nụ cười của cô dần méo xệch, rồi chuyển sang ngượng ngùng khi thứ mà cô đập xuống không phải là quả dưa\dots

\dots mà chính là nàng Cáo sa mạc xấu số kia, người đã bất tỉnh nhân sự, nằm sõng soài trên sàn với đôi mắt đang như nhìn thấy ngôi sao trên trời.

Fubuki lập tức quỳ một chân xuống cạnh Polka, vén tóc cô lên, kiểm tra hơi thở. Cô nàng ngẩng lên trời, rơi nước mắt (cá sấu) và nói một cách truyền cảm: ``Polka, đồng chí của chúng ta, đã hy sinh vì nghệ thuật\dots'' trong khi các thành viên xung quanh thì người há hốc miệng, người ngán ngẩm lắc đầu, người cũng chỉ có thể cười trừ cạn lời.

``Bỏ văn vở đi em ơi,'' Kuon thở dài. ``Ru-chan cũng trượt rồi đó.''

\ct

Người cuối cùng trong lượt này là Ayame, người có tinh thần năng lượng cực lớn.

Cô buộc khăn che mắt xong thì nhẹ nhàng rút cây gậy như thể đang chuẩn bị\dots\ hành lễ. Vẻ mặt bình tĩnh và tập trung khiến khán giả trông chờ một cú đập thần sầu.

Ayame quay ba vòng đầy dứt khoát, không hề có dấu hiệu chao đảo. Như một kiếm sĩ cảm nhận được khí trường, cô bắt đầu di chuyển chậm rãi, dò từng bước. Chiếc gậy gỗ hơi nâng lên, cô khẽ nhấc mũi chân, rồi bất ngờ dừng lại để phán đoán hướng.

Từ phía ngoài, Roboco lên tiếng động viên: ``Cố lên Ayame-chan! Tiến thêm hai bước nữa là tới quả dưa rồi đó!''

Haato cũng hô lớn: ``Bên trái một chút nữa đi! Đúng rồi, đúng rồi!''

Choco nhẹ nhàng nhắc: ``Cẩn thận nhé, phía trước là quả dưa, đừng mạnh quá kẻo văng mất phần ngon!''

Ayame khẽ gật đầu, tập trung lắng nghe, rồi điều chỉnh lại hướng đi theo lời chỉ dẫn. Cô hít một hơi thật sâu, gậy tre hơi nâng lên, rồi bất ngờ đập xuống!

*PẶP!*

Không trúng dưa.

Mà chiếc gậy suýt văng trúng Ryuuhou, người đang lấy thêm miếng bánh đi ngang phía sau.

``\textbf{Ê ê! Suýt nữa trúng vai em rồi đó!!}'' cô bé lùi lại mấy bước, mắt mở to, suýt làm rơi cả chiếc bánh kem dừa trên tay.

Ayame vội gỡ khăn, mặt tái mét: ``Á á\dots\ xin lỗi xin lỗi! Chị không cố ý!''

Ông anh trai của Ryuuhou đứng từ xa cũng chỉ nhếch miệng cười: ``Xém nữa là hại chết em gái của anh rồi đấy.''

Roboco vội chạy lại, vỗ vai Ayame: ``Không sao đâu, lần sau chị sẽ chỉ kỹ hơn cho em!''

Haato cười xòa: ``Đúng rồi, lần sau nhớ nghe kỹ hiệu lệnh nha!''

Choco cũng dịu dàng: ``May mà không ai bị thương, Ayame-sama. Đừng có vung gậy xuống đường đột thế chứ.''

Sau khi hội ý chớp nhoáng, ban giám khảo gồm Fubuki, Choco và Okayu đưa ra quyết định:
``Người thắng vòng này là Ayame(-chan/-sama)! Dù không trúng dưa, nhưng tinh thần kiếm sĩ và sự bình tĩnh được chấm điểm cao!''

Nàng Quỷ sừng reo lên: ``Yatta!'' hớn hở cắp quả dưa hấu lên nách rồi vui vẻ chạy ra phía Kuon.

Polka từ dưới đất rên rỉ, nửa tỉnh nửa mơ: ``U\dots\ tui\dots\ tui muốn đổi lượt\dots\ đập lại\dots\''

Tiếng cười lại vang lên khắp khoảng sân tràn ngập cát, khi tiết mục ``Đập dưa phiên bản \hll'' chỉ mới bắt đầu.

Cả nhóm cùng bật cười, không khí lại trở nên vui vẻ và hào hứng tiếp tục lượt chơi tiếp theo.

Sau màn đầu tiên đầy hỗn loạn nhưng cực kỳ vui nhộn, các cô gái bắt đầu có chút dè chừng, vừa lo lắng cho Polka, vừa sợ\dots\ đến lượt họ ``gây án.''

Thế nhưng, Fubuki chẳng để mọi người kịp chùng xuống, cô vỗ tay bôm bốp tập trung: ``Tiếp theo! Đội ngũ tinh anh bước ra! Nào nào, Roboco-senpai, Aki-chan và Botan-chan!''

Cả ba người được gọi tên, người lóc cóc chạy lên, người vươn tay ra sau khởi động, người không ngần ngại vác chày lên.

``Ủa, em tưởng `senpai' không chơi mấy trò trẻ con này chứ?'' Subaru chen ngang, nhưng bị Sora kéo ra một bên: ``Shhh! Coi đi rồi biết, Subaru-chan!''

``Trò này ai chơi chả được hả em,'' Kuon bật cười. ``Làm như thể chỉ mình em chơi trò ấy.''

\ct

Roboco là người mở màn cho vòng thứ hai. Với vẻ ngoài dễ thương pha lẫn vẻ bí hiểm của một trí tuệ nhân tạo, cô gái Rô-bốt tiến ra với nụ cười nhẹ và chiếc khăn đen được buộc chặt ngang mắt.

``Và không được chơi bẩn đâu đấy!''

``Rồi\~{}''

Roboco xoay đúng ba vòng và dừng lại, hai chân vẫn còn đứng vững trên nền cát.

``Cảm biến con quay hồi chuyển ổn định thật sự!'' Okayu buột miệng, khiến mọi người phá lên cười.

``Cơ mà để xem tên PON này làm được trò trống gì không đã\dots'' Kuon khoanh tay, dõi theo mọi hành động của cô nàng Người máy.

Roboco bắt đầu bước về phía trước, từng bước đi chính xác đến kỳ lạ. Chưa cần dùng đến sự trợ giúp của bất kỳ ai, cô dừng lại đúng trước quả dưa hấu, chỉnh lại tư thế tay, rồi đập xuống không một giây do dự.

*ĐOẸT!*

Một nhát trúng giữa tâm. Quả dưa vỡ đôi hoàn hảo, nước văng nhẹ ra hai bên, những miếng nhỏ bắn ra, đáp xuống hẳn bên ngoài.

Cú đập ``ngọt sớt'' đến mức làm cả Polka - người vừa chuếnh choáng tỉnh, trán chườm đá lạnh - phải bừng dậy, hét lên từ bên ngoài: ``\textbf{Chính xác đến từng mi-li-mét!!!}''

Tuy nhiên, ngay sau đó Subaru nheo mắt nghi ngờ: ``Khoan, có khi nào chị ấy dùng chức năng định vị gì không!? Định chơi bài gian lận hả!?''

Roboco vội gỡ khăn, tay giơ lên tỏ vẻ vô tội: ``Không hề! Chị tắt ra-đa đi rồi! Chị hoàn toàn dùng cảm tính luôn!''

``Cảm tính kiểu gì mà chính xác vậy\dots\ Robo-PON có cả mô hình thuật toán tới từ tương lai sao\dots?'' Kuon nhíu mày châm chọc, khiến cho khổ chủ phồng má lên, còn tất cả mọi người lại bật cười.

Cuối cùng, sau khi tham khảo "VAR" (xem lại), kết quả của Roboco được công nhận, và cô nàng được thưởng chính quả dưa hấu ấy.

``Mà Roboco-onee-chan ăn kiểu gì vậy? Chị là người máy mà nhỉ?'' Ryuuhou đang ăn miếng bánh cà phê vừa mới cắt kia.

``Chị có cách riêng của chị thôi\~{}''

\ct

Người thứ hai chuẩn bị cho lượt chơi của mình là Aki. Nàng Dị giới tóc vàng với nụ cười nhẹ nhàng. Cô buộc khăn lên, trông có chút lo lắng.

``Chị thực sự kém mấy trò này lắm\dots\ Nhưng mà, chị sẽ cố gắng hết sức!''

Korone vẫy đuôi động viên:``Aki-chan cố lên nha! Đừng đập bọn này là được\~{}''

``Em làm được mà!'' Sora tiếp lời, kèm một cái nắm tay cổ vũ.

Aki quay ba vòng hơi chậm, rồi bước từng bước nhỏ. ``Xem nào, quả dưa hấu chỗ nào nhỉ\dots\ Có phải là đằng trước không nhỉ?''

Kiên định với suy đoán của mình, Aki đã nhanh nhẹn tiến về phía trước, hoàn toàn lệch hướng sang phải. Vừa đi, cô vừa dò đường bằng chiếc gậy của mình.

*CỘC CỘC*

``A-Aki-senpai!?'' một giọng nói hốt hoảng vang lên.

Không ai khác, đó là Polka, người vẫn đang ngồi chườm đá sau cú vụt kinh hoàng vừa rồi. Cô nàng hét lên thất thanh: ``\textbf{Nooooo!! Là em đây, Polka đây! Đừng vụt em!}''

May mắn thay, Aki nghe thấy tiếng hét của nàng Cáo sa mạc, vội giật mình lùi ra sau: ``P-Polka-chan đấy à? Chị xin lỗi nhé!''

Fubuki và Rushia liền la lên cùng lúc: ``Bên trái đó! Xoay qua bên trái đi!!''

Aki lập tức chuyển hướng. Vừa lúc đó, chân cô đụng phải gì đó tròn tròn.

``A, đây rồi! Chắc chắn là nó!''

Cô dừng lại, chậm rãi vung tay\dots

*POẸT!*

Dưa hấu vỡ tan thành mảnh, phần ruột đỏ hồng văng ra đẹp mắt như cảnh phim quảng cáo nước ép.

Aki gỡ khăn ra, mắt sáng rỡ khi cô đập trúng mục tiêu: ``À-rế? Trúng thật rồi này!''

``Yeahhh! Chị làm được rồi đó, Aki-senpai!'' Ayame vỗ tay reo hò.

Trong khi đó, Polka trốn hẳn vào trong đám đông, thì thầm: ``Xíu nữa thì thành\dots\ dưa phụ\dots''

\ct

Cuối cùng là Botan, nàng Sư tử trắng lạnh lùng với kỹ năng siêu việt. Khi nghe tên mình, cô chỉ khẽ cười, cột khăn lên một cách gọn gàng như thể chuẩn bị vào đấu trường bắn súng.

``Đừng nói chị ấy dùng\dots\ cảm nhận âm thanh nha\dots'' Subaru lẩm bẩm, bắt đầu khiếp sợ sau lượt chơi của hai người trước đó. ``Đến em cũng không làm được như thế.''

Luna đưa gậy cho Botan, nhỏ giọng đùa: ``Níu chúng được đầu dưa, có thể em sẽ đựcc xưng nàm bà hoàng xạ thủ á-nora\~{}''

Botan chỉ đáp bằng một tiếng cười trầm: ``Vậy thì phải headshot rồi.''

Ba vòng xoay, không có gì quá khó với cô. Cũng giống như Roboco, lần này coo không cần ai hướng dẫn, Botan bước lên ba bước và dừng lại như đã tính toán từ trước.

``Trò này dễ ợt. Hây!''

Cô vung chiếc gậy đầy mạnh mẽ.

*TOẸT!*

Quả dưa hấu thứ ba vỡ tung trong tiếng hò reo của khán giả. Không chệch một hướng, cũng không bước thừa một nhịp.

``\textbf{Quả không ngoa khi nói Botan-chan là nữ hoàng FPS!}'' Fubuki hét lớn. ``Chúc mừng em!''

``Đây gọi là gì nhỉ\dots\ `One tap' đúng không?'' Flare hỏi hậu bối của mình. ``Chị để ý em hay làm thế suốt, bắn hạ bọn em chỉ trong một nốt nhạc thôi à.''

Nàng Sư tử trắng gỡ chiếc khăn, nói đùa: ``Có thể hiểu như vậy. Cái trò này dễ như ăn đạn vào đầu đối thủ thôi.''

Với hai cú đập chính xác liên tiếp từ Roboco và Botan, và một pha ``lật kèo'' đầy ngọt ngào từ Aki, vòng hai kết thúc trong tiếng cười rộn ràng.

Kuon nhìn số dưa còn lại, rồi sau sáu lượt chơi, cậu hỏi: ``Em có cách nào để đưa trò chơi khó hơn một chút không?''

Fubuki cười khanh khách: ``Yên tâm\~{} vòng ba sẽ vui hơn nhiều! \textit{Và đây cũng sẽ là dịp để em trả thù anh vụ đua xe ban nãy!}''

\ct

Trò chơi đập dưa bước sang lượt cuối cùng, và như mọi lần, Fubuki luôn biết cách làm mới cuộc vui.

Nàng Cáo tráng vỗ tay đầy phấn khích để gây chú ý tới mọi người, rồi cất giọng cao vút như đang công bố một sự kiện long trọng: ``Lượt này sẽ là Winner takes all (Ai thắng thì lấy tất)! Có năm quả dưa, ai đập trúng được ba trước thì giành hết nha!''

Mọi người cùng ``ồ'' lên đầy hứng khởi với cách chơi mới mẻ này.

Ngay lập tức, Miko giơ tay hét lớn: ``Ể!? Thắng thì được lấy hết á!? Vậy là ai thắng thì ăn hết dưa luôn hả, ne!?''

Watame ôm má, mắt sáng rỡ: ``Nếu vậy thì phải cố gắng hết sức thôi! Chúng ta sẽ ăn hết dưa luôn cho mà xem!''

Matsuri thì cười gian, vung gậy lên (dù biết rằng cô nàng chưa tới lượt): ``Hehe\~{} Winner takes all là đúng gu chị rồi! Chuẩn bị tinh thần đi nha, dưa này là của chị hết!''

Trong những tiếng hò reo đó, nàng Cáo rắng ngoắc Mio ra một góc, khẽ thì thầm bàn kế hoạch của mình: ``\hl{Này Mio, tôi sẽ đặt cả năm quả vào cùng một chỗ. Coi như là một món quà nhỏ cho Marine-senpai á\~{}}''

Mio nhíu mày, lắc đầu: ``Fubuki\dots\ đừng chấp nhặt chuyện Kuon-senpai thắng bà vụ M*r*o K**t ban nãy nữa. Trò chơi này đâu có liên quan gì?''

Fubuki nheo mắt, nhún vai: ``Yên tâm\~{} Không có ác ý gì đâu. Trái dưa vẫn còn nguyên đó, người vẫn nguyên vẹn mà!''

Mio chỉ biết ngán ngẩm đáp lại: ``Bà đúng là chỉ giỏi lý sự\dots\ Bà mà bị nghiệp quật là tôi bỏ mặc đấy nha.''

``Rồi\~{}'' nàng Cáo trắng đáp, giọng đầy ngây ngô nhưng nguy hiểm.

Khi tất cả đã chuẩn bị xong, Fubuki hô lớn, hoàn toàn tạm ngó lơ tới kế hoạch phục thù của mình. ``Rồi, rồi! Cặp đầu tiên sẽ là Marine-chan và\dots\ Ryuuhou-chan! Mời lên!''

Nàng Thuyền trưởng giơ tay vẫy đầy khí thế: ``Tuyệt quá! Em sẽ cho mọi người thấy thực lực của Senchou như thế nào!"'

Cô em gái của cậu Tổng thì tròn mắt nhìn xung quanh: ``Hể? E-Em đấy ạ? Tại sao lại là em?''

Nàng Cáo trắng nhắm mắt lại, cười khẩy rồi đột nhiên chỉ thẳng vào cô: ``Vì đó là em chứ còn sao! Chị thấy em mang sát khí tiềm ẩn đấy! Lên đi!''

``Ể? Nhưng em có phải là Rushia-onee-chan đâu-''

``Cứ lên đi!'' Fubuki tiếp tục luôn miệng, giọng nói của cô như áp đặt.

Ryuuhou chỉ biết thở dài, ngoan ngoãn tiến lên giữa ánh mắt đầy thương cảm nhưng phấn khích từ các thành viên.

Kaien đứng dậy khỏi chiếc ghế nghỉ, dặn dò: ``\vie{Con gái, cứ bình tĩnh mà chơi. Nhắm mắt cũng phải nghe bằng tai, hiểu chưa?}''

Kuon thì nhấc nhẹ tay vỗ vai em gái: ``\vie{Đừng để mình rơi vào bẫy mọi người. Tập trung cố gắng hết sức là được.}''

Ryuuhou gật đầu, lấy lại một chút tự tin: ``\vie{Em sẽ đập đến vỡ dưa cho hai người xem, Okaa-san, Onii-san!}''

Cùng lúc đó, Marine thì đang được các chị Gen 3 dốc toàn lực hướng dẫn:

``Bà nhớ duỗi thẳng lưng, rồi vung mạnh từ bên phải nha, peko!'' Pekora la lớn.

``Đừng đập sớm quá! Dò mặt đất trước đã!'' Noel nói, tay làm động tác mô phỏng cực kỳ nhiệt tình. ``Như tôi vừa làm đấy!''

``Ừm, biết rồi!'' nàng Thuyền trưởng gật đầu, khí chất bừng bừng tiến về khu vực thi đấu.

Cả hai cùng bịt mắt, xoay ba vòng, và trò chơi bắt đầu.

Fubuki tranh thủ lúc mọi người không để ý, cô nhảy vào khu vực thi đấu và bốc những quả dưa hấu được xếp trước đó lên, tính toán ước lượng để Senchou có được lợi thế.

``Dồn hết dưa về phía trái\dots\ Đúng ngay\~{} chỗ Marine-chan xoay mặt khi xong lượt quay.''

Cô nàng nhảy về khán đài, hí hửng xem thành quả của mình, trong khi Mio nhăn mặt, rõ ràng không đồng tình với hành vi đó.

``Ể? Ai đưa hết dưa hấu về một chỗ vậy?'' Towa thắc mắc, chỉ về lô dưa hấu được xếp gọn lệch bên trái so với Marine.

Chưa để ai nói thêm điều gì, Fubuki lớn giọng tuyên bố: ``Được rồi! Bắt đầu trò chơi nào!''

Ngay khi Marine hoàn thành vòng xoay cuối, cô chao đảo một chút, rồi cất tiếng rên khe khẽ: ``Ôi\dots\ chóng mặt quá\dots''

Ryuuhou cũng không kém cạnh hơn, cô bé hơi tập tễnh nhẹ ngay khi hoàn thành cú xoay.

Cả hai người rón rén bước về trước vài bước, đầu lắc lư như đang cố nghe ngóng về phía khán đài.

``Sang bên trái một chút đi Senchou ơi!'' Aqua nói, vẫy tay quét sang trái.

``Quay một vòng khoảng 30 độ nữa! Không! Quá rồi, sang phải một tí!'' Kuon cũng không kém cạnh.

Cả hai người vừa cầm chày, vừa dỏng tai nghe khán giả, vừa rê chày xuống bãi cát.

``Mấy quả dưa hấu đâu hết rồi!?'' nàng Thuyền trưởng rít lên, tỏ rõ khó chịu. ``Đi hết cả khu rồi mà chả thấy đâu cả!''

``Onii-san\dots\ Sora-onee-chan\dots\ Dưa hấu ở đâu vậy!?'' cô em gái của cậu Tổng hoàn toàn mất phương hướng, bắt đầu sợ sệt.

Cả hai người cứ lò mò đi trong bóng tối, cho đến khi\dots

*CỘP!*

``Có cái gì cộm cộm dưới chân này\dots''

Đó chính là kho dưa bí mật mà Fubuki đặt, và nó ngay dưới chân của Marine. Cô khẽ cười, nâng cây gậy lên\dots

``\textbf{Đống kho báu này là của ta!! Về đây với mẹ nào!!}''

*KHỰC!*

Một tiếng động mạnh vang lên cùng với cái chày bị rơi xuống. Không phải là tiếng dưa hấu được đập, càng không phải là tiếng vụt trượt.

Tất cả mọi người, trừ Ryuuhou, đều đồng loạt há hốc miệng với cảnh tượng trước mắt, kèm với những tiếng xì xào bàn tán bên dưới. Nhưng rồi Kuon và Kaien dần lấy lại bình tĩnh để tiếp tục chỉ đạo.

Marine bỗng chốc gập người xuống, hai tay ôm lưng, miệng rên rỉ: ``AAA!! Lưng tôi\dots\ Nguyền rủa cái thân già của tôi!!''

Fubuki chết đứng, hoàn toàn không ngờ được tình huống éo le này, tới độ không nói nên lời. Mio đứng bên kia lẩm bẩm, tặc lưỡi: ``Thấy chưa\dots\ gian lận không có hậu quả tốt đâu.''

Senchou không thể đứng dậy được, chỉ có thể rên khe khẽ trong đau đớn, buộc phải để Pekora, Noel và Flare chạy đến nâng cô dậy.

``Marine-chan!!'' Pekora thất thanh.

``Bà ổn chứ!?'' Flare sà tới hỏi han. ``Nào, bà có đứng dậy được không?''

Giữa lúc đó, Ryuuhou bước chậm rãi. Cô lắng nghe tiếng chỉ hướng từ xa của mẹ và anh trai. Tất cả mọi người sau vụ tai nạn bất ngờ đó đã\dots\ quay sang cổ vũ cho cô.

\begin{dialogue}
    \speak{Kuon} Đi thêm hai bước nữa, hơi lệch trái một chút!
    \speak{Kaien} Giờ quét gậy nhẹ qua đất, con sẽ cảm nhận được mặt dưa. Cứ bình tĩnh, đừng đập vội!
\end{dialogue}

Cô làm đúng từng lời chỉ giáo. Gậy bóng chày của cô va nhẹ vào một vật gì đó mịn màng và tròn trịa. Cô thậm chí còn dùng gậy soi xét thật kỹ.

``{\textbf{Chính nó!}}''

Khỏi cần nói thêm một lời nào, Ryuuhou  hít sâu, rồi giơ chiếc chày lên.

*BỐP!* Quả đầu tiên vỡ đôi.

Fubuki há hốc mồm với những gì cô thấy. Kế hoạch của cô đã phá sản.

``Lùi lại một bước rồi sang phải một chút\dots\ Rồi, đập đi!'' giọng Kuon vang lên.

Cô em gái của cậu nghe theo, di chuyển, và vụt lần nữa.

*BỐP!* Quả thứ hai cũng đã tan tành.

``Con giỏi lắm!'' Kaien nói to đầy tự hào.

``A! Còn quả nữa này!'' Ryuuhou di chiếc chày tới một quả tròn nữa.

Fubuki thấy tình hình nguy cấp, liền cố chen vào như đang đánh lạc hướng để câu thời gian: ``Ơ, hình như đó là\dots\ bóng đổ của một người đó nha? Cẩn thận nha, trông giống Watame lắm đó!''

Nhưng cô bé không hề dao động. ``Em chỉ tin những gì mẹ em và Onii-san nói thôi!'' câu nói đó như một đòn chí mạng về phía nàng Cáo trắng xảo quyệt đó.

Cô dò bằng đầu gậy, vỗ nhẹ trước khi giơ lên.

``Nó đó!'' Aki thốt lên.

*BỐP!* Quả thứ ba cũng đã được đập trúng mục tiêu. Ngay sau đó, tiếng mọi người hò reo chúc mừng vang lên cùng những tràng pháo tay rầm rộ.

Ryuuhou kéo khăn bịt mắt ra, mặt đỏ bừng nhưng ánh mắt rực sáng: ``Em thắng rồi hả?''

``3/5! Tuyệt đối luôn!'' Mio công bố. ``Chúc mừng em!''

``Tuyệt quá Ryuuhou-chan!'' Aqua nhảy vào hộp cát, ôm chầm lấy cô bé. ``Em đánh bại được Marine rồi!''

Marine vẫn đang được dìu ngồi bên cạnh, cười méo xệch tiếc nuối: ``Mình\dots\ chiến đấu bằng cả tấm lưng này\dots\ và tôi đã thất bại\dots''

``Không sao đâu Marine-chan! Ta sẽ thắng lần tới!'' Noel lên tiếng động viên.

``Để tôi lo cho bà ở lượt sau cho!'' Flare hòa âm.

Ở một góc khán đài, Fubuki thở dài một tiếng, dường như chưa chấp nhận được thất bại: ``Đúng là\dots\ đời không như là trong trò chơi thật\dots''

Mio vỗ vai cô bạn cáo: ``Thấy chưa? Mưu mô cũng không thắng nổi sự phối hợp gia đình đâu!''

``Ư\~{} Tức thật chứ! Hãy đợi đó, Kuon-senpai!''

\ct

Tưởng rằng trận đấu giữa Marine và Ryuuhou là màn hạ màn hoàn hảo, ai ngờ một tiếng hô dõng dạc lại vang lên từ mép sân: ``\textbf{Gượm đã!}''

Mọi ánh nhìn đều lập tức đổ dồn về cô nàng Cáo trắng của Thế hệ thứ Nhất. ``Có chuyện gì vậy, Fubuki-senpai?'' Subaru hỏi.

Fubuki khoanh tay, bản mặt hiện rõ một sự nghiêm túc. Đến mức đáng ngờ.

Và rồi, không cần chờ bất cứ ai lên tiếng, cô tuyên bố một cách hùng hồn: ``\textbf{Em muốn thách đấu với Kuon-senpai!}''

Mio suýt phun nước đang uống, thảng thốt gọi cô bạn: ``F-Fubuki! Bà vẫn chưa chịu từ bỏ à!?''

Kuon đứng phía sau, một tay đút túi quần, mắt nhắm hờ như thể đã thấy trước được chuyện này từ trước. ``Anh biết kiểu gì em cũng sẽ nói vậy mà. Định trả thù việc anh cướp mất chiến thắng trong cái trò M*r*o Kart ban nãy sao?''

Nàng Cáo tuyết giật mình khi bị cậu bắt bài: ``Ủ-Ủa!? Anh biết rồi sao?''

``Em tưởng anh mù chắc? Anh thấy em nhảy vào sân cát để đưa hết dưa hấu vào một cụm, xong rồi còn cố gắng đánh lạc hướng em gái anh còn gì. Ai mà chả đoán được.''

Fubuki vội xua tay: ``Ơ không không, cái đó là\dots! Em chỉ muốn thử tài năng thực chiến thôi mà\~{} Hoàn toàn không phải vì cay cú đâu\~{}!''

``Cái giọng của em lại cho anh điều ngược lại đấy\dots'' cậu thở dài. Từ phía xa, nàng Sói đen lắc đầu ngán ngẩm với cô bạn hiếu thắng kia: ``\textit{Tôi lạy bà đó Fubuki\dots\ Đừng có chấp nhặt mấy cái trò như vậy nữa!}''

\ct

Dẫu nói gì thì nói, lời thách đấu cũng đã được đưa ra, Kuon đành bước lên.

Nhưng lần này, trước khi cậu cầm chiếc gậy lên, cậu đã nhấc tay lên ra hiệu mọi người: ``Lần này anh không cần ai chỉ đường. Tai của anh coi như điếc rồi.''

``Ý anh là sao vậy, Kuon-senpai?'' Botan tò mò hỏi.

Cậu lấy ra cặp tai nghe đeo lên, rồi bật nhạc trò chơi quen thuộc ở mức âm lượng lớn đến mức đủ để bên ngoài nghe được.

Tiếng nhạc lớn như sấm nện vào màng nhĩ khiến cho tất cả mọi người trố mắt: ``Ê\dots\ Ê!? Thật đấy à?''

``Anh không thấy đau tai sao, Kuon-senpai!?'' Towa trố mắt, sững sờ hỏi.

Cậu đọc khẩu hình miệng của nàng rồi giơ ngón cái, ý rằng muốn nói: ``Không sao.'' Cậu lấy chiếc khăn và tự bịt mắt lại.

Fubuki nhìn cảnh tượng Kuon tự làm khó bản thân mình, nhếch mép đắc thắng: ``Hừm\dots\ Kuon-senpai chơi thế là dở rồi\~{}''

Cô chậm rãi đeo bịt mắt, nhưng chỉ kéo đến gần hết mắt phải. Một bên mắt khẽ hé ra, chỉ đủ để cô nhìn thấy bóng người xung quanh và. vị trí những quả dưa hấu nằm ở đâu.

``\textit{He he\dots\ Đây gọi là `cân bằng chiến thuật' đây\~{}}''

Trận đấu thứ hai của vòng thứ ba bắt đầu.

Cả hai xoay ba vòng, không hề nao núng, cũng không có điều gì bất trắc xảy ra. Quay xong, họ bắt đầu dùng những gì họ có để bắt đầu dò tìm những quả dưa hấu.

Với một bên mắt được lộ ra khéo léo, nàng Cáo tráng dễ dàng tìm được đường đi nước bước. Cô nàng vừa cười cành cạch trong lòng vừa nghĩ: ``\textit{Với ánh mắt của loài cáo này, chiến thắng sẽ về tay mình thôi!}''

Ở phía bên kia, Kuon hoàn toàn dựa vào cảm giác để xác định phương hướng. ``\textit{Chỗ này hơi nhô lên thì phải\dots\ Không, đó là cái cột nhà. Dưa hấu chắc ở đâu đó quanh đây\dots}''

Chẳng bao lâu sau khi trò chơi bắt đầu, Fubuki đã lần ra được một quả dưa hấu ngay trước mặt. Cô nàng hí hửng nghĩ thầm:
``\textit{Dễ quá! Thêm một bước nữa là tới rồi!}''

Nhưng ngay khi cô tưởng rằng quả dưa đầu tiên sẽ thuộc về mình\dots

*PHỤP!*

\dots một thứ gì đó trơn tuột nằm ngay dưới chân cô.

``\textbf{ÁÁÁÁÁÁÁÁÁÁÁÁÁÁÁÁÁ!!!}''

Fubuki trượt ngã ra sau, rơi cả gậy, chân tay giãy đành đạch trong tư thế hết sức kỳ quặc.

Trong lúc lơ mơ mở mắt, cô nhận ra một vỏ chuối (không biết từ đâu xuất hiện) đang nằm chềnh ềnh ngay đó.

Cô hét toáng lên: ``Ai để thứ này ở đây vậy!!?''

Mio nhìn chiếc vỏ rồi thở dài: ``Chắc là Luna ăn xong lại vứt bừa ra đấy\dots''

Luna ở đằng xa vừa nhai hoa quả vừa quay lưng lại khi nghe ai đó gọi: ``Hả? Gọi em zì thê-nora?''

Trong khi Fubuki còn đang lồm cồm bò dậy và quờ quạng tìm gậy, Kuon vẫn tiếp tục dò dẫm từng bước cẩn trọng. Không nghe thấy gì ngoài nhạc đập ầm ầm trong tai, cậu chỉ còn trông cậy vào cảm giác ở bàn chân và chiếc gậy trên tay để định vị mọi thứ.

Trông cậu chẳng khác gì một người mù vừa nghe nhạc vừa dò đường bằng cảm tính.

``Chỗ này cũng không có\dots\ Thế thì nó ở đâu nhỉ?''

Cây gậy sượt nhẹ qua mặt đất, rồi huơ lên không trung, dò xem có gì bất thường không.

*BỤP*

Một vật gì đó mềm mềm vừa chạm vào đầu gậy. ``Chắc là nó đây rồi.''

Không ngờ, đó lại chính là cái đầu của một cô gái vừa bị gậy bóng chày ghé thăm hai lần trước đó.

Là Polka.

Lại là Polka.

``Ể!? Không phải lần nữa chứ!!'' Cô hét toáng lên, định la lên cảnh báo\dots

``\textbf{Kuon-senpai! Em không phải dưa hấu! Không phải dưa hấu!!}''

\dots nhưng quên mất rằng hai tai của cậu đang bị nhạc bập bùng tra tấn, chẳng nghe thấy gì.

Polka đứng trân người, mắt mở to như sắp bị hành hình. Chiếc gậy bóng chày lại lơ lửng ngay trên đầu cô, thêm một lần nữa.

Mọi người nín thở khi nạn nhân xui xẻo chuẩn bị đón thêm một đòn trời giáng.

``Tội Polka-chan thật đó\dots\ Bị gậy ghé thăm tới ba lần!'' Sora vừa cười vừa mếu với tình cảnh trước mắt.

``Polka đã hy sinh vì chúng ta quá nhiều rồi\dots'' Nene chắp tay cầu nguyện, ``Chúng tôi sẽ nhớ ơn bà!''

``Ê này!? Đừng có nói gở như thế chứ!!'' nàng Cáo sa mạc hét lớn. ``Mau cứu tôi đi!!''

Ngay lúc này, chiếc gậy bóng chày được đưa lên cao.

``\textbf{Đừng mà!!!}''

Nhưng Kuon khựng lại, dừng ngay trước đỉnh đầu cô. Hai tay cậu giữ nguyên tư thế, rồi nhẹ nhàng dùng đầu gậy chạm vào bề mặt của vật kia.

Tròn, nhưng gồ ghề. Không giống hoa văn trên quả dưa.

Cậu lẩm bẩm một câu mà chẳng ai nghe rõ, nhưng nhìn khẩu hình có thể đoán là: ``Không phải dưa.''

Rồi Kuon quay đi, tiếp tục tìm mục tiêu khác.

Polka ngồi sụp xuống, thoát khỏi cú đánh trời giáng. Cô không hét, không khóc, chỉ còn khuôn mặt trắng bệch như bị hút cạn sinh khí, và\dots\ một dòng nước ấm lặng lẽ chảy xuống váy.

Watame thấy vậy, vội vàng kéo cô ra một bên: ``Thôi thôi, để chị dắt em đi thay đồ\dots\'' Trớ trêu thay, chính nàng Cừu cũng từng suýt thành nạn nhân của con dao cắt bánh lúc trước, và cũng đã ``ra quần'' vì sợ.

\ct

Sau vài phút dò xét cẩn thận cùng vài cú đánh hụt, Kuon cuối cùng cũng đập trúng ba quả dưa theo đúng luật, dù tốc độ không quá nhanh.

Fubuki thì\dots\ hoàn toàn trắng tay. Ngay khi định chơi gian, cô liên tiếp gặp xui xẻo: nào là trượt vỏ chuối, bị quạt thổi cát vào mắt, rồi hắt hơi đánh trật mục tiêu dù quả dưa ngay trước mặt, để rồi bị đối thủ đập mất. Nói cách khác, cô vừa mất gà lại vừa mất cả thóc.

Khi Kuon tháo tai nghe ra, không khí đã lặng đi, mọi người đều biết kết quả đã rõ ràng. Fubuki thì vẫn đang được gỡ chân khỏi tấm bạt dính, miệng không ngừng than vãn.

``Anh ấy\dots\ thắng rồi đúng không?'' Aqua nghiêng đầu hỏi nhỏ, mắt mở to lấp lánh.

``Ừm, mà còn không cần ai chỉ đường luôn\dots'' Aki gật đầu, vừa ngạc nhiên vừa thán phục.

``Đỉnh thật\dots\ Đeo tai nghe với nhạc xập xình thế kia mà vẫn không đập nhầm ai,'' Watame vừa nói vừa liếc nhìn Polka vẫn đang được dìu đi, ``mặc dù suýt thì\dots''

``Tại em là vật thể không xác định mà,'' Noel thì thầm. ``Chứ nếu là chị thì chắc là đập từ đầu rồi.''

Matsuri huých vai cô: ``Ý gì đó hả!?''

Kuon lại gần, chìa tay ra: ``Chơi tốt đấy,'' như một lời xã giao.

Fubuki hậm hực quay mặt đi, rõ ràng vẫn còn dỗi vì kết quả thảm hại.

``Đừng trẻ con nữa, Fubuki. Bà chẳng có tinh thần thể thao gì cả,'' Mio nhăn mặt trách.

``Đúng đó! Người ta đã có thành ý rồi thì phải đáp lại đi chứ!'' Suisei lớn tiếng.

Kuon đặt gậy lên đống dưa, cười nhẹ: ``Nói vậy thôi. Anh cũng chẳng cần đống dưa đó làm gì đâu. Nhường em phần thắng đấy.''

Hành động của cậu khiến mọi người bất ngờ. Chưa ai từng thấy người thắng lại nhường phần thưởng cho người thua, nhất là khi người thua lại còn là một kẻ chơi gian.

``\dots Cái gì?'' nàng Cáo trắng lắp bắp ``Anh\dots\ anh nói là anh nhường à?''

Câu nói ấy như một cú đánh vào lòng tự trọng của cô. Fubuki quay lại, mắt long lanh: ``Thật hả? Anh không cần phần thưởng à?''

Cậu gật đầu: ``Em cần nó hơn. Anh có thể ăn năm quả dưa hấu của Ryuuhou giành được ban nãy rồi. Đừng thù dai vì mấy chuyện nhỏ nhặt nữa.''

``Đáng yêu quá đáng luôn\dots'' Aki ôm má, nghiêng đầu mơ mộng. ``Một quý ông đúng nghĩa\dots''

Matsuri quay sang Suisei: ``Tôi hiểu rồi, gu mấy chị em chúng mình là đây này.''

Suisei khoanh tay, gật đầu trầm tư: ``Ừm\dots\ có tài, khiêm tốn, lại ưa nhìn, mỗi tội hơi béo. Bảo sao\dots''

Fubuki cúi đầu, tóc mái che đi ánh mắt đang long lanh. ``\dots Vậy tôi có thể gõ đầu cậu một cái không?''

Cả nhóm đồng loạt ``\textbf{HẢ???}''

``Ể!?''

Aqua suýt ngã ngửa: ``Từ nhường dưa chuyển sang gõ đầu là sao!?''

``Kiểu logic gì vậy trời\dots'' Mio ôm trán.

``Chị ấy vẫn không chịu thua thật mà\dots'' Subaru bật cười.

Kuon đứng đó, nửa muốn cười, nửa như đang đắn đo xem\dots\ có nên cho một cú không.

``Cứ việc.''

``Vậy thì em không khách sáo nữa nha\~{}'' Fubuki không do dự, cầm chày lên.

*BỐP!*

Một cú gõ ``nhẹ'' bằng tay, nhưng nghe rõ là\dots\ chẳng hề nhẹ chút nào.

``AAAA!'' cậu ôm đầu, lùi lại một bước.
``Đó mà là nhẹ hả!?''

Fubuki nhe răng cười:

``Ờ thì\dots\ cáo nhẹ tay đó mà\~{} Coi như huề nha. M*r*o K**t, dưa hấu, hết nợ rồi. Tha cho anh á\~{}''

``Có mà cáo lươn thì có,'' cậu bĩu môi, giả vờ tỏ vẻ mặt đáng thương, nhưng rồi cũng bụm miệng cười.

Sau đó, trò chơi đập dưa vẫn tiếp tục diễn ra với ba cặp đấu khác, nhưng lần này chỉ được kể ngắn gọn hơn do không có nhiều biến cố quá đà xảy ra.

\ct

Thay vì Fubuki, người đang đánh chén một ngon lành món dưa hấu, Mio sẽ là người cầm cây nảy mực điều khiển trận đấu (một phần vì mọi người không tin tưởng nàng Cáo trắng gian xảo kia, điều mà cô cũng chẳng để tâm).

Towa và Nene bước vào ``đấu trường'' đầy dưa hấu, mỗi người cầm một cây gậy nhỏ do Mio phát. Khác với vẻ ngoài ngầu lòi thường ngày, nàng Ác ma lại tỏ ra khá căng thẳng, liên tục quay sang hỏi: ``Đập nhẹ thôi có tính không vậy?'' còn Nene thì tràn đầy năng lượng, vừa bịt mắt vừa nhảy chân sáo quanh sân, tới mức đôi ba lần vấp vào quả dưa ngay khi trò chơi còn chưa khai màn.

``\textbf{Bắt đầu!}'' Mio hô lớn.

Towa vừa xoay đủ ba vòng đã loạng choạng như con rô-bốt mất thăng bằng, rồi vung gậy đập một cú trúng ngay vào hộp quạt hơi nước đặt ở góc tường.

Tiếng *BỤP!* vang lên, nước bắn tung tóe, làm một bên váy của Suisei đứng gần đó bị ướt một mảng.

``Towa-chan!!'' nàng Ca sĩ tránh ra một bên, giật mình với cú tai nạn vừa rồi.

``Xin lỗi chị mà!'' Towa cuống cuồng cúi đầu xin lỗi, dù cô còn đang quay ở sai hướng.

Còn Nene thì vừa xoay được một vòng rưỡi đã ngã bẹp xuống chiếu bãi cát vì chóng mặt. Khi được đỡ dậy, cô chỉ lẩm bẩm: ``Dưa nó quay, không phải em\dots''

Tới cuối cùng, Towa giành chiến thắng sau những cú đập đầy chật vật, còn đối thủ thì đập được một quả sau một lần ăn may.

Mio từ xa giơ tay: ``Và giờ thì hết trò rồi! Tất cả mọi người chuẩn bị dọn dẹp đi!''

\ct

Tiếp theo là cặp đôi Pekora và Miko, hai người vốn đã có ``ân oán'' từ những lần chí chóe nhau trong Mi**cr*ft.

Pekora bước vào trận với khí thế của một chiến binh thực thụ, tay cầm gậy như chuẩn bị ra trận. Miko thì vừa đi vừa vung gậy múa may, miệng không ngừng cười khúc khích. Cả hai bịt mắt, được dẫn ra giữa sân và xoay vòng dưới tiếng reo hò cổ vũ của mọi người.

Ngay khi tiếng ``\textbf{Bắt đầu!}'' vang lên, nàng Thỏ đã hét lớn: ``\textbf{Miko-chan, đừng có đập vào tôi đó, peko!}''

Nàng Vu nữ thì đáp lại tỉnh rụi: ``Chị đập là đập dưa mà\~{} không phải chị đâu, ne ha ha!''

Nhưng đến cuối cùng thì\dots\ chẳng ai nghe lời hướng dẫn cả. Hai người cứ mải đuổi theo tiếng chân nhau, vừa đi vừa vung gậy loạn xạ, hoàn toàn ngó lơ mục tiêu chính. Pekora suýt nữa thì đập trúng cái đèn treo trên trần nhà, còn Miko thì vung gậy trúng ngay cái gối ôm hình Korone mà ai đó để lạc chỗ, rồi tự tin hô to: ``\textbf{Trúng rồi!}''

Korone chỉ tay hét lớn: ``Đó là cái gối của em mà!!'', mặc cho nàng Vu nữ cầm chiếc chày đập liên tục vào chiếc gối.

Cả phòng cười ồ lên, còn Pekora thì vừa thở phào vừa càu nhàu: ``Đúng là Miko-chi, lúc nào cũng phá game peko\dots''

``Đến cuối cùng Miko-chi vẫn là PON thôi nhỉ?'' Sora cười trừ.

Miko thì vẫn cười toe, ôm lấy cái gối của nàng Khuyển như vừa lập được chiến công hiển hách, còn chủ nhân của nó thì chạy lại giành lại ``nạn nhân'', vừa cười vừa giả vờ giận dỗi.

Kết quả chung cuộc, không ai giành chiến thắng, nhưng tiếng cười thì vang khắp căn phòng.

\ct

Cặp đôi cuối cùng là cặp đôi Hành-Tỏi - Aqua và Shion, đôi bạn cùng Thế hệ nổi tiếng với những màn châm chọc và ``tán tỉnh'' nhau không hồi kết.

Aqua bước vào sân với vẻ mặt nghiêm túc hiếm thấy, tay nắm chặt cây gậy như chuẩn bị ra trận. Shion thì trái lại, vừa cười khúc khích vừa lắc lư cây chày gỗ, miệng ngân nga bài hát ``Mage of Violet'' như thể đang ở trong một cuộc dạo chơi chứ không phải chuẩn bị thi đấu.

Cả hai cùng bịt mắt, xoay vòng xong thì bắt đầu lao đi như thỏ rừng tránh lũ. Tiếng hướng dẫn từ khán giả vang lên loạn xạ: ``Trái!'', rồi ``Phải!'', ``Lùi lại!'', ``Chạy tới đây này!'', ``Dừng lại!'' và còn ``Đập đi!'', hết từ những thành viên cùng thế hệ lại tới các đàn chị đàn em khác, khiến cả hai càng thêm rối loạn, cứ xoay vòng vòng mà chẳng biết đâu là hướng đúng. Aqua dò dẫm tiến về phía trước, nghe theo tiếng hô của mọi người, bất ngờ vung gậy thật mạnh.

*BỐP!* Một quả dưa hấu vỡ toang, nước bắn tung tóe khắp sàn.

Nhưng ngay sau đó, Kaien từ phía xa hét lên: ``\textbf{Ấy, đó là quả dưa để sáng mai cô ép nước mà!}'' Aqua nhảy dựng lên, mặt méo xệch: ``Tôi lạnh quá trời ơi! Sao lại để dưa ở đây vậy trời!'' Đó hoàn toàn là quả dưa hấu bên ngoài khu vực thi đấu, không phải là bất kỳ quả nào trong sân cát kia.

Shion thì loay hoay mãi mới tìm được vị trí, nhưng khi chuẩn bị vung gậy thì trượt chân, ngã bệt xuống sàn, rên rỉ: ``Em đau mông rồi\dots\ Đúng là không có duyên với trò này mà\dots'' Với việc cô nàng không thể dùng phép thuật do hết ma lực, cô trông không khác gì một người bình thường cả.

Không khí trở nên náo nhiệt hơn bao giờ hết khi cả hai cùng lúc phát hiện ra một quả dưa còn sót lại ở giữa sân. Cả hai lao tới, gần như cùng lúc vung gậy xuống.

*BỐP!* Hai tiếng đập vang lên gần như đồng thời, nhưng chỉ có một người đập trúng tâm quả dưa, còn người kia chỉ sượt qua mép.

Sau một hồi kiểm tra, Mio - trọng tài bất đắc dĩ - ra tuyên bố: ``Người chiến thắng là\dots\ \textbf{Aqua!} Shion chỉ sượt qua thôi, còn Aqua đập trúng giữa luôn!''

Nàng Hấu gái nhảy cẫng lên ăn mừng, còn nàng Phù thủy thì phụng phịu bỏ chiếc bịt ra: ``Lần sau tôi sẽ phục thù! Hành thắng tỏi lần này thôi nhé!''

Trận chiến dưa hấu khép lại trong tiếng cười rộn rã, kết thúc một ngày sinh nhật đầy ắp kỷ niệm và niềm vui.

\ct

Đồng hồ đã chỉ về gần mười hai giờ đêm. Chỉ còn ít phút nữa là ngày sinh nhật của Kuon sẽ chính thức khép lại.

A-chan bước ra giữa phòng, hai tay đưa lên vỗ nhịp vài cái để thu hút sự chú ý của mọi người. ``Rồi, rồi! Mọi người nghe tôi nói đây! Trò đập dưa hấu chính thức kết thúc! Vậy là cả bốn hoạt động trong chuỗi sinh nhật của Hikawa-san cũng đã hoàn thành luôn rồi đó!''

Một tràng vỗ tay rộ lên, pha lẫn những tiếng hú hét vui vẻ và mừng rỡ từ tất cả mọi người vang lên khe khẽ, khi những nhà khác đã chìm vào giấc ngủ.

``Và nếu tính luôn cả hoạt động ăn uống lúc chiều tối khi Kuon-senpai tuyên bố khai màn\dots'' Sora tiếp lời, giọng đều đều như đang điểm danh, ``\dots thì tổng cộng chúng ta đã trải qua: một màn cắt bánh sinh nhật đầy bất ngờ và tai nạn\dots''

Tất cả mọi người đều đổ dồn về phía Watame và Suisei, những người vô tình dính vào vụ đó.

``\dots Một buổi hát karaoke sôi động đến vỡ trần nhà, một cuộc đua M*r*o K**t mà không ít người vẫn còn cay cú\dots\ nhìn Fubuki kìa\~{}'' Mio liếc nhẹ về phía hồ ly trắng đang giả vờ huýt sáo, dù rằng cô nàng đã hả giận.

``\dots tiếp theo là tiết mục ảo thuật của Polka-chan khiến mọi người khóc cười lẫn lộn, và cuối cùng là trò đập dưa hấu, hoạt động\dots\ bắn nước và gào thét ồn ào nhất hôm nay!'' nàng Thần tượng kết thúc tổng kết.

``Bà nói đúng ghê, ne\~{}'' Miko bật cười. ``Cái sàn phòng vẫn còn nước nè\~{}''

``Vẫn còn phảng phất mùi dưa hấu quanh đâu đây\dots'' Korone ngửi mũi, gật gù đồng tình.

``Vậy chúng ta cùng bắt đầu dọn dẹp lấy chỗ ngủ nào\~{}'' Okayu giơ nắm đấm lên cao, tỏ vẻ như đang là người chỉ huy.

``\textbf{Đồng ý!}''

Ngay sau lời tổng kết đó, đợt dọn dẹp cuối cùng được Mio và Sora tổ chức nhanh chóng. Các thành viên tự túc chia nhau các nhóm nhỏ: nhóm thu dọn, nhóm lau sàn, nhóm gấp đệm futon, và nhóm lo các thiết bị điện tử.

Sân cát tạm thời trải trong phòng tầng ba được gom gọn lại bằng chổi cứng và bàn tay đầy kiên nhẫn của Korone, Okayu và Ayame, những người gần như ``nằm vùng'' tại khu vực đó trong suốt hoạt động dưa hấu, dù hai trong số họ không tham gia trò chơi. Các mảnh vỡ vỏ dưa, khăn lót, khăn ướt đều được gom vào túi rác phân loại theo yêu cầu của Kaien.

Ở một góc khác, những món đạo cụ mà Polka dùng cho màn ảo thuật được Mio và Aki kiểm tra lại từng món: từ chiếc hộp gỗ, bóng bay trang trí, gậy phép màu giả (mà thực chất là bộ điều khiển) đến bộ bài đặc biệt. Polka cười ngượng khi thấy cái khăn lụa của mình vẫn đang quấn quanh chân bàn vì lúc diễn xong quên tháo.

Chiếc TV cỡ lớn và bộ dàn âm thanh được Shion, Noel và Towa - những người còn đủ sức - hợp lực khiêng lên từng bậc cầu thang về chỗ cũ, trong khi Suisei giữ ổ cắm và gỡ từng sợi dây nối. Aqua chạy theo hô: ``Cẩn thận đó, TV đó nặng lắm nhaaa\~{}!'' còn Ryujin thì theo sau, trực tiếp nối lại dây điện.

Phía bếp của các tầng, đồ ăn còn lại, gồm vài phần bánh sinh nhật còn nguyên lát đã được cắt gọn, dưa hấu chưa bổ, nước ép và vài xiên thịt nướng dự phòng thì được nhóm của Subaru, Choco và Nene chia vào các hộp đựng.

Kuon bước tới, nhẹ giọng nói với mẹ: ``\vie{Cho con giữ lại một ít cho nhà mình, mẹ nhé\dots\ để sáng mai có đồ ăn sẵn cho cả nhà ạ.}''

Kaien mỉm cười, khẽ gật đầu. Cả hai đều hiểu, dù không nói ra, phần đồ ăn ấy sẽ được Kuon lặng lẽ đem sang hai căn phòng nằm ở tầng hai và tầng bảy, nơi hai người khách đặc biệt đang ở ẩn theo chỉ thị từ YAGOO.

Dần dần, ánh sáng trong phòng được hạ thấp. Những tiếng bước chân vội vã đổi thành tiếng trò chuyện nhỏ nhẹ, tiếng gấp chăn, trải đệm. Đêm muộn đã buông xuống tầng ba của ngôi nhà Hikawa một cách yên bình và trọn vẹn.

\il

Và như thế, ngày Mười Bảy tháng Bảy năm đó đã khép lại với biết bao điều không thể nào quên.

Từ buổi sáng với cuộc họp mật lên kế hoạch chớp nhoáng nhưng đâu ra đấy, khi cả nhóm \hll\ cùng nhau với gia đình Hikawa đã bày mưu tính bí mật tổ chức một sinh nhật hoành tráng cho Kuon. Mọi người phân chia nhiệm vụ như một đội đặc nhiệm tinh nhuệ, với tốc độ và sự đồng lòng đáng kinh ngạc, tất cả đã mở màn cho những chuỗi sự kiện náo loạn kế tiếp.

Rồi là cuộc đột kích siêu thị JUSCO Ryuuhen trong buổi trưa không ai ngờ tới, nơi mà những thành viên như Suisei, Noel, Fubuki, bộ đôi Chó-Mèo\dots\ thi nhau tranh giành nguyên liệu làm bánh kem, trái cây, đồ trang trí, hay thậm chí là những hàng thịt tau. Tiếng hò hét, xe đẩy va vào nhau, và tiếng cười giòn tan cứ thế vang vọng khắp các lối đi, biến một buổi mua sắm thành trận ``càn quét thương trường'' có một không hai.

Tiếp đó là cuộc ``xâm lược Kawachi'' bất đắc dĩ nhưng lại kỳ công: những đường chạy vội vã, tiếng cười nói rôm rả, những màn ``cải trang'' của những Haato, Fubuki hay bộ ba Thế hệ thứ Nămm, rồi cả việc vô tình làm xáo trộn tuyến đường Thủ đô và thậm chí\dots\ vô tình giúp lực lượng an ninh khu vực chớp lấy thời cơ ra tay tóm sống đám xã hội đen khét tiếng của Thủ đô - những kẻ sau đó bị đánh cho tơi tả do những cơn thịnh nộ của người dân. Không ai có thể tưởng tượng rằng một ngày sinh nhật lại để lại ảnh hưởng đến cả trật tự an ninh đô thị như thế.

Buổi chiều tiếp diễn với những hoạt động hậu trường đầy hỗn loạn, từ việc bánh kem đa sắc màu của nhóm, đến việc trang trí đầy rực rỡ căn phòng sau một cuộc bàn bạc của nhóm, cả tập thể đã góp phần tô điểm cho bầu không khí đầy ắp sự sống động, tiếng cười và tình cảm chân thành giữa các thành viên. Ai cũng chạy ngược chạy xuôi, mồ hôi đẫm áo, nhưng không ai than phiền. Bởi vì, họ biết rằng mình đang làm điều gì đó đặc biệt cho một người trong lòng họ.

Và rồi, bữa tiệc sinh nhật chính thức bắt đầu. Màn karaoke tưng bừng, cuộc đua M*r*o K**t kịch tính, tiết mục ảo thuật ``ảo đến phát hoảng'' của Polka, và trò chơi đập dưa hấu vừa cười chảy nước mắt vừa đậm tính chiến thuật, tất cả đã kết thành một bữa tiệc khó quên, không chỉ với Kuon mà còn với mọi người có mặt tại đó.

Toàn bộ ba mươi hai cô gái và gia đình Hikawa đã tổ chức một ngày sinh nhật rất to lớn, rất quy mô, và cũng rất thành công. Mọi người đã cùng nhau chuẩn bị, cùng nhau tạo bất ngờ, và cũng cùng nhau hưởng trái ngọt từ đó.

Đó cũng là sinh nhật đáng nhớ nhất mà Hikawa Kuon từng được trải nghiệm trên đời, và cậu biết rằng, rất có thể, sẽ không bao giờ có cơ hội thứ hai cho một ngày như thế nữa.

Một ngày thứ bảy huyên náo đến lạ thường trong căn nhà của gia đình Hikawa đã diễn ra như thế đấy!
\end{document}